% Options for packages loaded elsewhere
\PassOptionsToPackage{unicode}{hyperref}
\PassOptionsToPackage{hyphens}{url}
\documentclass[
]{article}
\usepackage{xcolor}
\usepackage{amsmath,amssymb}
\setcounter{secnumdepth}{-\maxdimen} % remove section numbering
\usepackage{iftex}
\ifPDFTeX
  \usepackage[T1]{fontenc}
  \usepackage[utf8]{inputenc}
  \usepackage{textcomp} % provide euro and other symbols
\else % if luatex or xetex
  \usepackage{unicode-math} % this also loads fontspec
  \defaultfontfeatures{Scale=MatchLowercase}
  \defaultfontfeatures[\rmfamily]{Ligatures=TeX,Scale=1}
\fi
\usepackage{lmodern}
\ifPDFTeX\else
  % xetex/luatex font selection
  \setmainfont[]{Times New Roman}
\fi
% Use upquote if available, for straight quotes in verbatim environments
\IfFileExists{upquote.sty}{\usepackage{upquote}}{}
\IfFileExists{microtype.sty}{% use microtype if available
  \usepackage[]{microtype}
  \UseMicrotypeSet[protrusion]{basicmath} % disable protrusion for tt fonts
}{}
\makeatletter
\@ifundefined{KOMAClassName}{% if non-KOMA class
  \IfFileExists{parskip.sty}{%
    \usepackage{parskip}
  }{% else
    \setlength{\parindent}{0pt}
    \setlength{\parskip}{6pt plus 2pt minus 1pt}}
}{% if KOMA class
  \KOMAoptions{parskip=half}}
\makeatother
\usepackage{color}
\usepackage{fancyvrb}
\newcommand{\VerbBar}{|}
\newcommand{\VERB}{\Verb[commandchars=\\\{\}]}
\DefineVerbatimEnvironment{Highlighting}{Verbatim}{commandchars=\\\{\}}
% Add ',fontsize=\small' for more characters per line
\newenvironment{Shaded}{}{}
\newcommand{\AlertTok}[1]{\textcolor[rgb]{1.00,0.00,0.00}{\textbf{#1}}}
\newcommand{\AnnotationTok}[1]{\textcolor[rgb]{0.38,0.63,0.69}{\textbf{\textit{#1}}}}
\newcommand{\AttributeTok}[1]{\textcolor[rgb]{0.49,0.56,0.16}{#1}}
\newcommand{\BaseNTok}[1]{\textcolor[rgb]{0.25,0.63,0.44}{#1}}
\newcommand{\BuiltInTok}[1]{\textcolor[rgb]{0.00,0.50,0.00}{#1}}
\newcommand{\CharTok}[1]{\textcolor[rgb]{0.25,0.44,0.63}{#1}}
\newcommand{\CommentTok}[1]{\textcolor[rgb]{0.38,0.63,0.69}{\textit{#1}}}
\newcommand{\CommentVarTok}[1]{\textcolor[rgb]{0.38,0.63,0.69}{\textbf{\textit{#1}}}}
\newcommand{\ConstantTok}[1]{\textcolor[rgb]{0.53,0.00,0.00}{#1}}
\newcommand{\ControlFlowTok}[1]{\textcolor[rgb]{0.00,0.44,0.13}{\textbf{#1}}}
\newcommand{\DataTypeTok}[1]{\textcolor[rgb]{0.56,0.13,0.00}{#1}}
\newcommand{\DecValTok}[1]{\textcolor[rgb]{0.25,0.63,0.44}{#1}}
\newcommand{\DocumentationTok}[1]{\textcolor[rgb]{0.73,0.13,0.13}{\textit{#1}}}
\newcommand{\ErrorTok}[1]{\textcolor[rgb]{1.00,0.00,0.00}{\textbf{#1}}}
\newcommand{\ExtensionTok}[1]{#1}
\newcommand{\FloatTok}[1]{\textcolor[rgb]{0.25,0.63,0.44}{#1}}
\newcommand{\FunctionTok}[1]{\textcolor[rgb]{0.02,0.16,0.49}{#1}}
\newcommand{\ImportTok}[1]{\textcolor[rgb]{0.00,0.50,0.00}{\textbf{#1}}}
\newcommand{\InformationTok}[1]{\textcolor[rgb]{0.38,0.63,0.69}{\textbf{\textit{#1}}}}
\newcommand{\KeywordTok}[1]{\textcolor[rgb]{0.00,0.44,0.13}{\textbf{#1}}}
\newcommand{\NormalTok}[1]{#1}
\newcommand{\OperatorTok}[1]{\textcolor[rgb]{0.40,0.40,0.40}{#1}}
\newcommand{\OtherTok}[1]{\textcolor[rgb]{0.00,0.44,0.13}{#1}}
\newcommand{\PreprocessorTok}[1]{\textcolor[rgb]{0.74,0.48,0.00}{#1}}
\newcommand{\RegionMarkerTok}[1]{#1}
\newcommand{\SpecialCharTok}[1]{\textcolor[rgb]{0.25,0.44,0.63}{#1}}
\newcommand{\SpecialStringTok}[1]{\textcolor[rgb]{0.73,0.40,0.53}{#1}}
\newcommand{\StringTok}[1]{\textcolor[rgb]{0.25,0.44,0.63}{#1}}
\newcommand{\VariableTok}[1]{\textcolor[rgb]{0.10,0.09,0.49}{#1}}
\newcommand{\VerbatimStringTok}[1]{\textcolor[rgb]{0.25,0.44,0.63}{#1}}
\newcommand{\WarningTok}[1]{\textcolor[rgb]{0.38,0.63,0.69}{\textbf{\textit{#1}}}}
\usepackage{longtable,booktabs,array}
\usepackage{caption}
\captionsetup[table]{skip=6pt}
\newcounter{none} % for unnumbered tables
\usepackage{calc} % for calculating minipage widths
% Correct order of tables after \paragraph or \subparagraph
\usepackage{etoolbox}
\makeatletter
\patchcmd\longtable{\par}{\if@noskipsec\mbox{}\fi\par}{}{}
\makeatother
% Allow footnotes in longtable head/foot
\IfFileExists{footnotehyper.sty}{\usepackage{footnotehyper}}{\usepackage{footnote}}
\makesavenoteenv{longtable}
\usepackage{graphicx}
\makeatletter
\newsavebox\pandoc@box
\newcommand*\pandocbounded[1]{% scales image to fit in text height/width
  \sbox\pandoc@box{#1}%
  \Gscale@div\@tempa{\textheight}{\dimexpr\ht\pandoc@box+\dp\pandoc@box\relax}%
  \Gscale@div\@tempb{\linewidth}{\wd\pandoc@box}%
  \ifdim\@tempb\p@<\@tempa\p@\let\@tempa\@tempb\fi% select the smaller of both
  \ifdim\@tempa\p@<\p@\scalebox{\@tempa}{\usebox\pandoc@box}%
  \else\usebox{\pandoc@box}%
  \fi%
}
% Set default figure placement to htbp
\def\fps@figure{htbp}
\makeatother
\setlength{\emergencystretch}{3em} % prevent overfull lines
\providecommand{\tightlist}{%
  \setlength{\itemsep}{0pt}\setlength{\parskip}{0pt}}
\usepackage{bookmark}
\IfFileExists{xurl.sty}{\usepackage{xurl}}{} % add URL line breaks if available
\urlstyle{same}
% fallback for those not using the hyperref driver hyperxmp:
\makeatletter
\@ifundefined{xmpquote}{\newcommand{\xmpquote}[1]{#1}}{}
\makeatother
\hypersetup{
  hidelinks,
  pdfcreator={LaTeX via pandoc}}

\author{}
\date{}

\begin{document}

\section{HỆ THỐNG KIẾN THỨC KINH TẾ HỌC TOÀN DIỆN (LÝ THUYẾT \& ỨNG DỤNG
PHD-LEVEL)}\label{hux1ec7-thux1ed1ng-kiux1ebfn-thux1ee9c-kinh-tux1ebf-hux1ecdc-touxe0n-diux1ec7n-luxfd-thuyux1ebft-ux1ee9ng-dux1ee5ng-phd-level}

\begin{quote}
\textbf{Lời tựa \& Hướng dẫn sử dụng:} Đây là tài liệu giáo trình và tài
liệu tham khảo chuyên sâu cấp độ Tiến sĩ (PhD-Level) tổng hợp 17 Chương,
3 Phụ lục kỹ thuật và Ngân hàng câu hỏi tự kiểm tra. Mỗi chuyên đề tuân
thủ mô hình 4 lớp lý luận: (1) Cốt lõi Lý thuyết toán học LaTeX, (2)
Điểm mù nhận thức thực tế, (3) Ứng dụng thực tiễn Việt Nam \& Quốc tế,
và (4) Ví dụ minh họa tính toán chi tiết.
\end{quote}

\begin{center}\rule{0.5\linewidth}{0.5pt}\end{center}

\subsection{SƠ ĐỒ LỘ TRÌNH ĐỌC \& PHỤ THUỘC KIẾN THỨC (PREREQUISITE
ROADMAP)}\label{sux1a1-ux111ux1ed3-lux1ed9-truxecnh-ux111ux1ecdc-phux1ee5-thuux1ed9c-kiux1ebfn-thux1ee9c-prerequisite-roadmap}

\begin{verbatim}
[Nền tảng Vi mô & Vĩ mô]
  ├── CHƯƠNG 1 (Thị trường & Giá) & CHƯƠNG 3 (Ra quyết định bất định) ──> CHƯƠNG 8 (Kinh tế công) & CHƯƠNG 12 (Thông tin bất cân xứng)
  ├── CHƯƠNG 2 (Thất bại thị trường) & CHƯƠNG 5 (Trò chơi cơ bản)     ──> CHƯƠNG 11 (Trò chơi nâng cao) & CHƯƠNG 17 (Tổ chức ngành IO)
  ├── CHƯƠNG 4 (Thương mại Melitz) & CHƯƠNG 7 (Vĩ mô Solow/IS-LM)      ──> CHƯƠNG 13 (Vĩ mô RBC/DSGE/OLG) & CHƯƠNG 14 (Phát triển thể chế)
  ├── CHƯƠNG 6 (Hành vi cơ bản)                                        ──> CHƯƠNG 16 (Hành vi nâng cao & Thần kinh Neuroeconomics)
  └── CHƯƠNG 9 (Lao động) & CHƯƠNG 10 (Tiền tệ & Tài chính)            ──> (Cầu nối giữa Vi mô và Vĩ mô thực nghiệm)

[Công cụ Định lượng & Toán học Nền tảng]
  ├── PHỤ LỤC A (Nền tảng Toán & Đối ngẫu: Shephard, Roy, Slutsky, vNM) ─> Công cụ lý thuyết cho toàn bộ 17 Chương
  ├── PHỤ LỤC B (Quy hoạch Động & Phương trình Bellman)                 ─> Công cụ giải tích cho Chương 10 (Tài chính), Chương 13 (DSGE/OLG), Chương 15.8 (VAR/VECM)
  └── CHƯƠNG 15 (Kinh tế lượng & Suy luận Nhân quả DiD/RDD/Panel FE-RE) ─> Công cụ kiểm định thực nghiệm cho mọi mô hình lý thuyết
\end{verbatim}

\begin{center}\rule{0.5\linewidth}{0.5pt}\end{center}

\subsection{MỤC LỤC TỔNG HỢP (TABLE OF
CONTENTS)}\label{mux1ee5c-lux1ee5c-tux1ed5ng-hux1ee3p-table-of-contents}

\begin{itemize}
\tightlist
\item
  \hyperref[phux1ea7n-a--nux1ec1n-tux1ea3ng-kinh-tux1ebf-hux1ecdc]{PHẦN
  A --- NỀN TẢNG KINH TẾ HỌC}

  \begin{itemize}
  \tightlist
  \item
    \hyperref[chux1b0ux1a1ng-1-thux1ecb-trux1b0ux1eddng--giuxe1-markets--prices]{CHƯƠNG
    1: THỊ TRƯỜNG \& GIÁ (Markets \& Prices)}
  \item
    \hyperref[chux1b0ux1a1ng-2-thux1ea5t-bux1ea1i-thux1ecb-trux1b0ux1eddng-market-failures]{CHƯƠNG
    2: THẤT BẠI THỊ TRƯỜNG (Market Failures)}
  \item
    \hyperref[chux1b0ux1a1ng-3-ra-quyux1ebft-ux111ux1ecbnh-trong-houxe0n-cux1ea3nh-bux1ea5t-ux111ux1ecbnh--thux1eddi-gian-decision-making]{CHƯƠNG
    3: RA QUYẾT ĐỊNH TRONG HOÀN CẢNH BẤT ĐỊNH \& THỜI GIAN (Decision
    Making)}
  \item
    \hyperref[chux1b0ux1a1ng-4-thux1b0ux1a1ng-mux1ea1i--vux129-muxf4-quux1ed1c-tux1ebf-international-trade--macroeconomics]{CHƯƠNG
    4: THƯƠNG MẠI \& VĨ MÔ QUỐC TẾ (International Trade \&
    Macroeconomics)}
  \item
    \hyperref[chux1b0ux1a1ng-5-tux1b0ux1a1ng-tuxe1c-chiux1ebfn-lux1b0ux1ee3c--luxfd-thuyux1ebft-truxf2-chux1a1i-cux1a1-bux1ea3n-strategic-interaction--game-theory]{CHƯƠNG
    5: TƯƠNG TÁC CHIẾN LƯỢC \& LÝ THUYẾT TRÒ CHƠI CƠ BẢN (Strategic
    Interaction \& Game Theory)}
  \item
    \hyperref[chux1b0ux1a1ng-6-kinh-tux1ebf-hux1ecdc-huxe0nh-vi-cux1a1-bux1ea3n-basic-behavioral-economics]{CHƯƠNG
    6: KINH TẾ HỌC HÀNH VI CƠ BẢN (Basic Behavioral Economics)}
  \item
    \hyperref[chux1b0ux1a1ng-7-vux129-muxf4-tux1ed5ng-quuxe1t-general-macroeconomics]{CHƯƠNG
    7: VĨ MÔ TỔNG QUÁT (General Macroeconomics)}
  \end{itemize}
\item
  \hyperref[phux1ea7n-b--chuyuxean-ux111ux1ec1-nuxe2ng-cao-phd-level]{PHẦN
  B --- CHUYÊN ĐỀ NÂNG CAO (PHD LEVEL)}

  \begin{itemize}
  \tightlist
  \item
    \hyperref[chux1b0ux1a1ng-8-kinh-tux1ebf-hux1ecdc-cuxf4ng-cux1ed9ng--lux1ef1a-chux1ecdn-xuxe3-hux1ed9i-public-economics--social-choice]{CHƯƠNG
    8: KINH TẾ HỌC CÔNG CỘNG \& LỰA CHỌN XÃ HỘI (Public Economics \&
    Social Choice)}
  \item
    \hyperref[chux1b0ux1a1ng-9-kinh-tux1ebf-hux1ecdc-lao-ux111ux1ed9ng--vux1ed1n-nhuxe2n-lux1ef1c-labor-economics--human-capital]{CHƯƠNG
    9: KINH TẾ HỌC LAO ĐỘNG \& VỐN NHÂN LỰC (Labor Economics \& Human
    Capital)}
  \item
    \hyperref[chux1b0ux1a1ng-10-kinh-tux1ebf-hux1ecdc-tiux1ec1n-tux1ec7--tuxe0i-chuxednh-monetary--financial-economics]{CHƯƠNG
    10: KINH TẾ HỌC TIỀN TỆ \& TÀI CHÍNH (Monetary \& Financial
    Economics)}
  \item
    \hyperref[chux1b0ux1a1ng-11-luxfd-thuyux1ebft-truxf2-chux1a1i-nuxe2ng-cao--thiux1ebft-kux1ebf-cux1a1-chux1ebf]{CHƯƠNG
    11: LÝ THUYẾT TRÒ CHƠI NÂNG CAO \& THIẾT KẾ CƠ CHẾ (Advanced Game
    Theory \& Mechanism Design)}
  \item
    \hyperref[chux1b0ux1a1ng-12-thuxf4ng-tin-bux1ea5t-cuxe2n-xux1ee9ng--luxfd-thuyux1ebft-hux1ee3p-ux111ux1ed3ng-information-economics--contract-theory]{CHƯƠNG
    12: THÔNG TIN BẤT CÂN XỨNG \& LÝ THUYẾT HỢP ĐỒNG (Information
    Economics \& Contract Theory)}
  \item
    \hyperref[chux1b0ux1a1ng-13-kinh-tux1ebf-hux1ecdc-vux129-muxf4-nuxe2ng-cao-advanced-macroeconomics]{CHƯƠNG
    13: KINH TẾ HỌC VĨ MÔ NÂNG CAO (Advanced Macroeconomics)}
  \item
    \hyperref[chux1b0ux1a1ng-14-kinh-tux1ebf-hux1ecdc-phuxe1t-triux1ec3n--thux1ec3-chux1ebf-development--institutional-economics]{CHƯƠNG
    14: KINH TẾ HỌC PHÁT TRIỂN \& THỂ CHẾ (Development \& Institutional
    Economics)}
  \item
    \hyperref[chux1b0ux1a1ng-15-kinh-tux1ebf-lux1b0ux1ee3ng--suy-luux1eadn-nhuxe2n-quux1ea3-econometrics--causal-inference]{CHƯƠNG
    15: KINH TẾ LƯỢNG \& SUY LUẬN NHÂN QUẢ (Econometrics \& Causal
    Inference)}
  \item
    \hyperref[chux1b0ux1a1ng-16-kinh-tux1ebf-hux1ecdc-huxe0nh-vi-nuxe2ng-cao--thux1ea7n-kinh-advanced-behavioral--neuroeconomics]{CHƯƠNG
    16: KINH TẾ HỌC HÀNH VI NÂNG CAO \& THẦN KINH (Advanced Behavioral
    \& Neuroeconomics)}
  \item
    \hyperref[chux1b0ux1a1ng-17-kinh-tux1ebf-hux1ecdc-tux1ed5-chux1ee9c-nguxe0nh-nuxe2ng-cao-advanced-industrial-organization]{CHƯƠNG
    17: KINH TẾ HỌC TỔ CHỨC NGÀNH NÂNG CAO (Advanced Industrial
    Organization)}
  \end{itemize}
\item
  \hyperref[phux1ee5-lux1ee5c-a-nux1ec1n-tux1ea3ng-touxe1n-hux1ecdc--luxfd-thuyux1ebft-ux111ux1ed1i-ngux1eabu-trong-vi-muxf4vux129-muxf4-phd]{PHỤ
  LỤC VÀ TÀI LIỆU THAM KHẢO}

  \begin{itemize}
  \tightlist
  \item
    \hyperref[phux1ee5-lux1ee5c-a-nux1ec1n-tux1ea3ng-touxe1n-hux1ecdc--luxfd-thuyux1ebft-ux111ux1ed1i-ngux1eabu-trong-vi-muxf4vux129-muxf4-phd]{PHỤ
    LỤC A: NỀN TẢNG TOÁN HỌC \& LÝ THUYẾT ĐỐI NGẪU TRONG VI MÔ/VĨ MÔ
    PHD}
  \item
    \hyperref[phux1ee5-lux1ee5c-b-quy-houx1ea1ch-ux111ux1ed9ng--phux1b0ux1a1ng-truxecnh-bellman-dynamic-programming--bellman]{PHỤ
    LỤC B: QUY HOẠCH ĐỘNG \& PHƯƠNG TRÌNH BELLMAN (Dynamic Programming
    \& Bellman)}
  \item
    \hyperref[phux1ee5-lux1ee5c-c-nguxe2n-huxe0ng-cuxe2u-hux1ecfi-tux1ef1-kiux1ec3m-tra--uxf4n-luyux1ec7n-tux1eeb-ux111iux1ec3m-muxf9-nhux1eadn-thux1ee9c-self-assessment-bank]{PHỤ
    LỤC C: NGÂN HÀNG CÂU HỎI TỰ KIỂM TRA \& ÔN LUYỆN TỪ ĐIỂM MÙ NHẬN
    THỨC}
  \item
    \hyperref[danh-mux1ee5c-tuxe0i-liux1ec7u-tham-khux1ea3o-tux1ed5ng-hux1ee3p-comprehensive-bibliography]{DANH
    MỤC TÀI LIỆU THAM KHẢO TỔNG HỢP (COMPREHENSIVE BIBLIOGRAPHY)}
  \end{itemize}
\end{itemize}

\begin{center}\rule{0.5\linewidth}{0.5pt}\end{center}

\section{PHẦN A --- NỀN TẢNG KINH TẾ
HỌC}\label{phux1ea7n-a-nux1ec1n-tux1ea3ng-kinh-tux1ebf-hux1ecdc}

\section{HỆ THỐNG KIẾN THỨC KINH TẾ HỌC TOÀN DIỆN (LÝ THUYẾT \& ỨNG
DỤNG)}\label{hux1ec7-thux1ed1ng-kiux1ebfn-thux1ee9c-kinh-tux1ebf-hux1ecdc-touxe0n-diux1ec7n-luxfd-thuyux1ebft-ux1ee9ng-dux1ee5ng}

\textbf{Biên soạn theo tiêu chuẩn Tiến sĩ Kinh tế học (PhD Level)}
\emph{Phần A: Nền tảng và Ứng dụng · Phần B: Chuyên đề Nâng cao \& Hình
thức hóa Toán học}

\begin{center}\rule{0.5\linewidth}{0.5pt}\end{center}

\begin{center}\rule{0.5\linewidth}{0.5pt}\end{center}

\section{CHƯƠNG 1: THỊ TRƯỜNG \& GIÁ (Markets \&
Prices)}\label{chux1b0ux1a1ng-1-thux1ecb-trux1b0ux1eddng-giuxe1-markets-prices}

Nhóm nền tảng nhất: giải thích giá cả hình thành thế nào và vì sao thị
trường tự do thường phân bổ nguồn lực hiệu quả. Đây là bộ công cụ để đọc
mọi biến động giá trong đời sống --- từ giá xăng, giá nhà đến lương nhân
sự.

\begin{center}\rule{0.5\linewidth}{0.5pt}\end{center}

\subsection{Cung -- Cầu (Supply \&
Demand)}\label{cung-cux1ea7u-supply-demand}

\textbf{Cốt lõi:} Giá và sản lượng cân bằng của một thị trường được xác
định tại giao điểm giữa đường cung (lượng người bán sẵn sàng bán ở mỗi
mức giá) và đường cầu (lượng người mua sẵn sàng mua ở mỗi mức giá). Khi
một trong hai đường dịch chuyển do các yếu tố ngoài giá (thu nhập, thị
hiếu, chi phí đầu vào, kỳ vọng\ldots), điểm cân bằng mới sẽ hình thành.
Cơ chế điều chỉnh: nếu giá tạm thời cao hơn mức cân bằng, dư cung xuất
hiện, người bán buộc phải hạ giá để bán hết hàng tồn; nếu giá thấp hơn,
thiếu hụt xuất hiện, người mua sẵn sàng trả cao hơn để giành hàng --- cả
hai lực đều đẩy giá VỀ điểm cân bằng.

\textbf{Điểm mù thường gặp:} Khi đọc một tin biến động giá (xăng tăng,
nhà đất sốt\ldots), phản xạ đầu tiên thường là kết luận cảm tính ngay
(``chắc do nhu cầu tăng'') mà quên dừng lại đặt đúng câu hỏi phân tích:
yếu tố nào đang dịch chuyển --- đường cung hay đường cầu --- trước khi
suy ra nguyên nhân thực sự. Một điểm mù khác: khi CẢ hai đường cùng dịch
chuyển đồng thời (ví dụ vừa tăng thu nhập vừa tăng chi phí sản xuất),
chiều thay đổi của GIÁ có thể dự đoán chắc chắn nhưng chiều thay đổi của
LƯỢNG lại không xác định được nếu không biết mức độ dịch chuyển tương
đối của mỗi đường --- kết luận vội ``lượng chắc cũng tăng'' là sai lầm
phổ biến.

\textbf{Ứng dụng thực tế:} Khi giá xăng thế giới tăng đột biến, chi phí
vận chuyển (một đầu vào) tăng khiến đường cung nhiều loại hàng hóa dịch
trái --- giá bán lẻ tăng dù nhu cầu tiêu dùng không đổi. Phân biệt đúng
nguyên nhân giúp doanh nghiệp chọn đúng biện pháp ứng phó: đàm phán lại
nhà cung cấp thay vì cắt giảm marketing.

\textbf{Ví dụ minh hoạ:} Giả sử vé xem phim có phương trình cầu
\(Q_d = 100 - 2P\) và phương trình cung \(Q_s = -20 + 3P\) (\(P\) tính
bằng nghìn đồng/vé, \(Q\) tính bằng nghìn vé/tháng). Cân bằng:
\(Q_d = Q_s \Rightarrow 100 - 2P = -20 + 3P \Rightarrow 120 = 5P \Rightarrow P = 24\),
\(Q = 52\). Nếu thu nhập khán giả tăng khiến cầu dịch phải thành
\(Q_d' = 130 - 2P\) (mua nhiều hơn ở MỌI mức giá), cân bằng mới:
\(130 - 2P = -20 + 3P \Rightarrow 150 = 5P \Rightarrow P = 30\),
\(Q = 70\) --- cả giá lẫn lượng cân bằng đều tăng, đúng như dự đoán khi
cầu dịch phải trong lúc cung giữ nguyên.

\begin{center}\rule{0.5\linewidth}{0.5pt}\end{center}

\subsection{Độ co giãn
(Elasticity)}\label{ux111ux1ed9-co-giuxe3n-elasticity}

\textbf{Cốt lõi:} Đo mức độ nhạy cảm của lượng cầu (hoặc cung) khi giá
thay đổi 1\%, tính bằng công thức \(E = (\%\Delta Q) / (\%\Delta P)\).
Cầu co giãn nhiều (\(|E| > 1\)) nghĩa là người mua dễ dàng chuyển sang
lựa chọn khác khi giá tăng; cầu co giãn ít (\(|E| < 1\)) nghĩa là người
mua vẫn mua gần như cũ dù giá đổi --- thường gặp ở hàng thiết yếu, ít
hàng thay thế. Độ co giãn quyết định trực tiếp doanh thu phản ứng thế
nào khi đổi giá: nếu \(|E| > 1\), tăng giá làm TỔNG DOANH THU giảm (vì
lượng bán giảm mạnh hơn phần trăm giá tăng); nếu \(|E| < 1\), tăng giá
làm tổng doanh thu tăng.

\textbf{Điểm mù thường gặp:} Khi chi phí đầu vào tăng và cần điều chỉnh
giá bán, phản xạ thường là tăng giá đồng loạt mọi sản phẩm mà quên dừng
lại hỏi: sản phẩm nào khách hàng dễ tìm được lựa chọn thay thế? Bỏ qua
câu hỏi này khiến chiến lược giá phản tác dụng đúng ở những sản phẩm co
giãn nhiều nhất. Một điểm mù khác: độ co giãn không phải hằng số cố định
của một sản phẩm --- nó thay đổi dọc theo đường cầu (thường co giãn
nhiều hơn ở vùng giá cao, ít hơn ở vùng giá thấp) và thay đổi theo khung
thời gian (co giãn theo giá xăng gần như bằng 0 trong vài ngày nhưng lớn
hơn nhiều sau vài năm khi người tiêu dùng đổi sang xe tiết kiệm nhiên
liệu).

\textbf{Ứng dụng thực tế:} Một quán cà phê có thể tăng giá cà phê đặc
sản (cầu co giãn ít, khách trung thành, ít nơi thay thế) nhưng nên giữ
nguyên giá nước lọc/trà đá (cầu co giãn nhiều, dễ bị so sánh) để tối ưu
tổng doanh thu thay vì tăng đều mọi món.

\textbf{Ví dụ minh hoạ:} Một cửa hàng tăng giá một mặt hàng từ 20.000đ
lên 22.000đ (tăng 10\%), lượng bán giảm từ 1.000 xuống 850 sản
phẩm/tháng (giảm 15\%) \(\Rightarrow E = -15\% / 10\% = -1{,}5\) (cầu co
giãn nhiều). Doanh thu trước: \(20.000 \times 1.000 = 20\) triệu đồng;
doanh thu sau: \(22.000 \times 850 = 18{,}7\) triệu đồng --- doanh thu
GIẢM dù giá tăng, đúng dự đoán khi \(|E| > 1\). Nếu cửa hàng biết trước
độ co giãn này, quyết định đúng đắn hơn là KHÔNG tăng giá mặt hàng đó.

\begin{center}\rule{0.5\linewidth}{0.5pt}\end{center}

\subsection{Thặng dư tiêu dùng / sản xuất (Consumer \& Producer
Surplus)}\label{thux1eb7ng-dux1b0-tiuxeau-duxf9ng-sux1ea3n-xuux1ea5t-consumer-producer-surplus}

\textbf{Cốt lõi:} Thặng dư tiêu dùng là phần chênh lệch giữa mức giá
người mua sẵn sàng trả và mức giá thực trả (họ được lợi). Thặng dư sản
xuất là chênh lệch giữa giá bán thực nhận và chi phí tối thiểu người bán
chấp nhận. Tổng hai loại thặng dư đo phúc lợi xã hội tạo ra từ một giao
dịch --- về mặt hình học, đó là diện tích hai tam giác nằm giữa đường
cầu/cung và mức giá cân bằng, tính bằng công thức diện tích tam giác:
\(\frac{1}{2} \times \text{đáy} \times \text{chiều cao}\).

\textbf{Điểm mù thường gặp:} Khi đánh giá một chính sách (áp trần giá,
đánh thuế, trợ giá), phản xạ thường dừng lại ở câu hỏi ``ai được lợi''
mà quên đi tìm phần ``tổn thất vô ích'' (deadweight loss) --- phần thặng
dư biến mất hoàn toàn chứ không chuyển từ bên này sang bên kia. Một điểm
mù khác: thặng dư tiêu dùng đo bằng TIỀN không phải lúc nào cũng phản
ánh đúng mức độ ``hạnh phúc'' tăng thêm --- một khoản thặng dư 100.000đ
có ý nghĩa khác nhau với người thu nhập thấp so với người thu nhập cao
(độ thoả dụng biên của tiền giảm dần), điều mà mô hình thặng dư cổ điển
không tính đến.

\textbf{Ứng dụng thực tế:} Áp trần giá thuê nhà giúp người thuê hiện tại
được lợi (thặng dư tiêu dùng cao hơn) nhưng làm giảm nguồn cung nhà cho
thuê mới, tạo tổn thất vô ích cho cả xã hội --- đây là lý do nhiều nhà
kinh tế cảnh báo về hệ quả dài hạn của kiểm soát giá.

\textbf{Ví dụ minh hoạ:} Dùng lại ví dụ vé xem phim (\(Q_d = 100 - 2P\),
\(Q_s = -20 + 3P\), cân bằng \(P = 24\), \(Q = 52\)): tại \(Q = 0\),
đường cầu cho \(P = 50\) (mức giá cao nhất người mua sẵn sàng trả)
\(\Rightarrow\) thặng dư tiêu dùng
\(= \frac{1}{2} \times 52 \times (50 - 24) = \frac{1}{2} \times 52 \times 26 = 676\)
(nghìn đồng). Đường cung tại \(Q = 0\) cho \(P \approx 6{,}67\) (mức giá
thấp nhất người bán chấp nhận) \(\Rightarrow\) thặng dư sản xuất
\(= \frac{1}{2} \times 52 \times (24 - 6{,}67) \approx \frac{1}{2} \times 52 \times 17{,}33 \approx 450{,}7\)
(nghìn đồng). Tổng phúc lợi xã hội tại điểm cân bằng
\(\approx 1.126{,}7\) (nghìn đồng) --- đây chính là mức TỐI ĐA có thể
đạt được; bất kỳ can thiệp nào đẩy lượng giao dịch ra khỏi \(Q = 52\)
(dù bằng giá trần, giá sàn hay thuế) đều làm tổng này giảm.

\begin{center}\rule{0.5\linewidth}{0.5pt}\end{center}

\subsection{Giá trần \& giá sàn (Price ceiling \& Price
floor)}\label{giuxe1-trux1ea7n-giuxe1-suxe0n-price-ceiling-price-floor}

\textbf{Cốt lõi:} Giá trần là mức giá tối đa hợp pháp, đặt DƯỚI giá cân
bằng để bảo vệ người mua (ví dụ giá thuê nhà, giá xăng thời khủng hoảng)
--- gây ra THIẾU HỤT vì tại mức giá thấp đó người mua muốn mua nhiều hơn
người bán muốn bán. Giá sàn là mức giá tối thiểu hợp pháp, đặt TRÊN giá
cân bằng để bảo vệ người bán (ví dụ lương tối thiểu, giá thu mua nông
sản) --- gây ra DƯ THỪA vì lượng cung vượt lượng cầu ở mức giá đó. Cả
hai chỉ có hiệu lực (ràng buộc thị trường) khi đặt LỆCH khỏi giá cân
bằng tự nhiên --- một giá trần cao hơn giá cân bằng hoặc giá sàn thấp
hơn giá cân bằng hoàn toàn không có tác dụng gì.

\textbf{Điểm mù thường gặp:} Trực giác thường nghĩ ``áp giá trần thấp
hơn = người mua luôn có lợi'', nhưng bỏ qua rằng thiếu hụt buộc phải
phân bổ hàng hoá bằng cơ chế KHÁC ngoài giá --- xếp hàng, quen biết,
hoặc ``chợ đen'' --- người thực sự nhận được hàng ở giá thấp không nhất
thiết là người cần nhất, mà là người đến sớm nhất hoặc quen biết nhất.
Tương tự, giá sàn (như lương tối thiểu) không tự động có lợi cho TOÀN BỘ
người lao động --- những người vẫn giữ được việc làm có lợi, nhưng phần
dư thừa lao động (thất nghiệp tăng thêm) lại là chi phí của chính sách,
một đánh đổi hay bị bỏ qua trong tranh luận đơn giản hoá.

\textbf{Ứng dụng thực tế:} Kiểm soát giá thuê nhà (rent control) tại
nhiều thành phố lớn --- trong ngắn hạn người thuê hiện tại có lợi, nhưng
dài hạn chủ nhà giảm đầu tư bảo trì/xây mới, làm trầm trọng thêm khan
hiếm nhà cho người thuê MỚI, đây là lý do nhiều nhà kinh tế (không phân
biệt khuynh hướng chính trị) đồng thuận phản đối kiểm soát giá thuê nhà
cứng nhắc kéo dài.

\textbf{Ví dụ minh hoạ:} Với \(Q_d = 100 - 2P\), \(Q_s = -20 + 3P\) (cân
bằng \(P = 24\), \(Q = 52\)), nếu chính phủ áp giá trần \(P = 18\):
\(Q_d = 100 - 36 = 64\), \(Q_s = -20 + 54 = 34 \Rightarrow\) thiếu hụt
\(= 64 - 34 = 30\) (nghìn vé) --- chỉ 34 nghìn người mua được vé dù 64
nghìn người muốn mua ở mức giá đó, 30 nghìn người còn lại phải xếp hàng,
tìm quan hệ, hoặc mua chợ đen với giá cao hơn nhiều mức trần.

\begin{center}\rule{0.5\linewidth}{0.5pt}\end{center}

\subsection{Ai thực sự chịu thuế? (Tax
incidence)}\label{ai-thux1ef1c-sux1ef1-chux1ecbu-thuux1ebf-tax-incidence}

\textbf{Cốt lõi:} Khi chính phủ đánh thuế một hàng hoá, gánh nặng thuế
thực tế (ai trả bao nhiêu) KHÔNG phụ thuộc việc luật quy định người bán
hay người mua phải nộp thuế, mà phụ thuộc ĐỘ CO GIÃN TƯƠNG ĐỐI của cung
và cầu: bên nào co giãn ÍT hơn (khó thay đổi hành vi hơn khi giá đổi) sẽ
gánh phần thuế LỚN hơn, bất kể ai là người trực tiếp nộp tiền cho cơ
quan thuế. Đây là hệ quả trực tiếp của mô hình cung-cầu: thuế chèn một
``khoảng cách'' giữa giá người mua trả và giá người bán nhận, và thị
trường tự phân bổ khoảng cách đó theo độ co giãn.

\textbf{Điểm mù thường gặp:} Tranh luận chính sách thường tập trung vào
``đánh thuế vào ai'' (thuế VAT người bán nộp hay thuế thu nhập người lao
động nộp) như thể điều đó quyết định ai chịu thiệt --- bỏ qua rằng về
mặt kinh tế, việc luật quy định ai nộp thuế trên giấy tờ hầu như không
ảnh hưởng đến người thực sự chịu gánh nặng, chỉ ảnh hưởng ai phải làm
thủ tục hành chính nộp thuế. Đánh thuế người bán hay người mua, xét về
mặt toán học, cho ra CÙNG một kết quả cân bằng cuối cùng.

\textbf{Ứng dụng thực tế:} Thuế thuốc lá thường được thiết kế đánh vào
nhà sản xuất, nhưng vì cầu thuốc lá co giãn RẤT ÍT (người nghiện khó bỏ
dù giá tăng), gần như toàn bộ thuế được chuyển sang người tiêu dùng qua
giá bán cao hơn --- ngược lại, thuế đánh vào hàng xa xỉ dễ thay thế (cầu
co giãn nhiều) sẽ do người bán/nhà sản xuất gánh phần lớn vì tăng giá sẽ
mất khách nhanh chóng.

\textbf{Ví dụ minh hoạ:} Với \(Q_d = 100 - 2P\), \(Q_s = -20 + 3P\) (cân
bằng gốc \(P = 24\)), giả sử chính phủ đánh thuế 5 (nghìn đồng)/vé thu
từ người bán. Người bán chỉ còn giữ lại \((P - 5)\) mỗi vé nên cung mới:
\(Q_s = -20 + 3(P-5) = -35 + 3P\). Cân bằng mới:
\(100 - 2P = -35 + 3P \Rightarrow 135 = 5P \Rightarrow P = 27\) (giá
người mua trả); người bán thực nhận \(27 - 5 = 22\). So với giá gốc 24:
người mua trả THÊM 3, người bán nhận ÍT HƠN 2 (tổng \(3+2=5\) đúng bằng
mức thuế). Dù thuế thu từ người bán trên giấy tờ, người MUA lại chịu
phần lớn hơn (3 so với 2) --- vì tại vùng giá này, cầu co giãn ít hơn
cung một chút (độ co giãn cầu \(\approx -0{,}92\) so với cung
\(\approx 1{,}38\) tại điểm cân bằng gốc).

\begin{center}\rule{0.5\linewidth}{0.5pt}\end{center}

\subsection{Di chuyển trên đường vs dịch chuyển đường (Movement along vs
Shift)}\label{di-chuyux1ec3n-truxean-ux111ux1b0ux1eddng-vs-dux1ecbch-chuyux1ec3n-ux111ux1b0ux1eddng-movement-along-vs-shift}

\textbf{Cốt lõi:} Phân biệt hai loại thay đổi thường bị gộp làm một.
``Di chuyển DỌC đường cầu/cung'' (movement along the curve) xảy ra CHỈ
khi giá của chính hàng hoá đó thay đổi --- đường cầu/cung giữ nguyên,
chỉ điểm trên đường thay đổi. ``Dịch chuyển CẢ đường cầu/cung'' (shift
of the curve) xảy ra khi một yếu tố NGOÀI giá thay đổi (thu nhập, giá
hàng liên quan, thị hiếu, kỳ vọng, công nghệ, số người mua/bán trên thị
trường) --- toàn bộ đường di chuyển sang vị trí mới, nghĩa là ở CÙNG một
mức giá, lượng mua/bán mong muốn đã khác trước.

\textbf{Điểm mù thường gặp:} Cụm từ hàng ngày ``cầu tăng'' bị dùng lẫn
lộn cho hai trường hợp hoàn toàn khác nhau: giá giảm khiến lượng cầu
tăng (di chuyển dọc đường cầu CŨ, đường cầu không đổi) và thu nhập tăng
khiến toàn bộ đường cầu dịch phải (ở MỌI mức giá, người ta mua nhiều
hơn). Nhầm lẫn hai loại ``tăng'' này khiến phân tích nguyên nhân-kết quả
sai ngay từ bước đầu, vì mỗi loại đòi hỏi một cách diễn giải và một hành
động đáp trả hoàn toàn khác nhau.

\textbf{Ứng dụng thực tế:} Khi phân tích vì sao doanh số một sản phẩm
tăng, cần hỏi trước tiên: có phải do CHÍNH doanh nghiệp giảm giá (di
chuyển dọc đường cầu, doanh thu có thể tăng hoặc giảm tuỳ độ co giãn)
hay do yếu tố bên ngoài như xu hướng thị trường/thu nhập tăng (dịch
chuyển đường cầu, doanh số tăng ở MỌI mức giá kể cả khi không đổi giá)
--- hai nguyên nhân đòi hỏi chiến lược hoàn toàn khác: một bên là quyết
định giá của chính mình, bên kia là tín hiệu thị trường ngoài tầm kiểm
soát giá.

\textbf{Ví dụ minh hoạ:} Doanh nghiệp X giảm giá sản phẩm từ 100.000đ
xuống 90.000đ, doanh số tăng từ 500 lên 600 sản phẩm/tháng --- đây là DI
CHUYỂN dọc đường cầu (do chính X thay đổi giá, đường cầu không đổi).
Tháng sau, dù X giữ nguyên giá 90.000đ, doanh số vẫn tiếp tục tăng lên
750 vì đối thủ cạnh tranh chính đóng cửa --- đây là DỊCH CHUYỂN đường
cầu (nguyên nhân đến từ bên ngoài quyết định giá của X, đường cầu dịch
phải hoàn toàn).

\begin{center}\rule{0.5\linewidth}{0.5pt}\end{center}

\subsection{Sản lượng tối ưu \& quyết định có nên sản xuất (Optimal
Output \& Shutdown
Decision)}\label{sux1ea3n-lux1b0ux1ee3ng-tux1ed1i-ux1b0u-quyux1ebft-ux111ux1ecbnh-cuxf3-nuxean-sux1ea3n-xuux1ea5t-optimal-output-shutdown-decision}

\textbf{Cốt lõi:} Doanh nghiệp tối đa hoá lợi nhuận bằng cách sản xuất
đến mức sản lượng nơi DOANH THU BIÊN (\(MR\) --- doanh thu tăng thêm từ
một đơn vị bán thêm; trong thị trường cạnh tranh
\(MR = \text{giá bán } P\)) bằng CHI PHÍ BIÊN (\(MC\) --- chi phí tăng
thêm để sản xuất một đơn vị nữa): nếu \(MR > MC\), sản xuất thêm vẫn có
lợi; nếu \(MR < MC\), đơn vị đó đang lỗ --- quy tắc \(MR = MC\) luôn xác
định đúng điểm dừng tối ưu. Câu hỏi thứ hai, độc lập với câu hỏi trên:
có nên sản xuất chút nào không? Trong ngắn hạn, so giá bán \(P\) với CHI
PHÍ BIẾN ĐỔI BÌNH QUÂN (\(AVC\)) --- nếu \(P \ge AVC\), vẫn nên tiếp tục
sản xuất dù đang lỗ (vì phần chênh giúp bù bớt chi phí cố định); nếu
\(P < AVC\), nên đóng cửa ngay lập tức. Trong dài hạn, ngưỡng chuyển
thành so \(P\) với TỔNG CHI PHÍ BÌNH QUÂN (\(ATC\)) --- \(P < ATC\) kéo
dài là tín hiệu nên rút khỏi ngành.

\textbf{Điểm mù thường gặp:} Nhầm lẫn phổ biến nhất: coi ``đang lỗ''
đồng nghĩa với ``nên đóng cửa ngay''. Một doanh nghiệp có thể đang lỗ
(\(P < ATC\)) nhưng vẫn nên TIẾP TỤC sản xuất trong ngắn hạn nếu \(P\)
vẫn lớn hơn \(AVC\) --- vì đóng cửa hoàn toàn sẽ mất luôn phần bù đắp
chi phí cố định (khoản phải trả dù có sản xuất hay không, xem Chương 3
--- chi phí chìm). Chỉ khi \(P\) tụt xuống dưới \(AVC\), tiếp tục sản
xuất mới thực sự tệ hơn đóng cửa. Điểm mù thứ hai: nhầm mục tiêu ``tối
đa hoá lợi nhuận'' với ``tối thiểu hoá chi phí bình quân trên một đơn
vị'' --- sản xuất đến mức \(ATC\) thấp nhất thường đã VƯỢT QUA điểm
\(MR = MC\) từ lâu, tức đang sản xuất NHIỀU HƠN mức tối ưu.

\textbf{Ứng dụng thực tế:} Nhiều hãng hàng không trong giai đoạn khủng
hoảng vẫn duy trì một số chuyến bay dù bán vé dưới giá hoà vốn --- vì
giá vé vẫn đủ bù chi phí nhiên liệu/phi hành đoàn biến đổi (\(AVC\)) và
góp thêm một phần bù chi phí thuê máy bay cố định, tốt hơn để máy bay
nằm không hoàn toàn không sinh doanh thu. Một xưởng may cân nhắc nhận
thêm đơn hàng nên hỏi đơn hàng TIẾP THEO (chi phí biên để may thêm) có
đem lại doanh thu biên cao hơn chi phí biên hay không --- chứ không hỏi
giá trung bình các đơn trước đang lời bao nhiêu.

\textbf{Ví dụ minh hoạ:} Với chi phí biên \(MC = 3 + 0{,}8Q\) và giá thị
trường \(P = 12\): giải
\(3 + 0{,}8Q = 12 \Rightarrow Q^* \approx 11{,}25\) đơn vị --- sản lượng
tối đa hoá lợi nhuận. Riêng câu hỏi có nên sản xuất: giả sử chi phí cố
định \(= 20\); tại \(Q = 10\), \(AVC = 3 + 0{,}4 \times 10 = 7\) và
\(ATC = 20/10 + 3 + 0{,}4 \times 10 = 9\). Nếu giá thị trường tụt còn
\(P = 5\): vì \(AVC (7) > 5\), giá không đủ bù cả chi phí biến đổi
\(\Rightarrow\) nên đóng cửa ngay dù đã lỡ chi 20 cho chi phí cố định.
Nếu giá là \(P = 8\): \(AVC (7) < 8 < ATC (9) \Rightarrow\) vẫn lỗ trên
tổng thể nhưng NÊN tiếp tục sản xuất ngắn hạn vì đang bù được một phần
chi phí cố định.

\begin{center}\rule{0.5\linewidth}{0.5pt}\end{center}

\section{CHƯƠNG 2: THẤT BẠI THỊ TRƯỜNG (Market
Failures)}\label{chux1b0ux1a1ng-2-thux1ea5t-bux1ea1i-thux1ecb-trux1b0ux1eddng-market-failures}

Thị trường tự do không phải lúc nào cũng hiệu quả. Nhóm này hệ thống 5
lý do kinh điển khiến thị trường ``thất bại'' --- và cũng là lý do chính
đáng cho sự can thiệp của nhà nước hoặc thiết kế hợp đồng/thể chế.

\begin{center}\rule{0.5\linewidth}{0.5pt}\end{center}

\subsection{Ngoại tác
(Externalities)}\label{ngoux1ea1i-tuxe1c-externalities}

\textbf{Cốt lõi:} Xảy ra khi hành động của một bên tạo ra chi phí hoặc
lợi ích cho bên thứ ba không tham gia giao dịch, mà không được phản ánh
vào giá. Ngoại tác tiêu cực (ô nhiễm) khiến chi phí biên xã hội
(\(MSC = \text{MPC} + \text{MEC}\)) cao hơn chi phí biên tư nhân --- thị
trường tự do sản xuất tại điểm \(D = MPC\), NHIỀU hơn mức tối ưu xã hội
(nơi \(D = MSC\)). Ngoại tác tích cực (giáo dục, tiêm chủng) khiến lợi
ích biên xã hội cao hơn lợi ích biên tư nhân, khiến thị trường cung cấp
ÍT hơn mức tối ưu xã hội.

\textbf{Điểm mù thường gặp:} Chỉ nghĩ ngay đến ngoại tác môi trường mà
quên nhận ra ngoại tác cũng xảy ra trong công việc hằng ngày --- một
nhân viên trễ deadline tạo ngoại tác tiêu cực lên cả nhóm, nhưng vì
không mang tên ``ô nhiễm'' nên hiếm khi được nhận diện và tính chi phí
đúng như bản chất của nó. Điểm mù thứ hai: thuế Pigou đúng bằng \(MEC\)
chỉ khôi phục hiệu quả nếu \(MEC\) không đổi theo \(Q\) hoặc theo NGƯỜI
gây ra --- trong thực tế các nguồn ô nhiễm khác nhau (nhà máy cũ/mới) có
\(MEC\) khác nhau, nên một mức thuế đồng nhất cho cả ngành chỉ là gần
đúng, không tối ưu tuyệt đối cho từng doanh nghiệp.

\textbf{Ứng dụng thực tế:} Thuế carbon là công cụ kinh điển để ``nội
hoá'' ngoại tác tiêu cực --- buộc người gây ô nhiễm trả đúng chi phí xã
hội họ tạo ra, thay vì đẩy chi phí đó cho cộng đồng gánh chịu miễn phí.
Cơ chế ``cap-and-trade'' (mua bán hạn ngạch phát thải) là cách khác đạt
cùng mục tiêu: thay vì ấn định GIÁ (thuế), nhà nước ấn định LƯỢNG phát
thải tối đa rồi để thị trường tự tìm ra giá.

\textbf{Ví dụ minh hoạ:} Giả sử một nhà máy có chi phí biên tư nhân
\(MPC = 2 + 0{,}5Q\) và gây ô nhiễm với chi phí ngoại tác biên không đổi
\(MEC = 6\) (nghìn đồng/đơn vị)
\(\Rightarrow MSC = MPC + MEC = 8 + 0{,}5Q\). Cầu thị trường
\(P = 30 - 0{,}5Q\). Thị trường tự do sản xuất tại \(D = MPC\):
\(2+0{,}5Q = 30-0{,}5Q \Rightarrow Q_{\text{thị trường}} = 28\). Mức tối
ưu xã hội tại \(D = MSC\):
\(8+0{,}5Q = 30-0{,}5Q \Rightarrow Q^* = 22\). Sản xuất thừa 6 đơn vị so
với tối ưu --- thuế Pigou đúng bằng \(MEC = 6\)/đơn vị sẽ đẩy chi phí tư
nhân lên ngang chi phí xã hội, buộc thị trường tự điều chỉnh đúng về
\(Q^* = 22\).

\begin{center}\rule{0.5\linewidth}{0.5pt}\end{center}

\subsection{Thông tin bất cân xứng (Asymmetric
Information)}\label{thuxf4ng-tin-bux1ea5t-cuxe2n-xux1ee9ng-asymmetric-information}

\textbf{Cốt lõi:} Một bên trong giao dịch có nhiều thông tin hơn bên
kia. Akerlof (1970) hình thức hoá qua mô hình ``thị trường chanh''
(market for lemons): nếu người mua chỉ biết chất lượng TRUNG BÌNH của
thị trường, họ chỉ sẵn sàng trả mức giá tương ứng chất lượng trung bình
--- đẩy người bán hàng TỐT (đòi giá cao hơn mức trung bình) rời thị
trường trước, kéo chất lượng trung bình còn lại xuống thấp hơn nữa. Lặp
lại vòng xoáy này (adverse selection death spiral) có thể khiến cả thị
trường sụp đổ dù đôi bên đều có thể có lợi nếu thông tin đầy đủ. Phân
biệt hai thời điểm: lựa chọn ngược (adverse selection) xảy ra TRƯỚC giao
dịch do thông tin ẩn giấu; rủi ro đạo đức (moral hazard) xảy ra SAU giao
dịch do hành động ẩn giấu.

\textbf{Điểm mù thường gặp:} Trong đàm phán M\&A hoặc mua nhà đất, người
mua thường tin lời người bán mà quên nhận ra một sự thật cấu trúc: bên
bán LUÔN có động cơ và thông tin để che giấu nhược điểm --- đây chính là
lý do due diligence, bảo hành (warranty), và tài khoản ký quỹ (escrow)
tồn tại như công cụ bắt buộc chứ không phải thủ tục hình thức. Điểm mù
thứ hai: nhầm lẫn hai loại thường xảy ra --- cho rằng mọi vấn đề thông
tin đều giải quyết được bằng cách ``minh bạch hoá'' (chỉ đúng với lựa
chọn ngược, vì vấn đề là thông tin ẩn), trong khi rủi ro đạo đức là vấn
đề HÀNH VI xảy ra sau, minh bạch hoá thông tin ban đầu không giải quyết
được --- cần giám sát hoặc thiết kế khuyến khích (như mức khấu trừ bảo
hiểm) thay vì chỉ công khai thông tin.

\textbf{Ứng dụng thực tế:} Bảo hiểm y tế tự nguyện dễ gặp lựa chọn
ngược: người biết mình có nguy cơ sức khỏe cao có xu hướng mua bảo hiểm
nhiều hơn người khỏe mạnh, đẩy chi phí trung bình của quỹ bảo hiểm lên
cao --- buộc công ty bảo hiểm phải sàng lọc hoặc tăng phí. Rủi ro đạo
đức xuất hiện riêng biệt: một khi đã mua bảo hiểm ô tô toàn diện, người
lái có xu hướng cẩn trọng kém hơn (vì hậu quả tai nạn phần lớn do bên
bảo hiểm gánh) --- đây là lý do hợp đồng bảo hiểm luôn có mức khấu trừ
(deductible) bắt người được bảo hiểm chịu một phần chi phí.

\textbf{Ví dụ minh hoạ:} Thị trường xe cũ có 2 loại: xe tốt (50\% thị
trường, người bán sẵn lòng bán từ 200 triệu, người mua sẵn lòng trả đến
240 triệu) và xe kém (50\%, người bán sẵn lòng bán từ 80 triệu, người
mua sẵn lòng trả đến 100 triệu). Nếu người mua KHÔNG phân biệt được, họ
chỉ trả mức kỳ vọng trung bình:
\(0{,}5 \times 240 + 0{,}5 \times 100 = 170\) triệu. Nhưng 170 triệu
THẤP hơn mức 200 triệu người bán xe tốt đòi \(\Rightarrow\) người bán xe
tốt RÚT khỏi thị trường. Người mua nhận ra chỉ còn xe kém, hạ kỳ vọng
xuống 100 triệu --- thị trường xe tốt biến mất hoàn toàn dù cả người mua
lẫn người bán xe tốt đều có thể có lợi nếu giao dịch với đầy đủ thông
tin.

\begin{center}\rule{0.5\linewidth}{0.5pt}\end{center}

\subsection{Hàng hóa công (Public
goods)}\label{huxe0ng-huxf3a-cuxf4ng-public-goods}

\textbf{Cốt lõi:} Hàng hóa có hai đặc tính độc lập nhau: cạnh tranh
trong tiêu dùng hay không (người này dùng có làm giảm phần người khác
không), và loại trừ được hay không (có thể ngăn người không trả tiền vẫn
hưởng lợi không). Kết hợp hai trục này cho ra bốn loại hàng hoá kinh
điển: hàng hoá tư nhân (cạnh tranh + loại trừ được), hàng hoá câu lạc bộ
(không cạnh tranh + loại trừ được), tài nguyên chung (cạnh tranh + không
loại trừ được --- xem khái niệm tiếp theo), và hàng hoá công thuần tuý
(không cạnh tranh + không loại trừ được). Với hàng hoá công thuần tuý,
thị trường tư nhân có xu hướng cung cấp dưới mức cần thiết --- vấn đề
``kẻ ăn theo'' (free-rider): mỗi người có động cơ chờ người khác trả
tiền trong khi mình vẫn hưởng lợi.

\textbf{Điểm mù thường gặp:} Trước một vấn đề tập thể (ô nhiễm không
khí, an ninh khu phố, nghiên cứu cơ bản), phản xạ ``để thị trường/tư
nhân tự lo'' thường xuất hiện trước khi kịp nhận ra bản chất hàng hoá
công của vấn đề --- không cạnh tranh và không loại trừ được --- khiến cơ
chế giá tự nó không thể hoạt động. Điểm mù thứ hai: gộp chung ``hàng hoá
công'' (public goods, khái niệm kinh tế học chính xác) với ``hàng hoá do
nhà nước cung cấp'' (khái niệm hành chính) --- nhiều thứ nhà nước cung
cấp thực chất là hàng hoá tư nhân hoặc câu lạc bộ (giáo dục, y tế có thể
loại trừ được bằng học phí/viện phí) được nhà nước trợ giá vì lý do công
bằng, không phải vì đặc tính kinh tế học của chúng đòi hỏi vậy.

\textbf{Ứng dụng thực tế:} Nghiên cứu khoa học cơ bản (không tạo sản
phẩm bán được ngay) thường do nhà nước hoặc quỹ phi lợi nhuận tài trợ,
vì doanh nghiệp tư nhân không thể thu phí từ tất cả những ai hưởng lợi
từ tri thức đó sau này. Radio/truyền hình phát sóng miễn phí (không loại
trừ được về mặt kỹ thuật truyền thống) thường được tài trợ gián tiếp qua
quảng cáo thay vì thu phí trực tiếp --- một cách né vấn đề free-rider mà
không cần chính phủ can thiệp.

\textbf{Ví dụ minh hoạ:} Pháo hoa lễ hội: chi phí tổ chức 500 triệu
đồng, thành phố có 20.000 người thực sự sẵn lòng trả trung bình 30.000đ
để xem (nhưng không thể thu phí vì không loại trừ được ai xem từ xa)
\(\Rightarrow\) tổng lợi ích xã hội \(= 20.000 \times 30.000 = 600\)
triệu \(> 500\) triệu chi phí --- về mặt xã hội ĐÁNG tổ chức. Nhưng vì
không thể thu phí trực tiếp từng người xem, không doanh nghiệp tư nhân
nào tự nguyện đứng ra tổ chức dù có lợi cho xã hội --- cần ngân sách
công hoặc nhà tài trợ (được bù lại bằng gắn thương hiệu) đứng ra chi trả
thay cơ chế giá thông thường.

\begin{center}\rule{0.5\linewidth}{0.5pt}\end{center}

\subsection{Sức mạnh độc quyền (Market
power)}\label{sux1ee9c-mux1ea1nh-ux111ux1ed9c-quyux1ec1n-market-power}

\textbf{Cốt lõi:} Khi một doanh nghiệp (độc quyền) hoặc một nhóm nhỏ
doanh nghiệp (độc quyền nhóm) có khả năng định giá cao hơn chi phí biên
mà không mất hết khách hàng, sản lượng thị trường sẽ thấp hơn và giá cao
hơn mức cạnh tranh --- tạo tổn thất vô ích. Khác biệt cốt lõi so với
cạnh tranh hoàn hảo: doanh nghiệp độc quyền đối diện TOÀN BỘ đường cầu
thị trường (dốc xuống) thay vì đường cầu nằm ngang tại giá thị trường,
nên để bán thêm phải HẠ giá cho TẤT CẢ đơn vị đã bán trước --- khiến
doanh thu biên (\(MR\)) luôn THẤP HƠN giá bán (\(MR < P\)), khác cạnh
tranh hoàn hảo nơi \(MR = P\). Độc quyền vẫn áp dụng quy tắc \(MR = MC\)
để tìm sản lượng tối ưu (xem Chương 1), nhưng vì \(MR < P\), mức giá bán
cuối cùng cao hơn chi phí biên --- khoảng cách này là nguồn gốc tổn thất
vô ích.

\textbf{Điểm mù thường gặp:} Khi đánh giá một ngành hay một thương vụ
M\&A, dễ chỉ nhìn vào QUY MÔ doanh nghiệp mà quên kiểm tra đúng hai yếu
tố quyết định sức mạnh định giá thực sự: rào cản gia nhập ngành và mức
độ thay thế của sản phẩm --- một công ty nhỏ trong thị trường ngách
không đối thủ vẫn có thể có sức mạnh định giá cao. Điểm mù thứ hai: nhầm
``độc quyền có lợi nhuận cao'' với ``độc quyền luôn xấu cho xã hội'' ---
lợi nhuận độc quyền TỰ THÂN chỉ là chuyển giao từ người tiêu dùng sang
doanh nghiệp (không mất mát ròng); phần thực sự gây hại là TỔN THẤT VÔ
ÍCH từ sản lượng bị cắt giảm, một phân biệt quan trọng khi thiết kế
chính sách (đánh thuế lợi nhuận độc quyền khác về bản chất với buộc tăng
sản lượng).

\textbf{Ứng dụng thực tế:} Cơ quan quản lý cạnh tranh thường chặn các
thương vụ M\&A giữa hai đối thủ trực tiếp lớn trong cùng ngành hẹp (ví
dụ hai hãng viễn thông) vì lo ngại giảm số người chơi sẽ làm tăng sức
mạnh định giá chung của ngành, dù mỗi công ty trước sáp nhập không hẳn
đã ``độc quyền'' theo nghĩa tuyệt đối.

\textbf{Ví dụ minh hoạ:} Cầu thị trường \(P = 100 - Q\), chi phí biên
không đổi \(MC = 20\). Độc quyền: doanh thu
\(TR = PQ = 100Q - Q^2 \Rightarrow MR = 100 - 2Q\). Đặt \(MR = MC\):
\(100 - 2Q = 20 \Rightarrow Q_{\text{độc quyền}} = 40\),
\(P_{\text{độc quyền}} = 100 - 40 = 60\). Nếu là cạnh tranh hoàn hảo
(\(P = MC\)): \(100 - Q = 20 \Rightarrow Q_{\text{cạnh tranh}} = 80\),
\(P_{\text{cạnh tranh}} = 20\). Độc quyền bán ÍT hơn một nửa (40 so với
80) với giá CAO gấp 3 lần (60 so với 20) --- khoảng cách này tạo ra tổn
thất vô ích đáng kể so với cạnh tranh hoàn hảo.

\begin{center}\rule{0.5\linewidth}{0.5pt}\end{center}

\subsection{Bi kịch của cái chung (Tragedy of the
Commons)}\label{bi-kux1ecbch-cux1ee7a-cuxe1i-chung-tragedy-of-the-commons}

\textbf{Cốt lõi:} Khi một tài nguyên có tính CẠNH TRANH trong tiêu dùng
(một người khai thác làm giảm phần còn lại cho người khác) nhưng KHÔNG
LOẠI TRỪ ĐƯỢC (không thể ngăn ai khai thác) --- góc phần tư ``tài nguyên
chung'' trong ma trận ở khái niệm trước --- mỗi cá nhân chỉ tính chi
phí/lợi ích RIÊNG của mình khi quyết định khai thác, bỏ qua chi phí họ
áp lên người khác (về bản chất là một dạng ngoại tác tiêu cực, xem khái
niệm đầu chương). Hardin (1968) đặt tên hiện tượng khai thác vượt mức
bền vững này là ``bi kịch của cái chung'': mỗi cá nhân hành động hoàn
toàn duy lý theo lợi ích riêng, nhưng tổng hợp lại gây hại cho tất cả.

\textbf{Điểm mù thường gặp:} Thường bị nhầm là vấn đề ``lòng tham'' hay
``thiếu đạo đức'' của người khai thác --- nhưng về bản chất đây là vấn
đề CẤU TRÚC KHUYẾN KHÍCH: ngay cả người có đạo đức cao nhất cũng đối
diện cùng động cơ khai thác thêm, vì phần họ TỪ CHỐI khai thác không tự
động được giữ lại mà thường bị người khác khai thác thay --- ``nhường
nhịn đơn phương'' không giải quyết được vấn đề nếu không đi kèm cơ chế
ràng buộc TẤT CẢ mọi người tham gia.

\textbf{Ứng dụng thực tế:} Đánh bắt cá quá mức tại các ngư trường quốc
tế (không thuộc chủ quyền riêng ai) là ví dụ kinh điển --- mỗi tàu cá
đánh bắt tối đa có lợi cho riêng họ, nhưng tổng hợp hành vi này làm suy
giảm quần thể cá cho TẤT CẢ trong tương lai. Hai hướng giải pháp chính:
(1) hạn ngạch đánh bắt kết hợp giấy phép có thể chuyển nhượng (ITQ ---
biến tài nguyên chung thành gần giống hàng hoá tư nhân có thể sở
hữu/giao dịch), hoặc (2) Ostrom (giải Nobel 2009) chỉ ra cộng đồng địa
phương tự quản lý qua quy tắc tự thiết lập (không cần nhà nước hay tư
nhân hoá hoàn toàn) cũng có thể thành công trong nhiều trường hợp thực
tế được nghiên cứu.

\textbf{Ví dụ minh hoạ:} Một hồ đánh cá chung có lợi ích ròng biên KHÔNG
đổi cho mỗi ngư dân khi tự đánh bắt thêm là 10 (họ chỉ tính riêng phần
mình), nhưng mỗi đơn vị đánh bắt thêm cũng làm giảm sản lượng chung
tương lai, gây tổng chi phí ngoại tác biên quy đổi \(= 4\)/đơn vị mà
người đánh bắt không tính đến (lợi ích biên xã hội thực
\(= 10 - 4 = 6\)). Với chi phí khai thác biên \(MC = 2 + 0{,}3Q\): mức
khai thác cá nhân tự phát (bỏ qua ngoại tác):
\(10 = 2 + 0{,}3Q \Rightarrow Q \approx 26{,}7\)/người. Mức khai thác
bền vững xã hội phải trừ ngoại tác:
\(6 = 2 + 0{,}3Q \Rightarrow Q^* \approx 13{,}3\)/người --- chỉ bằng MỘT
NỬA mức khai thác tự phát, minh hoạ mức độ khai thác quá mức khi không
có cơ chế điều phối.

\begin{center}\rule{0.5\linewidth}{0.5pt}\end{center}

\section{CHƯƠNG 3: RA QUYẾT ĐỊNH TRONG HOÀN CẢNH BẤT ĐỊNH \& THỜI GIAN
(Decision
Making)}\label{chux1b0ux1a1ng-3-ra-quyux1ebft-ux111ux1ecbnh-trong-houxe0n-cux1ea3nh-bux1ea5t-ux111ux1ecbnh-thux1eddi-gian-decision-making}

Nhóm công cụ tư duy cá nhân --- cách một nhà kinh tế ``lý tính'' cân
nhắc lựa chọn. Đây là nhóm dễ nhớ về mặt lý thuyết nhưng khó áp dụng
nhất trong thực tế, vì nó đi ngược lại nhiều bản năng tâm lý tự nhiên.

\begin{center}\rule{0.5\linewidth}{0.5pt}\end{center}

\subsection{Chi phí cơ hội (Opportunity
cost)}\label{chi-phuxed-cux1a1-hux1ed9i-opportunity-cost}

\textbf{Cốt lõi:} Giá trị của lựa chọn tốt nhất bị bỏ qua khi ta chọn
một phương án --- đây là chi phí thực sự cần cân nhắc khi ra quyết định,
không chỉ là tiền bỏ ra mà cả thời gian, công sức, và những gì lẽ ra có
thể đạt được nếu chọn hướng khác. Về mặt hình thức, nếu nguồn lực khan
hiếm (thời gian, vốn, lao động) có thể tạo ra hai kết quả A hoặc B,
đường giới hạn khả năng sản xuất (PPF --- production possibility
frontier) mô tả mọi tổ hợp (A, B) tối đa có thể đạt được; độ dốc của PPF
tại một điểm chính là chi phí cơ hội biên của A tính theo B. PPF thường
lõm (bowed outward) vì nguồn lực không đồng nhất --- khi chuyển ngày
càng nhiều nguồn lực từ B sang A, những nguồn lực kém phù hợp nhất với A
(nhưng từng rất hiệu quả cho B) mới bị chuyển sau cùng, khiến chi phí cơ
hội của mỗi đơn vị A tăng dần (quy luật chi phí cơ hội tăng dần). Cần
phân biệt chi phí kế toán (accounting cost --- chỉ tính chi tiêu bằng
tiền thực trả) với chi phí kinh tế (economic cost = chi phí kế toán +
chi phí cơ hội của các nguồn lực tự có, như vốn tự bỏ hay thời gian chủ
sở hữu); lợi nhuận kinh tế do đó luôn nhỏ hơn lợi nhuận kế toán, và một
doanh nghiệp có lãi kế toán vẫn có thể đang lỗ về kinh tế nếu vốn và
công sức đó sinh lời cao hơn ở phương án khác. Đây cũng chính là nền
tảng của lợi thế so sánh (comparative advantage --- xem Chương 4): một
quốc gia hay cá nhân nên chuyên môn hoá vào hoạt động có chi phí cơ hội
thấp nhất, không nhất thiết là hoạt động mình giỏi nhất tuyệt đối.

\textbf{Điểm mù thường gặp:} Chỉ tính chi phí bằng tiền mặt bỏ ra mà
quên tính thời gian hoặc nguồn lực đã dùng vào việc khác --- ví dụ tự
làm một việc ``để tiết kiệm tiền thuê ngoài'' trong khi thời gian đó lẽ
ra tạo ra giá trị lớn hơn nếu dùng vào chuyên môn chính. Một điểm mù
khác: nhầm chi phí cơ hội với tổng giá trị của MỌI phương án bị từ chối
--- đúng ra chỉ tính phương án tốt nhất trong số bị bỏ qua (nếu có ba
lựa chọn thay thế, chi phí cơ hội bằng giá trị của lựa chọn tốt nhất
trong ba, không phải tổng cả ba). Điều này áp dụng cả cho lựa chọn
``không làm gì'': giữ tiền mặt nhàn rỗi hoặc trì hoãn quyết định vẫn có
chi phí cơ hội --- đúng bằng lợi nhuận kỳ vọng của phương án đầu tư tốt
nhất bị bỏ lỡ trong thời gian đó, dù không có dòng tiền nào thực sự
``mất đi''.

\textbf{Ứng dụng thực tế:} Một nhà sáng lập tự làm kế toán để ``tiết
kiệm chi phí'' --- nhưng nếu 5 giờ mỗi tuần đó dùng để gặp khách hàng có
thể mang lại doanh thu lớn hơn nhiều lần phí thuê kế toán, thì việc tự
làm đang tốn chi phí cơ hội cao hơn lợi ích tưởng tượng. Tương tự, một
chủ tiệm báo cáo lãi kế toán 200 triệu đồng/năm nhưng nếu số vốn 2 tỷ
đồng bỏ vào tiệm có thể sinh lời 15\%/năm (300 triệu đồng) nếu gửi tiết
kiệm hoặc đầu tư nơi khác với mức rủi ro tương đương, thì về bản chất
kinh tế, tiệm đang lỗ cơ hội 100 triệu đồng mỗi năm dù sổ sách vẫn ghi
nhận lãi.

\textbf{Ví dụ minh hoạ:} Một xưởng may có 120 giờ công/tuần. May một áo
sơ mi tốn 2 giờ công; may một quần jean tốn 4 giờ công. Nếu dồn toàn bộ
giờ công cho áo sơ mi: \(120 \div 2 = 60\) áo/tuần (0 quần); nếu dồn
toàn bộ cho quần jean: \(120 \div 4 = 30\) quần/tuần (0 áo) --- đây là
hai điểm đầu mút của PPF. Vì công nghệ đơn giản (một loại lao động, năng
suất không đổi), PPF ở đây là đoạn thẳng chứ không lõm --- chi phí cơ
hội không đổi dọc đường: mỗi 1 quần jean sản xuất thêm luôn ``tốn'' đúng
2 áo sơ mi bị bỏ qua (vì 4 giờ dùng cho 1 quần đúng bằng giờ có thể may
2 áo), bất kể đang sản xuất bao nhiêu quần. Nếu xưởng hiện sản xuất tại
điểm (40 áo, 10 quần) --- dùng đúng
\(40 \times 2 + 10 \times 4 = 80 + 40 = 120\) giờ, đúng năng lực --- và
cân nhắc nhận thêm một đơn hàng 5 quần jean nữa, họ phải hiểu ngay: chấp
nhận đơn này đồng nghĩa từ bỏ \(5 \times 2 = 10\) áo sơ mi, bất kể giá
bán quần jean trông ``hấp dẫn'' thế nào trên hoá đơn --- quyết định đúng
chỉ nên nhận đơn nếu lợi nhuận từ 5 quần jean vượt lợi nhuận có thể có
từ 10 áo sơ mi bị bỏ lỡ.

\begin{center}\rule{0.5\linewidth}{0.5pt}\end{center}

\subsection{Chi phí chìm (Sunk
cost)}\label{chi-phuxed-chuxecm-sunk-cost}

\textbf{Cốt lõi:} Chi phí đã bỏ ra và không thể thu hồi, dù có tiếp tục
hay dừng lại. Nguyên tắc ra quyết định hợp lý (chiết khấu dòng tiền,
relevant costing) là chỉ dòng tiền và lợi ích trong tương lai mới nên
ảnh hưởng đến quyết định --- chi phí chìm không nên xuất hiện trong phép
tính. Về mặt hình thức, nếu bài toán tối ưu là chọn \(x\) để tối đa hoá
lợi ích ròng tương lai \(V(x) - K(x)\), trong đó \(K(x)\) là chi phí
phải chi THÊM tuỳ theo lựa chọn \(x\), thì một khoản chi phí \(C\) đã
trả trong quá khứ --- không phụ thuộc \(x\) --- chỉ làm dịch toàn bộ hàm
mục tiêu xuống một hằng số không đổi (\(V(x) - K(x) - C\)), không hề làm
thay đổi giá trị \(x\) tối ưu. Đây là lý do các chuẩn mực kế toán quản
trị (relevant costing) và các mô hình tối ưu hoá động (dynamic
programming, nguyên lý Bellman) đều loại chi phí chìm khỏi hàm mục tiêu
ngay từ đầu --- quyết định chỉ nên là hàm số của trạng thái hiện tại và
các dòng tiền còn lại phía trước, không phải hàm số của lịch sử đã qua.
Về mặt hành vi, xu hướng vẫn tiếp tục đầu tư vào một phương án chỉ vì đã
đầu tư nhiều --- dù bằng chứng mới cho thấy nên dừng --- được gọi là
ngụy biện chi phí chìm (sunk cost fallacy), một trong những thiên lệch
được ghi nhận thực nghiệm rõ ràng nhất trong kinh tế học hành vi (Arkes
\& Blumer, 1985; xem thêm Nhóm 6).

\textbf{Điểm mù thường gặp:} Ngay tại thời điểm cân nhắc có nên tiếp tục
một dự án đang lỗ, ta hiếm khi dừng lại để nhận ra đây chính là một tình
huống ``chi phí chìm'' --- cảm giác ``đã đầu tư quá nhiều để bỏ giữa
chừng'' xuất hiện trước, và phép tính lý tính chỉ được gọi ra (nếu có)
sau khi quyết định cảm tính đã hình thành. Một điểm mù khác: đánh đồng
``chi phí chìm'' với ``chi phí cố định'' (fixed cost) --- hai khái niệm
nằm trên hai trục phân loại khác nhau. Chi phí cố định là chi phí không
đổi theo sản lượng trong ngắn hạn (ví dụ tiền thuê mặt bằng theo hợp
đồng); chi phí chìm là chi phí không thể thu hồi bất kể quyết định nào
được chọn. Một khoản chi phí có thể vừa cố định vừa chìm (tiền thuê nhà
xưởng đã trả trước, không hoàn lại), nhưng cũng có thể cố định mà không
chìm --- một máy móc đã mua nhưng còn giá trị thanh lý đáng kể; nếu dừng
dự án, phần giá trị thu hồi được từ việc bán máy KHÔNG phải chi phí chìm
và cần đưa vào phép so sánh (xem thêm Chương 1 --- sản lượng tối ưu và
quyết định có nên sản xuất).

\textbf{Ứng dụng thực tế:} Một công ty đã chi 2 tỷ đồng phát triển một
sản phẩm nhưng nghiên cứu thị trường mới cho thấy sản phẩm sẽ không bán
được --- quyết định đúng là dừng lại ngay, đánh giá dựa trên chi phí
hoàn thiện còn lại so với doanh thu kỳ vọng, không phải dựa trên 2 tỷ đã
mất --- trừ phần nào trong 2 tỷ đó còn có thể thu hồi được (ví dụ bán
lại bản quyền công nghệ hoặc nguyên vật liệu dở dang cho đối tác khác);
phần thu hồi này cần đưa vào bài toán vì nó là một dòng tiền tương lai
thực sự, không phải chi phí chìm.

\textbf{Ví dụ minh hoạ:} Mở rộng ví dụ trên bằng số liệu cụ thể: công ty
đã chi 2.000 triệu đồng phát triển sản phẩm (chi phí chìm --- không đưa
vào phép tính tiếp theo). Hai phương án còn lại: (1) Tiếp tục hoàn thiện
--- cần chi thêm 500 triệu đồng, dự kiến thu về 700 triệu đồng
\(\Rightarrow\) lợi ích ròng tương lai \(= 700 - 500 = +200\) triệu
đồng. (2) Dừng lại và thanh lý tài sản dở dang (bản quyền, nguyên vật
liệu) --- thu về 150 triệu đồng \(\Rightarrow\) lợi ích ròng tương lai
\(= +150\) triệu đồng. So sánh đúng chỉ dựa trên hai con số tương lai
này (200 so với 150): nên chọn TIẾP TỤC vì \(200 > 150\) --- kết luận
này hoàn toàn không phụ thuộc việc 2.000 triệu đã chi là 1.000 hay 5.000
triệu, vì con số đó không xuất hiện trong phép so sánh. Lối tư duy sai
(thường gặp) sẽ so sánh ``đã đầu tư 2.000 triệu, nếu dừng coi như mất
trắng, còn tiếp tục ít ra cũng thu về gần đủ'' --- cách nghĩ này lẫn chi
phí chìm vào bài toán và dễ dẫn đến việc tiếp tục rót vốn dù phương án
dừng lại thực ra tốt hơn. Phép thử: nếu chi phí hoàn thiện là 600 triệu
đồng thay vì 500 (mọi thứ khác giữ nguyên), lợi ích ròng của phương án
tiếp tục giảm còn \(700 - 600 = 100 < 150\) --- quyết định đúng khi đó
ĐẢO CHIỀU thành DỪNG LẠI, điều mà lối tư duy ``tiếc của đã bỏ ra''
thường bỏ lỡ vì vẫn neo vào con số 2.000 triệu đã chi.

\begin{center}\rule{0.5\linewidth}{0.5pt}\end{center}

\subsection{Tư duy biên (Marginal
thinking)}\label{tux1b0-duy-biuxean-marginal-thinking}

\textbf{Cốt lõi:} So sánh lợi ích tăng thêm (marginal benefit ---
\(MB\)) với chi phí tăng thêm (marginal cost --- \(MC\)) của MỘT đơn vị
quyết định tiếp theo, thay vì so sánh tổng lợi ích với tổng chi phí.
Quyết định tối ưu là tiếp tục cho đến khi lợi ích biên bằng chi phí
biên. Về mặt toán học, nếu lợi ích ròng là \(NB(Q) = TB(Q) - TC(Q)\)
(tổng lợi ích trừ tổng chi phí theo biến quyết định \(Q\)), điều kiện
cần bậc một để \(Q\) tối đa hoá \(NB\) là đạo hàm
\(NB'(Q) = TB'(Q) - TC'(Q) = MB(Q) - MC(Q) = 0\), tức \(MB = MC\) ---
``tư duy biên'' vì vậy không phải một mẹo kinh nghiệm mà là hệ quả trực
tiếp của phép tính vi phân tối ưu hoá, áp dụng cho mọi bài toán ra quyết
định có biến liên tục. Đây chính là nguyên lý nền đứng sau lựa chọn tối
ưu của người tiêu dùng, quyết định sản lượng của doanh nghiệp
(\(MR = MC\), xem Chương 1 --- sản lượng tối ưu), và phần lớn các bài
toán tối ưu hoá có ràng buộc trong kinh tế học. Điều kiện bậc một
\(MB = MC\) chỉ đảm bảo một điểm dừng (stationary point); điều kiện bậc
hai đòi hỏi \(MC\) phải tăng nhanh hơn \(MB\) (hay \(MB\) giảm nhanh hơn
\(MC\)) tại điểm đó để đảm bảo đây thực sự là điểm tối đa chứ không phải
tối thiểu --- về trực giác, đây là lý do quy tắc ``tiếp tục miễn \(MB\)
lớn hơn \(MC\), dừng khi \(MB\) nhỏ hơn \(MC\)'' luôn đúng: nó tự động
thoả điều kiện bậc hai đó.

\textbf{Điểm mù thường gặp:} Khi cân nhắc có nên nhận thêm một đơn hàng,
khách hàng, hay ca làm việc, phản xạ thường là nhìn vào chi phí hoặc lợi
nhuận TRUNG BÌNH đã biết, mà quên dừng lại tính đúng chi phí và lợi ích
BIÊN của riêng quyết định tiếp theo đó. Một điểm mù khác: ngộ nhận chi
phí biên (\(MC\)) và chi phí bình quân (\(AC\)) luôn di chuyển cùng
chiều. Về mặt toán học, \(MC\) cắt \(AC\) đúng tại điểm \(AC\) nhỏ nhất:
khi \(MC\) nhỏ hơn \(AC\), mỗi đơn vị thêm kéo bình quân giảm xuống
(\(AC\) đang giảm); khi \(MC\) lớn hơn \(AC\), mỗi đơn vị thêm kéo bình
quân tăng lên (\(AC\) đang tăng). Điều này nghĩa là một doanh nghiệp
hoàn toàn có thể đang ở vùng chi phí bình quân giảm dần (trông có vẻ
``ngày càng hiệu quả'') nhưng chi phí biên đã bắt đầu tăng --- quyết
định sản lượng tối ưu (\(MR = MC\)) không hề chờ đến khi \(AC\) ngừng
giảm mới cần quan tâm.

\textbf{Ứng dụng thực tế:} Một hãng hàng không sẵn sàng bán vé giờ chót
với giá rất thấp --- vì chi phí biên của một hành khách thêm trên chuyến
bay đã lên lịch gần như bằng 0 (ghế trống vẫn bay), nên bất kỳ giá nào
cao hơn chi phí biên đó đều có lợi, dù thấp hơn nhiều giá vé trung bình.
Đây cũng là minh hoạ trực tiếp cho điểm mù thứ hai ở trên: chi phí BÌNH
QUÂN mỗi ghế (đã gánh cả chi phí nhiên liệu, phi hành đoàn, khấu hao máy
bay) có thể là 2 triệu đồng, nhưng chi phí BIÊN của ghế trống cuối cùng
gần bằng 0 --- bán vé 500 nghìn đồng cho ghế đó vẫn có lợi dù thấp hơn
nhiều lần mức bình quân.

\textbf{Ví dụ minh hoạ:} Một quán cà phê cân nhắc thuê thêm nhân viên
theo giờ vào cuối tuần. Lợi ích biên (doanh thu tăng thêm nhờ phục vụ
nhanh hơn, ít khách bỏ đi) của giờ nhân viên thứ \(Q\) ước tính
\(MB = 200 - 8Q\) (nghìn đồng/giờ, \(Q\) tính bằng số giờ nhân viên
thêm). Chi phí biên (lương giờ và phụ cấp) không đổi \(MC = 40\) (nghìn
đồng/giờ). Điểm tối ưu:
\(MB = MC \Rightarrow 200 - 8Q = 40 \Rightarrow Q^* = 20\) giờ nhân viên
thêm mỗi cuối tuần. Nếu quán chỉ nhìn vào lợi nhuận trung bình mỗi giờ
nhân viên đã thuê trước đó (giả sử đang dương, ví dụ 60 nghìn đồng/giờ)
mà quyết định thuê thêm ồ ạt đến 30 giờ, thực tế tại \(Q = 30\):
\(MB = 200 - 240 = -40\), nhỏ hơn \(MC = 40\) rất nhiều --- giờ nhân
viên thứ 30 đang lỗ biên nặng (lợi ích biên thực ra đã ÂM, có thể do quá
đông nhân viên gây vướng víu, chi phí phối hợp tăng) dù con số ``trung
bình'' nhìn vẫn có vẻ ổn. Đây là ví dụ cụ thể cho thấy bình quân và biên
có thể lệch pha hoàn toàn.

\begin{center}\rule{0.5\linewidth}{0.5pt}\end{center}

\subsection{Lợi ích biên giảm dần (Diminishing Marginal
Utility)}\label{lux1ee3i-uxedch-biuxean-giux1ea3m-dux1ea7n-diminishing-marginal-utility}

\textbf{Cốt lõi:} Mỗi đơn vị tiêu dùng thêm của cùng một hàng hóa mang
lại lợi ích (thỏa dụng) tăng thêm ít hơn đơn vị trước đó. Đây là lý do
đường cầu dốc xuống, và cũng lý giải vì sao đa dạng hóa thường tốt hơn
dồn hết vào một thứ. Về mặt hình thức, nếu hàm thỏa dụng theo số lượng
tiêu dùng là \(U(Q)\), lợi ích biên \(MU(Q) = U'(Q)\); ``lợi ích biên
giảm dần'' tương đương \(U''(Q) < 0\) --- hàm thỏa dụng lõm (concave)
theo \(Q\). Đây cũng là lý do đường cầu cá nhân dốc xuống: người tiêu
dùng chỉ sẵn sàng trả thêm tiền cho một đơn vị hàng hoá đúng bằng lợi
ích biên nó mang lại, và vì lợi ích biên giảm dần theo số lượng đã có,
mức giá sẵn trả cho đơn vị tiếp theo cũng giảm dần. Khái niệm này còn là
nền tảng của lý thuyết e ngại rủi ro (risk aversion): nếu thỏa dụng đối
với TIỀN cũng lõm, một người sẽ luôn thích một khoản tiền chắc chắn hơn
một canh bạc có cùng giá trị kỳ vọng, vì phần thỏa dụng mất đi khi thua
lớn hơn phần thỏa dụng có thêm khi thắng cùng số tiền. Cần phân biệt
``lợi ích biên giảm dần'' (\(MU\) giảm nhưng vẫn DƯƠNG --- tổng lợi ích
\(TU\) vẫn tăng, chỉ tăng chậm lại) với trường hợp cực đoan hơn là lợi
ích biên ÂM (tiêu dùng thêm khiến tổng lợi ích giảm thực sự, ví dụ ăn
quá no) --- phần lớn phân tích kinh tế chỉ cần giả định \(MU\) dương và
giảm dần, không cần đến vùng \(MU\) âm.

\textbf{Điểm mù thường gặp:} Khi có thêm ngân sách hoặc nguồn lực, phản
xạ ``nhiều hơn luôn tốt hơn'' thường xuất hiện trước khi kịp nhận ra lợi
ích biên đang giảm dần rất nhanh --- ví dụ tăng gấp đôi ngân sách quảng
cáo không có nghĩa doanh thu tăng gấp đôi, vì các kênh và đối tượng hiệu
quả nhất đã được khai thác ở lượt chi tiêu đầu. Một điểm mù khác: nhầm
``lợi ích biên giảm dần'' với ``tổng lợi ích giảm dần'' --- hai điều
hoàn toàn khác nhau. Lợi ích biên giảm dần chỉ có nghĩa TỔNG lợi ích vẫn
tăng nhưng tăng chậm lại, không phải tổng lợi ích đang đi xuống. Nhầm
lẫn này khiến nhiều người kết luận sai rằng nên ngừng tiêu dùng hoặc đầu
tư ngay khi lợi ích biên bắt đầu giảm --- trong khi quyết định đúng (xem
mục Tư duy biên ở trên) là tiếp tục miễn lợi ích biên còn lớn hơn chi
phí biên, bất kể lợi ích biên đã giảm bao nhiêu so với đơn vị đầu tiên.

\textbf{Ứng dụng thực tế:} Trong đầu tư tài chính, lợi ích biên giảm dần
của tài sản là nền tảng lý thuyết cho việc đa dạng hóa danh mục thay vì
dồn toàn bộ vốn vào một cổ phiếu, dù cổ phiếu đó có kỳ vọng sinh lời cao
nhất --- về bản chất, đây là ứng dụng trực tiếp của thỏa dụng biên giảm
dần đối với của cải: mất một khoản tiền X gây đau đớn (giảm thỏa dụng)
nhiều hơn mức vui sướng (tăng thỏa dụng) khi được thêm đúng khoản X đó,
nên dàn trải rủi ro luôn tốt hơn về mặt thỏa dụng kỳ vọng dù không thay
đổi giá trị kỳ vọng bằng tiền.

\textbf{Ví dụ minh hoạ:} Một người tiêu dùng ly trà sữa với thỏa dụng
(tính bằng đơn vị thỏa dụng quy ước) như sau: ly 1 mang lại 50 đơn vị;
ly 2 mang thêm 34 đơn vị (tổng 84); ly 3 mang thêm 20 đơn vị (tổng 104);
ly 4 mang thêm 8 đơn vị (tổng 112); ly 5 mang thêm 2 đơn vị (tổng 114)
--- lợi ích BIÊN (50, 34, 20, 8, 2) giảm dần liên tục dù tổng lợi ích
(50, 84, 104, 112, 114) vẫn tăng ở mọi ly, chỉ tăng chậm dần. Giả sử 1
đơn vị thỏa dụng quy đổi tương đương 2.000 đồng và giá một ly trà sữa là
30.000 đồng (tương đương 15 đơn vị thỏa dụng), áp dụng nguyên tắc tư duy
biên (\(MB = MC\)) ở mục trên: nên mua đến ly mà lợi ích biên còn LỚN
HƠN 15 đơn vị --- ly 3 mang thêm 20 đơn vị (còn lời 5 đơn vị, nên mua),
ly 4 chỉ thêm 8 đơn vị (lỗ 7 đơn vị so với giá, không nên mua)
\(\Rightarrow\) lựa chọn tối ưu là dừng lại ở 3 ly, đúng tại điểm lợi
ích biên vừa vượt qua rồi sắp rơi xuống dưới chi phí biên quy đổi.

\begin{center}\rule{0.5\linewidth}{0.5pt}\end{center}

\section{CHƯƠNG 4: THƯƠNG MẠI \& VĨ MÔ QUỐC TẾ (International Trade \&
Macroeconomics)}\label{chux1b0ux1a1ng-4-thux1b0ux1a1ng-mux1ea1i-vux129-muxf4-quux1ed1c-tux1ebf-international-trade-macroeconomics}

Nhóm giải thích vì sao các quốc gia giao thương với nhau, điều gì xảy ra
khi dòng hàng hóa, vốn và tiền tệ di chuyển qua biên giới --- và tại sao
cùng một chính sách có thể là ``tốt'' cho bên này nhưng ``xấu'' cho bên
kia trong cùng một nền kinh tế. (Xem thêm Phần B, Chương 12 để có bản
hình thức hóa đầy đủ với mô hình Heckscher--Ohlin, Mundell--Fleming và
Melitz.)

\begin{center}\rule{0.5\linewidth}{0.5pt}\end{center}

\subsection{Lợi thế so sánh (Comparative
advantage)}\label{lux1ee3i-thux1ebf-so-suxe1nh-comparative-advantage}

\textbf{Cốt lõi:} Một quốc gia (hoặc cá nhân, doanh nghiệp) nên chuyên
môn hóa vào sản xuất và xuất khẩu hàng hóa mà mình có \textbf{chi phí cơ
hội tương đối thấp hơn} so với đối tác --- ngay cả khi kém hiệu quả hơn
đối tác ở \emph{mọi} mặt về năng suất tuyệt đối. Ricardo (1817) hình
thức hóa ý tưởng này: với 2 quốc gia và 2 hàng hóa, quốc gia H có lợi
thế so sánh về hàng hóa X nếu chi phí cơ hội sản xuất X tính theo Y thấp
hơn ở quốc gia F: \(a_{XL}/a_{YL} < a^*_{XL}/a^*_{YL}\), trong đó
\(a_{iL}\) là số giờ công cần để sản xuất một đơn vị hàng \(i\). Khi
điều kiện đó đúng và tỷ giá thương mại quốc tế \(P_X/P_Y\) nằm giữa hai
mức chi phí cơ hội nội địa, cả hai quốc gia đều tiêu dùng được điểm
\textbf{nằm ngoài} đường giới hạn khả năng sản xuất (PPF) của mình ---
tức đạt tổ hợp hàng hóa không thể có trong trường hợp tự cấp tự túc. Lưu
ý quan trọng: lý thuyết cổ điển giả định công nghệ cố định và lao động
là yếu tố duy nhất; mô hình Heckscher--Ohlin (Phần B) mở rộng thêm vốn
và đất đai, dẫn đến kết quả sâu hơn về phân phối thu nhập theo yếu tố
sản xuất.

\textbf{Điểm mù thường gặp:} Trong tranh luận về mở cửa thương mại, lập
luận phổ biến nhất thường là so sánh năng suất \textbf{tuyệt đối} ---
``nước đó sản xuất gì cũng nhanh và rẻ hơn, mở cửa thì ngành của mình
chết hết'' --- mà quên rằng giao thương là bài toán phân bổ nguồn lực
\textbf{khan hiếm} vào nơi chi phí cơ hội \textbf{tương đối} thấp nhất,
không phải cuộc thi tuyệt đối. Ngay cả khi một bên thua kém tuyệt đối ở
mọi thứ, họ vẫn có lợi khi chuyên môn hóa vào điểm yếu kém \textbf{ít
nhất} của mình.

\textbf{Điểm mù thứ hai:} Lý thuyết cổ điển giả định \textbf{công nghệ
tĩnh} và yếu tố sản xuất không di chuyển xuyên biên giới --- hai điều
hiếm khi đúng dài hạn. Nếu một quốc gia đang phát triển chuyên môn hóa
vào khâu giá trị thấp (gia công, lắp ráp) dựa theo lợi thế so sánh tĩnh,
có nguy cơ bị \textbf{``khóa chặt năng lực''} (capability lock-in): theo
thời gian, khi tiền lương tăng, lợi thế chi phí thấp biến mất nhưng năng
lực công nghệ nội địa không được tích lũy đủ để chuyển sang khâu giá trị
cao hơn --- đây là phản biện học thuật có cơ sở (không phải lập luận bảo
hộ đơn thuần) cần được xem xét nghiêm túc trong thiết kế chính sách công
nghiệp.

\textbf{Ứng dụng thực tế:} Trong tranh luận về CPTPP/EVFTA, câu hỏi
không chỉ là ``Việt Nam có lợi thế so sánh ở đâu hiện tại'' (thường là
gia công, chi phí lao động thấp) mà còn là liệu chuyên môn hóa dài hạn
vào các khâu giá trị thấp có làm chậm quá trình nâng cấp công nghệ và
năng lực nội địa hay không --- và liệu chính sách công nghiệp đủ mạnh để
chủ động dịch chuyển lợi thế so sánh sang các khâu cao hơn theo thời
gian hay không.

\textbf{Ví dụ minh hoạ:} Giả sử nước Home (H) và Foreign (F) sản xuất Áo
(A) và Máy tính (M). Giờ công cần thiết:

{\def\LTcaptype{none} % do not increment counter
\begin{longtable}[]{@{}
  >{\raggedright\arraybackslash}p{(\linewidth - 8\tabcolsep) * \real{0.2000}}
  >{\raggedright\arraybackslash}p{(\linewidth - 8\tabcolsep) * \real{0.2000}}
  >{\raggedright\arraybackslash}p{(\linewidth - 8\tabcolsep) * \real{0.2000}}
  >{\raggedright\arraybackslash}p{(\linewidth - 8\tabcolsep) * \real{0.2000}}
  >{\raggedright\arraybackslash}p{(\linewidth - 8\tabcolsep) * \real{0.2000}}@{}}
\toprule\noalign{}
\begin{minipage}[b]{\linewidth}\raggedright
Quốc gia
\end{minipage} & \begin{minipage}[b]{\linewidth}\raggedright
Áo (\(a_{AL}\))
\end{minipage} & \begin{minipage}[b]{\linewidth}\raggedright
Máy tính (\(a_{ML}\))
\end{minipage} & \begin{minipage}[b]{\linewidth}\raggedright
Chi phí cơ hội của 1 Áo
\end{minipage} & \begin{minipage}[b]{\linewidth}\raggedright
Chi phí cơ hội của 1 Máy tính
\end{minipage} \\
\midrule\noalign{}
\endhead
\bottomrule\noalign{}
\endlastfoot
\textbf{Home (H)} & 2 giờ & 8 giờ & \textbf{0,25 Máy tính} & 4 Áo \\
\textbf{Foreign (F)} & 1 giờ & 2 giờ & \textbf{0,50 Máy tính} & 2 Áo \\
\end{longtable}
}

\begin{itemize}
\tightlist
\item
  F giỏi hơn H ở \emph{cả hai} sản phẩm (lợi thế tuyệt đối thuộc về F).
\item
  Nhưng chi phí cơ hội sản xuất 1 Áo ở H là \(2/8 = 0{,}25\) Máy tính,
  nhỏ hơn ở F (\(1/2 = 0{,}50\) Máy tính) \(\Rightarrow\) \textbf{H có
  lợi thế so sánh về Áo}.
\item
  Chi phí cơ hội sản xuất 1 Máy tính ở F là \(2/1 = 2\) Áo, nhỏ hơn ở H
  (\(8/2 = 4\) Áo) \(\Rightarrow\) \textbf{F có lợi thế so sánh về Máy
  tính}.
\item
  Tỷ giá thương mại cân bằng thỏa mãn: \(0{,}25 < P_A/P_M < 0{,}50\).
  Giả sử \(P_A/P_M = 0{,}33\) (1 Máy tính = 3 Áo): H xuất 3 Áo đổi 1 Máy
  tính (tốn 6 giờ công, thay vì tự làm tốn 8 giờ --- tiết kiệm 2 giờ); F
  xuất 1 Máy tính đổi 3 Áo (tốn 2 giờ công, thay vì tự làm tốn 3 giờ ---
  tiết kiệm 1 giờ). \textbf{Cả hai cùng lợi.}
\end{itemize}

\begin{center}\rule{0.5\linewidth}{0.5pt}\end{center}

\subsection{Bộ ba bất khả thi (Impossible Trinity /
Trilemma)}\label{bux1ed9-ba-bux1ea5t-khux1ea3-thi-impossible-trinity-trilemma}

\textbf{Cốt lõi:} Mundell (1963) và Fleming (1962) chứng minh rằng một
nền kinh tế mở \textbf{không thể đồng thời đạt được cả ba mục tiêu} sau:
\textbf{(1)} Tỷ giá hối đoái cố định (Fixed exchange rate); \textbf{(2)}
Tự do di chuyển dòng vốn (Free capital mobility); \textbf{(3)} Chính
sách tiền tệ độc lập (Independent monetary policy). Một quốc gia bắt
buộc phải chọn tối đa \textbf{hai trong ba} mục tiêu, và chấp nhận từ bỏ
mục tiêu thứ ba.

Cơ sở cơ chế: Nếu vốn di chuyển tự do và tỷ giá cố định, mọi chênh lệch
lãi suất nội địa so với quốc tế sẽ gây ra dòng vốn chảy ra/vào khổng lồ,
ép Ngân hàng Trung ương (NHTW) phải can thiệp bằng dự trữ ngoại hối để
giữ giá đồng tiền --- nghĩa là cung tiền nội địa bị bịt chặt bởi hành vi
can thiệp tỷ giá, NHTW \textbf{mất hoàn toàn khả năng đặt lãi suất độc
lập}.

Ba góc của Tam giác Mundell: - \textbf{Góc A (Mỹ, EU):} Tự do dòng vốn +
Tiền tệ độc lập \(\Rightarrow\) \textbf{Chấp nhận tỷ giá thả nổi}. -
\textbf{Góc B (Hồng Kông, Đan Mạch):} Tự do dòng vốn + Tỷ giá cố định
\(\Rightarrow\) \textbf{Từ bỏ chính sách tiền tệ độc lập} (lãi suất Hồng
Kông tự động chạy theo Fed Mỹ). - \textbf{Góc C (Trung Quốc trước 2015,
Việt Nam giai đoạn trước):} Tỷ giá cố định/quản lý + Tiền tệ độc lập
\(\Rightarrow\) \textbf{Phải kiểm soát dòng vốn} (Capital controls).

\textbf{Điểm mù thường gặp:} Trực giác chính trị thường muốn ``tất cả'':
vừa muốn giữ tỷ giá ổn định (để tránh lạm phát nhập khẩu), vừa muốn hạ
lãi suất nội địa để kích thích kinh tế, lại vừa muốn thu hút dòng vốn
đầu tư nước ngoài tự do. Bỏ qua Bộ ba bất khả thi khiến chính sách tiền
tệ rơi vào tình trạng \textbf{mâu thuẫn cấu trúc} --- can thiệp tỷ giá
đốt sạch dự trữ ngoại hối khi dòng vốn đảo chiều vội vã.

\textbf{Điểm mù thứ hai:} Nhầm lẫn giữa ``kiểm soát dòng vốn trên giấy
tờ'' và ``kiểm soát dòng vốn thực tế''. Khi quy mô thương mại và tài
chính mở rộng, các doanh nghiệp và nhà đầu tư có vô số kênh lách kiểm
soát vốn (chuyển giá, kết hối chậm, tiền mã hóa) --- làm cho góc C thực
tế trượt dần về góc A hoặc B mà chính phủ không chủ động kiểm soát được.

\textbf{Ứng dụng thực tế:} Cuộc khủng hoảng tài chính châu Á (1997):
Thái Lan, Indonesia, Hàn Quốc cố gắng duy trì tỷ giá cố định neo vào USD
+ mở cửa dòng vốn, nhưng đồng thời muốn giữ lãi suất nội địa cao để kiềm
chế lạm phát. Khi Fed tăng lãi suất và USD tăng giá, dòng vốn nóng tháo
chạy; các NHTW đốt hết dự trữ ngoại hối để bảo vệ tỷ giá và thất bại sụp
đổ hoàn toàn trong vài tuần --- minh chứng lịch sử kinh điển về cái giá
phải trả khi vi phạm Bộ ba bất khả thi.

\textbf{Ví dụ minh hoạ:} Giả sử lãi suất Mỹ là \(i^* = 5\%\). NHTW nước
X muốn cố định tỷ giá \(E = 25.000\text{ VND/USD}\) và cho phép dòng vốn
tự do lưu thông. Theo Điều kiện Ngang giá Lãi suất Chưa Bù đắp (UIP):
\(i = i^* + \frac{\Delta E^e}{E}\). Vì tỷ giá cố định kỳ vọng
\(\Delta E^e = 0\), bắt buộc lãi suất nội địa \(i = i^* = 5\%\). Nếu
NHTW X cố tình hạ lãi suất xuống \(i = 3\%\) để kích thích kinh tế: -
Nhà đầu tư rút tiền khỏi X (chuyển VND sang USD) để gửi sang Mỹ hưởng
5\% \(\Rightarrow\) Cầu USD tăng vọt. - Để giữ tỷ giá 25.000, NHTW X
phải bán USD ra từ dự trữ ngoại hối. - Nếu dự trữ ngoại hối là 30 tỷ USD
và dòng vốn tháo chạy là 5 tỷ USD/tuần, dự trữ sẽ \textbf{cạn sạch trong
6 tuần} --- buộc phải phá giá đồng tiền rầm rộ hoặc nâng lãi suất trở
lại 5\%.

\begin{center}\rule{0.5\linewidth}{0.5pt}\end{center}

\subsection{Cán cân thanh toán (Balance of Payments -
BoP)}\label{cuxe1n-cuxe2n-thanh-touxe1n-balance-of-payments---bop}

\textbf{Cốt lõi:} Báo cáo thống kê tổng hợp mọi giao dịch kinh tế giữa
người cư trú (residents) và người không cư trú (non-residents) trong một
thời kỳ. Cấu trúc cơ bản:
\[ \text{Cán cân vãng lai (CA)} + \text{Cán cân Tài chính/Vốn (FA)} + \text{Lỗi và Sơ sót (EO)} = \Delta \text{Dự trữ ngoại hối chính thức (RAM)} \]
Trong đó: \textbf{CA} (Current Account) bao gồm Cán cân thương mại (Xuất
- Nhập khẩu hàng hóa \& dịch vụ), Thu nhập sơ cấp (tiền lương, lợi nhuận
đầu tư chuyển về), và Thu nhập thứ cấp (kiều hối, viện trợ). \textbf{FA}
(Financial Account) bao gồm FDI, FII (vốn đầu tư gián tiếp/chứng khoán),
và các dòng vay nợ.

Đồng nhất thức vĩ mô cốt lõi: \(CA = S - I = (S_{p} - I) + (T - G)\),
trong đó \(S\) là tổng tiết kiệm quốc gia, \(I\) là tổng đầu tư. Thâm
hụt cán cân vãng lai (\(CA < 0\)) tương đương với việc \textbf{tiết kiệm
quốc gia không đủ tài trợ cho đầu tư nội địa} (\(S < I\)), buộc quốc gia
đó phải vay nợ hoặc bán tài sản cho nước ngoài (thặng dư cán cân tài
chính \(FA > 0\)).

\textbf{Điểm mù thường gặp:} Nhìn thấy ``thâm hụt thương mại'' hay
``thâm hụt vãng lai'' là kết luận ngay nền kinh tế ``yếu kém'' hay ``mất
sức cạnh tranh'' --- mà không phân tích vế \(S - I\). Với một nền kinh
tế đang phát triển có nhiều cơ hội đầu tư sinh lời cao (\(I\) lớn), thâm
hụt CA để nhập khẩu máy móc, công nghệ (được tài trợ bằng dòng vốn FDI
dài hạn) là hoàn toàn lành mạnh và cần thiết cho tăng trưởng.

\textbf{Điểm mù thứ hai:} Đánh đồng ``thặng dư thương mại'' với ``phúc
lợi quốc gia cao''. Một quốc gia thặng dư CA liên tục lớn có thể là biểu
hiện của \textbf{tiết kiệm quá mức / tiêu dùng nội địa yếu kém} (như Đức
hoặc Nhật Bản trong nhiều giai đoạn), dẫn đến việc phải tích lũy tài sản
nước ngoài với lợi suất thấp thay vì cải thiện đời sống người dân trong
nước.

\textbf{Ứng dụng thực tế:} Phân tích mô hình tăng trưởng của Việt Nam:
Việt Nam thường xuyên thâm hụt thương mại trong giai đoạn 2000--2011 do
nhập khẩu tư liệu sản xuất mạnh phục vụ công nghiệp hóa; từ 2012 đến nay
chuyển sang thặng dư thương mại nhờ khối FDI (Samsung, LG\ldots) xuất
khẩu điện tử. Tuy nhiên, vế Thu nhập Sơ cấp (Primary Income) thâm hụt
ngày càng lớn do doanh nghiệp FDI chuyển lợi nhuận về nước --- cho thấy
sự chênh lệch giữa GDP (tạo ra trên lãnh thổ) và GNI (thu nhập thực tế
của người dân Việt Nam).

\textbf{Ví dụ minh hoạ:} Xét nền kinh tế A năm 2025: - Xuất khẩu hàng
hóa \& dịch vụ: 300 tỷ USD; Nhập khẩu: 310 tỷ USD
\(\Rightarrow \text{Cán cân thương mại} = -10\) tỷ USD. - Chi trả lợi
nhuận FDI ra nước ngoài: −12 tỷ USD; Kiều hối nhận về: +8 tỷ USD. -
\(\Rightarrow CA = -10 - 12 + 8 = -14\) tỷ USD. - Dòng vốn FDI ròng giải
ngân: +18 tỷ USD; Dòng vốn đầu tư gián tiếp (FII) ròng: +2 tỷ USD. -
\(\Rightarrow FA = +20\) tỷ USD. - Giả sử \(EO = -1\) tỷ USD. - Thay vào
công thức: \(\Delta RAM = CA + FA + EO = -14 + 20 - 1 = +5\) tỷ USD. -
Kết luận: Mặc dù thâm hụt vãng lai 14 tỷ USD, nền kinh tế vẫn
\textbf{tăng dự trữ ngoại hối 5 tỷ USD} nhờ thu hút dòng vốn FDI và FII
vượt trội.

\begin{center}\rule{0.5\linewidth}{0.5pt}\end{center}

\subsection{Tỷ giá hối đoái \& Điều kiện Marshall--Lerner (Exchange Rate
\&
Marshall--Lerner)}\label{tux1ef7-giuxe1-hux1ed1i-ux111ouxe1i-ux111iux1ec1u-kiux1ec7n-marshalllerner-exchange-rate-marshalllerner}

\textbf{Cốt lõi:} Tỷ giá hối đoái danh nghĩa (\(E\)) là giá của một đơn
vị ngoại tệ tính bằng nội tệ (ví dụ 25.000 VND/USD; \(E\) tăng = nội tệ
mất giá). Tỷ giá hối đoái thực (\(e\) hoặc \(REER\)) đo lường giá cả
tương đối của hàng hóa nước ngoài so với hàng hóa nội địa:
\[ e = \frac{E \cdot P^*}{P} \] trong đó \(P^*\) là mức giá nước ngoài,
\(P\) là mức giá nội địa. Khi nội tệ mất giá thực (\(e\) tăng), hàng nội
địa trở nên rẻ hơn tương đối so với hàng ngoại, kích thích xuất khẩu
(\(X\)) và giảm nhập khẩu (\(M\)).

Tuy nhiên, phá giá/mất giá nội tệ \textbf{chỉ cải thiện cán cân thương
mại (\(CA\)) nếu Điều kiện Marshall--Lerner thỏa mãn}:
\[ |\epsilon_x| + |\epsilon_m| > 1 \] trong đó \(\epsilon_x\) là độ co
giãn của cầu xuất khẩu theo giá, \(\epsilon_m\) là độ co giãn của cầu
nhập khẩu theo giá. Nếu tổng độ co giãn nhỏ hơn 1, việc phá giá sẽ làm
\textbf{tồi tệ thêm} cán cân thương mại vì khối lượng hàng hóa chưa kịp
thay đổi nhưng giá trị tiền tệ nhập khẩu đã tăng ngay lập tức.

\textbf{Hiệu ứng Đường cong J (J-Curve Effect):} Trong ngắn hạn (vài
tháng đầu sau khi phá giá), các hợp đồng thương mại đã ký kết cố định số
lượng, độ co giãn rất thấp (\(|\epsilon_x| + |\epsilon_m| < 1\)), nên
cán cân thương mại \textbf{xấu đi trước}; chỉ trong trung và dài hạn khi
các doanh nghiệp điều chỉnh được hợp đồng và chuỗi cung ứng, khối lượng
hàng hóa mới phản ứng (\(|\epsilon_x| + |\epsilon_m| > 1\)), đưa cán cân
thương mại \textbf{cải thiện trở lại} --- tạo thành hình đường cong chữ
J.

\textbf{Điểm mù thường gặp:} Phản xạ đơn giản cho rằng ``muốn tăng xuất
khẩu chỉ cần phá giá đồng tiền''. Bỏ qua hai nguy cơ: \textbf{(1)} Nếu
nền kinh tế phụ thuộc nặng nề vào nguyên liệu nhập khẩu để sản xuất hàng
xuất khẩu, phá giá làm chi phí sản xuất tăng tương ứng, triệt tiêu lợi
thế giá; \textbf{(2)} Phá giá gây ra nhập khẩu lạm phát (\(P\) tăng),
đẩy tỷ giá thực \(e = E P^* / P\) trở lại vị trí cũ.

\textbf{Điểm mù thứ hai:} Bỏ qua cấu trúc đồng tiền thanh toán
(Invoicing Currency). Nếu toàn bộ xuất nhập khẩu đều định danh bằng USD
(Dominant Currency Pricing --- Gopinath 2015), việc phá giá đồng nội tệ
so với USD không làm giảm giá hàng xuất khẩu trên thị trường quốc tế
ngay lập tức, làm yếu đi kênh truyền dẫn giá truyền thống.

\textbf{Ứng dụng thực tế:} Khi Việt Nam điều chỉnh tăng tỷ giá USD/VND:
các ngành nông sản, thủy sản, dệt may (dùng nhiều nguyên liệu nội địa)
hưởng lợi rõ rệt. Nhưng ngành sản xuất linh kiện điện tử, thép (nhập
khẩu nguyên liệu đầu vào lớn bằng USD) chịu chi phí tăng vọt. Tổng thể,
việc giữ tỷ giá VND ổn định tương đối giúp kiểm soát lạm phát và tạo môi
trường vĩ mô dự đoán được cho dòng vốn FDI.

\textbf{Ví dụ minh hoạ:} Giả sử Việt Nam đang thâm hụt thương mại:
\(X = 100\) tỷ USD, \(M = 110\) tỷ USD. Ngân hàng Nhà nước phá giá VND
10\% so với USD. - \textbf{Trường hợp 1 --- Ngắn hạn (Độ co giãn thấp):}
\(\epsilon_x = 0{,}2\), \(\epsilon_m = 0{,}3\). Tổng \(= 0{,}5 < 1\). -
Lượng xuất khẩu chỉ tăng
\(10\% \times 0{,}2 = +2\% \Rightarrow X_{\text{mới}} = 102\) tỷ USD. -
Lượng nhập khẩu giảm \(10\% \times 0{,}3 = -3\%\), nhưng giá tính bằng
nội tệ/USD tăng 10\% \(\Rightarrow\) giá trị nhập khẩu thực tế tính theo
đồng tiền quy đổi tăng
\(\Rightarrow M_{\text{mới}} = 110 \times (1 - 0{,}03) \times 1{,}10 / 1{,}10 = \dots\)
Cán cân thương mại tính theo nội tệ bị \textbf{thâm hụt sâu hơn}. -
\textbf{Trường hợp 2 --- Dài hạn (Độ co giãn cao):}
\(\epsilon_x = 0{,}8\), \(\epsilon_m = 0{,}7\). Tổng \(= 1{,}5 > 1\). -
Lượng xuất khẩu tăng
\(10\% \times 0{,}8 = +8\% \Rightarrow X_{\text{mới}} = 108\) tỷ USD. -
Lượng nhập khẩu giảm
\(10\% \times 0{,}7 = -7\% \Rightarrow M_{\text{mới}} = 110 \times 0{,}93 = 102{,}3\)
tỷ USD. - Cán cân thương mại chuyển từ \textbf{thâm hụt −10 tỷ sang
thặng dư +5,7 tỷ USD} --- Marshall--Lerner nghiệm đúng.

\begin{center}\rule{0.5\linewidth}{0.5pt}\end{center}

\subsection{Định lý Stolper--Samuelson \& Phân phối thu nhập
(Stolper--Samuelson
Theorem)}\label{ux111ux1ecbnh-luxfd-stolpersamuelson-phuxe2n-phux1ed1i-thu-nhux1eadp-stolpersamuelson-theorem}

\textbf{Cốt lõi:} Nằm trong khung mô hình Heckscher--Ohlin (2x2x2), Định
lý Stolper--Samuelson (1941) phát biểu về tác động của thương mại tự do
lên phân phối thu nhập nội địa: \textbf{Khi mở cửa thương mại tự do, thu
nhập thực tế của chủ sở hữu yếu tố sản xuất dồi dào tương đối trong quốc
gia sẽ TĂNG LÊN, trong khi thu nhập thực tế của chủ sở hữu yếu tố sản
xuất khan hiếm tương đối sẽ GIẢM ĐI --- bất kể họ làm việc trong ngành
nào.}

Cơ chếMagnification Effect (Jones 1965): Sự tăng giá của hàng hóa xuất
khẩu đẩy giá yếu tố được sử dụng thâm dụng trong ngành đó tăng với
\textbf{tỷ lệ lớn hơn} tỷ lệ tăng giá hàng hóa:
\[ \hat{w} > \hat{P}_X > \hat{P}_Y > \hat{r} \quad (\text{khi ngành } X \text{ thâm dụng lao động } L) \]
Hệ quả trực tiếp: Thương mại tự do tạo ra \textbf{người thắng và người
thua tuyệt đối} trong nội bộ một quốc gia, dù tổng phúc lợi toàn quốc
gia đó tăng lên (Pareto gains potentially, không phải Pareto improvement
tự động).

\textbf{Điểm mù thường gặp:} Giả định rằng ``thương mại tự do tốt cho
tất cả mọi người'' và coi những nhóm phản đối mở cửa thương mại (như
công đoàn công nhân nhà máy tại các nước phát triển) là ``thiếu hiểu
biết''. Lý thuyết Stolper--Samuelson chứng minh tổn thất của họ là
\textbf{thực tế và có cơ sở lý thuyết chuẩn xác} --- lao động phổ thông
ở Mỹ thực sự chịu giảm thu nhập thực tế khi thương mại mở rộng với các
nước dồi dào lao động giá rẻ.

\textbf{Điểm mù thứ hai:} Cho rằng sự phản đối thương mại đến từ ``mâu
thuẫn giữa các ngành'' (ví dụ ngành dệt may vs ngành công nghệ).
Stolper--Samuelson chỉ ra mâu thuẫn thực sự nằm ở \textbf{giai cấp/yếu
tố sản xuất} (Lao động vs Tư bản/Kỹ năng), vì lao động và vốn có thể di
chuyển giữa các ngành trong dài hạn.

\textbf{Ứng dụng thực tế:} Hiệu ứng ``China Shock'' (Autor, Dorn \&
Hanson 2013): Việc Trung Quốc gia nhập WTO (2001) làm tăng tổng phúc lợi
kinh tế Mỹ (hàng hóa rẻ hơn cho người tiêu dùng), nhưng gây tổn thất
nặng nề cho công nhân sản xuất tại các vùng công nghiệp truyền thống
(Rust Belt) của Mỹ. Việc thiếu các chính sách tái phân phối hoặc đào tạo
lại hiệu quả đã dẫn đến sự gia tăng của chủ nghĩa bảo hộ và biến động
chính trị tại Mỹ sau đó.

\textbf{Ví dụ minh hoạ:} Xét nước phát triển H dồi dào Vốn (\(K\)), khan
hiếm Lao động (\(L\)). Nước đang phát triển F dồi dào Lao động (\(L\)),
khan hiếm Vốn (\(K\)). - Ngành X sản xuất Máy tính (thâm dụng Vốn
\(K\)). - Ngành Y sản xuất Giày da (thâm dụng Lao động \(L\)).

Khi H và F ký hiệp định thương mại tự do: - H chuyên môn hóa vào Máy
tính (giá Máy tính \(P_X\) tăng tương đối). - Theo Stolper--Samuelson,
tại nước H: - Lợi tức của Vốn (\(r\)) tăng \textbf{nhiều hơn} mức tăng
của \(P_X\) \(\Rightarrow\) Chủ sở hữu tư bản ở H \textbf{giàu lên thực
tế}. - Tiền lương của Lao động (\(w\)) \textbf{giảm thực tế} (mua được
ít cả Máy tính lẫn Giày da hơn trước). - Tác động ngược lại xảy ra tại
nước F: Tiền lương lao động phổ thông tăng thực tế, lợi tức vốn giảm
tương đối.

\begin{center}\rule{0.5\linewidth}{0.5pt}\end{center}

\subsection{FDI \& Năng lực hấp thụ (FDI \& Absorptive
Capacity)}\label{fdi-nux103ng-lux1ef1c-hux1ea5p-thux1ee5-fdi-absorptive-capacity}

\textbf{Cốt lõi:} Đầu tư trực tiếp nước ngoài (FDI) không chỉ mang theo
nguồn vốn tài chính mà quan trọng hơn là mang theo \textbf{công nghệ, kỹ
năng quản trị, và quyền truy cập thị trường toàn cầu}. Tuy nhiên, sự tồn
tại của FDI \textbf{không tự động đảm bảo} sự lan tỏa công nghệ
(technology spillovers) cho nền kinh tế nhận đầu tư.

Hàm lan tỏa công nghệ (Spillover Function):
\[ \Delta A_{\text{dom}} = f\left(\text{FDI}, \text{Absorptive Capacity}\right) = f\left(\text{FDI}, H, \text{R\&D}_{\text{dom}}, \text{TechGap}\right) \]
Trong đó \textbf{Năng lực hấp thụ} (Absorptive Capacity --- Cohen \&
Levinthal 1990) phụ thuộc vào: \textbf{(1)} Vốn con người (\(H\));
\textbf{(2)} Đầu tư R\&D nội địa; \textbf{(3)} Khoảng cách công nghệ
(Tech Gap). Nếu khoảng cách công nghệ quá lớn hoặc năng lực hấp thụ nội
địa quá yếu, FDI sẽ hoạt động dưới dạng \textbf{``ốc đảo công nghệ''}
(enclave economics) --- nhập khẩu toàn bộ linh kiện, lắp ráp bằng lao
động giá rẻ, rồi xuất khẩu, không tạo liên kết ngược/xuôi
(backward/forward linkages) với doanh nghiệp nội địa.

\textbf{Điểm mù thường gặp:} Đánh giá thành công của chính sách thu hút
FDI thuần túy qua \textbf{tổng số vốn đăng ký/giải ngân} hoặc tổng kim
ngạch xuất khẩu của khối FDI --- mà bỏ qua tỷ lệ giá trị gia tăng nội
địa (Domestic Value Added - DVA) và số lượng doanh nghiệp nội địa tham
gia được vào chuỗi cung ứng Tier-1, Tier-2.

\textbf{Điểm mù thứ hai:} Nhầm lẫn giữa ``chuyển giao công nghệ tự
động'' và ``tích lũy năng lực công nghệ chủ động''. Doanh nghiệp FDI
không có động cơ kinh tế để chuyển giao bí quyết công nghệ cốt lõi cho
đối thủ nội địa trừ khi bị buộc bởi áp lực cạnh tranh hoặc được thu hút
bởi năng lực đối ứng của đối tác nội địa.

\textbf{Ứng dụng thực tế:} Bài học đối lập giữa Hàn Quốc (thập niên
1960--1980) và một số nước Đông Nam Á: Hàn Quốc hạn chế FDI thuần túy,
ưu tiên mua bản quyền công nghệ (licensing) và vay nợ nước ngoài để các
Chaebol tự đầu tư R\&D tích lũy năng lực nội địa. Ngược lại, nhiều nước
phụ thuộc hoàn toàn vào FDI mà không có chính sách nâng cao năng lực hấp
thụ nội địa, dẫn đến nguy cơ bẫy thu nhập trung bình khi chi phí lao
động tăng.

\textbf{Ví dụ minh hoạ:} Xét 2 kịch bản của một ngành công nghiệp tại
quốc gia N nhận 10 tỷ USD FDI từ tập đoàn công nghệ X:

{\def\LTcaptype{none} % do not increment counter
\begin{longtable}[]{@{}
  >{\raggedright\arraybackslash}p{(\linewidth - 4\tabcolsep) * \real{0.3333}}
  >{\raggedright\arraybackslash}p{(\linewidth - 4\tabcolsep) * \real{0.3333}}
  >{\raggedright\arraybackslash}p{(\linewidth - 4\tabcolsep) * \real{0.3333}}@{}}
\toprule\noalign{}
\begin{minipage}[b]{\linewidth}\raggedright
Chỉ số
\end{minipage} & \begin{minipage}[b]{\linewidth}\raggedright
Kịch bản A (Năng lực hấp thụ yếu)
\end{minipage} & \begin{minipage}[b]{\linewidth}\raggedright
Kịch bản B (Năng lực hấp thụ mạnh)
\end{minipage} \\
\midrule\noalign{}
\endhead
\bottomrule\noalign{}
\endlastfoot
\textbf{Vốn con người nội địa (\(H\))} & Thấp (chỉ đáp ứng lao động phổ
thông) & Cao (kỹ sư đào tạo bài bản, chuẩn quốc tế) \\
\textbf{R\&D nội địa (\% GDP)} & 0,2\% & 1,5\% \\
\textbf{Tỷ lệ cung ứng nội địa (Localization)} & \textbf{5\%} (chỉ có
bao bì, suất ăn, bảo vệ) & \textbf{45\%} (linh kiện cơ khí chính xác,
phần mềm) \\
\textbf{Spillover công nghệ} & Gần bằng 0 (ốc đảo) & Rất cao (kỹ sư tách
ra lập công ty riêng) \\
\textbf{Giá trị giữ lại nội địa (DVA)} & \textbf{10 USD} trên 100 USD
xuất khẩu & \textbf{50 USD} trên 100 USD xuất khẩu \\
\end{longtable}
}

Kịch bản B cho thấy cùng một lượng vốn FDI nhưng tạo ra giá trị gia tăng
nội địa gấp 5 lần nhờ năng lực hấp thụ cao.

\begin{center}\rule{0.5\linewidth}{0.5pt}\end{center}

\subsection{Rào cản thương mại \& Lý thuyết hành động tập thể (Trade
Barriers \& Collective
Action)}\label{ruxe0o-cux1ea3n-thux1b0ux1a1ng-mux1ea1i-luxfd-thuyux1ebft-huxe0nh-ux111ux1ed9ng-tux1eadp-thux1ec3-trade-barriers-collective-action}

\textbf{Cốt lõi:} Thuế quan (Tariff), hạn ngạch (Quota), và rào cản phi
thuế quan (Non-Tariff Barriers - NTBs) gây ra \textbf{tổn thất vô ích
(Deadweight Loss - DWL)} cho tổng phúc lợi xã hội của một nền kinh tế
mở. Phân tích tổn thất phúc lợi từ thuế quan tại nước nhập khẩu nhỏ
(Small Open Economy): - Thặng dư tiêu dùng giảm: \(-(A + B + C + D)\) -
Thặng dư sản xuất tăng: \(+A\) - Thu ngân sách nhà nước tăng: \(+C\) -
\textbf{Tổn thất vô ích ròng (DWL):} \(-(B + D)\), trong đó \(B\) là tổn
thất do sản xuất không hiệu quả (production distortion), \(D\) là tổn
thất do tiêu dùng bị thu hẹp (consumption distortion).

Vậy tại sao chính phủ vẫn áp đặt rào cản thương mại? Lý thuyết Hành động
Tập thể (Olson 1965, Grossman \& Helpman 1994 --- ``Protection for
Sale'') giải thích: Group size \& Concentration: Nhóm ngành sản xuất nội
địa được hưởng lợi từ bảo hộ có \textbf{số lượng nhỏ, lợi ích tập trung
cao} \(\Rightarrow\) động cơ mạnh và chi phí phối hợp thấp để vận động
hành lang (lobby). Trong khi đó, người tiêu dùng gánh chịu chi phí bảo
hộ có \textbf{số lượng rất đông, lợi ích bị phân tán nhỏ lẻ}
\(\Rightarrow\) không có động cơ và chi phí tổ chức quá cao để tập hợp
lực lượng chống lại bảo hộ.

\textbf{Điểm mù thường gặp:} Lập luận ``Bảo hộ ngành công nghiệp non
trẻ'' (Infant Industry Argument --- List 1841) thường được dùng để biện
minh cho thuế quan. Tuy nhiên, điểm mù là: một khi thuế quan được ban
hành, ngành được bảo hộ sẽ dùng nguồn lực lobby để \textbf{vĩnh cửu hóa}
rào cản đó thay vì chịu áp lực nâng cao năng lực cạnh tranh để tháo bỏ
rào cản --- ``ngành công nghiệp non trẻ không bao giờ chịu lớn''.

\textbf{Điểm mù thứ hai:} Nhầm lẫn giữa ``bảo vệ việc làm trong một
ngành'' và ``tổng việc làm của nền kinh tế''. Thuế quan bảo vệ việc làm
trong ngành được bảo hộ, nhưng làm tăng chi phí đầu vào của các ngành hạ
nguồn và gây tăng giá cho người tiêu dùng, triệt tiêu số lượng việc làm
lớn hơn ở các ngành khác.

\textbf{Ứng dụng thực tế:} Thuế quan thép và nhôm của Mỹ (Section 232):
Áp thuế 25\% lên thép nhập khẩu giúp bảo vệ khoảng 1.000 việc làm trong
ngành sản xuất thép Mỹ, nhưng làm tăng chi phí đầu vào của các ngành
tiêu thụ thép (ô tô, xây dựng, máy móc) làm mất đi ước tính khoảng
75.000 việc làm trong các ngành này (Schrecker et al.~2019) --- đúng
theo dự đoán tổn thất ròng của lý thuyết kinh tế.

\textbf{Ví dụ minh hoạ:} Thị trường nội địa sản phẩm X. Giá thế giới
\(P_w = 100\). - Phương trình cầu nội địa: \(Q_d = 400 - P\) - Phương
trình cung nội địa: \(Q_s = -50 + P\)

\textbf{Khi thương mại tự do (bỏ thuế):} - \(P = P_w = 100\). -
\(Q_d = 400 - 100 = 300\); \(Q_s = -50 + 100 = 50\). - Nhập khẩu
\(= 300 - 50 = 250\).

\textbf{Khi chính phủ áp thuế quan \(t = 20\) / đơn vị
\(\Rightarrow P_{\text{mới}} = 120\):} - \(Q_d' = 400 - 120 = 280\);
\(Q_s' = -50 + 120 = 70\). - Nhập khẩu mới \(= 280 - 70 = 210\). - Thặng
dư tiêu dùng giảm:
\(\Delta CS = -\left[280 \times 20 + \frac{1}{2}(300-280) \times 20\right] = -5.600 - 200 = -5.800\).
- Thặng dư sản xuất tăng:
\(\Delta PS = +\left[50 \times 20 + \frac{1}{2}(70-50) \times 20\right] = +1.000 + 200 = +1.200\).
- Thu thuế chính phủ: \(T = 210 \times 20 = 4.200\). - \textbf{Tổn thất
vô ích ròng (DWL):}
\(\Delta CS + \Delta PS + T = -5.800 + 1.200 + 4.200 = -400\). - Phân
tích: Tổn thất sản xuất không hiệu quả
\(B = \frac{1}{2}(70-50) \times 20 = 200\); Tổn thất tiêu dùng bị thu
hẹp \(D = \frac{1}{2}(300-280) \times 20 = 200\). Tổng
\(DWL = B + D = 400\). Xã hội \textbf{mất trắng 400 đơn vị tiền tệ} từ
chính sách này.

\begin{center}\rule{0.5\linewidth}{0.5pt}\end{center}

\subsection{REER \& Cạnh tranh (Real Effective Exchange
Rate)}\label{reer-cux1ea1nh-tranh-real-effective-exchange-rate}

\textbf{Cốt lõi:} Tỷ giá hối đoái thực hiệu dụng (REER --- Real
Effective Exchange Rate) là chỉ số đo lường sức cạnh tranh giá cả tổng
hợp của một quốc gia so với rổ các đối tác thương mại chính, đã điều
chỉnh theo lạm phát:
\[ \text{REER} = \sum_{i=1}^{k} w_i \cdot \left( \frac{E_i \cdot P^*_i}{P} \right) \]
trong đó \(w_i\) là trọng số thương mại của đối tác \(i\), \(E_i\) là tỷ
giá song phương, \(P^*_i\) là chỉ số giá của đối tác \(i\), và \(P\) là
chỉ số giá nội địa. - \textbf{REER tăng (Real Appreciation):} Hàng hóa
nội địa trở nên \textbf{đắt hơn tương đối} so với đối tác
\(\Rightarrow\) \textbf{Sức cạnh tranh giá giảm}. - \textbf{REER giảm
(Real Depreciation):} Hàng hóa nội địa trở nên \textbf{rẻ hơn tương đối}
\(\Rightarrow\) \textbf{Sức cạnh tranh giá tăng}.

Khái niệm \textbf{Phá giá nội địa (Internal Devaluation)}: Khi một nước
nằm trong liên minh tiền tệ (như Hy Lạp trong Eurozone) không thể tự phá
giá đồng tiền danh nghĩa (\(E\) cố định), cách duy nhất để giảm REER
nhằm khôi phục sức cạnh tranh là \textbf{nội giảm giá (\(P\) giảm)} qua
việc cắt giảm tiền lương danh nghĩa và chi tiêu công --- một quá trình
vô cùng đau đớn và gây suy thoái kinh tế nghiêm trọng.

\textbf{Điểm mù thường gặp:} Nhìn vào sự ổn định của tỷ giá danh nghĩa
song phương (ví dụ USD/VND) và kết luận rằng ``sức cạnh tranh tỷ giá
đang ổn định'' --- bỏ qua biến động lạm phát tương đối và biến động của
các đồng tiền khác trong rổ thương mại (EUR, JPY, CNY). Nếu các đối thủ
cạnh tranh (như Nhân dân tệ hay Taka Bangladesh) mất giá mạnh so với USD
trong khi VND giữ nguyên giá so với USD, \textbf{REER của VND thực chất
đã tăng} \(\Rightarrow\) Việt Nam bị mất sức cạnh tranh giá.

\textbf{Điểm mù thứ hai:} Coi REER giảm là mục tiêu duy nhất. REER thấp
dựa trên tiền lương rẻ không phải là chiến lược cạnh tranh bền vững; cải
thiện năng suất lao động (TFP) để nâng cao sức cạnh tranh phi giá cả
(chất lượng, thương hiệu) mới là mục tiêu dài hạn.

\textbf{Ứng dụng thực tế:} Cuộc khủng hoảng nợ công Hy Lạp (2010--2015):
Hy Lạp bị tăng REER mạnh trong thập niên 2000 do lạm phát nội địa cao
hơn mức trung bình Eurozone. Vì không thể phá giá đồng Euro, Hy Lạp phải
thực hiện ``internal devaluation'' theo yêu cầu của Troika --- cắt giảm
20--30\% lương danh nghĩa, dẫn đến thất nghiệp vọt lên 27\% và GDP sụt
giảm 25\%.

\textbf{Ví dụ minh hoạ:} So sánh sức cạnh tranh xuất khẩu dệt may giữa
Việt Nam và Bangladesh sang thị trường Mỹ năm \(t\) và \(t+1\): - Năm
\(t\): Tỷ giá USD/VND = 24.000; Tỷ giá USD/BDT = 100. Chỉ số giá cả ba
nước = 100. - Năm \(t+1\): - Lạm phát Mỹ: 2\%; Lạm phát Việt Nam: 4\%;
Lạm phát Bangladesh: 6\%. - Tỷ giá USD/VND giữ nguyên 24.000 (VND cố
định). - Bangladesh cho đồng BDT mất giá 10\% so với USD \(\Rightarrow\)
USD/BDT = 110.

\textbf{Tính biến động REER tương đối so với Mỹ:} - Đối với Việt Nam:
\(REER_{\text{VN}} = \frac{E_{\text{mới}}}{E_{\text{gốc}}} \times \frac{P_{\text{Mỹ}}}{P_{\text{VN}}} = 1 \times \frac{102}{104} = 0{,}9807\)
(REER giảm nhẹ 1,93\% nhờ lạm phát Mỹ, tăng sức cạnh tranh nhẹ so với
Mỹ). - Đối với Bangladesh:
\(REER_{\text{BDT}} = \frac{110}{100} \times \frac{102}{106} = 1{,}10 \times 0{,}9622 = 1{,}0584\)
\(\Rightarrow\) BDT mất giá thực \textbf{5,84\%} so với USD. -
\textbf{Kết quả so sánh trực tiếp VN vs Bangladesh:} Hàng dệt may
Bangladesh trở nên \textbf{rẻ hơn tương đối khoảng 7,7\%} so với hàng
Việt Nam tại Mỹ. Mặc dù lạm phát Bangladesh cao hơn, việc chủ động phá
giá đồng BDT đã giúp Bangladesh giành thế mạnh cạnh tranh giá so với
Việt Nam.

\begin{center}\rule{0.5\linewidth}{0.5pt}\end{center}

\subsubsection{4.8 Mô hình Melitz (2003): Thương mại \& Doanh nghiệp Dị
biệt (Heterogeneous Firms \&
Trade)}\label{muxf4-huxecnh-melitz-2003-thux1b0ux1a1ng-mux1ea1i-doanh-nghiux1ec7p-dux1ecb-biux1ec7t-heterogeneous-firms-trade}

\subsubsection{4.8.1 Cốt lõi Lý thuyết (Core
Theory)}\label{cux1ed1t-luxf5i-luxfd-thuyux1ebft-core-theory}

Trong khi các mô hình thương mại truyền thống như Ricardo hay
Heckscher-Ohlin giải thích thương mại quốc tế dựa trên lợi thế so sánh ở
cấp độ quốc gia hoặc ngành, và mô hình Krugman (1980) sử dụng cạnh tranh
độc quyền để giải thích thương mại nội ngành, chúng đều không thể lý
giải một hiện tượng thực nghiệm quan trọng: trong cùng một ngành, chỉ có
một số ít doanh nghiệp tham gia xuất khẩu, và chúng thường có năng suất
cao hơn các doanh nghiệp chỉ phục vụ thị trường nội địa. Mô hình Melitz
(2003) giải quyết vấn đề này bằng cách đưa sự dị biệt của doanh nghiệp
(firm heterogeneity) vào mô hình cạnh tranh độc quyền.

\textbf{Cấu trúc mô hình cơ bản:} Xét một nền kinh tế với lao động là
yếu tố sản xuất duy nhất, được chuẩn hóa sao cho mức lương \$ w = 1 \$.
Phía cầu được đặc trưng bởi hàm thỏa dụng CES (Constant Elasticity of
Substitution):
\[ U = \left[ \int_{\omega \in \Omega} q(\omega)^{\frac{\sigma-1}{\sigma}} d\omega \right]^{\frac{\sigma}{\sigma-1}} \]
trong đó \$ \sigma \textgreater{} 1 \$ là độ co giãn thay thế giữa các
hàng hóa, và \$ \Omega \$ là tập hợp các loại hàng hóa hiện có. Cầu cho
mỗi sản phẩm là:
\[ q(\omega) = Q \left( \frac{p(\omega)}{P} \right)^{-\sigma} \]

Ở phía cung, các doanh nghiệp phải chịu một chi phí chìm \$ f\_e \$ để
gia nhập ngành. Sau khi trả \$ f\_e \$, doanh nghiệp bốc thăm một mức
năng suất (productivity) \$ \varphi \$ từ một hàm phân phối xác suất \$
g(\varphi) \$ với hàm phân phối tích lũy \$ G(\varphi) \$. Hàm chi phí
của một doanh nghiệp có năng suất \$ \varphi \$ là:
\[ C(q) = f + \frac{q}{\varphi} \] trong đó \$ f \$ là chi phí cố định
để duy trì sản xuất tại thị trường nội địa. Doanh nghiệp định giá theo
quy tắc mark-up trên chi phí biên:
\[ p(\varphi) = \frac{\sigma}{\sigma-1} \frac{1}{\varphi} \]

Lợi nhuận từ thị trường nội địa là một hàm tăng theo năng suất \$
\varphi \$:
\[ \pi_d(\varphi) = \frac{R}{\sigma} \left( \frac{P \rho \varphi}{1} \right)^{\sigma-1} - f \]
với \$ \rho = \frac{\sigma-1}{\sigma} \$ và \$ R \$ là tổng chi tiêu.

\textbf{Điều kiện Tồn tại và Xuất khẩu (Cutoff Conditions):} Do có chi
phí cố định \$ f \$, không phải doanh nghiệp nào cũng tồn tại. Chỉ những
doanh nghiệp có năng suất đủ lớn, \$ \varphi \geq \varphi\^{}* \$, mới
có lợi nhuận không âm. Ngưỡng năng suất nội địa \$ \varphi\^{}* \$ được
xác định bởi điều kiện lợi nhuận bằng 0 (ZCP - Zero Cutoff Profit):
\[ \pi_d(\varphi^*) = 0 \iff \frac{R}{\sigma} (P \rho \varphi^*)^{\sigma-1} = f \]

Nếu doanh nghiệp muốn xuất khẩu, họ phải chịu thêm chi phí cố định xuất
khẩu \$ f\_x \$ và chi phí thương mại dạng tảng băng (iceberg trade
costs) \$ \tau \textgreater{} 1 \$. Lợi nhuận từ xuất khẩu là:
\[ \pi_x(\varphi) = \frac{R^*}{\sigma} \left( \frac{P^* \rho \varphi}{\tau} \right)^{\sigma-1} - f_x \]
Tương tự, ngưỡng năng suất xuất khẩu \$ \varphi\_x\^{}* \$ được xác định
bởi: \[ \pi_x(\varphi_x^*) = 0 \] Do \$ f\_x \$ và \$ \tau \$ thường đủ
lớn so với quy mô thị trường, ta có kết quả cốt lõi: \$ \varphi\_x\^{}*
\textgreater{} \varphi\^{}* \$. Điều này có nghĩa là sự tự chọn
(self-selection): chỉ những doanh nghiệp năng suất nhất mới xuất khẩu.

\textbf{Điều kiện Gia nhập Tự do (Free Entry Condition):} Lợi nhuận kỳ
vọng trước khi gia nhập phải bằng chi phí chìm \$ f\_e \$:
\[ \int_{\varphi^*}^\infty \pi(\varphi) g(\varphi) d\varphi = f_e \]

Khi chi phí thương mại \$ \tau \$ giảm (tự do hóa thương mại), cơ hội
xuất khẩu tăng lên đẩy nhu cầu lao động lên cao, làm tăng lương thực tế.
Điều này đẩy ngưỡng năng suất nội địa \$ \varphi\^{}* \$ lên cao, loại
bỏ các doanh nghiệp kém hiệu quả nhất (Melitz effect), đồng thời làm
giảm \$ \varphi\_x\^{}* \$. Kết quả là năng suất trung bình của toàn
ngành tăng lên, tạo ra một nguồn lợi ích từ thương mại mới: sự phân bổ
lại nguồn lực (reallocation).

\subsubsection{4.8.2 Điểm mù thường gặp (Cognitive Traps \&
Limitations)}\label{ux111iux1ec3m-muxf9-thux1b0ux1eddng-gux1eb7p-cognitive-traps-limitations}

\begin{enumerate}
\def\labelenumi{\arabic{enumi}.}
\tightlist
\item
  \textbf{Đồng nhất hóa Tiến bộ công nghệ và Phân bổ lại:} Sinh viên
  thường nhầm lẫn giữa việc năng suất trung bình tăng trong mô hình
  Melitz với tiến bộ công nghệ (technological progress). Trong Melitz,
  năng suất của \emph{từng doanh nghiệp} không hề thay đổi sau khi bốc
  thăm. Năng suất tổng thể tăng lên hoàn toàn là do sự phân bổ lại lao
  động từ các doanh nghiệp năng suất thấp (bị phá sản) sang các doanh
  nghiệp năng suất cao (mở rộng xuất khẩu).
\item
  \textbf{Vai trò của chi phí cố định xuất khẩu (\(f_x\)):} Nhiều người
  cho rằng chỉ cần chi phí vận chuyển (\$ \tau \$) là đủ để giải thích
  tại sao không phải ai cũng xuất khẩu. Tuy nhiên, nếu \$ f\_x = 0
  \(, mọi doanh nghiệp đang sản xuất (\) \varphi \ge \varphi\^{}* \$)
  đều sẽ xuất khẩu một lượng nhỏ. Sự hiện diện của \$ f\_x \$ là điểm
  cốt lõi tạo ra việc chia rẽ thị trường thành hai nhóm (người xuất khẩu
  và không xuất khẩu).
\end{enumerate}

\subsubsection{4.8.3 Ứng dụng thực tế (Practical
Applications)}\label{ux1ee9ng-dux1ee5ng-thux1ef1c-tux1ebf-practical-applications}

\subsubsection{4.8.3 Ứng dụng thực tế (Practical
Applications)}\label{ux1ee9ng-dux1ee5ng-thux1ef1c-tux1ebf-practical-applications-1}

\subsubsection{4.8.4 Ví dụ minh họa tính toán (Worked Numerical
Examples)}\label{vuxed-dux1ee5-minh-hux1ecda-tuxednh-touxe1n-worked-numerical-examples}

Giả sử năng suất \$ \varphi \$ tuân theo phân phối Pareto \$ G(\varphi)
= 1 - \left(\frac{b}{\varphi}\right)\^{}k \$ với \$ b=1, k=4 \$. Các
thông số kinh tế: \$ \sigma = 3 \$, \$ f = 10 \$, \$ f\_x = 20 \$, \$
\tau = 1.5 \$. Ta thiết lập sự cân bằng tĩnh: Theo định lý ZCP và tỷ lệ
giữa lợi nhuận nội địa và xuất khẩu, ta có tỷ lệ giữa hai ngưỡng năng
suất trong điều kiện đối xứng:
\[ \varphi_x^* = \tau \left( \frac{f_x}{f} \right)^{\frac{1}{\sigma-1}} \varphi^* \]
Thay số: \$ \varphi\_x\^{}* = 1.5 \times \left( \frac{20}{10}
\right)\^{}\{\frac{1}{2}\} \varphi\^{}* = 1.5 \times \sqrt{2}
\varphi\^{}* \approx 2.12 \varphi\^{}* \$. Điều này minh họa rõ rằng
ngưỡng năng suất để xuất khẩu cao hơn gấp đôi ngưỡng để tồn tại ở thị
trường nội địa. Doanh nghiệp phải hiệu quả gấp đôi mới bù đắp được chi
phí vận chuyển và chi phí cố định thâm nhập thị trường nước ngoài.

\begin{center}\rule{0.5\linewidth}{0.5pt}\end{center}

\begin{center}\rule{0.5\linewidth}{0.5pt}\end{center}

\section{CHƯƠNG 5: TƯƠNG TÁC CHIẾN LƯỢC \& LÝ THUYẾT TRÒ CHƠI CƠ BẢN
(Strategic Interaction \& Game
Theory)}\label{chux1b0ux1a1ng-5-tux1b0ux1a1ng-tuxe1c-chiux1ebfn-lux1b0ux1ee3c-luxfd-thuyux1ebft-truxf2-chux1a1i-cux1a1-bux1ea3n-strategic-interaction-game-theory}

Khi kết quả của một quyết định phụ thuộc vào việc \textbf{người khác}
quyết định gì --- đây là địa hạt của lý thuyết trò chơi. Không giống các
bài toán tối ưu đơn giản ở Chương 3 (tối đa hóa lợi ích trong môi trường
thụ động), tương tác chiến lược đòi hỏi ta mô hình hóa \emph{tư duy của
đối thủ} để tìm ra hành động tốt nhất của mình --- và hành động tốt nhất
của mình lại phụ thuộc vào dự đoán của ta về đối thủ, vòng lặp đó có thể
tiếp tục vô tận. Nhóm này cung cấp công cụ để cắt vòng lặp đó và tìm ra
điểm ổn định --- cân bằng --- cũng như thiết kế các cơ chế thay đổi điểm
ổn định đó theo hướng có lợi hơn. (Xem thêm Phần B, Chương 11 để có bản
hình thức hóa đầy đủ: Cân bằng Bayes--Nash, thiết kế cơ chế, đấu giá và
định lý doanh thu tương đương.)

\begin{center}\rule{0.5\linewidth}{0.5pt}\end{center}

\subsection{5.1 Cân bằng Nash (Nash
Equilibrium)}\label{cuxe2n-bux1eb1ng-nash-nash-equilibrium}

\subsubsection{5.1.1 Cốt lõi Lý thuyết (Core
Theory)}\label{cux1ed1t-luxf5i-luxfd-thuyux1ebft-core-theory-1}

Về mặt hình thức, trong một trò chơi \(n\) người chơi, cân bằng Nash là
bộ chiến lược \((s_1^*, s_2^*, \ldots, s_n^*)\) sao cho với mọi \(i\):
u\_i(s\_i\^{}\emph{, s\_\{-i\}\^{}}) \ge u\_i(s\_i, s\_\{-i\}\^{}*)
\quad \forall s\_i \in S\_i tức \(s_i^*\) là \textbf{phản ứng tốt nhất}
(best response) của người chơi \(i\) trước chiến lược \(s_{-i}^*\) của
tất cả người khác. Điều này đúng cho cả \textbf{chiến lược thuần túy}
(pure strategy --- chọn một hành động xác định) và \textbf{chiến lược
hỗn hợp} (mixed strategy --- xác suất hóa trên tập chiến lược thuần
túy). Định lý Nash đảm bảo cân bằng hỗn hợp luôn tồn tại khi cân bằng
thuần túy không tồn tại.

Phân biệt hai khái niệm hay bị gộp làm một: \textbf{trạng thái ổn định}
(không ai muốn lệch đơn phương --- định nghĩa Nash) và \textbf{kết quả
tốt nhất cho tất cả} (tối ưu Pareto). Cân bằng Nash không đảm bảo hiệu
quả Pareto --- thực tế, nhiều cân bằng Nash nổi tiếng nhất lại
\textbf{kém hơn} tất cả đối với mọi bên so với kết quả hợp tác có thể
đạt được. Đây là điểm cốt lõi của Thế lưỡng nan người tù ngay bên dưới.

\subsubsection{5.1.2 Điểm mù thường gặp (Cognitive Traps \&
Limitations)}\label{ux111iux1ec3m-muxf9-thux1b0ux1eddng-gux1eb7p-cognitive-traps-limitations-1}

\subsubsection{5.1.2 Điểm mù thường gặp (Cognitive Traps \&
Limitations)}\label{ux111iux1ec3m-muxf9-thux1b0ux1eddng-gux1eb7p-cognitive-traps-limitations-2}

Đa cân bằng (Multiple equilibria) và vấn đề chọn lựa cân bằng
(Equilibrium selection). Sinh viên thường cho rằng cân bằng Nash là duy
nhất hoặc có tính dự báo cao, nhưng theo Định lý Dân gian (Folk
Theorem), nhiều trò chơi lặp lại có vô số cân bằng. Điều này đòi hỏi các
khái niệm tinh chỉnh (refinements) như cân bằng Focal point (Schelling).

\subsubsection{5.1.3 Ứng dụng thực tế (Practical
Applications)}\label{ux1ee9ng-dux1ee5ng-thux1ef1c-tux1ebf-practical-applications-2}

Dự báo kết quả đàm phán thương mại (như chiến tranh thuế quan Mỹ-Trung)
hoặc phân tích sự ổn định của các hiệp ước vũ khí hạt nhân quốc tế.

\subsubsection{5.1.4 Ví dụ minh họa tính toán (Worked Numerical
Examples)}\label{vuxed-dux1ee5-minh-hux1ecda-tuxednh-touxe1n-worked-numerical-examples-1}

{\def\LTcaptype{none} % do not increment counter
\begin{longtable}[]{@{}lll@{}}
\toprule\noalign{}
& \textbf{B chọn C} & \textbf{B chọn T} \\
\midrule\noalign{}
\endhead
\bottomrule\noalign{}
\endlastfoot
\textbf{A chọn C} & A=10, B=10 & A=2, B=14 \\
\textbf{A chọn T} & A=14, B=2 & A=5, B=5 \\
\end{longtable}
}

\textbf{Tìm cân bằng Nash bằng phương pháp phản ứng tốt nhất:} - Nếu B
chọn C: A so sánh 10 (chọn C) với 14 (chọn T) \(\Rightarrow\) A chọn T.
- Nếu B chọn T: A so sánh 2 (chọn C) với 5 (chọn T) \(\Rightarrow\) A
chọn T. - \(\Rightarrow\) Chiến lược trội của A là \textbf{luôn chọn T}
(không phụ thuộc vào B). - Tương tự, chiến lược trội của B cũng là
\textbf{luôn chọn T}.

\textbf{Cân bằng Nash duy nhất: (T, T) với lợi nhuận (5, 5).} Nhưng cả
hai đều thích (C, C) hơn vì đem lại (10, 10) --- đây chính là cấu trúc
thế lưỡng nan người tù: cân bằng Nash không phải kết quả tốt nhất cho
mọi người.

\begin{center}\rule{0.5\linewidth}{0.5pt}\end{center}

\subsection{5.2 Thế lưỡng nan Người tù \& Trò chơi lặp lại (Prisoner's
Dilemma \& Repeated
Games)}\label{thux1ebf-lux1b0ux1ee1ng-nan-ngux1b0ux1eddi-tuxf9-truxf2-chux1a1i-lux1eb7p-lux1ea1i-prisoners-dilemma-repeated-games}

\subsubsection{5.2.1 Cốt lõi Lý thuyết (Core
Theory)}\label{cux1ed1t-luxf5i-luxfd-thuyux1ebft-core-theory-2}

Nhưng điều gì xảy ra khi trò chơi \textbf{lặp lại vô hạn} (hay số lần
chưa biết trước)? \textbf{Định lý Dân gian (Folk Theorem --- Aumann \&
Shapley, 1976; Fudenberg \& Maskin, 1986)} trả lời: bất kỳ lợi ích trung
bình chiết khấu nào nằm trong tập hợp \emph{cá nhân hợp lý}
(individually rational payoffs --- ít nhất bằng lợi ích minimax của mỗi
bên) đều có thể được duy trì như một cân bằng Nash trong trò chơi lặp
lại vô hạn, với hệ số chiết khấu \(\delta\) đủ cao. Về trực giác: khi
bóng tối của tương lai đủ dài và đủ nặng, đe dọa ``nếu anh phản bội hôm
nay, tôi sẽ không bao giờ hợp tác lại'' trở nên đáng tin ---
\textbf{chiến lược ăn miếng trả miếng} (tit-for-tat --- hợp tác ở vòng
đầu, sau đó sao chép hành động của đối phương ở vòng trước) được Axelrod
(1984) chứng minh rất hiệu quả trong thực nghiệm.

Điều kiện hợp tác bền vững trong trò chơi lặp lại: lợi ích từ hợp tác
dài hạn phải vượt cám dỗ phản bội một lần rồi chịu trừng phạt. Bất đẳng
thức điều kiện: \(\frac{R}{1-\delta} \ge T + \frac{P}{1-\delta}\), suy
ra \(\delta \ge \frac{T-R}{T-P}\). Tức hệ số chiết khấu (coi trọng tương
lai) phải đủ cao --- khi \(\delta\) thấp (tương lai ít quan trọng, hoặc
quan hệ dự kiến ngắn hạn), lý thuyết dự đoán hợp tác sụp đổ, ngay cả khi
các bên đều muốn hợp tác.

\subsubsection{5.2.2 Điểm mù thường gặp (Cognitive Traps \&
Limitations)}\label{ux111iux1ec3m-muxf9-thux1b0ux1eddng-gux1eb7p-cognitive-traps-limitations-3}

\subsubsection{5.2.2 Điểm mù thường gặp (Cognitive Traps \&
Limitations)}\label{ux111iux1ec3m-muxf9-thux1b0ux1eddng-gux1eb7p-cognitive-traps-limitations-4}

Cạm bẫy tư duy ngắn hạn và bỏ qua giá trị của yếu tố chiết khấu
(\(\delta\)). Hợp tác chỉ bền vững nếu các bên đủ kiên nhẫn (chiết khấu
dòng tiền tương lai không quá mạnh). Nếu \(\delta\) quá thấp (tương lai
ít ý nghĩa), mọi cơ chế hợp tác đều sụp đổ.

\subsubsection{5.2.3 Ứng dụng thực tế (Practical
Applications)}\label{ux1ee9ng-dux1ee5ng-thux1ef1c-tux1ebf-practical-applications-3}

Giải thích sự ổn định và sụp đổ luân phiên của Cartel dầu mỏ (OPEC) và
các cuộc chiến tranh giá (price wars) trong ngành hàng không.

\subsubsection{5.2.4 Ví dụ minh họa tính toán (Worked Numerical
Examples)}\label{vuxed-dux1ee5-minh-hux1ecda-tuxednh-touxe1n-worked-numerical-examples-2}

{\def\LTcaptype{none} % do not increment counter
\begin{longtable}[]{@{}lll@{}}
\toprule\noalign{}
& \textbf{B: Không QC} & \textbf{B: Quảng cáo} \\
\midrule\noalign{}
\endhead
\bottomrule\noalign{}
\endlastfoot
\textbf{A: Không QC} & A=20, B=20 & A=8, B=25 \\
\textbf{A: Quảng cáo} & A=25, B=8 & A=12, B=12 \\
\end{longtable}
}

Cân bằng Nash một lần: (Quảng cáo, Quảng cáo) = (12, 12). Kết quả tốt
hơn: (Không QC, Không QC) = (20, 20). \(T=25, R=20, P=12, S=8\). Điều
kiện hợp tác:
\(\delta \ge \frac{T-R}{T-P} = \frac{25-20}{25-12} = \frac{5}{13} \approx 0{,}385\).
Nếu hệ số chiết khấu theo quý của mỗi hãng là \(\delta = 0{,}9\) (tương
đương lãi suất kỳ vọng \textasciitilde11\%/năm, quan hệ dài hạn), thì
\(0{,}9 > 0{,}385\) --- điều kiện hợp tác thỏa mãn. Chiến lược
tit-for-tat dưới dạng ``không quảng cáo ở quý đầu; nếu đối phương quảng
cáo quý trước thì quảng cáo quý này'' tạo ra cân bằng hợp tác bền vững
với lợi nhuận kỳ vọng mỗi bên: \(20/(1-0{,}9) = 200\) triệu (theo giá
trị hiện tại tổng các quý tương lai), so với \(12/(1-0{,}9) = 120\)
triệu nếu cùng quảng cáo --- hợp tác mang lại thêm \textbf{80 triệu} giá
trị hiện tại, đủ động lực để duy trì.

\begin{center}\rule{0.5\linewidth}{0.5pt}\end{center}

\subsection{5.3 Trò chơi tuần tự \& Cam kết không thể đảo ngược
(Sequential games \&
Commitment)}\label{truxf2-chux1a1i-tuux1ea7n-tux1ef1-cam-kux1ebft-khuxf4ng-thux1ec3-ux111ux1ea3o-ngux1b0ux1ee3c-sequential-games-commitment}

\subsubsection{5.3.1 Cốt lõi Lý thuyết (Core
Theory)}\label{cux1ed1t-luxf5i-luxfd-thuyux1ebft-core-theory-3}

Tại sao trò chơi tuần tự tạo ra kết quả khác trò chơi đồng thời? Người
đi \textbf{trước} có \textbf{lợi thế cam kết}: bằng cách hành động
trước, họ thay đổi ràng buộc và cấu trúc khuyến khích của toàn bộ trò
chơi. Lợi thế này đặc biệt mạnh khi hành động đó là \textbf{không thể
đảo ngược} --- đầu tư vào nhà máy lớn, ký hợp đồng dài hạn, công bố công
khai mức giá --- vì tính không thể đảo ngược biến cam kết thành
\textbf{đáng tin cậy} (credible). Ngược lại, một đe dọa hoặc cam kết
\textbf{có thể đảo ngược} dễ dàng không đáng tin theo SPE, vì người đe
dọa biết rằng khi đến thời điểm thực hiện, họ có thể tốt hơn khi không
làm vậy --- đây là ``đe dọa rỗng'' (empty threat).

Mô hình Stackelberg (1934) là ứng dụng kinh điển nhất: trong cạnh tranh
sản lượng, người đi trước (Stackelberg leader) công bố sản lượng của
mình, buộc người đi sau phản ứng. Leader chọn sản lượng cao hơn Cournot
để cam kết không co lại, kéo follower về điểm sản lượng thấp hơn theo
hàm phản ứng. Kết quả: leader có lợi nhuận cao hơn follower và cao hơn
mức Cournot --- \textbf{lợi thế đi trước} (first-mover advantage) thực
sự tồn tại trong cạnh tranh sản lượng.

\subsubsection{5.3.2 Điểm mù thường gặp (Cognitive Traps \&
Limitations)}\label{ux111iux1ec3m-muxf9-thux1b0ux1eddng-gux1eb7p-cognitive-traps-limitations-5}

\subsubsection{5.3.2 Điểm mù thường gặp (Cognitive Traps \&
Limitations)}\label{ux111iux1ec3m-muxf9-thux1b0ux1eddng-gux1eb7p-cognitive-traps-limitations-6}

Đe dọa không đáng tin cậy (Incredible threats). Lỗi phổ biến là tin vào
những lời đe dọa không tối ưu hóa lợi ích của người đe dọa tại thời điểm
họ phải thực thi nó. Phải dùng quy nạp ngược (Backward Induction) để
loại bỏ các cân bằng Nash không hoàn hảo cục bộ (non-SPNE).

\subsubsection{5.3.3 Ứng dụng thực tế (Practical
Applications)}\label{ux1ee9ng-dux1ee5ng-thux1ef1c-tux1ebf-practical-applications-4}

Chiến lược ``đốt thuyền'' (Burning bridges) trong quân sự hoặc cam kết
công suất không thể hoàn tác của các hãng sản xuất để răn đe đối thủ
mới.

\subsubsection{5.3.4 Ví dụ minh họa tính toán (Worked Numerical
Examples)}\label{vuxed-dux1ee5-minh-hux1ecda-tuxednh-touxe1n-worked-numerical-examples-3}

\textbf{Bước 1 --- Tìm hàm phản ứng của Follower B:} B tối đa hóa
\(\pi_B = (P - 10)q_B = (90 - q_A - q_B)q_B\). Đặt
\(\partial \pi_B / \partial q_B = 0\):
\(90 - q_A - 2q_B = 0 \Rightarrow q_B^*(q_A) = (90 - q_A)/2\) --- đây là
hàm phản ứng tốt nhất của B.

\textbf{Bước 2 --- Leader A quy nạp ngược, thế vào hàm phản ứng của B:}
\(\pi_A = (P-10)q_A = (90 - q_A - q_B^*(q_A))q_A = (90 - q_A - (90-q_A)/2)q_A = ((90-q_A)/2)q_A\).
Tối đa hóa:
\(\partial \pi_A / \partial q_A = (90 - 2q_A)/2 = 0 \Rightarrow q_A^* = 45\).

\textbf{Bước 3 --- Kết quả:} \(q_B^* = (90-45)/2 = 22{,}5\);
\(Q = 67{,}5\); \(P = 32{,}5\).
\(\pi_A = (32{,}5-10) \times 45 = 1.012{,}5\);
\(\pi_B = (32{,}5-10) \times 22{,}5 = 506{,}25\).

\textbf{So sánh với Cournot (cùng công nghệ, đồng thời):} Cân bằng
Cournot cho \(q_A = q_B = 26,\overline{6}\), \(P = 36,\overline{6}\),
\(\pi_A = \pi_B = \approx 711\). Leader Stackelberg có lợi nhuận
\textbf{1.012,5 \textgreater{} 711} (lợi thế đi trước), follower chỉ đạt
\textbf{506,25 \textless{} 711} (bất lợi đi sau). Người tiêu dùng hưởng
lợi: \(P_{Stack} = 32{,}5 < P_{Cournot} = 36,\overline{6}\) --- sản
lượng tổng cao hơn, giá thấp hơn so với Cournot.

\begin{center}\rule{0.5\linewidth}{0.5pt}\end{center}

\subsection{5.4 Cạnh tranh Cournot vs Bertrand --- Nghịch lý giá
(Cournot vs Bertrand --- The
Paradox)}\label{cux1ea1nh-tranh-cournot-vs-bertrand-nghux1ecbch-luxfd-giuxe1-cournot-vs-bertrand-the-paradox}

\subsubsection{5.4.1 Cốt lõi Lý thuyết (Core
Theory)}\label{cux1ed1t-luxf5i-luxfd-thuyux1ebft-core-theory-4}

Đây là \textbf{Nghịch lý Bertrand (Bertrand Paradox)}: trực giác thông
thường nói ``nhiều hãng hơn = cạnh tranh hơn'', nhưng Bertrand chứng
minh chỉ cần \textbf{2 hãng} là đủ tạo ra kết quả như thể có vô số hãng
--- miễn là họ cạnh tranh bằng giá với sản phẩm đồng nhất. Tại sao? Vì
nếu cả hai đang đặt \(P > MC\), mỗi bên có thể cắt giá xuống
\(P - \varepsilon\) để giành toàn bộ thị trường, quá trình này chỉ dừng
khi \(P = MC\).

Các giải pháp thoát khỏi Nghịch lý Bertrand --- tại sao các ngành thực
tế với ít hãng vẫn duy trì lợi nhuận: \textbf{(a) Năng lực sản xuất có
hạn} (Edgeworth, 1897): nếu mỗi hãng không thể phục vụ toàn bộ thị
trường, chiến lược cắt giá mạnh không còn hấp dẫn vì bị giới hạn sản
lượng; \textbf{(b) Dị biệt hóa sản phẩm}: khi sản phẩm không đồng nhất,
mất giá không còn mất toàn bộ thị trường ngay lập tức --- cầu co giãn
theo giá hữu hạn, tạo không gian cho \(P > MC\); \textbf{(c) Chi phí
chuyển đổi} (switching costs): khách hàng hiện tại không phản ứng hoàn
hảo với giá vì có chi phí chuyển sang đối thủ.

\subsubsection{5.4.2 Điểm mù thường gặp (Cognitive Traps \&
Limitations)}\label{ux111iux1ec3m-muxf9-thux1b0ux1eddng-gux1eb7p-cognitive-traps-limitations-7}

\subsubsection{5.4.2 Điểm mù thường gặp (Cognitive Traps \&
Limitations)}\label{ux111iux1ec3m-muxf9-thux1b0ux1eddng-gux1eb7p-cognitive-traps-limitations-8}

Giả định sản phẩm đồng nhất dẫn đến Nghịch lý Bertrand (chỉ cần 2 hãng
là đủ để giá giảm bằng chi phí biên). Trong thực tế, nghịch lý này bị
phá vỡ dễ dàng nếu có ràng buộc về công suất (Nghịch lý Edgeworth) hoặc
sự khác biệt hóa sản phẩm (Hotelling).

\subsubsection{5.4.3 Ứng dụng thực tế (Practical
Applications)}\label{ux1ee9ng-dux1ee5ng-thux1ef1c-tux1ebf-practical-applications-5}

Phân tích sáp nhập (Merger analysis) trong các vụ án chống độc quyền, dự
báo biến động giá khi cấu trúc thị trường thay đổi (vd: Uber và Grab).

\subsubsection{5.4.4 Ví dụ minh họa tính toán (Worked Numerical
Examples)}\label{vuxed-dux1ee5-minh-hux1ecda-tuxednh-touxe1n-worked-numerical-examples-4}

\textbf{Cân bằng Cournot:} Hàm phản ứng tốt nhất mỗi hãng:
\(q_i^* = (80 - q_j)/2\). Giải hệ:
\(q_1^* = q_2^* = 80/3 \approx 26{,}7\); \(Q = 160/3 \approx 53{,}3\);
\(P = 100 - 160/3 \approx 46{,}7\). Lợi nhuận mỗi hãng:
\(\pi = (46{,}7-20) \times 26{,}7 \approx 712\).

\textbf{Cân bằng Bertrand (sản phẩm đồng nhất, không giới hạn năng
lực):} \(P^* = MC = 20\); \(Q = 80\); lợi nhuận mỗi hãng = \textbf{0}.

\textbf{Độc quyền (so sánh):}
\(MR = 100 - 2Q = 20 \Rightarrow Q_{mono} = 40\); \(P_{mono} = 60\);
\(\pi_{mono} = (60-20) \times 40 = 1.600\).

{\def\LTcaptype{none} % do not increment counter
\begin{longtable}[]{@{}llll@{}}
\toprule\noalign{}
Cấu trúc & Giá & Sản lượng tổng & Lợi nhuận/hãng \\
\midrule\noalign{}
\endhead
\bottomrule\noalign{}
\endlastfoot
Độc quyền & 60 & 40 & 1.600 (cho 1 hãng) \\
Cournot (2 hãng) & ≈46,7 & ≈53,3 & ≈712 \\
Bertrand (2 hãng) & 20 (= MC) & 80 & \textbf{0} \\
Cạnh tranh hoàn hảo & 20 & 80 & 0 \\
\end{longtable}
}

Nghịch lý Bertrand rõ ràng: thêm từ 1 lên 2 hãng dưới cạnh tranh
Bertrand tạo ra kết quả y hệt như cạnh tranh hoàn hảo với vô số hãng.
Nhưng thêm từ 1 lên 2 hãng dưới cạnh tranh Cournot chỉ giảm giá từ 60
xuống \textasciitilde46,7, không về đến mức cạnh tranh hoàn hảo.

\begin{center}\rule{0.5\linewidth}{0.5pt}\end{center}

\subsection{5.5 Rào cản gia nhập chiến lược (Strategic entry
deterrence)}\label{ruxe0o-cux1ea3n-gia-nhux1eadp-chiux1ebfn-lux1b0ux1ee3c-strategic-entry-deterrence}

\subsubsection{5.5.1 Cốt lõi Lý thuyết (Core
Theory)}\label{cux1ed1t-luxf5i-luxfd-thuyux1ebft-core-theory-5}

Ba chiến lược ngăn chặn gia nhập chính:

\textbf{(a) Dư thừa năng lực (excess capacity):} Đầu tư xây nhà máy lớn
hơn cần thiết ở trạng thái độc quyền. Điều này làm đáng tin đe dọa ``nếu
anh vào, tôi sẽ mở hết công suất và giá sụt mạnh'' --- vì incumbent đã
chi phí chìm cho năng lực đó, chi phí biên thực của việc mở rộng sản
lượng rất thấp. Dixit (1980) hình thức hóa: incumbent chọn năng lực
\(k\) trước khi entrant quyết định. Nếu \(k\) đủ lớn, lợi nhuận kỳ vọng
của entrant (sau khi tính phản ứng mở rộng sản lượng của incumbent) trở
thành âm.

\textbf{(b) Định giá giới hạn (limit pricing):} Đặt giá \textbf{thấp
hơn} mức độc quyền tối ưu ngắn hạn để hai mục tiêu: tạo tín hiệu về chi
phí thấp (thông tin bất cân xứng) và/hoặc thu hẹp không gian lợi nhuận
cho entrant tiềm năng. Bain (1956) đề xuất nguyên tắc; Milgrom--Roberts
(1982) hình thức hóa trong khuôn khổ lý thuyết tín hiệu --- giá thấp như
tín hiệu về chi phí biên thấp, làm entrant tin rằng thị trường không đủ
hấp dẫn.

\textbf{(c) Chi phí chìm chiến lược (strategic sunk costs):} Đầu tư vào
thương hiệu, R\&D, quan hệ khách hàng, hoặc hợp đồng dài hạn --- tất cả
đều tạo ra \textbf{chi phí chìm cho người muốn gia nhập} để cạnh tranh
hiệu quả. Đây là cơ sở lý thuyết của khái niệm ``rào cản gia nhập tự
nhiên'' theo Bain (1956) được định nghĩa lại bởi Stigler (1968) và hình
thức hóa hiện đại hơn bởi Sutton (1991).

\subsubsection{5.5.2 Điểm mù thường gặp (Cognitive Traps \&
Limitations)}\label{ux111iux1ec3m-muxf9-thux1b0ux1eddng-gux1eb7p-cognitive-traps-limitations-9}

\subsubsection{5.5.2 Điểm mù thường gặp (Cognitive Traps \&
Limitations)}\label{ux111iux1ec3m-muxf9-thux1b0ux1eddng-gux1eb7p-cognitive-traps-limitations-10}

Đầu tư quá mức vào công suất dư thừa (Excess capacity) hoặc quảng cáo
chỉ nhằm mục đích đe dọa đối thủ, dẫn đến phân bổ nguồn lực phi hiệu quả
theo góc nhìn phúc lợi xã hội dù là tối ưu với doanh nghiệp.

\subsubsection{5.5.3 Ứng dụng thực tế (Practical
Applications)}\label{ux1ee9ng-dux1ee5ng-thux1ef1c-tux1ebf-practical-applications-6}

Phân tích cạnh tranh trong các ngành có chi phí chìm lớn như viễn thông,
dược phẩm, hoặc vận tải biển.

\subsubsection{5.5.4 Ví dụ minh họa tính toán (Worked Numerical
Examples)}\label{vuxed-dux1ee5-minh-hux1ecda-tuxednh-touxe1n-worked-numerical-examples-5}

Chiến lược ngăn chặn: I đầu tư trước vào năng lực \(K = 60\) (lớn hơn
mức Cournot), chi phí đầu tư năng lực \(= 300\). Với năng lực đó, nếu E
gia nhập, I mở hết công suất, sản lượng thị trường tăng mạnh, giá sụt về
\(P' \approx 30\). Lợi nhuận E sau gia nhập:
\(\pi_E' \approx (30-20) \times 20 - 400 = 200 - 400 = -200 < 0\) --- E
không gia nhập. I giữ độc quyền với chi phí đầu tư năng lực \(300\)
nhưng duy trì \(\pi_I = 1.600\) \(\Rightarrow\) \textbf{lợi nhuận ròng =
1.300 \textgreater{} 712} (lợi nhuận Cournot nếu không đầu tư). Đầu tư
dư thừa năng lực 300 đơn vị tạo ra \textbf{890 đơn vị lợi nhuận tăng
thêm} --- sinh lời rõ ràng.

\begin{center}\rule{0.5\linewidth}{0.5pt}\end{center}

\subsection{5.6 Phát tín hiệu \& Sàng lọc (Signaling \&
Screening)}\label{phuxe1t-tuxedn-hiux1ec7u-suxe0ng-lux1ecdc-signaling-screening}

\subsubsection{5.6.1 Cốt lõi Lý thuyết (Core
Theory)}\label{cux1ed1t-luxf5i-luxfd-thuyux1ebft-core-theory-6}

Phân biệt \textbf{phát tín hiệu} (signaling --- bên có thông tin chủ
động) với \textbf{sàng lọc} (screening --- bên không có thông tin thiết
kế cơ chế để bên kia tự tiết lộ). Ví dụ: ứng viên tốt chủ động nộp bằng
cấp (phát tín hiệu); nhà tuyển dụng đưa ra hai gói hợp đồng với mức
lương và yêu cầu khác nhau để ứng viên tự phân loại mình (sàng lọc). Về
mặt kết quả phúc lợi: sàng lọc thường kém hiệu quả hơn thông tin đầy đủ
(dẫn đến bóp méo/distortion để phân biệt các loại), nhưng vẫn tốt hơn
thị trường sụp đổ hoàn toàn do adverse selection.

Một cân bằng phân tách không phải luôn là kết quả duy nhất ---
\textbf{cân bằng gộp} (pooling equilibrium) cũng có thể tồn tại: khi chi
phí phát tín hiệu cao đến mức người chất lượng cao cũng không muốn phát,
cả hai loại gộp chung tại cùng một hành động. Lý thuyết tinh chế
(Cho--Kreps Intuitive Criterion, 1987) thường được dùng để loại bỏ các
cân bằng gộp ``không hợp lý'' và giữ lại cân bằng phân tách.

\subsubsection{5.6.2 Điểm mù thường gặp (Cognitive Traps \&
Limitations)}\label{ux111iux1ec3m-muxf9-thux1b0ux1eddng-gux1eb7p-cognitive-traps-limitations-11}

\subsubsection{5.6.2 Điểm mù thường gặp (Cognitive Traps \&
Limitations)}\label{ux111iux1ec3m-muxf9-thux1b0ux1eddng-gux1eb7p-cognitive-traps-limitations-12}

Sự nhầm lẫn về chi phí tín hiệu. Tín hiệu chỉ đáng tin nếu nó tốn kém để
làm giả (Costly to fake). Nếu chi phí phát tín hiệu của người giỏi bằng
người dở, tín hiệu trở nên vô tác dụng và thị trường rơi vào cân bằng
gộp (Pooling equilibrium).

\subsubsection{5.6.3 Ứng dụng thực tế (Practical
Applications)}\label{ux1ee9ng-dux1ee5ng-thux1ef1c-tux1ebf-practical-applications-7}

Bằng cấp đại học như một công cụ phát tín hiệu (Spence 1973), chính sách
bảo hành dài hạn của xe hơi hạng sang, hoặc cổ tức báo hiệu sức khỏe tài
chính.

\subsubsection{5.6.4 Ví dụ minh họa tính toán (Worked Numerical
Examples)}\label{vuxed-dux1ee5-minh-hux1ecda-tuxednh-touxe1n-worked-numerical-examples-6}

\textbf{Tìm cân bằng phân tách:} Giả sử nhà tuyển dụng trả
\(w(e \ge e^*) = 200\) và \(w(e < e^*) = 100\). Để loại L không muốn bắt
chước: \(200 - e^* < 100 \Rightarrow e^* > 100\). Để loại H muốn phát
tín hiệu: \(200 - e^*/2 > 100 \Rightarrow e^* < 200\). Phạm vi \(e^*\)
hợp lệ: \(100 < e^* < 200\). Chọn \(e^* = 120\): - Loại H: Thu nhập ròng
= \(200 - 60 = 140\) (phát tín hiệu) vs \(100 - 0 = 100\) (không phát)
\(\Rightarrow\) \textbf{Phát tín hiệu (140 \textgreater{} 100)}. - Loại
L: Thu nhập ròng = \(200 - 120 = 80\) (bắt chước) vs \(100 - 0 = 100\)
(không bắt chước) \(\Rightarrow\) \textbf{Không bắt chước (100
\textgreater{} 80)}.

Cân bằng phân tách tồn tại. \textbf{Chi phí xã hội}: chi phí giáo dục
của loại H là 60 --- nếu giáo dục không tạo năng suất thêm, đây là
\textbf{60 đơn vị lãng phí xã hội thuần túy} cho mỗi ứng viên loại H,
chỉ để phân biệt với loại L.

\begin{center}\rule{0.5\linewidth}{0.5pt}\end{center}

\subsection{5.7 Bài toán Người ủy quyền--Đại diện (Principal--Agent
Problem)}\label{buxe0i-touxe1n-ngux1b0ux1eddi-ux1ee7y-quyux1ec1nux111ux1ea1i-diux1ec7n-principalagent-problem}

\subsubsection{5.7.1 Cốt lõi Lý thuyết (Core
Theory)}\label{cux1ed1t-luxf5i-luxfd-thuyux1ebft-core-theory-7}

Về mặt hình thức: kết quả \(y = e + \varepsilon\), trong đó \(e\) là nỗ
lực không quan sát được và \(\varepsilon \sim \mathcal{N}(0, \sigma^2)\)
là nhiễu ngẫu nhiên. Principal thiết kế hợp đồng lương
\(w(y) = \alpha + \beta y\) (\(\alpha\) là lương cơ bản, \(\beta\) là hệ
số chia sẻ rủi ro). \textbf{Đánh đổi cốt lõi giữa khuyến khích và chia
sẻ rủi ro:} - \(\beta\) cao \(\rightarrow\) agent có khuyến khích mạnh
hơn (lương phụ thuộc nhiều vào kết quả), nhưng agent gánh nhiều rủi ro
hơn (ngay cả khi nỗ lực cao, kết quả xấu do \(\varepsilon\) âm cũng làm
giảm lương) --- agent không thích rủi ro nên đòi \textbf{phần bù rủi ro}
(risk premium) cao hơn, tốn kém cho principal. - \(\beta\) thấp
\(\rightarrow\) agent chịu ít rủi ro, không cần phần bù rủi ro, nhưng nỗ
lực thấp vì lương gần như cố định.

Điều kiện tối ưu: \(\beta^* = \frac{1}{1 + r\sigma^2 c''}\), trong đó
\(r\) là mức ngại rủi ro của agent, \(\sigma^2\) là phương sai kết quả,
và \(c''\) đo độ tốn kém biên của nỗ lực. Khi rủi ro lớn (\(\sigma^2\)
cao), agent ngại rủi ro nhiều (\(r\) cao), hoặc nỗ lực thêm rất tốn kém
(\(c''\) lớn) --- hệ số khuyến khích tối ưu \(\beta^*\) thấp hơn.

\subsubsection{5.7.2 Điểm mù thường gặp (Cognitive Traps \&
Limitations)}\label{ux111iux1ec3m-muxf9-thux1b0ux1eddng-gux1eb7p-cognitive-traps-limitations-13}

\subsubsection{5.7.2 Điểm mù thường gặp (Cognitive Traps \&
Limitations)}\label{ux111iux1ec3m-muxf9-thux1b0ux1eddng-gux1eb7p-cognitive-traps-limitations-14}

Đánh đổi giữa rủi ro và khích lệ (Risk-Incentive tradeoff). Trả lương
100\% theo hiệu quả sẽ tạo động lực tối đa, nhưng lại chuyển toàn bộ rủi
ro biến động sang người lao động (người vốn e ngại rủi ro hơn giới chủ),
dẫn đến mất hiệu quả Pareto.

\subsubsection{5.7.3 Ứng dụng thực tế (Practical
Applications)}\label{ux1ee9ng-dux1ee5ng-thux1ef1c-tux1ebf-practical-applications-8}

Thiết kế gói lương thưởng cho CEO (Stock options), hợp đồng chia sẻ rủi
ro nông nghiệp (Sharecropping), và cơ chế trả tiền theo hiệu quả trong y
tế.

\subsubsection{5.7.4 Ví dụ minh họa tính toán (Worked Numerical
Examples)}\label{vuxed-dux1ee5-minh-hux1ecda-tuxednh-touxe1n-worked-numerical-examples-7}

\textbf{Bước 1 --- Nỗ lực tối ưu của agent:} Agent tối đa hóa
\(\mathbb{E}[w] - c(e) = \alpha + \beta e - e^2/2\). Điều kiện bậc nhất:
\(\beta - e = 0 \Rightarrow e^*(\beta) = \beta\) --- nỗ lực agent chọn
bằng đúng hệ số khuyến khích.

\textbf{Bước 2 --- Chi phí rủi ro agent gánh:} Agent nhận lương kỳ vọng
\(\mu_w = \alpha + \beta^2\) nhưng chịu rủi ro lương. Chi phí rủi ro
\(= \frac{1}{2} r \beta^2 \sigma^2 = \frac{1}{2}(0{,}5)\beta^2(4) = \beta^2\).
Để agent chấp nhận hợp đồng:
\(\alpha + \beta^2 - \beta^2/2 - \beta^2 \ge \bar{u} = 0\) (mức thỏa
dụng để lại) \(\Rightarrow \alpha = \beta^2/2\).

\textbf{Bước 3 --- Principal tối ưu \(\beta\):} Lợi nhuận principal
\(= \mathbb{E}[y] - \mathbb{E}[w] = \beta - (\alpha + \beta^2) = \beta - \beta^2/2 - \beta^2 = \beta - 3\beta^2/2\).
Tối đa hóa: \(1 - 3\beta = 0 \Rightarrow \beta^* = 1/3\).

\textbf{Kết quả:} Hệ số khuyến khích tối ưu \(\beta^* = 1/3\); nỗ lực
agent \(e^* = 1/3\); lương cơ bản \(\alpha = 1/18\); lương kỳ vọng
\(\mathbb{E}[w] = 1/18 + (1/3)^2 = 1/18 + 1/9 = 1/6\). So sánh với hợp
đồng lý tưởng (thông tin đầy đủ, không nhiễu): \(e_{FB} = 1\) --- nỗ lực
thực tế chỉ bằng 1/3 nỗ lực tối ưu xã hội vì \(\sigma^2\) cao và agent
ngại rủi ro, buộc principal phải nhượng bộ về mức khuyến khích để tiết
kiệm phần bù rủi ro. \textbf{Chi phí thông tin bất cân xứng} (deadweight
loss do moral hazard) = lợi nhuận tối ưu xã hội trừ lợi nhuận đạt được
với \(\beta^*\): một phần giá trị bị mất chuyển thành phần bù rủi ro và
nỗ lực dưới mức.

\begin{center}\rule{0.5\linewidth}{0.5pt}\end{center}

\section{CHƯƠNG 6: KINH TẾ HỌC HÀNH VI CƠ BẢN (Basic Behavioral
Economics)}\label{chux1b0ux1a1ng-6-kinh-tux1ebf-hux1ecdc-huxe0nh-vi-cux1a1-bux1ea3n-basic-behavioral-economics}

Con người thực tế không luôn ra quyết định ``duy lý'' như mô hình kinh
tế cổ điển giả định. Nhưng kinh tế học hành vi không chỉ đơn giản là
``liệt kê những sai lầm của con người'' --- mục tiêu sâu hơn là xây dựng
mô hình \emph{dự đoán được} và \emph{có cơ sở tâm lý học} về cách con
người thực tế hành xử, để từ đó thiết kế chính sách và cơ chế tốt hơn.
Nhóm này hệ thống các sai lệch hành vi có hệ thống và có thể dự đoán ---
không phải ngẫu nhiên, mà là những thiên kiến nhất quán xuất hiện ở phần
lớn người trong những tình huống nhất định. (Xem thêm Phần B, Chương 16
để có bản hình thức hóa đầy đủ của lý thuyết triển vọng, mô hình β-δ và
kinh tế học thần kinh.)

\begin{center}\rule{0.5\linewidth}{0.5pt}\end{center}

\subsection{6.1 Lý thuyết triển vọng \& Ác cảm mất mát (Prospect Theory
\& Loss
Aversion)}\label{luxfd-thuyux1ebft-triux1ec3n-vux1ecdng-uxe1c-cux1ea3m-mux1ea5t-muxe1t-prospect-theory-loss-aversion}

\subsubsection{6.1.1 Cốt lõi Lý thuyết (Core
Theory)}\label{cux1ed1t-luxf5i-luxfd-thuyux1ebft-core-theory-8}

\textbf{(1) Hàm giá trị} \(v(\cdot)\): thay vì đánh giá tài sản theo mức
tuyệt đối, con người đánh giá theo \textbf{thay đổi so với điểm tham
chiếu} \(r\) (reference point --- thường là hiện trạng hoặc kỳ vọng).
Hàm giá trị có ba đặc điểm: - Lõm ở vùng lãi (\(x > 0\), tương tự ngại
rủi ro khi được lợi): \(v''(x) < 0\) với \(x > 0\) - Lồi ở vùng lỗ
(\(x < 0\), tương tự ưa thích rủi ro khi thua lỗ): \(v''(x) > 0\) với
\(x < 0\) - Dốc hơn ở vùng lỗ: \(|v'(-x)| > |v'(x)|\) --- \textbf{ác cảm
mất mát} (loss aversion), với hệ số \(\lambda \approx 2{,}25\) (Tversky
\& Kahneman 1992)

\textbf{(2) Hàm trọng số xác suất} \(\pi(p)\): thay vì dùng xác suất
khách quan, con người dùng trọng số quyết định biến dạng xác suất ---
\textbf{phóng đại} các xác suất nhỏ (\(\pi(p) > p\) với \(p\) nhỏ, giải
thích tại sao người ta mua vé số và bảo hiểm thảm họa) và \textbf{thu
nhỏ} các xác suất trung bình đến lớn (\(\pi(p) < p\)).

Hàm giá trị trong lý thuyết triển vọng thường được đặc tả dạng lũy thừa:
\[ v(x) = \begin{cases} x^\alpha & \text{nếu } x \ge 0 \\ -\lambda(-x)^\alpha & \text{nếu } x < 0 \end{cases} \]
với \(\alpha \approx 0{,}88\) và \(\lambda \approx 2{,}25\). Giá trị kỳ
vọng dưới lý thuyết triển vọng: \(V = \sum_i \pi(p_i) v(x_i - r)\).

\textbf{Mẫu bộ tứ thái độ rủi ro} (fourfold pattern): lý thuyết triển
vọng dự đoán bốn thái độ rủi ro khác nhau tùy vùng lãi/lỗ và xác suất
cao/thấp --- ví dụ trong vùng lãi + xác suất cao: ngại rủi ro (chốt lời
sớm); trong vùng lỗ + xác suất cao: ưa thích rủi ro (giữ tài sản thua lỗ
chờ phục hồi). Cả bốn ô này đều quan sát được nhất quán trong thực
nghiệm.

\subsubsection{6.1.2 Điểm mù thường gặp (Cognitive Traps \&
Limitations)}\label{ux111iux1ec3m-muxf9-thux1b0ux1eddng-gux1eb7p-cognitive-traps-limitations-15}

\subsubsection{6.1.2 Điểm mù thường gặp (Cognitive Traps \&
Limitations)}\label{ux111iux1ec3m-muxf9-thux1b0ux1eddng-gux1eb7p-cognitive-traps-limitations-16}

Không phải lúc nào con người cũng e ngại rủi ro. Hàm giá trị phi tuyến
của Lý thuyết Triển vọng chỉ ra rằng con người tìm kiếm rủi ro
(risk-seeking) trong miền thua lỗ để ``gỡ gạc'', và e ngại rủi ro trong
miền lợi nhuận.

\subsubsection{6.1.3 Ứng dụng thực tế (Practical
Applications)}\label{ux1ee9ng-dux1ee5ng-thux1ef1c-tux1ebf-practical-applications-9}

Hiệu ứng ngược (Disposition effect) trên thị trường chứng khoán: nhà đầu
tư có xu hướng chốt lời quá sớm (khi có lãi) và gồng lỗ quá lâu (khi
đang lỗ).

\subsubsection{6.1.4 Ví dụ minh họa tính toán (Worked Numerical
Examples)}\label{vuxed-dux1ee5-minh-hux1ecda-tuxednh-touxe1n-worked-numerical-examples-8}

\textbf{Tình huống A (Vùng lãi):} Phương án 1: nhận chắc 500 (lợi
nhuận). Phương án 2: 50\% được 1.000, 50\% được 0.

\(V_1 = v(500) = 500^{0,88} \approx 261{,}6\).
\(V_2 = \pi(0{,}5) \cdot v(1.000) + \pi(0{,}5) \cdot v(0)\). Dùng
\(\pi(0{,}5) \approx 0{,}42\) (trọng số xác suất):
\(V_2 \approx 0{,}42 \times 1000^{0,88} + 0{,}42 \times 0 = 0{,}42 \times 527{,}5 \approx 221{,}6\).
Vì \(261{,}6 > 221{,}6\), chọn \textbf{chắc 500} --- \textbf{ngại rủi ro
trong vùng lãi} (dù kỳ vọng toán học bằng nhau = 500).

\textbf{Tình huống B (Vùng lỗ):} Phương án 1: mất chắc −500. Phương án
2: 50\% mất −1.000, 50\% mất 0.

\(V_1 = -\lambda \cdot 500^{0,88} = -2{,}25 \times 261{,}6 = -588{,}6\).
\(V_2 \approx 0{,}42 \times (-2{,}25 \times 527{,}5) + 0{,}42 \times 0 = 0{,}42 \times (-1.186{,}9) \approx -498{,}5\).
Vì \(-498{,}5 > -588{,}6\), chọn \textbf{đánh cược} --- \textbf{ưa thích
rủi ro trong vùng lỗ} (dù kỳ vọng toán học bằng nhau = −500). Đây chính
là ``Mẫu bộ tứ'': cùng một cấu trúc xác suất nhưng đổi dấu (lãi ↔ lỗ)
tạo ra thái độ rủi ro hoàn toàn ngược chiều.

\begin{center}\rule{0.5\linewidth}{0.5pt}\end{center}

\subsection{6.2 Chiết khấu hyperbolic \& Tự kiểm soát (Hyperbolic
discounting \&
Self-control)}\label{chiux1ebft-khux1ea5u-hyperbolic-tux1ef1-kiux1ec3m-souxe1t-hyperbolic-discounting-self-control}

\subsubsection{6.2.1 Cốt lõi Lý thuyết (Core
Theory)}\label{cux1ed1t-luxf5i-luxfd-thuyux1ebft-core-theory-9}

Thực nghiệm (Thaler 1981; Loewenstein \& Prelec 1992) cho thấy con người
thực tế dùng tỷ lệ chiết khấu \textbf{rất cao cho tương lai gần} và thấp
hơn nhiều cho tương lai xa --- hành vi này được mô hình hóa tốt hơn bởi
\textbf{chiết khấu hyperbolic}: \(D(t) = 1/(1 + \alpha t)\), tỷ lệ chiết
khấu tức thời giảm dần theo thời gian. Hệ quả: \textbf{thiên kiến hiện
tại} (present bias) và \textbf{hành vi không nhất quán theo thời gian}
(time-inconsistent behavior) --- kế hoạch ngày hôm nay cho ngày mai
không được thực hiện khi ngày mai đến.

Laibson (1997) đề xuất mô hình β-δ đơn giản hóa chiết khấu
quasi-hyperbolic:
\[ U_t = u(c_t) + \beta \sum_{\tau=1}^{T} \delta^\tau u(c_{t+\tau}) \]
trong đó \(0 < \beta < 1\) là \textbf{hệ số thiên kiến hiện tại}
(present bias factor) và \(\delta\) là hệ số chiết khấu dài hạn bình
thường. Khi \(\beta = 1\): chiết khấu hàm mũ chuẩn. Khi \(\beta < 1\):
hiện tại được đánh giá cao hơn tất cả tương lai bằng nhân tử
\(1/\beta\), nhưng bên trong tương lai thì vẫn chiết khấu theo
\(\delta\) bình thường. Phân biệt hai loại tác nhân:
\textbf{``Sophisticates''} --- biết mình bị thiên kiến hiện tại trong
tương lai, nên chủ động ràng buộc trước (commitment devices);
\textbf{``Naifs''} --- mỗi kỳ đều nghĩ ``lần sau mình sẽ kỷ luật hơn''
nhưng không bao giờ thực hiện.

\subsubsection{6.2.2 Điểm mù thường gặp (Cognitive Traps \&
Limitations)}\label{ux111iux1ec3m-muxf9-thux1b0ux1eddng-gux1eb7p-cognitive-traps-limitations-17}

\subsubsection{6.2.2 Điểm mù thường gặp (Cognitive Traps \&
Limitations)}\label{ux111iux1ec3m-muxf9-thux1b0ux1eddng-gux1eb7p-cognitive-traps-limitations-18}

Sai lầm trong dự báo tương lai của chính bản thân (Naivete vs
Sophistication). Một người có chiết khấu hyperbolic ngây thơ sẽ liên tục
trì hoãn việc tập thể dục vì tin rằng ``ngày mai mình sẽ làm'', không
nhận ra ngày mai họ vẫn sẽ chọn trì hoãn.

\subsubsection{6.2.3 Ứng dụng thực tế (Practical
Applications)}\label{ux1ee9ng-dux1ee5ng-thux1ef1c-tux1ebf-practical-applications-10}

Thiết kế các cơ chế cam kết (Commitment devices) như tài khoản tiết kiệm
hưu trí khóa chặt, hoặc các ứng dụng chặn mạng xã hội.

\subsubsection{6.2.4 Ví dụ minh họa tính toán (Worked Numerical
Examples)}\label{vuxed-dux1ee5-minh-hux1ecda-tuxednh-touxe1n-worked-numerical-examples-9}

\textbf{Đánh giá hôm nay (t=0):} Tập hôm nay:
\(V_0^{today} = -5 + \beta\delta \cdot 8 = -5 + 0{,}7 \times 0{,}95 \times 8 = -5 + 5{,}32 = +0{,}32 > 0\)
\(\Rightarrow\) Đáng tập. Trì hoãn (tập ngày mai nhìn từ hôm nay):
\(V_0^{tomorrow} = 0 + \beta\delta^2(-5) + \beta\delta^2 \cdot \delta \cdot 8 = 0 + 0{,}7\times0{,}9025\times(-5) + 0{,}7\times0{,}9025\times0{,}95\times8\).
So sánh cho thấy tập hôm nay có lợi hơn.

\textbf{Đánh giá ngày mai (t=1, khi ``ngày mai'' đến):} ``Tập hôm nay''
(vẫn là ngày mai so với hôm qua):
\(V_1 = -5 + \beta\delta \cdot 8 = -5 + 5{,}32 = +0{,}32\). Nhưng bây
giờ ``trì hoãn sang ngày kia'' trở thành:
\(V_1^{delay} = 0 + \beta\delta^2 \cdot (-5) + \beta\delta^3 \cdot 8 \approx 0 + 0{,}7\times0{,}9025\times(-5) + 0{,}7\times0{,}857\times8 = -3{,}16 + 4{,}8 = +1{,}64\).
\textbf{Trì hoãn sang ngày kia nhìn từ ngày mai có giá trị +1,64
\textgreater{} +0,32 (tập hôm nay)} \(\Rightarrow\) người này lại chọn
trì hoãn --- và vòng lặp này lặp lại mỗi ngày. Đây chính là cơ chế trì
hoãn mãn tính (chronic procrastination) được mô hình β-δ dự đoán:
\textbf{mỗi ngày nhìn tới ngày mai đều có vẻ hợp lý để trì hoãn}.

\begin{center}\rule{0.5\linewidth}{0.5pt}\end{center}

\subsection{6.3 Sở thích xã hội: Công bằng, Tương hỗ \& Vị tha (Social
Preferences)}\label{sux1edf-thuxedch-xuxe3-hux1ed9i-cuxf4ng-bux1eb1ng-tux1b0ux1a1ng-hux1ed7-vux1ecb-tha-social-preferences}

\subsubsection{6.3.1 Cốt lõi Lý thuyết (Core
Theory)}\label{cux1ed1t-luxf5i-luxfd-thuyux1ebft-core-theory-10}

\textbf{Ultimatum Game (Güth et al.~1982):} Người chơi A chia 100 đơn vị
với B; nếu B từ chối, cả hai nhận 0. Dự đoán homo economicus: A đề xuất
giữ 99, cho B 1; B chấp nhận vì 1 \textgreater{} 0. Thực nghiệm: các đề
xuất dưới 30\% thường bị từ chối --- B sẵn sàng nhận 0 để trừng phạt đề
xuất ``bất công''. Kết quả trên nhất quán qua hàng trăm thực nghiệm,
hàng chục quốc gia, mọi trình độ học vấn.

\textbf{Mô hình Fehr--Schmidt (1999)} hình thức hóa sở thích xã hội qua
\textbf{ác cảm bất bình đẳng} (inequality aversion): mỗi cá nhân \(i\)
không chỉ quan tâm đến lợi ích cá nhân \(x_i\) mà còn quan tâm đến bất
bình đẳng giữa mình và người khác:
\[ U_i = x_i - \alpha_i \max(x_j - x_i, 0) - \beta_i \max(x_i - x_j, 0) \]
trong đó \(\alpha_i\) đo \textbf{ác cảm với bất lợi} (envy --- không
thích khi người khác được nhiều hơn mình) và \(\beta_i\) đo \textbf{ác
cảm với lợi thế} (guilt/altruism --- không thích khi mình được nhiều hơn
người khác). Kết quả thực nghiệm: \(\alpha_i > \beta_i\) với phần lớn
mọi người (đố kỵ mạnh hơn tội lỗi), và \(\beta_i < 1\) (người ta vẫn
thích được nhiều hơn nhưng chịu một ``chi phí tâm lý'' khi bất bình đẳng
quá lớn).

\subsubsection{6.3.2 Điểm mù thường gặp (Cognitive Traps \&
Limitations)}\label{ux111iux1ec3m-muxf9-thux1b0ux1eddng-gux1eb7p-cognitive-traps-limitations-19}

\subsubsection{6.3.2 Điểm mù thường gặp (Cognitive Traps \&
Limitations)}\label{ux111iux1ec3m-muxf9-thux1b0ux1eddng-gux1eb7p-cognitive-traps-limitations-20}

Động lực tài chính có thể lấn át (Crowding out) động lực đạo đức/xã hội.
Việc áp dụng tiền phạt hoặc tiền thưởng đôi khi làm hỏng các chuẩn mực
xã hội sẵn có (như nghiên cứu của Gneezy \& Rustichini 2000 về việc phạt
đón muộn ở nhà trẻ).

\subsubsection{6.3.3 Ứng dụng thực tế (Practical
Applications)}\label{ux1ee9ng-dux1ee5ng-thux1ef1c-tux1ebf-practical-applications-11}

Thiết kế chính sách hiến máu nhân đạo (không trả tiền mặt để tránh thu
hút máu bẩn), hoặc các chiến dịch nâng cao ý thức môi trường.

\subsubsection{6.3.4 Ví dụ minh họa tính toán (Worked Numerical
Examples)}\label{vuxed-dux1ee5-minh-hux1ecda-tuxednh-touxe1n-worked-numerical-examples-10}

B đánh giá đề xuất theo Fehr--Schmidt: \(x_B = 30\), \(x_A = 70\). Vì
\(x_A > x_B\):
\(U_B = 30 - 1{,}5 \times (70 - 30) - 0 = 30 - 60 = -30 < 0\). \textbf{B
từ chối} (dù từ chối mất 30 đơn vị tiền tệ).

A muốn tối đa hóa \(x_A\) nhưng biết B sẽ từ chối nếu đề xuất quá bất
công. A cần tìm mức cho B tối thiểu \(x_B^*\) sao cho B chấp nhận:
\(x_B - \alpha_B(x_A - x_B) \ge 0 \Rightarrow x_B - 1{,}5(100 - 2x_B) \ge 0 \Rightarrow x_B - 150 + 3x_B \ge 0 \Rightarrow 4x_B \ge 150 \Rightarrow x_B^* \ge 37{,}5\).
A tối ưu: đề xuất \textbf{37,5 cho B, giữ 62,5} --- cao hơn nhiều so với
1 đơn vị mà mô hình homo economicus dự đoán, hoàn toàn phù hợp với kết
quả thực nghiệm thường quan sát được trong phạm vi 35--45\%.

\begin{center}\rule{0.5\linewidth}{0.5pt}\end{center}

\subsection{6.4 Thiên kiến neo (Anchoring
effect)}\label{thiuxean-kiux1ebfn-neo-anchoring-effect}

\subsubsection{6.4.1 Cốt lõi Lý thuyết (Core
Theory)}\label{cux1ed1t-luxf5i-luxfd-thuyux1ebft-core-theory-11}

Cơ chế tâm lý: sau khi tiếp xúc mỏ neo, quá trình ước lượng diễn ra bằng
cách \textbf{điều chỉnh từ mỏ neo ra} (anchoring-and-adjustment) --- và
con người thường điều chỉnh \textbf{không đủ xa}, nên kết quả cuối cùng
vẫn bị kéo về phía mỏ neo. Đây không phải vì thiếu thông tin mà là vì
quá trình nhận thức sử dụng mỏ neo như điểm khởi đầu của quá trình tìm
kiếm câu trả lời.

Hiệu ứng neo xuất hiện trong ba loại tình huống chính: \textbf{(a)} Đàm
phán: giá đề xuất đầu tiên neo toàn bộ kết quả thương lượng;
\textbf{(b)} Định giá tài sản: giá mua ban đầu neo kỳ vọng khi bán lại;
\textbf{(c)} Dự báo và kế hoạch: con số kế hoạch năm trước neo kế hoạch
năm nay (kể cả khi điều kiện thị trường thay đổi hoàn toàn).

\subsubsection{6.4.2 Điểm mù thường gặp (Cognitive Traps \&
Limitations)}\label{ux111iux1ec3m-muxf9-thux1b0ux1eddng-gux1eb7p-cognitive-traps-limitations-21}

\subsubsection{6.4.2 Điểm mù thường gặp (Cognitive Traps \&
Limitations)}\label{ux111iux1ec3m-muxf9-thux1b0ux1eddng-gux1eb7p-cognitive-traps-limitations-22}

Mỏ neo hoàn toàn vô lý (như số an sinh xã hội ngẫu nhiên) vẫn ảnh hưởng
mạnh mẽ đến định giá (Ariely). Người ta điều chỉnh từ mỏ neo nhưng
thường điều chỉnh không đủ (insufficient adjustment).

\subsubsection{6.4.3 Ứng dụng thực tế (Practical
Applications)}\label{ux1ee9ng-dux1ee5ng-thux1ef1c-tux1ebf-practical-applications-12}

Chiến lược định giá bán lẻ (để giá gốc cực cao rồi gạch đi giảm giá),
thương lượng lương khởi điểm, hoặc yêu cầu bồi thường tại tòa án.

\subsubsection{6.4.4 Ví dụ minh họa tính toán (Worked Numerical
Examples)}\label{vuxed-dux1ee5-minh-hux1ecda-tuxednh-touxe1n-worked-numerical-examples-11}

\textbf{Kịch bản 1 (Không có neo):} Hai bên thảo luận tự do, không bên
nào đề xuất trước. Sau đàm phán, kết quả thường hội tụ về mức giữa:
\textbf{khoảng 27--28 triệu}.

\textbf{Kịch bản 2 (Neo từ ứng viên):} Ứng viên đề xuất trước: ``Tôi kỳ
vọng mức lương 35 triệu.'' Nhà tuyển dụng biết đây là cao hơn thị trường
nhưng không phủ nhận hoàn toàn --- bắt đầu đàm phán từ 35 triệu xuống.
Kết quả thực nghiệm: thường hội tụ cao hơn kịch bản 1 khoảng
\textbf{2--5 triệu}, tức \textbf{29--30 triệu}. Thử nghiệm ngược: nhà
tuyển dụng đề xuất trước ``ngân sách chúng tôi là 22 triệu'' --- kết quả
thường hội tụ \textbf{22--24 triệu}, thấp hơn kịch bản 1. \textbf{Ai đề
xuất trước và đề xuất bao nhiêu quyết định phần lớn kết quả đàm phán}
--- không phải ``giá trị thực'' của vị trí.

\begin{center}\rule{0.5\linewidth}{0.5pt}\end{center}

\subsection{6.5 Hiệu ứng sở hữu \& Sự mất mát của nguyên trạng
(Endowment effect \& Status quo
bias)}\label{hiux1ec7u-ux1ee9ng-sux1edf-hux1eefu-sux1ef1-mux1ea5t-muxe1t-cux1ee7a-nguyuxean-trux1ea1ng-endowment-effect-status-quo-bias}

\subsubsection{6.5.1 Cốt lõi Lý thuyết (Core
Theory)}\label{cux1ed1t-luxf5i-luxfd-thuyux1ebft-core-theory-12}

\textbf{Thiên kiến nguyên trạng} (status quo bias --- Samuelson \&
Zeckhauser 1988) là dạng rộng hơn: xu hướng giữ nguyên hiện trạng và
tránh thay đổi, ngay cả khi thay đổi có lợi ích ròng dương. Cơ chế kép:
ác cảm mất mát (từ bỏ lợi ích của hiện trạng cảm giác như mất) +
heuristics ``tránh hối tiếc'' (nếu thay đổi tệ đi, tự trách mình; nếu
không thay đổi mà tệ đi, đổ thừa ngoại cảnh). Kết quả: ngưỡng để thay
đổi hành vi cao hơn nhiều so với lý thuyết tối ưu dự đoán.

Quan hệ với thiết kế cơ chế: hiệu ứng sở hữu và thiên kiến nguyên trạng
tạo ra \textbf{độ dính} (stickiness) trong hành vi --- lựa chọn mặc định
(default option) có ảnh hưởng phi lý rất lớn vì hầu hết người không chủ
động thoát khỏi nó. Điều này được khai thác có chủ ý trong thiết kế
``Kiến trúc lựa chọn'' (xem mục tiếp theo).

\subsubsection{6.5.2 Điểm mù thường gặp (Cognitive Traps \&
Limitations)}\label{ux111iux1ec3m-muxf9-thux1b0ux1eddng-gux1eb7p-cognitive-traps-limitations-23}

\subsubsection{6.5.2 Điểm mù thường gặp (Cognitive Traps \&
Limitations)}\label{ux111iux1ec3m-muxf9-thux1b0ux1eddng-gux1eb7p-cognitive-traps-limitations-24}

Vi phạm định lý Coase. Hiệu ứng sở hữu làm cho Giá sẵn sàng chấp nhận
(WTA) lớn hơn nhiều so với Giá sẵn sàng trả (WTP). Ngay cả khi chi phí
giao dịch bằng không, thương lượng vẫn có thể bế tắc do các bên định giá
quá cao tài sản họ đang cầm.

\subsubsection{6.5.3 Ứng dụng thực tế (Practical
Applications)}\label{ux1ee9ng-dux1ee5ng-thux1ef1c-tux1ebf-practical-applications-13}

Chính sách dùng thử miễn phí (Free trial) hoặc hoàn tiền 100\%
(Money-back guarantee) của các nhà bán lẻ để tạo cảm giác sở hữu.

\subsubsection{6.5.4 Ví dụ minh họa tính toán (Worked Numerical
Examples)}\label{vuxed-dux1ee5-minh-hux1ecda-tuxednh-touxe1n-worked-numerical-examples-12}

\textbf{Tính toán lý thuyết:} NPV với lãi suất chiết khấu 10\%/năm, phần
mềm dùng thêm 7 năm:
\(NPV = -500 + \sum_{t=1}^{7} \frac{200}{1{,}1^t} = -500 + 200 \times \frac{1-(1/1{,}1)^7}{0{,}1} = -500 + 200 \times 4{,}868 = -500 + 973{,}6 = +473{,}6\)
triệu \(\Rightarrow\) Đáng chuyển đổi rõ ràng.

\textbf{Nhưng thực tế:} Nhóm ra quyết định cảm nhận 500 triệu đầu tư như
``mất'' (chi phí chìm tinh thần dù chưa phải chi phí chìm thực), 200
triệu tiết kiệm mỗi năm như ``được'' --- do ác cảm mất mát, ``mất ngay''
được đánh giá nặng hơn ``được dần'' theo hệ số
\(\lambda \approx 2{,}25\). Họ cảm nhận NPV tâm lý:
\(\approx -500 \times 2{,}25 + 973{,}6 = -1.125 + 973{,}6 = -151 < 0\)
--- \textbf{cảm giác ``không đáng''} dù phân tích tài chính khách quan
rõ ràng đáng đầu tư. Kết quả: quyết định bị trì hoãn hoặc từ chối vì
thiên kiến nguyên trạng, tổ chức tiếp tục dùng phần mềm kém hiệu quả
thêm nhiều năm.

\begin{center}\rule{0.5\linewidth}{0.5pt}\end{center}

\subsection{6.6 Cú hích \& Kiến trúc lựa chọn (Nudge \& Choice
Architecture)}\label{cuxfa-huxedch-kiux1ebfn-truxfac-lux1ef1a-chux1ecdn-nudge-choice-architecture}

\subsubsection{6.6.1 Cốt lõi Lý thuyết (Core
Theory)}\label{cux1ed1t-luxf5i-luxfd-thuyux1ebft-core-theory-13}

Các công cụ kiến trúc lựa chọn chính: \textbf{(a) Lựa chọn mặc định}
(default rules): khai thác thiên kiến nguyên trạng --- hầu hết người giữ
nguyên lựa chọn mặc định dù có thể thay đổi dễ dàng. Cùng một chính
sách, opt-out (mặc định là có tham gia, muốn không tham gia phải chủ
động hủy) tạo tỷ lệ tham gia cao hơn nhiều so với opt-in (mặc định không
tham gia). \textbf{(b) Đơn giản hóa} (simplification): giảm chi phí nhận
thức của quyết định đúng đắn bằng cách hiển thị thông tin quan trọng rõ
ràng hơn. \textbf{(c) Phản hồi xã hội} (social norms): ``90\% hàng xóm
của bạn đã nộp thuế đúng hạn'' hiệu quả hơn nhắc nhở pháp lý về nghĩa vụ
nộp thuế. \textbf{(d) Sắp xếp không gian} (salience): đặt thực phẩm lành
mạnh ở tầm mắt trong căng-tin trường học làm tăng lựa chọn lành mạnh mà
không cấm thực phẩm không lành mạnh.

Thaler \& Sunstein gọi triết lý này là \textbf{``Chủ nghĩa gia trưởng tự
do''} (libertarian paternalism): can thiệp vào kiến trúc lựa chọn để dẫn
dắt hành vi theo hướng tốt hơn cho người chọn, trong khi vẫn hoàn toàn
tôn trọng quyền tự do lựa chọn --- khác với chủ nghĩa gia trưởng truyền
thống (cấm đoán) và thị trường tự do thuần túy (không can thiệp dù thị
trường dẫn đến kết quả tệ cho cá nhân).

\subsubsection{6.6.2 Điểm mù thường gặp (Cognitive Traps \&
Limitations)}\label{ux111iux1ec3m-muxf9-thux1b0ux1eddng-gux1eb7p-cognitive-traps-limitations-25}

\subsubsection{6.6.2 Điểm mù thường gặp (Cognitive Traps \&
Limitations)}\label{ux111iux1ec3m-muxf9-thux1b0ux1eddng-gux1eb7p-cognitive-traps-limitations-26}

Ranh giới đạo đức của Chủ nghĩa Gia trưởng Tự do (Libertarian
Paternalism). Đôi khi Nudge có thể bị lạm dụng bởi doanh nghiệp thành
Sludge (gây rắc rối, làm khó người dùng khi muốn hủy đăng ký dịch vụ).

\subsubsection{6.6.3 Ứng dụng thực tế (Practical
Applications)}\label{ux1ee9ng-dux1ee5ng-thux1ef1c-tux1ebf-practical-applications-14}

Thiết kế mặc định (Opt-out thay vì Opt-in) cho chương trình hiến tạng
hoặc tiết kiệm hưu trí tự động (Save More Tomorrow).

\subsubsection{6.6.4 Ví dụ minh họa tính toán (Worked Numerical
Examples)}\label{vuxed-dux1ee5-minh-hux1ecda-tuxednh-touxe1n-worked-numerical-examples-13}

Dữ liệu thực nghiệm (Madrian \& Shea 2001): Công ty A (opt-in): sau 3
tháng, 37\% nhân viên đã đăng ký; sau 36 tháng, 65\%. Công ty B
(opt-out): sau 3 tháng, 86\% nhân viên tham gia; sau 36 tháng, 98\%.
\textbf{Tỷ lệ tham gia dài hạn: 65\% vs 98\% --- chênh lệch 33 điểm phần
trăm.} Chi phí thay đổi: gần như bằng 0 (chỉ thay đổi từ ``cần điền form
đăng ký'' thành ``cần điền form hủy''). Không có sự cưỡng bức, không có
thay đổi khuyến khích tài chính --- toàn bộ sự khác biệt đến từ
\textbf{kiến trúc lựa chọn đơn thuần}, khai thác thiên kiến nguyên trạng
và thiên kiến hiện tại.

\begin{center}\rule{0.5\linewidth}{0.5pt}\end{center}

\subsection{6.7 Tài chính hành vi (Behavioral
Finance)}\label{tuxe0i-chuxednh-huxe0nh-vi-behavioral-finance}

\subsubsection{6.7.1 Cốt lõi Lý thuyết (Core
Theory)}\label{cux1ed1t-luxf5i-luxfd-thuyux1ebft-core-theory-14}

Ba dòng bằng chứng nền tảng:

\textbf{(a) Shiller (1981) --- Biến động giá quá mức:} Thị trường cổ
phiếu Mỹ biến động mạnh hơn \emph{nhiều} so với những gì mô hình giá trị
hiện tại thuần túy dự đoán từ biến động của cổ tức thực. Bằng chứng:
phương sai của giá cổ phiếu lớn hơn cận trên lý thuyết cho phép bởi hàng
chục lần. Hàm ý: giá tài sản không chỉ phản ánh giá trị cơ bản
(fundamental value) mà còn bị ảnh hưởng bởi tâm lý đám đông, kỳ vọng
ngoại suy, và bong bóng đầu cơ.

\textbf{(b) Anomalies --- Các hiệu ứng bất thường có hệ thống:} Hiệu ứng
Tháng Một (January Effect --- giá tăng bất thường); Hiệu ứng Thứ Hai
(Monday Effect --- giá giảm bất thường); Hiệu ứng Quy mô (Size Effect
--- công ty nhỏ có lợi suất cao hơn sau điều chỉnh rủi ro --- Fama \&
French 1992); Hiệu ứng Đà (Momentum --- cổ phiếu đang tăng tiếp tục tăng
trong 3--12 tháng --- Jegadeesh \& Titman 1993). Nếu thị trường hoàn
toàn hiệu quả, các mẫu này không nên tồn tại dưới dạng có thể khai thác
bền vững.

\textbf{(c) Thaler --- Kế toán tâm lý và Disposition Effect:} Con người
không đánh giá tổng tài sản theo giá trị tổng hợp mà phân loại tiền vào
các ``tài khoản tinh thần'' (mental accounts) riêng biệt --- tiền
thưởng, tiền lương, tiền thừa kế được đối xử khác nhau dù về mặt kinh tế
là như nhau. \textbf{Disposition Effect} (Shefrin \& Statman 1985): nhà
đầu tư bán cổ phiếu lãi quá sớm và giữ cổ phiếu lỗ quá lâu --- hành vi
ngược hoàn toàn với tối ưu thuế (nên cắt lỗ sớm để được hoàn thuế).

\subsubsection{6.7.2 Điểm mù thường gặp (Cognitive Traps \&
Limitations)}\label{ux111iux1ec3m-muxf9-thux1b0ux1eddng-gux1eb7p-cognitive-traps-limitations-27}

\subsubsection{6.7.2 Điểm mù thường gặp (Cognitive Traps \&
Limitations)}\label{ux111iux1ec3m-muxf9-thux1b0ux1eddng-gux1eb7p-cognitive-traps-limitations-28}

Noise traders có thể không bị thị trường đào thải như giả định của Fama.
De Long et al.~(1990) chứng minh rằng những nhà đầu tư nhiễu này gánh
chịu nhiều rủi ro hơn, nên trung bình họ lại có tỷ suất sinh lời cao
hơn, giúp họ tồn tại lâu dài trên thị trường.

\subsubsection{6.7.3 Ứng dụng thực tế (Practical
Applications)}\label{ux1ee9ng-dux1ee5ng-thux1ef1c-tux1ebf-practical-applications-15}

Giải thích bong bóng tài sản (Dot-com bubble, Crypto), và chiến lược
giao dịch đà (Momentum trading).

\subsubsection{6.7.4 Ví dụ minh họa tính toán (Worked Numerical
Examples)}\label{vuxed-dux1ee5-minh-hux1ecda-tuxednh-touxe1n-worked-numerical-examples-14}

\textbf{Quyết định tối ưu thuế (duy lý):} Bán cổ phiếu B đang lỗ --- ghi
nhận lỗ vốn 30.000đ/cp (10 cp \(\rightarrow\) lỗ 300.000đ) được khấu trừ
vào lợi nhuận vốn khác, tiết kiệm thuế. Giữ cổ phiếu A tiếp tục tăng
trưởng kỳ vọng.

\textbf{Thực tế --- Disposition Effect:} Phần lớn nhà đầu tư chọn bán cổ
phiếu A (đang lãi) để ``chốt lời'' và giữ cổ phiếu B (đang lỗ) để ``đợi
phục hồi''. Lý do theo lý thuyết triển vọng: trong vùng lãi
\(\rightarrow\) ngại rủi ro, muốn khóa lợi nhuận chắc; trong vùng lỗ
\(\rightarrow\) ưa thích rủi ro, không muốn chấp nhận lỗ thực. Kết quả:
giữ cổ phiếu lỗ thường tiếp tục lỗ (các nghiên cứu cho thấy cổ phiếu bị
bán và cổ phiếu được giữ từ các quyết định Disposition Effect cho thấy
cổ phiếu bị giữ underperform sau đó trung bình 3,4\% trong 1 năm). Chi
phí ẩn: không tiết kiệm được thuế + giữ cổ phiếu xấu = \textbf{tổn thất
kép} từ một thiên kiến duy nhất.

\begin{center}\rule{0.5\linewidth}{0.5pt}\end{center}

\subsection{6.8 Kinh tế học thực nghiệm \& Thiết kế thực nghiệm ngẫu
nhiên (Experimental Economics \&
RCT)}\label{kinh-tux1ebf-hux1ecdc-thux1ef1c-nghiux1ec7m-thiux1ebft-kux1ebf-thux1ef1c-nghiux1ec7m-ngux1eabu-nhiuxean-experimental-economics-rct}

\subsubsection{6.8.1 Cốt lõi Lý thuyết (Core
Theory)}\label{cux1ed1t-luxf5i-luxfd-thuyux1ebft-core-theory-15}

\textbf{RCT (Randomized Controlled Trial):} Phân ngẫu nhiên đối tượng
vào nhóm thực nghiệm (treatment) và nhóm đối chứng (control) --- đảm bảo
\emph{không có sự khác biệt có hệ thống} giữa hai nhóm ngoài chính
sách/can thiệp được kiểm tra. Hệ quả: \textbf{ước lượng nhân quả}
(causal identification) thay vì chỉ tương quan. Ước lượng hiệu quả bình
quân xử lý (Average Treatment Effect --- ATE):
\(\hat{\tau}_{ATE} = \bar{Y}_{treated} - \bar{Y}_{control}\). Dưới phân
ngẫu nhiên thuần túy, ước lượng này không chệch.

Phân biệt ba cấp độ bằng chứng: \textbf{(a) Tương quan đơn giản} ---
``các trường có nhiều sách hơn thì điểm thi cao hơn''; có thể do trường
giàu mua nhiều sách hơn, không phải do sách. \textbf{(b) Hồi quy kiểm
soát biến} --- kiểm soát thu nhập trường, khu vực, v.v.; vẫn có thể còn
biến nội sinh không quan sát được. \textbf{(c) RCT} --- phân ngẫu nhiên
phát sách cho một nhóm trường, không phát cho nhóm đối chứng; chênh lệch
điểm thi là ảnh hưởng nhân quả của sách. Kết quả RCT của Kremer et al.~ở
Kenya: phát sách giáo khoa KHÔNG cải thiện điểm thi của hầu hết học sinh
--- trái hoàn toàn với kỳ vọng trực giác và chính sách giáo dục của
nhiều quốc gia.

\subsubsection{6.8.2 Điểm mù thường gặp (Cognitive Traps \&
Limitations)}\label{ux111iux1ec3m-muxf9-thux1b0ux1eddng-gux1eb7p-cognitive-traps-limitations-29}

\subsubsection{6.8.2 Điểm mù thường gặp (Cognitive Traps \&
Limitations)}\label{ux111iux1ec3m-muxf9-thux1b0ux1eddng-gux1eb7p-cognitive-traps-limitations-30}

Giá trị ngoại suy (External validity). Một kết quả thí nghiệm trong
phòng lab với sinh viên đại học chơi trò chơi Ultimatum với số tiền nhỏ
có thể không phản ánh đúng hành vi của CEO khi chốt giao dịch hàng tỷ đô
trên thị trường.

\subsubsection{6.8.3 Ứng dụng thực tế (Practical
Applications)}\label{ux1ee9ng-dux1ee5ng-thux1ef1c-tux1ebf-practical-applications-16}

Đánh giá hiệu quả chính sách công thông qua Thử nghiệm ngẫu nhiên có đối
chứng (RCT) - công cụ mang về Nobel 2019 cho Duflo, Banerjee và Kremer.

\subsubsection{6.8.4 Ví dụ minh họa tính toán (Worked Numerical
Examples)}\label{vuxed-dux1ee5-minh-hux1ecda-tuxednh-touxe1n-worked-numerical-examples-15}

{\def\LTcaptype{none} % do not increment counter
\begin{longtable}[]{@{}lll@{}}
\toprule\noalign{}
Nhóm & Thu nhập bình quân & Tỷ lệ có việc làm phi nông nghiệp \\
\midrule\noalign{}
\endhead
\bottomrule\noalign{}
\endlastfoot
Control (n=2.000 hộ) & 15,2 triệu/năm & 32\% \\
Treatment (n=2.000 hộ) & 17,8 triệu/năm & 41\% \\
\end{longtable}
}

\(\hat{\tau}_{ATE}^{income} = 17{,}8 - 15{,}2 = +2{,}6\) triệu/năm (tăng
17,1\%); \(\hat{\tau}_{ATE}^{employment} = 41\% - 32\% = +9\) điểm phần
trăm. Kiểm định t: nếu p-value \textless{} 0,05 ở cả hai chỉ số
\(\Rightarrow\) bằng chứng nhân quả mạnh rằng chương trình có hiệu quả.

\textbf{Lưu ý quan trọng:} Trước khi khẳng định ``chương trình cho vay
siêu nhỏ nên triển khai toàn quốc'', phải kiểm tra: (a) SUTVA (Stable
Unit Treatment Value Assumption) --- xã trong treatment không làm thay
đổi cơ hội của xã trong control (ví dụ không thu hút lao động từ xã đối
chứng); (b) Compliance --- bao nhiêu phần trăm hộ trong treatment thực
sự vay và sử dụng khoản vay; (c) External validity --- 200 xã thí điểm
có đại diện cho cả nước về điều kiện địa lý, thị trường, và thể chế địa
phương không? Thiếu bất kỳ kiểm tra nào trong ba điểm đó, kết quả RCT
vẫn có thể dẫn đến chính sách sai.

\begin{center}\rule{0.5\linewidth}{0.5pt}\end{center}

\section{CHƯƠNG 7: VĨ MÔ TỔNG QUÁT (General
Macroeconomics)}\label{chux1b0ux1a1ng-7-vux129-muxf4-tux1ed5ng-quuxe1t-general-macroeconomics}

Kinh tế vi mô giải thích hành vi của các đơn vị riêng lẻ; kinh tế vĩ mô
giải thích cách toàn bộ nền kinh tế vận hành như một hệ thống --- tại
sao sản lượng tăng hay thu hẹp, tại sao lạm phát xuất hiện, tại sao thất
nghiệp tăng trong suy thoái, và chính phủ có thể (hoặc nên) làm gì. Sự
khác biệt trọng yếu: ở cấp độ vĩ mô, \textbf{ngụy biện hợp thành}
(fallacy of composition) rất phổ biến --- những gì đúng với một hộ gia
đình hay một công ty đơn lẻ có thể sai hoàn toàn khi áp dụng cho toàn
nền kinh tế cùng lúc (Nghịch lý của Tiết kiệm, Chương 1 đã đề cập). Nhóm
này cung cấp khung phân tích các quyết định chính sách lớn nhất và các
cú sốc kinh tế toàn cầu. (Xem Phần B, Chương 13 (RBC/DSGE/Solow), DSGE,
Tăng trưởng Solow và nội sinh.)

\begin{center}\rule{0.5\linewidth}{0.5pt}\end{center}

\subsection{7.1 Mô hình AD-AS \& Các cú sốc kinh tế (Aggregate
Demand--Aggregate
Supply)}\label{muxf4-huxecnh-ad-as-cuxe1c-cuxfa-sux1ed1c-kinh-tux1ebf-aggregate-demandaggregate-supply}

\subsubsection{7.1.1 Cốt lõi Lý thuyết (Core
Theory)}\label{cux1ed1t-luxf5i-luxfd-thuyux1ebft-core-theory-16}

\textbf{Tổng cung ngắn hạn} (Short-Run AS --- SRAS): sản lượng mà toàn
bộ nền kinh tế sẵn sàng cung ứng ở mỗi mức giá, khi tiền lương danh
nghĩa và các đầu vào danh nghĩa \textbf{chưa điều chỉnh}. Đường SRAS dốc
lên vì giá cao hơn tạo lợi nhuận cho nhà sản xuất trong ngắn hạn khi chi
phí chưa tăng tương ứng. \textbf{Tổng cung dài hạn} (Long-Run AS ---
LRAS): sản lượng tiềm năng \(Y^*\) mà nền kinh tế tạo ra khi tất cả thị
trường --- kể cả thị trường lao động --- cân bằng, tức tất cả kỳ vọng
được thực hiện; LRAS thẳng đứng vì sản lượng dài hạn chỉ phụ thuộc vào
nguồn lực thực (vốn, lao động, công nghệ), không phụ thuộc vào mức giá
chung.

Phân tích cú sốc trong mô hình AD-AS: - \textbf{Cú sốc cầu dương} (tăng
chi tiêu chính phủ, giảm lãi suất, tăng tiêu dùng tư nhân): AD dịch phải
\(\rightarrow\) sản lượng tăng và giá tăng trong ngắn hạn; dài hạn SRAS
dịch trái (lương điều chỉnh tăng) \(\rightarrow\) trở về \(Y^*\) nhưng
giá cao hơn. - \textbf{Cú sốc cung tiêu cực} (tăng giá dầu, gián đoạn
chuỗi cung ứng): SRAS dịch trái \(\rightarrow\) \textbf{đình trệ lạm
phát} (stagflation --- sản lượng giảm và giá tăng cùng lúc), tình huống
khó nhất cho chính sách vì kích thích cầu chỉ chữa một vế.

\subsubsection{7.1.2 Điểm mù thường gặp (Cognitive Traps \&
Limitations)}\label{ux111iux1ec3m-muxf9-thux1b0ux1eddng-gux1eb7p-cognitive-traps-limitations-31}

\subsubsection{7.1.2 Điểm mù thường gặp (Cognitive Traps \&
Limitations)}\label{ux111iux1ec3m-muxf9-thux1b0ux1eddng-gux1eb7p-cognitive-traps-limitations-32}

Tính cứng nhắc của giá (Sticky prices) không đối xứng --- giá dễ tăng
nhưng khó giảm (Downward nominal rigidity). Điều này làm cho các cú sốc
cung âm đặc biệt nguy hiểm.

\subsubsection{7.1.3 Ứng dụng thực tế (Practical
Applications)}\label{ux1ee9ng-dux1ee5ng-thux1ef1c-tux1ebf-practical-applications-17}

Phân tích hiện tượng lạm phát đình trệ (Stagflation) và đánh giá độ trễ
của chính sách phục hồi kinh tế.

\subsubsection{7.1.4 Ví dụ minh họa tính toán (Worked Numerical
Examples)}\label{vuxed-dux1ee5-minh-hux1ecda-tuxednh-touxe1n-worked-numerical-examples-16}

{\def\LTcaptype{none} % do not increment counter
\begin{longtable}[]{@{}
  >{\raggedright\arraybackslash}p{(\linewidth - 6\tabcolsep) * \real{0.2500}}
  >{\raggedright\arraybackslash}p{(\linewidth - 6\tabcolsep) * \real{0.2500}}
  >{\raggedright\arraybackslash}p{(\linewidth - 6\tabcolsep) * \real{0.2500}}
  >{\raggedright\arraybackslash}p{(\linewidth - 6\tabcolsep) * \real{0.2500}}@{}}
\toprule\noalign{}
\begin{minipage}[b]{\linewidth}\raggedright
Lựa chọn chính sách
\end{minipage} & \begin{minipage}[b]{\linewidth}\raggedright
Tác động lên Y
\end{minipage} & \begin{minipage}[b]{\linewidth}\raggedright
Tác động lên P
\end{minipage} & \begin{minipage}[b]{\linewidth}\raggedright
Ý nghĩa
\end{minipage} \\
\midrule\noalign{}
\endhead
\bottomrule\noalign{}
\endlastfoot
Không làm gì & Dần tự điều chỉnh về 100 & Tự điều chỉnh về
\textasciitilde108 & Tốn thời gian, đau đớn ngắn hạn \\
Kích thích cầu (AD phải) & Phục hồi Y về 100 nhanh hơn & P tăng thêm
\(\rightarrow\) có thể 115--120 & Chấp nhận lạm phát để cứu tăng
trưởng \\
Thắt chặt cầu (AD trái) & Y giảm thêm \(\rightarrow\) suy thoái sâu hơn
& Chống lạm phát, giữ P thấp & Chấp nhận suy thoái sâu để kiểm soát
giá \\
Chính sách cung (đầu tư R\&D, năng lượng tái tạo) & SRAS dịch phải, Y
phục hồi & P giảm & Giải quyết đúng nguồn gốc, nhưng chậm \\
\end{longtable}
}

Không có lựa chọn ``vừa giữ sản lượng vừa kiểm soát giá'' --- đây là
đánh đổi chính sách cơ bản trong stagflation.

\begin{center}\rule{0.5\linewidth}{0.5pt}\end{center}

\subsection{7.2 Mô hình IS-LM \& Chính sách tài khóa-tiền tệ phối hợp
(IS-LM \& Policy
Mix)}\label{muxf4-huxecnh-is-lm-chuxednh-suxe1ch-tuxe0i-khuxf3a-tiux1ec1n-tux1ec7-phux1ed1i-hux1ee3p-is-lm-policy-mix}

\subsubsection{7.2.1 Cốt lõi Lý thuyết (Core
Theory)}\label{cux1ed1t-luxf5i-luxfd-thuyux1ebft-core-theory-17}

Phương trình cơ bản IS: \(Y = C(Y-T) + I(r) + G + NX\). Đạo hàm ngầm:
\(\frac{dr}{dY}\bigg|_{IS} = \frac{1-c_1(1-t)}{-b} < 0\) (dốc xuống),
trong đó \(c_1\) là xu hướng tiêu dùng biên, \(t\) là thuế suất,
\(b > 0\) là độ nhạy đầu tư theo lãi suất. Số nhân (multiplier) của chi
tiêu chính phủ:
\(\frac{\partial Y}{\partial G}\bigg|_{IS} = \frac{1}{1-c_1(1-t)} > 1\)
--- mỗi đơn vị tăng G làm tổng cầu tăng nhiều hơn 1 đơn vị qua vòng lặp
tiêu dùng.

\textbf{Đường LM} (Liquidity preference = Money supply): tập hợp các cặp
(Y, r) mà thị trường tiền tệ cân bằng (\(M/P = L(Y, r)\), với
\(L_Y > 0\) và \(L_r < 0\)). Dốc lên: sản lượng cao hơn \(\rightarrow\)
cầu tiền giao dịch tăng \(\rightarrow\) lãi suất phải tăng để cân bằng
thị trường tiền tệ.

\textbf{Hiệu ứng lấn át (Crowding-out):} Khi chính phủ tăng chi tiêu G
(IS dịch phải), lãi suất tăng theo đường LM \(\rightarrow\) giảm đầu tư
tư nhân I. Nếu LM thẳng đứng (cầu tiền không co giãn theo lãi suất):
crowding-out hoàn toàn --- G tăng không tăng Y. Nếu LM nằm ngang (bẫy
thanh khoản, cầu tiền co giãn vô hạn theo lãi suất): không có
crowding-out --- số nhân tài khóa đầy đủ phát huy tác dụng.

\textbf{Phối hợp chính sách (Policy mix):} Mô hình IS-LM cho phép phân
tích tổ hợp chính sách tài khóa và tiền tệ. Để tăng Y mà không thay đổi
lãi suất: kích thích tài khóa (IS phải) + nới lỏng tiền tệ vừa đủ (LM
phải) để giữ r cố định. Để tăng Y mà không tăng thâm hụt ngân sách: nới
lỏng tiền tệ mạnh (LM phải), chấp nhận lãi suất giảm.

\subsubsection{7.2.2 Điểm mù thường gặp (Cognitive Traps \&
Limitations)}\label{ux111iux1ec3m-muxf9-thux1b0ux1eddng-gux1eb7p-cognitive-traps-limitations-33}

\subsubsection{7.2.2 Điểm mù thường gặp (Cognitive Traps \&
Limitations)}\label{ux111iux1ec3m-muxf9-thux1b0ux1eddng-gux1eb7p-cognitive-traps-limitations-34}

Bẫy thanh khoản (Liquidity trap). Khi lãi suất chạm mốc 0 (Zero Lower
Bound - ZLB), đường LM nằm ngang và chính sách tiền tệ truyền thống hoàn
toàn bất lực trong việc kích cầu.

\subsubsection{7.2.3 Ứng dụng thực tế (Practical
Applications)}\label{ux1ee9ng-dux1ee5ng-thux1ef1c-tux1ebf-practical-applications-18}

Chính sách nới lỏng định lượng (QE) và lãi suất âm (ZIRP) của Nhật Bản
(thập niên 1990) và Mỹ (hậu 2008).

\subsubsection{7.2.4 Ví dụ minh họa tính toán (Worked Numerical
Examples)}\label{vuxed-dux1ee5-minh-hux1ecda-tuxednh-touxe1n-worked-numerical-examples-17}

\textbf{Phương trình IS:}
\(Y = C + I + G = 200 + 0{,}6(Y-200) + 300 - 40r + 400\). Rút gọn:
\(Y = 200 + 0{,}6Y - 120 + 300 - 40r + 400 = 780 + 0{,}6Y - 40r\). Suy
ra: \(0{,}4Y = 780 - 40r \Rightarrow Y = 1.950 - 100r\). \textbf{Đây là
phương trình IS.}

\textbf{Phương trình LM:}
\(M/P = L(Y, r) \Rightarrow 500 = 0{,}5Y - 30r \Rightarrow Y = 1.000 + 60r\).
\textbf{Đây là phương trình LM.}

\textbf{Cân bằng IS-LM:}
\(1.950 - 100r = 1.000 + 60r \Rightarrow 950 = 160r \Rightarrow r^* = 5{,}94\%\).
\(Y^* = 1.000 + 60 \times 5{,}94 = 1.356\).

\textbf{Tác động khi G tăng thêm 50 (IS dịch phải):} IS mới:
\(Y = 2.075 - 100r\). Cân bằng mới:
\(2.075 - 100r = 1.000 + 60r \Rightarrow r^{**} = 6{,}72\%\),
\(Y^{**} = 1.403\). \(\Delta Y = 1.403 - 1.356 = 47\). Số nhân IS-LM:
\(\Delta Y / \Delta G = 47/50 = 0{,}94 < 1\) --- nhỏ hơn số nhân
Keynesian thuần túy (= \(1/(1-0{,}6) = 2{,}5\)) vì \textbf{crowding-out
một phần}:
\(\Delta r = 0{,}78\% \Rightarrow \Delta I = -40 \times 0{,}78 = -31\)
--- đầu tư giảm 31 đơn vị, ``ăn'' bớt hiệu quả của kích thích tài khóa.

\begin{center}\rule{0.5\linewidth}{0.5pt}\end{center}

\subsection{7.3 Đường Phillips \& Kỳ vọng lạm phát (Phillips Curve \&
Inflation
Expectations)}\label{ux111ux1b0ux1eddng-phillips-kux1ef3-vux1ecdng-lux1ea1m-phuxe1t-phillips-curve-inflation-expectations}

\begin{figure}
\centering
\pandocbounded{\includegraphics[keepaspectratio,alt={Minh họa Đường Phillips}]{C:/Users/Admin/.gemini/antigravity/brain/e5791226-a607-45b8-984e-554eb13f4da1/duong_phillips_1785202072852.jpg}}
\caption{Minh họa Đường Phillips}
\end{figure}

\subsubsection{7.3.1 Cốt lõi Lý thuyết (Core
Theory)}\label{cux1ed1t-luxf5i-luxfd-thuyux1ebft-core-theory-18}

Nhưng Friedman (1968) và Phelps (1968) đã phản biện căn bản: đường
Phillips cổ điển bỏ qua \textbf{kỳ vọng lạm phát}. Dạng đúng đắn hơn là
\textbf{Đường Phillips kỳ vọng tăng cường} (expectations-augmented
Phillips Curve):
\[ \pi_t = \pi_t^e - \alpha(u_t - u^n) + \varepsilon_t \] trong đó
\(\pi_t^e\) là lạm phát kỳ vọng tại kỳ \(t\). Dưới kỳ vọng thích ứng
(\(\pi_t^e = \pi_{t-1}\)): đường Phillips dịch chuyển theo lạm phát thực
từ kỳ trước --- không có đánh đổi dài hạn, chỉ có đường Phillips ngắn
hạn. Dưới kỳ vọng duy lý: nếu NHTW không tạo bất ngờ, không có đánh đổi
cả ngắn hạn. Dài hạn: LRPC thẳng đứng tại \(u = u^n\) --- xác nhận
Friedman-Phelps và giải thích thất bại của chính sách kích thích cầu
liên tục thập niên 1970.

Phiên bản hiện đại --- \textbf{Đường Phillips Mới Keynesian (NKPC)}:
\[ \pi_t = \beta \mathbb{E}_t[\pi_{t+1}] + \kappa (y_t - y_t^*) \] trong
đó \(\beta \approx 1\) là hệ số chiết khấu, \(\kappa > 0\) là hệ số dốc
(phụ thuộc độ cứng danh nghĩa), và \((y_t - y_t^*)\) là khoảng cách sản
lượng (output gap). Đây là cơ sở vi mô của đường Phillips từ mô hình
định giá Calvo (1983): các hãng điều chỉnh giá theo xác suất ngẫu nhiên
mỗi kỳ, không phải tất cả cùng lúc --- tạo ra độ cứng danh nghĩa trong
nền kinh tế.

\subsubsection{7.3.2 Điểm mù thường gặp (Cognitive Traps \&
Limitations)}\label{ux111iux1ec3m-muxf9-thux1b0ux1eddng-gux1eb7p-cognitive-traps-limitations-35}

\subsubsection{7.3.2 Điểm mù thường gặp (Cognitive Traps \&
Limitations)}\label{ux111iux1ec3m-muxf9-thux1b0ux1eddng-gux1eb7p-cognitive-traps-limitations-36}

Bỏ qua vai trò của kỳ vọng lạm phát (Inflation expectations). Việc cố
gắng khai thác đường Phillips ngắn hạn để giảm thất nghiệp sẽ làm kỳ
vọng lạm phát tăng lên, dịch chuyển toàn bộ đường Phillips lên trên
trong dài hạn (Friedman-Phelps).

\subsubsection{7.3.3 Ứng dụng thực tế (Practical
Applications)}\label{ux1ee9ng-dux1ee5ng-thux1ef1c-tux1ebf-practical-applications-19}

Sự sụp đổ của hệ thống Bretton Woods và thời kỳ Đại lạm phát (Great
Inflation) thập niên 1970 tại Mỹ.

\subsubsection{7.3.4 Ví dụ minh họa tính toán (Worked Numerical
Examples)}\label{vuxed-dux1ee5-minh-hux1ecda-tuxednh-touxe1n-worked-numerical-examples-18}

\textbf{Trạng thái ban đầu:} \(u_0 = 5\%\), \(\pi_0 = 2\%\). Ổn định ---
không có áp lực lạm phát.

\textbf{Chính phủ kích thích cầu, thất nghiệp giảm xuống 3\%:}

{\def\LTcaptype{none} % do not increment counter
\begin{longtable}[]{@{}llll@{}}
\toprule\noalign{}
Kỳ & \(u_t\) & \(\pi_t^e = \pi_{t-1}\) & \(\pi_t\) \\
\midrule\noalign{}
\endhead
\bottomrule\noalign{}
\endlastfoot
1 & 3\% & 2\% & 2\% + 0,5×(5−3) = 3\% \\
2 & 3\% & 3\% & 3\% + 0,5×2 = 4\% \\
3 & 3\% & 4\% & 4\% + 1 = 5\% \\
4 & 3\% & 5\% & 6\% \\
\end{longtable}
}

Lạm phát tăng dần và \textbf{không có giới hạn} khi thất nghiệp duy trì
dưới \(u^n\). Muốn kiểm soát lạm phát, phải chấp nhận \(u > u^n\) trong
một thời gian để ``kéo'' kỳ vọng xuống --- đây là chi phí thực tế của
sai lầm chính sách tiền tệ. Sacrifice ratio ước tính: cần duy trì
\(u = 7\%\) (tức \(u - u^n = 2\%\)) trong 2 năm để giảm lạm phát 4 điểm
phần trăm
(\(2\times2\times0{,}5 = 2\%/\text{năm} \times 2 \text{ năm}\)) ---
tương đương thiệt hại sản lượng đáng kể.

\begin{center}\rule{0.5\linewidth}{0.5pt}\end{center}

\subsection{7.4 Chính sách tiền tệ hiện đại \& Quy tắc Taylor (Modern
Monetary Policy \& Taylor
Rule)}\label{chuxednh-suxe1ch-tiux1ec1n-tux1ec7-hiux1ec7n-ux111ux1ea1i-quy-tux1eafc-taylor-modern-monetary-policy-taylor-rule}

\subsubsection{7.4.1 Cốt lõi Lý thuyết (Core
Theory)}\label{cux1ed1t-luxf5i-luxfd-thuyux1ebft-core-theory-19}

\textbf{Quy tắc Taylor (1993)} là mô tả thực nghiệm về cách các NHTW
thực tế (đặc biệt Fed) đặt lãi suất: i\_t = r\^{}* + \pi\^{}* +
\phi\_\pi(\pi\_t - \pi\^{}*) + \phi\_y(y\_t - y\_t\^{}\emph{) trong đó
\(r^*\) là lãi suất thực cân bằng dài hạn (r-star), \(\pi^*\) là mục
tiêu lạm phát (thường 2\%), \(\pi_t\) là lạm phát thực tế,
\(y_t - y_t^*\) là khoảng cách sản lượng, và \(\phi_\pi > 1\),
\(\phi_y > 0\) là trọng số phản ứng. \textbf{Điều kiện Taylor (Taylor
Principle)}: \(\phi_\pi > 1\) --- lãi suất danh nghĩa phải tăng }nhiều
hơn 1:1* so với lạm phát tăng; tức \textbf{lãi suất thực phải tăng} khi
lạm phát vượt mục tiêu. Nếu \(\phi_\pi < 1\), lãi suất thực giảm khi lạm
phát tăng \(\rightarrow\) phản hồi ổn định hóa sai chiều \(\rightarrow\)
có thể dẫn đến vòng xoáy lạm phát tự thực hiện (indeterminacy).

Taylor ước lượng nguyên gốc: \(\phi_\pi = 1{,}5\), \(\phi_y = 0{,}5\),
với \(r^* = \pi^* = 2\%\): i\_t = 4 + 1\{,\}5(\pi\_t - 2) + 0\{,\}5(y\_t
- y\_t\^{}*)

\subsubsection{7.4.2 Điểm mù thường gặp (Cognitive Traps \&
Limitations)}\label{ux111iux1ec3m-muxf9-thux1b0ux1eddng-gux1eb7p-cognitive-traps-limitations-37}

\subsubsection{7.4.2 Điểm mù thường gặp (Cognitive Traps \&
Limitations)}\label{ux111iux1ec3m-muxf9-thux1b0ux1eddng-gux1eb7p-cognitive-traps-limitations-38}

Việc đo lường ``Khoảng trống sản lượng'' (Output gap) và ``Lãi suất thực
tự nhiên'' theo thời gian thực là cực kỳ sai số và bị điều chỉnh lại
nhiều lần (Data revisions), khiến việc áp dụng quy tắc Taylor một cách
cơ học trở nên nguy hiểm.

\subsubsection{7.4.3 Ứng dụng thực tế (Practical
Applications)}\label{ux1ee9ng-dux1ee5ng-thux1ef1c-tux1ebf-practical-applications-20}

Giải thích chu kỳ thắt chặt/nới lỏng của FED và các sai lầm dự báo chính
sách trước các cú sốc bất ngờ.

\subsubsection{7.4.4 Ví dụ minh họa tính toán (Worked Numerical
Examples)}\label{vuxed-dux1ee5-minh-hux1ecda-tuxednh-touxe1n-worked-numerical-examples-19}

i = 2 + 4 + 1\{,\}5(6-4) + 0\{,\}5(1) = 2 + 4 + 3 + 0\{,\}5 = 9\{,\}5\%
Nếu lãi suất chính sách hiện tại là 7\%: lãi suất đang \textbf{thấp hơn
mức Taylor khoảng 2,5\%} --- hàm ý chính sách tiền tệ đang nới lỏng so
với mức cần thiết để kiềm chế lạm phát theo khung Taylor. Cảnh báo:
không nên tăng thẳng lên 9,5\% ngay lập tức vì: (a) r-star có thể thấp
hơn giả định; (b) độ trễ chính sách tiền tệ 6--18 tháng; (c) rủi ro
landing cứng nếu thắt chặt quá nhanh. Quy tắc Taylor là \textbf{điểm
tham chiếu}, không phải đơn thuốc cứng nhắc.

\begin{center}\rule{0.5\linewidth}{0.5pt}\end{center}

\subsection{7.5 Bài toán bất nhất quán theo thời gian \& Tính độc lập
của NHTW (Time
Inconsistency)}\label{buxe0i-touxe1n-bux1ea5t-nhux1ea5t-quuxe1n-theo-thux1eddi-gian-tuxednh-ux111ux1ed9c-lux1eadp-cux1ee7a-nhtw-time-inconsistency}

\subsubsection{7.5.1 Cốt lõi Lý thuyết (Core
Theory)}\label{cux1ed1t-luxf5i-luxfd-thuyux1ebft-core-theory-20}

Về mặt hình thức: hàm mất mát (loss function) của NHTW:
\(\mathcal{L} = (y - y^*)^2 + \lambda \pi^2\), với \(y^* > y^n\) (mục
tiêu sản lượng tham vọng hơn tiềm năng). Dưới tự quyết (discretion): lạm
phát cân bằng \(\bar{\pi} = \frac{\kappa(y^* - y^n)}{\lambda} > 0\) ---
thiên kiến lạm phát dương, tỷ lệ nghịch với mức ngại lạm phát
\(\lambda\). Giải pháp: \textbf{ủy quyền} (delegation) cho NHTW độc lập
với nhà chính trị và có mức ngại lạm phát
\(\lambda_{CB} > \lambda_{govt}\) (thống đốc ``bảo thủ'' --- Rogoff
1985), hay \textbf{hợp đồng khuyến khích} (Walsh 1995), hay \textbf{mục
tiêu lạm phát rõ ràng} (inflation targeting) tạo cam kết đáng tin.

\subsubsection{7.5.2 Điểm mù thường gặp (Cognitive Traps \&
Limitations)}\label{ux111iux1ec3m-muxf9-thux1b0ux1eddng-gux1eb7p-cognitive-traps-limitations-39}

\subsubsection{7.5.2 Điểm mù thường gặp (Cognitive Traps \&
Limitations)}\label{ux111iux1ec3m-muxf9-thux1b0ux1eddng-gux1eb7p-cognitive-traps-limitations-40}

Chính phủ luôn hứa duy trì lạm phát thấp nhưng luôn có động cơ
(incentive) lén tạo ra lạm phát bất ngờ để giảm thất nghiệp hoặc làm tan
chảy nợ công. Công chúng biết điều này nên sẽ không tin vào lời hứa.

\subsubsection{7.5.3 Ứng dụng thực tế (Practical
Applications)}\label{ux1ee9ng-dux1ee5ng-thux1ef1c-tux1ebf-practical-applications-21}

Thiết kế thể chế: Cấp sự độc lập hoàn toàn cho Ngân hàng Trung ương
(Central Bank Independence) và neo chính sách (Policy anchors).

\subsubsection{7.5.4 Ví dụ minh họa tính toán (Worked Numerical
Examples)}\label{vuxed-dux1ee5-minh-hux1ecda-tuxednh-touxe1n-worked-numerical-examples-20}

\begin{center}\rule{0.5\linewidth}{0.5pt}\end{center}

\subsection{7.6 Tăng trưởng kinh tế: Solow, Hội tụ \& Tăng trưởng nội
sinh (Economic
Growth)}\label{tux103ng-trux1b0ux1edfng-kinh-tux1ebf-solow-hux1ed9i-tux1ee5-tux103ng-trux1b0ux1edfng-nux1ed9i-sinh-economic-growth}

\subsubsection{7.6.1 Cốt lõi Lý thuyết (Core
Theory)}\label{cux1ed1t-luxf5i-luxfd-thuyux1ebft-core-theory-21}

\[ \dot{k} = s \cdot f(k) - (n + \delta)k \] trong đó \(s\) là tỷ lệ
tiết kiệm, \(n\) là tốc độ tăng dân số, \(\delta\) là tỷ lệ khấu hao.
Tại \textbf{trạng thái ổn định} (steady state): \(\dot{k} = 0\), tức đầu
tư (sf(k)) bằng đầu tư cần để duy trì tỷ lệ vốn/lao động
\((n+\delta)k\). Trong trạng thái ổn định: sản lượng bình quân
\(y^* = f(k^*)\) không tăng, chỉ có \textbf{tiến bộ công nghệ ngoại
sinh} (A tăng theo thời gian) mới tạo ra tăng trưởng bền vững trong dài
hạn.

Hai dự đoán kiểm định được: \textbf{(a) Hội tụ} (convergence): các nước
nghèo hơn (có k thấp hơn, MPK cao hơn) nên tăng trưởng nhanh hơn nước
giàu --- hội tụ về trạng thái ổn định chung (nếu các nước có thông số
giống nhau: hội tụ tuyệt đối; nếu khác nhau: hội tụ có điều kiện). Bằng
chứng thực nghiệm: hội tụ tuyệt đối yếu (không thấy nước nghèo nhất luôn
tăng nhanh nhất toàn cầu); hội tụ có điều kiện mạnh hơn (khi kiểm soát
các đặc điểm quốc gia --- Mankiw, Romer \& Weil 1992). \textbf{(b) Tỷ lệ
tiết kiệm}: nước có tỷ lệ tiết kiệm cao hơn có \(k^*\) cao hơn
\(\rightarrow\) thu nhập cao hơn trong trạng thái ổn định, nhưng không
cao hơn tốc độ tăng trưởng dài hạn.

\subsubsection{7.6.2 Điểm mù thường gặp (Cognitive Traps \&
Limitations)}\label{ux111iux1ec3m-muxf9-thux1b0ux1eddng-gux1eb7p-cognitive-traps-limitations-41}

\subsubsection{7.6.2 Điểm mù thường gặp (Cognitive Traps \&
Limitations)}\label{ux111iux1ec3m-muxf9-thux1b0ux1eddng-gux1eb7p-cognitive-traps-limitations-42}

Tích lũy vốn vật chất không thể thúc đẩy tăng trưởng vô hạn do quy luật
sinh lời giảm dần (Diminishing returns). Tăng trưởng dài hạn chỉ có thể
đến từ tiến bộ công nghệ (A).

\subsubsection{7.6.3 Ứng dụng thực tế (Practical
Applications)}\label{ux1ee9ng-dux1ee5ng-thux1ef1c-tux1ebf-practical-applications-22}

Phân tích Bẫy thu nhập trung bình (Middle-income trap) khi các quốc gia
đang phát triển cạn kiệt lợi thế vốn rẻ và dân số trẻ.

\subsubsection{7.6.4 Ví dụ minh họa tính toán (Worked Numerical
Examples)}\label{vuxed-dux1ee5-minh-hux1ecda-tuxednh-touxe1n-worked-numerical-examples-21}

\textbf{Đóng góp của tích lũy nhân tố:} - Đóng góp của K:
\(\alpha \times \dot{K}/K = \frac{1}{3} \times 4{,}5\% = 1{,}5\%\) -
Đóng góp của L:
\((1-\alpha) \times \dot{L}/L = \frac{2}{3} \times 1{,}5\% = 1{,}0\%\) -
Tổng đóng góp nhân tố: 2,5\%

\textbf{Solow Residual (TFP growth):}
\(\dot{A}/A = \dot{Y}/Y - \alpha\dot{K}/K - (1-\alpha)\dot{L}/L = 6\% - 2{,}5\% = 3{,}5\%\)

\textbf{Kết luận:} Tăng trưởng GDP 6\% bao gồm: 1,5\% từ tích lũy vốn,
1,0\% từ tăng lao động, và \textbf{3,5\% từ TFP} --- tức hơn một nửa
tăng trưởng đến từ hiệu quả và công nghệ (Solow residual), không chỉ từ
đầu vào. Điều này nhấn mạnh tầm quan trọng của \textbf{cải thiện TFP}
(cải cách thể chế, R\&D, giáo dục) so với chỉ tăng đầu tư tư bản đơn
thuần trong chiến lược tăng trưởng dài hạn.

\begin{center}\rule{0.5\linewidth}{0.5pt}\end{center}

\subsection{7.7 Thể chế, Quản trị \& Tăng trưởng (Institutions,
Governance \&
Growth)}\label{thux1ec3-chux1ebf-quux1ea3n-trux1ecb-tux103ng-trux1b0ux1edfng-institutions-governance-growth}

\subsubsection{7.7.1 Cốt lõi Lý thuyết (Core
Theory)}\label{cux1ed1t-luxf5i-luxfd-thuyux1ebft-core-theory-22}

Phân biệt hai loại thể chế kinh tế trong ``Why Nations Fail'' (Acemoglu
\& Robinson 2012): \textbf{Thể chế bao gồm} (inclusive) --- bảo vệ quyền
tài sản rộng rãi, pháp quyền, thị trường cạnh tranh, môi trường kinh
doanh bình đẳng --- tạo khuyến khích đầu tư và đổi mới. \textbf{Thể chế
trích lợi} (extractive) --- quyền tài sản chỉ cho nhóm nhỏ cầm quyền,
pháp lý bất bình đẳng, rào cản gia nhập thị trường bảo hộ nhóm lợi ích
--- tạo khuyến khích tìm kiếm địa tô (rent-seeking) thay vì sản xuất.

\subsubsection{7.7.2 Điểm mù thường gặp (Cognitive Traps \&
Limitations)}\label{ux111iux1ec3m-muxf9-thux1b0ux1eddng-gux1eb7p-cognitive-traps-limitations-43}

\subsubsection{7.7.2 Điểm mù thường gặp (Cognitive Traps \&
Limitations)}\label{ux111iux1ec3m-muxf9-thux1b0ux1eddng-gux1eb7p-cognitive-traps-limitations-44}

Thể chế khai thác (Extractive) có thể tạo ra tăng trưởng nhanh trong
giai đoạn đầu (bằng cách ép buộc di chuyển nguồn lực từ nông nghiệp sang
công nghiệp), nhưng sẽ đình trệ vì chúng kìm hãm ``hủy diệt sáng tạo''
(Creative destruction).

\subsubsection{7.7.3 Ứng dụng thực tế (Practical
Applications)}\label{ux1ee9ng-dux1ee5ng-thux1ef1c-tux1ebf-practical-applications-23}

Lý giải sự sụp đổ của mô hình kinh tế Liên Xô và sự khác biệt phát triển
giữa Nam Hàn - Bắc Hàn (Acemoglu \& Robinson).

\subsubsection{7.7.4 Ví dụ minh họa tính toán (Worked Numerical
Examples)}\label{vuxed-dux1ee5-minh-hux1ecda-tuxednh-touxe1n-worked-numerical-examples-22}

\begin{center}\rule{0.5\linewidth}{0.5pt}\end{center}

\subsection{7.8 Khủng hoảng tài chính \& Hệ thống ngân hàng (Financial
Crises \& Banking
System)}\label{khux1ee7ng-houx1ea3ng-tuxe0i-chuxednh-hux1ec7-thux1ed1ng-nguxe2n-huxe0ng-financial-crises-banking-system}

\subsubsection{7.8.1 Cốt lõi Lý thuyết (Core
Theory)}\label{cux1ed1t-luxf5i-luxfd-thuyux1ebft-core-theory-23}

\textbf{(a) Mô hình Diamond--Dybvig (1983):} Ngân hàng thực hiện chức
năng \textbf{biến đổi kỳ hạn} (maturity transformation) --- nhận tiền
gửi ngắn hạn (linh hoạt rút) và cho vay dài hạn (kém thanh khoản). Điều
này tạo ra \textbf{hai cân bằng Nash}: cân bằng ``tốt'' --- người gửi
tin tưởng hệ thống, chỉ rút khi thực sự cần tiêu dùng sớm (loại 1); cân
bằng ``xấu'' --- nếu đủ nhiều người nghĩ những người khác sẽ rút, ai
cũng vội rút trước (bank run). Cân bằng xấu là \textbf{cân bằng Nash tự
thực hiện}: mỗi người rút vì kỳ vọng người khác rút --- và kỳ vọng đó
trở thành hiện thực khi đủ nhiều người hành động theo. Bảo hiểm tiền gửi
(deposit insurance) loại bỏ cân bằng xấu vì triệt tiêu động cơ rút sớm.

\textbf{(b) Mô hình Mishkin về Khuếch đại Tài chính:} Tài sản thế chấp
sụt giá \(\rightarrow\) giá trị thực của tài sản ngân hàng giảm
\(\rightarrow\) hạn chế tín dụng (credit crunch) \(\rightarrow\) đầu tư
và tiêu dùng giảm \(\rightarrow\) giá tài sản tiếp tục giảm
\(\rightarrow\) vòng xoáy \textbf{bộ khuếch đại tài chính} (financial
accelerator --- Bernanke, Gertler \& Gilchrist 1999). Vòng xoáy này biến
một cú sốc tài chính ban đầu có thể nhỏ thành khủng hoảng kinh tế thực
nghiêm trọng.

\textbf{(c) Mô hình Kỳ vọng đầu cơ (Krugman 1979, 1998):} Khủng hoảng
tiền tệ và ngân hàng xảy ra khi thị trường kỳ vọng chính phủ không thể
duy trì cam kết (tỷ giá cố định, ngân hàng trung ương không đủ dự trữ)
--- hành động đầu cơ phối hợp của thị trường tạo ra chính khủng hoảng mà
họ kỳ vọng, ngay cả khi ban đầu các căn bản (fundamentals) còn ổn.

\subsubsection{7.8.2 Điểm mù thường gặp (Cognitive Traps \&
Limitations)}\label{ux111iux1ec3m-muxf9-thux1b0ux1eddng-gux1eb7p-cognitive-traps-limitations-45}

\subsubsection{7.8.2 Điểm mù thường gặp (Cognitive Traps \&
Limitations)}\label{ux111iux1ec3m-muxf9-thux1b0ux1eddng-gux1eb7p-cognitive-traps-limitations-46}

Rủi ro đạo đức (Moral hazard) do tâm lý ``Quá lớn để đổ vỡ'' (Too Big To
Fail). Việc cứu trợ ngân hàng hôm nay gieo mầm cho các hành vi chấp nhận
rủi ro liều lĩnh hơn vào ngày mai.

\subsubsection{7.8.3 Ứng dụng thực tế (Practical
Applications)}\label{ux1ee9ng-dux1ee5ng-thux1ef1c-tux1ebf-practical-applications-24}

Các gói giải cứu (Bailout) năm 2008, Đạo luật Dodd-Frank, và yêu cầu vốn
đệm (Basel III) để chống rủi ro hệ thống.

\subsubsection{7.8.4 Ví dụ minh họa tính toán (Worked Numerical
Examples)}\label{vuxed-dux1ee5-minh-hux1ecda-tuxednh-touxe1n-worked-numerical-examples-23}

\textbf{Cân bằng tốt (không có bank run):} Chỉ 40\% (Loại 1) rút ở T=1.
Ngân hàng thanh lý \(0{,}8 \times 400 = 320\) đơn vị từ dự án sớm, trả
đủ 400 đơn vị cho Loại 1 (cần tài trợ chéo nhỏ từ phần khác). 60\% còn
lại chờ đến T=2, nhận \(1{,}3 \times 600 / 600 = 1{,}3\) mỗi người. Mọi
người đều được phục vụ.

\textbf{Cân bằng xấu (bank run toàn bộ):} 100\% người gửi đổ xô rút ở
T=1. Ngân hàng phải thanh lý toàn bộ dự án sớm: thu về
\(0{,}8 \times 1.000 = 800\). Chia đều: 800/1.000 = \textbf{0,8 mỗi
người}. Cả loại 1 lẫn loại 2 đều chỉ nhận 0,8 (tệ hơn 1,3 đáng lẽ loại 2
nhận được). \textbf{Bank run tự thực hiện}: nếu Loại 2 tin rằng Loại 2
khác cũng rút, họ tốt hơn nên rút (nhận 0,8 \textgreater{} 0 nếu ngân
hàng vỡ hoàn toàn và họ đến muộn). Bảo hiểm tiền gửi đầy đủ loại bỏ lý
do để Loại 2 rút sớm --- cân bằng xấu biến mất.

\begin{center}\rule{0.5\linewidth}{0.5pt}\end{center}

\begin{center}\rule{0.5\linewidth}{0.5pt}\end{center}

\section{PHẦN B --- CHUYÊN ĐỀ NÂNG CAO (PHD
LEVEL)}\label{phux1ea7n-b-chuyuxean-ux111ux1ec1-nuxe2ng-cao-phd-level}

\section{CHƯƠNG 8: KINH TẾ HỌC CÔNG CỘNG \& LỰA CHỌN XÃ HỘI (Public
Economics \& Social
Choice)}\label{chux1b0ux1a1ng-8-kinh-tux1ebf-hux1ecdc-cuxf4ng-cux1ed9ng-lux1ef1a-chux1ecdn-xuxe3-hux1ed9i-public-economics-social-choice}

Chương này cung cấp nền tảng lý thuyết nghiêm ngặt cho vai trò của nhà
nước trong nền kinh tế, phân tích các thất bại của thị trường và các cơ
chế tổng hợp sở thích cá nhân thành quyết định tập thể. Chúng ta sẽ sử
dụng công cụ tối ưu hóa tĩnh và lý thuyết trò chơi để đánh giá hiệu quả,
công bằng và các hệ quả chính sách.

\begin{center}\rule{0.5\linewidth}{0.5pt}\end{center}

\subsection{8.1 Kinh tế học Phúc lợi \& Các Định lý Cơ bản (Welfare
Economics \& Fundamental
Theorems)}\label{kinh-tux1ebf-hux1ecdc-phuxfac-lux1ee3i-cuxe1c-ux111ux1ecbnh-luxfd-cux1a1-bux1ea3n-welfare-economics-fundamental-theorems}

\subsubsection{8.1.1 Cốt lõi Lý thuyết (Core
Theory)}\label{cux1ed1t-luxf5i-luxfd-thuyux1ebft-core-theory-24}

Kinh tế học phúc lợi nghiên cứu các điều kiện để phân bổ nguồn lực đạt
hiệu quả Pareto. Giả sử nền kinh tế có \(I\) cá nhân
(\(i = 1, \dots, I\)) và \(J\) doanh nghiệp (\(j = 1, \dots, J\)), với
\(L\) hàng hóa. Mỗi cá nhân có hàm thỏa dụng \(u^i(x^i)\) liên tục, đơn
điệu ngặt và tựa lõm, trong đó \(x^i \in \mathbb{R}^L_+\). Nguồn lực ban
đầu là \(\omega = \sum_i \omega^i\). Tập hợp các phân bổ khả thi được
xác định bởi:
\[ \sum_{i=1}^I x^i = \sum_{i=1}^I \omega^i + \sum_{j=1}^J y^j \] trong
đó \(y^j \in Y^j\) là vector sản xuất của doanh nghiệp \(j\), \(Y^j\) là
tập khả năng sản xuất.

\textbf{Định lý Cơ bản Thứ nhất của Kinh tế học Phúc lợi:} Nếu sở thích
là đơn điệu cục bộ và có hệ thống thị trường đầy đủ, mọi cân bằng cạnh
tranh Walras đều là tối ưu Pareto. \emph{Chứng minh (phác thảo):} Giả sử
cân bằng \((x^*, y^*, p^*)\) không là tối ưu Pareto. Khi đó tồn tại phân
bổ khả thi \((x, y)\) sao cho \(u^i(x^i) \geq u^i(x^{i*})\) với mọi
\(i\) và ngặt với ít nhất một \(i\). Sự tối đa hóa thỏa dụng ngụ ý
\(p^* \cdot x^i \geq p^* \cdot x^{i*}\), tổng lại ta có
\(p^* \cdot \sum x^i > p^* \cdot \sum x^{i*}\). Thay điều kiện khả thi
vào, ta nhận được mâu thuẫn với giả định tối đa hóa lợi nhuận của doanh
nghiệp tại giá \(p^*\).

\textbf{Định lý Cơ bản Thứ hai:} Bất kỳ phân bổ tối ưu Pareto nào cũng
có thể đạt được thông qua một cân bằng thị trường cạnh tranh, với điều
kiện có sự phân phối lại nguồn lực ban đầu bằng các khoản chuyển giao
gộp (lump-sum transfers), giả định rằng tập tiêu dùng và tập sản xuất là
lồi.

Việc tìm phân bổ Pareto có thể được hình thức hóa bằng bài toán quy
hoạch tối đa hóa hàm phúc lợi xã hội Bergson-Samuelson
\(W(u^1, \dots, u^I)\):
\[ \max_{\{x^i, y^j\}} W(u^1(x^1), \dots, u^I(x^I)) \quad \text{s.t.} \quad \sum x^i \leq \omega + \sum y^j \]
Điều kiện bậc nhất (FOCs) yêu cầu tỷ lệ thay thế biên (MRS) giữa hai
hàng hóa phải bằng nhau đối với tất cả người tiêu dùng và bằng tỷ lệ
chuyển đổi biên (MRT) trong sản xuất:
\[ MRS_{l,k}^i = MRS_{l,k}^{i'} = MRT_{l,k}^j \quad \forall i, i', j \]

\subsubsection{8.1.2 Điểm mù thường gặp (Cognitive Traps \&
Limitations)}\label{ux111iux1ec3m-muxf9-thux1b0ux1eddng-gux1eb7p-cognitive-traps-limitations-47}

\subsubsection{8.1.2 Điểm mù thường gặp (Cognitive Traps \&
Limitations)}\label{ux111iux1ec3m-muxf9-thux1b0ux1eddng-gux1eb7p-cognitive-traps-limitations-48}

\subsubsection{8.1.2 Điểm mù thường gặp (Cognitive Traps \&
Limitations)}\label{ux111iux1ec3m-muxf9-thux1b0ux1eddng-gux1eb7p-cognitive-traps-limitations-49}

Nhầm lẫn giữa Hiệu quả (Efficiency) và Công bằng (Equity). Định lý Phúc
lợi 1 chỉ đảm bảo thị trường đạt hiệu quả Pareto, nhưng một phân bổ mà
một người có tất cả tài sản cũng là hiệu quả Pareto. Không thể dùng định
lý này để ngụy biện rằng thị trường tự do luôn mang lại kết quả xã hội
tốt nhất.

\subsubsection{8.1.3 Ứng dụng thực tế (Practical
Applications)}\label{ux1ee9ng-dux1ee5ng-thux1ef1c-tux1ebf-practical-applications-25}

Các định lý này thiết lập khuôn khổ lý thuyết cho cuộc tranh luận giữa
kinh tế thị trường tự do và sự can thiệp của nhà nước. Ở Việt Nam, quá
trình Đổi mới năm 1986 là sự thừa nhận thực tế của Định lý 1: cơ chế thị
trường và giá cả tự do phân bổ nguồn lực hiệu quả hơn kế hoạch hóa tập
trung. Đồng thời, các chương trình xóa đói giảm nghèo hay trợ cấp cho
các vùng khó khăn là sự cố gắng hướng tới các mục tiêu phân phối theo
tinh thần của Định lý 2, dù phải đối mặt với bài toán đánh đổi giữa hiệu
quả và công bằng (equity-efficiency tradeoff).

\subsubsection{8.1.4 Ví dụ minh họa tính toán (Worked Numerical
Examples)}\label{vuxed-dux1ee5-minh-hux1ecda-tuxednh-touxe1n-worked-numerical-examples-24}

Giả sử nền kinh tế trao đổi thuần túy (Edgeworth box) có 2 cá nhân A, B
và 2 hàng hóa \(X, Y\). Nguồn lực ban đầu \(\omega^A = (8, 2)\) và
\(\omega^B = (2, 8)\). Hàm thỏa dụng:
\(u^A(x_A, y_A) = x_A^{0.5} y_A^{0.5}\),
\(u^B(x_B, y_B) = x_B^{0.5} y_B^{0.5}\). Tổng nguồn lực:
\(\bar{X} = 10, \bar{Y} = 10\). \textbf{Bước 1: Tìm đường hợp đồng
(Contract Curve)} Điều kiện tối ưu Pareto:
\(MRS^A_{x,y} = MRS^B_{x,y}\). Ta có:
\(MRS^A_{x,y} = \frac{\partial u^A / \partial x_A}{\partial u^A / \partial y_A} = \frac{y_A}{x_A}\).
Tương tự: \(MRS^B_{x,y} = \frac{y_B}{x_B} = \frac{10 - y_A}{10 - x_A}\).
Cân bằng hai MRS:
\(\frac{y_A}{x_A} = \frac{10 - y_A}{10 - x_A} \Rightarrow 10y_A - x_A y_A = 10x_A - x_A y_A \Rightarrow y_A = x_A\).
Đường hợp đồng là đường chéo \(y_A = x_A\) trong hộp Edgeworth.
\textbf{Bước 2: Tìm cân bằng cạnh tranh Walras} Chọn \(Y\) làm chuẩn
(numeraire) \(p_y = 1\), giá \(p_x = p\). Thu nhập của A:
\(I_A = 8p + 2\). Thu nhập của B: \(I_B = 2p + 8\). Hàm cầu của A đối
với X:
\(x_A = \frac{0.5 I_A}{p} = \frac{0.5(8p + 2)}{p} = 4 + \frac{1}{p}\).
Hàm cầu của B đối với X:
\(x_B = \frac{0.5 I_B}{p} = \frac{0.5(2p + 8)}{p} = 1 + \frac{4}{p}\).
Điều kiện thanh toán thị trường cho X:
\(x_A + x_B = 10 \Rightarrow 4 + \frac{1}{p} + 1 + \frac{4}{p} = 10 \Rightarrow 5 + \frac{5}{p} = 10 \Rightarrow p = 1\).
Vậy \(p_x = 1, p_y = 1\). Phân bổ cân bằng:
\(x_A = 5, y_A = 5, x_B = 5, y_B = 5\). Phân bổ này nằm trên đường hợp
đồng (\(5 = 5\)), minh họa chính xác Định lý Cơ bản Thứ nhất.

\begin{center}\rule{0.5\linewidth}{0.5pt}\end{center}

\subsection{8.2 Thuế Tối ưu \& Nguyên tắc Ramsey (Optimal Taxation \&
Ramsey
Rule)}\label{thuux1ebf-tux1ed1i-ux1b0u-nguyuxean-tux1eafc-ramsey-optimal-taxation-ramsey-rule}

\subsubsection{8.2.1 Cốt lõi Lý thuyết (Core
Theory)}\label{cux1ed1t-luxf5i-luxfd-thuyux1ebft-core-theory-25}

Lý thuyết thuế tối ưu xem xét cách thức chính phủ có thể thu một lượng
ngân sách \(\bar{R}\) với tổn thất vô ích (deadweight loss) nhỏ nhất.
Trong mô hình đánh thuế hàng hóa của Ramsey (1927), giả sử có \(K\) hàng
hóa với giá sản xuất không đổi \(q_k\). Chính phủ áp thuế \(t_k\), giá
tiêu dùng là \(p_k = q_k + t_k\). Bài toán tối ưu của chính phủ là tối
đa hóa thỏa dụng gián tiếp \(V(p_1, \dots, p_K, I)\) với điều kiện ràng
buộc ngân sách:
\[ \max_{\{t_k\}} V(q_1+t_1, \dots, q_K+t_K, I) \quad \text{s.t.} \quad \sum_{k=1}^K t_k x_k(p) \geq \bar{R} \]
Hàm Lagrange:
\(\mathcal{L} = V(p, I) + \lambda \left[ \sum t_k x_k(p) - \bar{R} \right]\).
Sử dụng hằng đẳng thức Roy
(\(\frac{\partial V}{\partial p_k} = -\alpha x_k\) với \(\alpha\) là
thỏa dụng biên của thu nhập), FOC đối với \(t_k\) là:
\[ -\alpha x_k + \lambda \left[ x_k + \sum_{j} t_j \frac{\partial x_j}{\partial p_k} \right] = 0 \]
Nếu giả định các hàm cầu chéo bằng 0
(\(\frac{\partial x_j}{\partial p_k} = 0\) với \(j \neq k\)), ta dẫn
xuất được \textbf{Quy tắc Độ co giãn nghịch đảo (Inverse Elasticity
Rule)}:
\[ \frac{t_k}{p_k} = \left( \frac{\lambda - \alpha}{\lambda} \right) \frac{1}{\varepsilon_k} \]
trong đó
\(\varepsilon_k = - \frac{\partial x_k}{\partial p_k} \frac{p_k}{x_k}\)
là độ co giãn của cầu theo giá. Thuế suất tỷ lệ nghịch với độ co giãn:
hàng hóa có cầu càng kém co giãn thì thuế suất càng cao.

Trong trường hợp tổng quát, với phương trình Slutsky, nguyên tắc Ramsey
tương đương với việc thuế làm giảm tỷ lệ phần trăm lượng tiêu dùng của
mọi hàng hóa đi cùng một mức (Quy tắc giảm cân xứng). Đối với thuế thu
nhập tối ưu, mô hình Mirrlees (1971) tối đa hóa hàm phúc lợi xã hội theo
ràng buộc tương thích khuyến khích (incentive compatibility), cân bằng
giữa việc tái phân phối và việc duy trì động lực làm việc.

\subsubsection{8.2.2 Điểm mù thường gặp (Cognitive Traps \&
Limitations)}\label{ux111iux1ec3m-muxf9-thux1b0ux1eddng-gux1eb7p-cognitive-traps-limitations-50}

\subsubsection{8.2.2 Điểm mù thường gặp (Cognitive Traps \&
Limitations)}\label{ux111iux1ec3m-muxf9-thux1b0ux1eddng-gux1eb7p-cognitive-traps-limitations-51}

\subsubsection{8.2.2 Điểm mù thường gặp (Cognitive Traps \&
Limitations)}\label{ux111iux1ec3m-muxf9-thux1b0ux1eddng-gux1eb7p-cognitive-traps-limitations-52}

\subsubsection{8.2.3 Ứng dụng thực tế (Practical
Applications)}\label{ux1ee9ng-dux1ee5ng-thux1ef1c-tux1ebf-practical-applications-26}

Tại Việt Nam, Luật Thuế Giá trị Gia tăng (VAT) chia thành nhiều mức thuế
suất (0\%, 5\%, 10\%). Các mặt hàng thiết yếu, y tế, giáo dục, hoặc hàng
hóa phục vụ sản xuất nông nghiệp thường chịu mức 0\% hoặc 5\%. Điều này
vi phạm nguyên tắc Ramsey thuần túy về hiệu quả (vì chúng kém co giãn),
nhưng lại hợp lý về mặt phân phối và công bằng xã hội. Các loại thuế
tiêu thụ đặc biệt (rượu bia, thuốc lá) áp dụng thuế suất rất cao phần
lớn là do ngoại tác tiêu cực, nhưng cũng vì cầu đối với chúng tương đối
không co giãn (phù hợp Ramsey).

\subsubsection{8.2.4 Ví dụ minh họa tính toán (Worked Numerical
Examples)}\label{vuxed-dux1ee5-minh-hux1ecda-tuxednh-touxe1n-worked-numerical-examples-25}

Giả sử có hai hàng hóa 1 và 2. Hàm cầu độc lập:
\(x_1(p_1) = 100 - p_1\), \(x_2(p_2) = 120 - 2p_2\). Giá cung hoàn toàn
co giãn (chi phí biên không đổi): \(q_1 = 20, q_2 = 20\). Chính phủ cần
huy động \(\bar{R} = 100\). Thuế đơn giá là \(t_1, t_2\), giá tiêu dùng
\(p_1 = 20+t_1, p_2 = 20+t_2\). Tổn thất vô ích (DWL) ở mỗi thị trường
là diện tích tam giác:
\(DWL_k = \frac{1}{2} t_k \Delta x_k = \frac{1}{2} t_k^2 |\frac{dx_k}{dp_k}|\).
Bài toán tối thiểu hóa tổng DWL:
\[ \min_{t_1, t_2} \left[ \frac{1}{2} t_1^2 (1) + \frac{1}{2} t_2^2 (2) \right] \quad \text{s.t.} \quad t_1(100 - (20+t_1)) + t_2(120 - (20+t_2)) = 100 \]
Ràng buộc: \(t_1(80 - t_1) + t_2(100 - 2t_2) = 100\). Giả định thuế suất
nhỏ, thu ngân sách xấp xỉ \(80t_1 + 100t_2 = 100\). Hàm Lagrange (tuyến
tính hóa ràng buộc để minh họa Ramsey rule):
\[ \mathcal{L} = 0.5 t_1^2 + t_2^2 - \lambda (80t_1 + 100t_2 - 100) \]
FOC: \(t_1 - 80\lambda = 0 \Rightarrow t_1 = 80\lambda\). Và
\(2t_2 - 100\lambda = 0 \Rightarrow t_2 = 50\lambda\). Tỷ lệ thuế tối
ưu: \(\frac{t_1}{t_2} = \frac{80}{50} = 1.6\). Độ co giãn tại điểm ban
đầu (\(p=20\)): \(\varepsilon_1 = 1 \cdot \frac{20}{80} = 0.25\).
\(\varepsilon_2 = 2 \cdot \frac{20}{100} = 0.4\). Kiểm tra quy tắc co
giãn nghịch đảo: \(\frac{t_1 / 20}{t_2 / 20} = \frac{t_1}{t_2} = 1.6\).
Trùng khớp với tỷ lệ độ co giãn nghịch đảo:
\(\frac{\varepsilon_2}{\varepsilon_1} = \frac{0.4}{0.25} = 1.6\). Chính
phủ áp thuế cao hơn (tương đối) vào hàng hóa 1 do cầu của nó kém co giãn
hơn.

\begin{center}\rule{0.5\linewidth}{0.5pt}\end{center}

\paragraph{8.2.5 Mô hình Mirrlees (1971) về Thuế Thu nhập Phi tuyến Tối
ưu (Nonlinear Optimal Income
Taxation)}\label{muxf4-huxecnh-mirrlees-1971-vux1ec1-thuux1ebf-thu-nhux1eadp-phi-tuyux1ebfn-tux1ed1i-ux1b0u-nonlinear-optimal-income-taxation}

\paragraph{8.2.5.1 Cốt lõi Lý thuyết (Core
Theory)}\label{cux1ed1t-luxf5i-luxfd-thuyux1ebft-core-theory-26}

Trong khi quy tắc Ramsey (1927) giải quyết bài toán thuế gián thu trên
hàng hóa, \textbf{Mô hình Mirrlees (1971)} đặt nền móng cho lý thuyết
thuế trực thu phi tuyến tối ưu trên thu nhập lao động. Bài toán của
chính phủ là thiết kế một biểu thuế thu nhập \(T(z)\), trong đó
\(z = w \cdot l\) là thu nhập trước thuế (tích của mức lương phản ánh
năng suất cá nhân \(w\) và số giờ lao động \(l\)), nhằm tối đa hóa hàm
phúc lợi xã hội phi tập trung (Social Welfare Function - SWF) dưới ràng
buộc thông tin bất cân xứng: chính phủ chỉ quan sát được thu nhập \(z\),
không quan sát được năng suất thực tế \(w\) hay nỗ lực \(l\).

Hàm phúc lợi xã hội có dạng tích phân theo phân phối năng suất \(F(w)\)
với hàm mật độ \(f(w)\) trên đoạn \([w_0, w_1]\):
\[\max_{c(w), z(w)} \int_{w_0}^{w_1} G(u(c(w), z(w)/w)) f(w) dw\] trong
đó \(c(w) = z(w) - T(z(w))\) là tiêu dùng (thu nhập sau thuế),
\(u(c, l)\) là hàm thỏa dụng cá nhân với \(u_c > 0, u_l < 0\), và
\(G(\cdot)\) biểu thị mức độ quan tâm của xã hội đối với sự công bằng
(nguyên tắc bàng quan hay Rawlsian cực đại hóa người nghèo nhất).

\textbf{Điều kiện Đơn điệu \& Khớp nối (Spence-Mirrlees Single Crossing
Condition):} Để cá nhân có năng suất cao không có động lực giả vờ làm
người có năng suất thấp (Ràng buộc Khuyến khích - Incentive
Compatibility Constraint \(IC\)), tỷ lệ thay thế biên giữa tiêu dùng và
thu nhập phải giảm theo năng suất:
\[-\frac{\partial}{\partial w} \left( \frac{u_l(c, z/w) / w}{u_c(c, z/w)} \right) < 0\]
Điều kiện Euler-Lagrange cho bài toán tối ưu hóa biến phân dẫn đến
\textbf{Công thức Saez (2001)} cho thuế suất biên tối ưu
\(\tau(z) = T'(z)\) tại mức thu nhập \(z\):
\[\frac{\tau(z)}{1 - \tau(z)} = \frac{1 - F(w)}{w f(w)} \cdot \frac{1 + \frac{1}{\varepsilon_u}}{1} \cdot \bar{g}(w)\]
trong đó: 1. \(\frac{1 - F(w)}{w f(w)}\) là tỷ lệ phân phối nghịch đảo
Hazard rate (độ dày của đuôi phân phối thu nhập). 2. \(\varepsilon_u\)
là độ co giãn cung lao động theo tiền lương bù đắp (Compensated
elasticity of labor supply). Khi cung lao động càng co giãn
(\(\varepsilon_u\) lớn), thuế suất biên tối ưu càng nhỏ để tránh tổn
thất vô ích. 3. \(\bar{g}(w)\) là trọng số phúc lợi xã hội trung bình
của những cá nhân có năng suất từ \(w\) trở lên. \textbf{Định lý Thuế
suất Biên bằng 0 ở đỉnh (Zero Marginal Tax Rate at the Top):} Nếu phân
phối năng suất có cận trên hữu hạn \(w_1\), thì tại đỉnh cao nhất,
\(1 - F(w_1) = 0 \implies \tau(z(w_1)) = 0\). Thuế suất biên bằng 0 tại
người giàu nhất ngăn chặn sự méo mó cung lao động mà không làm giảm
nguồn thu từ bất kỳ ai có thu nhập cao hơn.

\paragraph{8.2.5.2 Điểm mù thường gặp (Cognitive Traps \&
Limitations)}\label{ux111iux1ec3m-muxf9-thux1b0ux1eddng-gux1eb7p-cognitive-traps-limitations-53}

\begin{itemize}
\tightlist
\item
  \textbf{Điểm mù thứ nhất (Ảo tưởng về thuế suất biên cao tuyệt đối cho
  giới siêu giàu):} Nhiều nhà làm chính sách theo đuổi chủ nghĩa bình
  quân tin rằng áp đặt thuế suất biên 80-90\% lên nhóm 1\% người giàu
  nhất là cách tốt nhất để giảm bất bình đẳng. Tuy nhiên, theo công thức
  Mirrlees-Saez, nếu cung lao động hoặc hành vi trốn/tránh thuế có độ co
  giãn cao (\(\varepsilon_u > 0,5\)), việc tăng thuế vượt qua ngưỡng tối
  đa Laffer (\(\tau^* = \frac{1}{1 + a \cdot \varepsilon_u}\) với \(a\)
  là hệ số Pareto) sẽ phản tác dụng, làm giảm tổng thu ngân sách để tái
  phân phối cho người nghèo.
\item
  \textbf{Điểm mù thứ hai (Bỏ qua cấu trúc hộ gia đình và trốn thuế sang
  phần thưởng phi tiền tệ):} Mô hình Mirrlees giả định thu nhập \(z\) là
  biến duy nhất bị đánh thuế. Thực tế, khi thuế suất biên thu nhập cá
  nhân tăng cao, các giám đốc điều hành (CEO) và chuyên gia cao cấp lập
  tức chuyển dịch cấu trúc thù lao sang các hình thức phi tiền tệ (đặc
  quyền công ty, quyền chọn cổ phiếu, quỹ hưu trí phi thuế) hoặc chuyển
  giá sang công ty gia đình. Nếu không kiểm soát được độ co giãn thu
  nhập chịu thuế (Elasticity of Taxable Income - ETI), hệ thống thuế
  Mirrlees trên giấy sẽ thất bại hoàn toàn trong thực tiễn.
\end{itemize}

\paragraph{8.2.5.3 Ứng dụng thực tế (Practical
Applications)}\label{ux1ee9ng-dux1ee5ng-thux1ef1c-tux1ebf-practical-applications-27}

\begin{itemize}
\tightlist
\item
  \textbf{Thiết kế Biểu thuế Thu nhập Cá nhân (PIT) tại Việt Nam:} Hiện
  nay, biểu thuế lũy tiến từng phần của Việt Nam gồm 7 bậc (từ 5\% đến
  35\%). Khung thuế suất cao nhất 35\% áp dụng cho phần thu nhập trên 80
  triệu đồng/tháng. Vận dụng lý thuyết Mirrlees, Bộ Tài chính đang xem
  xét thu gọn biểu thuế xuống 5 bậc, đồng thời mở rộng khoảng cách giữa
  các bậc để giảm thiểu ``hiệu ứng nhảy bậc'' (bracket creep) gây ức chế
  động lực tăng năng suất của tầng lớp trung lưu chuyên nghiệp.
\item
  \textbf{Tín dụng Thuế Thu nhập Kiếm được (EITC - Earned Income Tax
  Credit):} Thay vì trợ cấp tiền mặt không điều kiện (gây bẫy nghèo do
  thuế suất biên ngầm lên tới 100\% khi đi làm lại), mô hình Mirrlees mở
  rộng cho thấy thuế suất biên âm (trợ cấp lương - negative income tax)
  cho nhóm lao động thu nhập thấp nhất là công cụ tối ưu để khuyến khích
  tham gia thị trường lao động (extensive margin).
\end{itemize}

\paragraph{8.2.5.4 Ví dụ minh họa tính toán (Worked Numerical
Examples)}\label{vuxed-dux1ee5-minh-hux1ecda-tuxednh-touxe1n-worked-numerical-examples-26}

Giả sử một nền kinh tế có hàm phân phối thu nhập tuân theo phân phối
Pareto ở phần đuôi cao với hệ số Pareto \(a = 2\). Độ co giãn của thu
nhập chịu thuế theo tỷ lệ giữ lại sau thuế (Elasticity of Taxable
Income) của nhóm siêu giàu được ước lượng thực nghiệm là
\(\varepsilon = 0,25\). Chính phủ muốn áp dụng mức thuế suất biên tối ưu
cho nhóm siêu giàu (thuế suất cao nhất tối đa hóa thu ngân sách từ nhóm
này để tái phân phối, tức \(\bar{g} = 0\)). 1. \textbf{Áp dụng công thức
Saez cho đuôi Pareto:} \[\tau^* = \frac{1}{1 + a \cdot \varepsilon}\] 2.
\textbf{Thay số tham số vào công thức:}
\[\tau^* = \frac{1}{1 + 2 \times 0,25} = \frac{1}{1 + 0,5} = \frac{1}{1,5} \approx 66,67\%\]
\textbf{Biện luận:} Mức thuế suất biên tối ưu để tối đa hóa thu ngân
sách từ nhóm có thu nhập cao nhất trong điều kiện này là \(66,67\%\).
Nếu độ co giãn trốn/tránh thuế tăng lên \(\varepsilon = 0,5\) (do hội
nhập tài chính toàn cầu dễ chuyển tiền ra nước ngoài), mức thuế tối ưu
sẽ giảm xuống chỉ còn \(\tau^* = \frac{1}{1 + 2 \times 0,5} = 50\%\).
Điều này chứng minh độ co giãn hành vi là thông số quyết định trần thuế
suất quốc gia.

\subsection{8.3 Cung cấp Hàng hóa Công Tối ưu \& Điều kiện Samuelson
(Public Goods Provision \& Samuelson
Condition)}\label{cung-cux1ea5p-huxe0ng-huxf3a-cuxf4ng-tux1ed1i-ux1b0u-ux111iux1ec1u-kiux1ec7n-samuelson-public-goods-provision-samuelson-condition}

\subsubsection{8.3.1 Cốt lõi Lý thuyết (Core
Theory)}\label{cux1ed1t-luxf5i-luxfd-thuyux1ebft-core-theory-27}

Hàng hóa công thuần túy (pure public goods) được đặc trưng bởi hai tính
chất: không có tính cạnh tranh (non-rivalry) và không có tính loại trừ
(non-excludability). Giả sử nền kinh tế có một hàng hóa tư \(x\) và một
hàng hóa công \(G\). Hàm thỏa dụng của cá nhân \(i\) là \(u^i(x^i, G)\).
Ràng buộc nguồn lực của nền kinh tế là:
\[ \sum_{i=1}^I x^i + c(G) = \omega \] trong đó \(c(G)\) là hàm chi phí
sản xuất hàng hóa công, \(\omega\) là tổng của cải. Để tìm phân bổ tối
ưu Pareto, ta tối đa hóa thỏa dụng của cá nhân 1, với ràng buộc giữ
nguyên mức thỏa dụng của các cá nhân khác ở mức \(\bar{u}^i\):
\[ \max_{x^i, G} u^1(x^1, G) \quad \text{s.t.} \quad u^i(x^i, G) \geq \bar{u}^i \ (\forall i \neq 1), \quad \sum x^i + c(G) = \omega \]
Hàm Lagrange:
\[ \mathcal{L} = u^1(x^1, G) + \sum_{i=2}^I \mu_i [u^i(x^i, G) - \bar{u}^i] + \lambda [\omega - \sum x^i - c(G)] \]
Điều kiện bậc nhất: Đối với \(x^1\):
\(\frac{\partial u^1}{\partial x^1} - \lambda = 0 \Rightarrow \lambda = \frac{\partial u^1}{\partial x^1}\)
Đối với \(x^i (i \neq 1)\):
\(\mu_i \frac{\partial u^i}{\partial x^i} - \lambda = 0 \Rightarrow \mu_i = \frac{\lambda}{\partial u^i / \partial x^i}\)
Đối với \(G\):
\(\frac{\partial u^1}{\partial G} + \sum_{i=2}^I \mu_i \frac{\partial u^i}{\partial G} - \lambda c'(G) = 0\)
Chia phương trình \(G\) cho \(\lambda\) và thay các biểu thức \(\mu_i\),
ta nhận được \textbf{Điều kiện Samuelson}:
\[ \sum_{i=1}^I \frac{\partial u^i / \partial G}{\partial u^i / \partial x^i} = c'(G) \iff \sum_{i=1}^I MRS_{G,x}^i = MRT_{G,x} \]
Ý nghĩa: Hàng hóa công nên được cung cấp cho đến khi tổng giá trị sẵn
lòng trả biên (tổng các MRS) bằng với chi phí biên để cung cấp thêm một
đơn vị (MRT).

\subsubsection{8.3.2 Điểm mù thường gặp (Cognitive Traps \&
Limitations)}\label{ux111iux1ec3m-muxf9-thux1b0ux1eddng-gux1eb7p-cognitive-traps-limitations-54}

\subsubsection{8.2.2 Điểm mù thường gặp (Cognitive Traps \&
Limitations)}\label{ux111iux1ec3m-muxf9-thux1b0ux1eddng-gux1eb7p-cognitive-traps-limitations-55}

\subsubsection{8.2.2 Điểm mù thường gặp (Cognitive Traps \&
Limitations)}\label{ux111iux1ec3m-muxf9-thux1b0ux1eddng-gux1eb7p-cognitive-traps-limitations-56}

\subsubsection{8.3.3 Ứng dụng thực tế (Practical
Applications)}\label{ux1ee9ng-dux1ee5ng-thux1ef1c-tux1ebf-practical-applications-28}

Xây dựng quốc phòng, hải đăng, hoặc một cây cầu công cộng là các ví dụ
điển hình. Tại các đô thị lớn ở Việt Nam, việc đóng góp xây dựng camera
an ninh khu phố hay hệ thống chiếu sáng ngõ hẻm thường đối mặt với bài
toán free-rider: một số hộ gia đình từ chối đóng góp nhưng vẫn hưởng
lợi. Do đó, việc cung cấp các hàng hóa này đòi hỏi cưỡng chế thông qua
thu thuế hoặc phí tổ dân phố thay vì đóng góp tự nguyện, nhằm ép buộc
nền kinh tế di chuyển đến điểm tối ưu Samuelson.

\subsubsection{8.3.4 Ví dụ minh họa tính toán (Worked Numerical
Examples)}\label{vuxed-dux1ee5-minh-hux1ecda-tuxednh-touxe1n-worked-numerical-examples-27}

Một cộng đồng có 3 người. Hàm thỏa dụng của người \(i\) (\(i=1,2,3\))
là: \(u^i(x^i, G) = x^i + \alpha_i \ln(G)\) với
\(\alpha_1 = 10, \alpha_2 = 20, \alpha_3 = 30\). Ngân sách ban đầu của
mỗi người là \(\omega_i = 100\). Chi phí sản xuất hàng hóa công là
\(c(G) = 5G\). \textbf{Bước 1: Tìm mức cung cấp tối ưu xã hội (Samuelson
Condition)} Tính
\(MRS^i_{G,x} = \frac{\partial u^i / \partial G}{\partial u^i / \partial x^i} = \frac{\alpha_i / G}{1} = \frac{\alpha_i}{G}\).
Điều kiện Samuelson: \(\sum_{i=1}^3 MRS^i_{G,x} = c'(G)\)
\(\frac{10}{G} + \frac{20}{G} + \frac{30}{G} = 5 \Rightarrow \frac{60}{G} = 5 \Rightarrow G^* = 12\).
\textbf{Bước 2: Tìm mức cung cấp tự nguyện (Cân bằng Nash)} Giả sử đóng
góp của mỗi người là \(g_i\), \(G = g_1 + g_2 + g_3\). Ràng buộc:
\(x^i + 5g_i = 100\). Người \(i\) tối đa hóa
\(100 - 5g_i + \alpha_i \ln(g_i + G_{-i})\). FOC cho \(g_i > 0\):
\(\frac{\alpha_i}{G} = 5 \Rightarrow G = \frac{\alpha_i}{5}\). Vì
\(\alpha_3 = 30\) là lớn nhất, người 3 muốn cung cấp \(G = 6\). Tại mức
này, người 1 và 2 có \(MRS < 5\), nên họ đóng góp \(g_1 = 0, g_2 = 0\).
Kết quả: \(G_{Nash} = 6\). \textbf{Diễn giải:} Cung cấp tự nguyện
(\(G=6\)) thấp hơn nhiều so với mức tối ưu Pareto (\(G=12\)) do hiện
tượng kẻ đi xe không vé. Người 1 và 2 tận hưởng miễn phí sự đóng góp của
người 3.

\begin{center}\rule{0.5\linewidth}{0.5pt}\end{center}

\subsection{8.4 Định lý Arrow \& Bỏ phiếu (Arrow's Impossibility Theorem
\& Voting
Theory)}\label{ux111ux1ecbnh-luxfd-arrow-bux1ecf-phiux1ebfu-arrows-impossibility-theorem-voting-theory}

\subsubsection{8.4.1 Cốt lõi Lý thuyết (Core
Theory)}\label{cux1ed1t-luxf5i-luxfd-thuyux1ebft-core-theory-28}

Lý thuyết Lựa chọn Xã hội (Social Choice Theory) nghiên cứu cách thức
tổng hợp các sở thích cá nhân thành một thứ tự ưu tiên chung của xã hội.
Gọi \(X\) là tập các lựa chọn thay thế. Mỗi cá nhân \(i\) có một quan hệ
ưu tiên (preference relation) \(\succeq_i\) trên \(X\) (đầy đủ và bắc
cầu). Hàm phúc lợi xã hội (SWF) \(F\) ánh xạ các hồ sơ sở thích
\(( \succeq_1, \dots, \succeq_I )\) thành một thứ tự sở thích xã hội
\(\succeq_S\).

\textbf{Định lý Bất khả thi của Arrow (1951):} Nếu \(|X| \geq 3\), không
tồn tại một hàm SWF nào thỏa mãn đồng thời 4 tiên đề sau đây (ngoại trừ
chế độ độc tài): 1. \textbf{Miền xác định phổ quát (Universal Domain -
U):} SWF phải xác định với mọi cấu hình sở thích cá nhân có thể có. 2.
\textbf{Nguyên lý Pareto yếu (Weak Pareto - P):} Nếu mọi cá nhân đều
thích \(x\) hơn \(y\) (\(x \succ_i y \ \forall i\)), thì xã hội cũng
phải thích \(x\) hơn \(y\) (\(x \succ_S y\)). 3. \textbf{Tính độc lập
của các lựa chọn không liên quan (Independence of Irrelevant
Alternatives - IIA):} Sở thích xã hội giữa \(x\) và \(y\) chỉ phụ thuộc
vào thứ tự xếp hạng của các cá nhân giữa \(x\) và \(y\), không phụ thuộc
vào một giải pháp \(z\) nào khác. 4. \textbf{Không độc tài
(Non-dictatorship - D):} Không tồn tại cá nhân \(h\) sao cho
\(x \succ_h y \implies x \succ_S y\) bất chấp sở thích của những người
khác.

Khi miền sở thích được thu hẹp, chẳng hạn sở thích đỉnh đơn
(single-peaked preferences) trên không gian một chiều, \textbf{Định lý
Cử tri Trung vị (Median Voter Theorem - MVT)} của Duncan Black cho biết:
Quyết định theo đa số sẽ dẫn đến kết cục là điểm lý tưởng của cử tri
trung vị, và đây là một điểm cân bằng Condorcet ổn định.

\subsubsection{8.4.2 Điểm mù thường gặp (Cognitive Traps \&
Limitations)}\label{ux111iux1ec3m-muxf9-thux1b0ux1eddng-gux1eb7p-cognitive-traps-limitations-57}

\subsubsection{8.4.2 Điểm mù thường gặp (Cognitive Traps \&
Limitations)}\label{ux111iux1ec3m-muxf9-thux1b0ux1eddng-gux1eb7p-cognitive-traps-limitations-58}

\subsubsection{8.4.2 Điểm mù thường gặp (Cognitive Traps \&
Limitations)}\label{ux111iux1ec3m-muxf9-thux1b0ux1eddng-gux1eb7p-cognitive-traps-limitations-59}

Cố gắng tìm kiếm một hệ thống bầu cử hoàn hảo. Định lý Arrow đã chứng
minh toán học rằng không có hệ thống bầu cử nào là hoàn hảo (thỏa mãn cả
4 tiên đề). Mọi hệ thống đều là sự đánh đổi (trade-off) và có thể bị
thao túng chiến lược.

\subsubsection{8.4.3 Ứng dụng thực tế (Practical
Applications)}\label{ux1ee9ng-dux1ee5ng-thux1ef1c-tux1ebf-practical-applications-29}

Trong các cuộc bỏ phiếu tại Quốc hội hoặc hội đồng cổ đông doanh nghiệp,
kết quả ngân sách hay đầu tư phụ thuộc rất lớn vào người chủ tọa sắp xếp
thứ tự biểu quyết. Theo Định lý Cử tri trung vị, trong các cuộc bầu cử
tổng thống Mỹ, hai đảng lớn có xu hướng điều chỉnh chính sách tiến về
phía trung dung để thu hút lá phiếu của cử tri dao động (cử tri trung
vị), làm cho các cương lĩnh tranh cử trở nên ít cực đoan hơn so với kỳ
bầu cử sơ bộ nội bộ đảng.

\subsubsection{8.4.4 Ví dụ minh họa tính toán (Worked Numerical
Examples)}\label{vuxed-dux1ee5-minh-hux1ecda-tuxednh-touxe1n-worked-numerical-examples-28}

\textbf{Minh họa Nghịch lý Condorcet \& Bầu cử.} Có 3 giải pháp chính
sách \(A, B, C\) và 3 nhóm cử tri với số lượng bằng nhau: - Nhóm 1:
\(A \succ_1 B \succ_1 C\) - Nhóm 2: \(B \succ_2 C \succ_2 A\) - Nhóm 3:
\(C \succ_3 A \succ_3 B\) \textbf{Biểu quyết cặp đôi (Majority rule
pairwise voting):} - Trận 1 (A vs B): Nhóm 1 và 3 thích A hơn B
\(\Rightarrow\) A thắng. - Trận 2 (B vs C): Nhóm 1 và 2 thích B hơn C
\(\Rightarrow\) B thắng. - Trận 3 (C vs A): Nhóm 2 và 3 thích C hơn A
\(\Rightarrow\) C thắng. Xã hội rơi vào chu trình:
\(A \succ_S B \succ_S C \succ_S A\). Không có hệ thống nhất quán. Thất
bại Tiên đề IIA và tính bắc cầu. \textbf{Áp dụng Định lý Cử tri trung
vị:} Giả sử chính sách là mức chi tiêu công \(G \in [0, 100]\). Điểm lý
tưởng (peak) của 3 người là: \(G_1 = 20, G_2 = 50, G_3 = 80\). Sở thích
của họ giảm dần khi \(G\) rời xa điểm lý tưởng. Sắp xếp: \(20, 50, 80\).
Cử tri trung vị là người thứ 2 với \(G_2 = 50\). Nếu đưa ra bỏ phiếu
giữa 50 và bất kỳ mức nào khác (vd 60). Với 60, người 3 ủng hộ 60 (gần
80 hơn), nhưng người 1 và 2 sẽ ủng hộ 50 (vì 50 gần điểm lý tưởng của họ
hơn). Do đó 50 thắng. \(G=50\) là kẻ chiến thắng Condorcet bất khả chiến
bại.

\begin{center}\rule{0.5\linewidth}{0.5pt}\end{center}

\subsection{8.5 Ngoại tác \& Giải pháp (Externalities Revisited:
Pigouvian vs
Coasean)}\label{ngoux1ea1i-tuxe1c-giux1ea3i-phuxe1p-externalities-revisited-pigouvian-vs-coasean}

\subsubsection{8.5.1 Cốt lõi Lý thuyết (Core
Theory)}\label{cux1ed1t-luxf5i-luxfd-thuyux1ebft-core-theory-29}

Ngoại tác xảy ra khi hành động của một chủ thể ảnh hưởng trực tiếp đến
hàm thỏa dụng hoặc hàm sản xuất của chủ thể khác mà không thông qua cơ
chế giá thị trường. Ký hiệu nhà máy hóa chất (sản xuất thép \(s\)) xả
thải gây hại cho ngư dân. Lợi nhuận nhà máy:
\(\pi_S(s) = p_s s - c_s(s)\). Tổn thất của ngư dân: \(D(s)\) với
\(D'(s) > 0\). - \textbf{Giải pháp Pigou:} Tối ưu xã hội đòi hỏi
\(\max_s \pi_S(s) - D(s)\). FOC: \(p_s - c_s'(s) - D'(s) = 0\). Tuy
nhiên, thị trường tự do sẽ chọn mức \(s\) tại \(p_s = c_s'(s)\). Để khắc
phục, chính phủ áp thuế Pigou \(t = D'(s^*)\) tính trên mỗi đơn vị xả
thải tại điểm tối ưu. Nhà máy tối đa hóa
\(p_s s - c_s(s) - ts \Rightarrow p_s - c_s'(s) - t = 0\), khôi phục
hiệu quả Pareto. - \textbf{Định lý Coase:} Nếu quyền sở hữu tài sản
(property rights) được xác định rõ ràng và chi phí giao dịch
(transaction costs) bằng 0, quá trình thương lượng tự do giữa các bên sẽ
dẫn đến mức sản xuất (hay ô nhiễm) tối ưu Pareto, bất kể quyền sở hữu
ban đầu được trao cho ai. \emph{Chứng minh (phác thảo):} Nếu nhà máy có
quyền ô nhiễm, ngư dân sẽ trả tiền hối lộ \(B\) để nhà máy giảm \(s\).
Tổng thặng dư: \((\pi_S(s) + B) - (D(s) + B) = \pi_S(s) - D(s)\). Tối đa
hóa tổng thặng dư đạt được \(s^*\). Nếu ngư dân có quyền nước sạch, nhà
máy sẽ trả tiền bồi thường \(C\) để được xả thải, tổng thặng dư vẫn là
\(\pi_S(s) - D(s)\). Mức ô nhiễm hiệu quả \(s^*\) luôn đạt được độc lập
với việc phân bổ quyền tài sản.

\subsubsection{8.5.2 Điểm mù thường gặp (Cognitive Traps \&
Limitations)}\label{ux111iux1ec3m-muxf9-thux1b0ux1eddng-gux1eb7p-cognitive-traps-limitations-60}

\subsubsection{8.5.2 Điểm mù thường gặp (Cognitive Traps \&
Limitations)}\label{ux111iux1ec3m-muxf9-thux1b0ux1eddng-gux1eb7p-cognitive-traps-limitations-61}

\subsubsection{8.5.2 Điểm mù thường gặp (Cognitive Traps \&
Limitations)}\label{ux111iux1ec3m-muxf9-thux1b0ux1eddng-gux1eb7p-cognitive-traps-limitations-62}

Định lý Coase bị hiểu lầm là thị trường luôn tự giải quyết được ngoại
tác. Điều kiện tối quan trọng là chi phí giao dịch (Transaction costs)
phải bằng không và quyền sở hữu được phân định rõ. Trong thực tế (như
biến đổi khí hậu), chi phí giao dịch là cực kỳ lớn và các bên đàm phán
rất đông.

\subsubsection{8.5.3 Ứng dụng thực tế (Practical
Applications)}\label{ux1ee9ng-dux1ee5ng-thux1ef1c-tux1ebf-practical-applications-30}

Thuế bảo vệ môi trường đối với xăng dầu tại Việt Nam là một nỗ lực áp
dụng lý thuyết Thuế Pigou để nội hóa chi phí biên xã hội (khí thải, hao
mòn đường sá). Ngược lại, Hệ thống mua bán phát thải châu Âu (EU ETS) là
sự kết hợp giữa thiết lập tổng lượng (cap) và trao quyền tài sản có thể
mua bán (trade), một ứng dụng rõ ràng của hướng tiếp cận thị trường kiểu
Coase để tối thiểu hóa tổng chi phí giảm thải đối với các nhà máy trên
toàn châu lục.

\subsubsection{8.5.4 Ví dụ minh họa tính toán (Worked Numerical
Examples)}\label{vuxed-dux1ee5-minh-hux1ecda-tuxednh-touxe1n-worked-numerical-examples-29}

Nhà máy A xả thải \(E\). Lợi ích biên từ việc xả thải (Marginal
Benefit): \(MB = 200 - 2E\). Chi phí biên gây ra cho xã hội (Marginal
Damage): \(MD = E\). \textbf{Bước 1: Tìm cân bằng tự do và Tối ưu xã
hội} Tự do: Nhà máy phát thải đến khi
\(MB = 0 \Rightarrow 200 - 2E = 0 \Rightarrow E_0 = 100\). Tối ưu xã
hội:
\(MB = MD \Rightarrow 200 - 2E = E \Rightarrow 3E = 200 \Rightarrow E^* = 66.67\).
\textbf{Bước 2: Tính Thuế Pigou} Tại mức tối ưu \(E^* = 66.67\),
\(MD(66.67) = 66.67\). Thuế Pigou \(t = 66.67\) trên mỗi đơn vị E. Nhà
máy đối diện bài toán
\(MB = t \Rightarrow 200 - 2E = 66.67 \Rightarrow E = 66.67\).
\textbf{Bước 3: Thương lượng Coase} Giả sử chi phí giao dịch bằng 0. Nạn
nhân có quyền cấm xả thải (\(E_{ban đầu} = 0\)). Nhà máy muốn tăng \(E\)
từ 0. Tại \(E=0\), \(MB=200 > MD=0\). Nhà máy có thể hối lộ nạn nhân một
khoản tiền \(P\) sao cho \(MD < P < MB\) để được xả 1 đơn vị. Thương
lượng tiếp tục mở rộng cho đến khi \(MB = MD \Rightarrow E^* = 66.67\).
Tại đây, nhà máy không sẵn sàng trả mức giá nạn nhân yêu cầu nữa. Cân
bằng đạt được hiệu quả Pareto.

\begin{center}\rule{0.5\linewidth}{0.5pt}\end{center}

\subsection{8.6 Phân tích Chi phí-Lợi ích \& Tiêu chí Kaldor-Hicks
(Cost-Benefit Analysis \&
Kaldor-Hicks)}\label{phuxe2n-tuxedch-chi-phuxed-lux1ee3i-uxedch-tiuxeau-chuxed-kaldor-hicks-cost-benefit-analysis-kaldor-hicks}

\subsubsection{8.6.1 Cốt lõi Lý thuyết (Core
Theory)}\label{cux1ed1t-luxf5i-luxfd-thuyux1ebft-core-theory-30}

Khi đánh giá một dự án công, chính sách hiếm khi thỏa mãn tối ưu Pareto
chặt chẽ (vì luôn có người chịu thiệt). Phân tích chi phí - lợi ích
(CBA) dựa trên \textbf{Tiêu chí bồi thường Kaldor-Hicks}: Một chính sách
được xem là nâng cao phúc lợi nếu những người được lợi có thể (về mặt
giả thuyết) bồi thường đầy đủ cho những người bị thiệt, và vẫn còn dư
lợi ích. Nói cách khác, tổng thặng dư kinh tế tăng lên.

Để đo lường sự thay đổi phúc lợi của người tiêu dùng do sự thay đổi giá
hoặc cung cấp hàng hóa công, ta sử dụng \textbf{Hàm chi tiêu
(Expenditure Function)} \(e(p, u)\) và hai khái niệm: 1. \textbf{Biến
thiên Tương đương (Equivalent Variation - EV):} Số tiền mà một người sẵn
sàng trả (hoặc nhận) TRƯỚC khi thay đổi để tránh (hoặc đạt được) sự thay
đổi đó. \[ EV = e(p_0, u_1) - e(p_0, u_0) = e(p_0, u_1) - I \] 2.
\textbf{Biến thiên Bù đắp (Compensating Variation - CV):} Số tiền cần
bồi thường SAU khi thay đổi xảy ra để cá nhân đó trở về mức thỏa dụng
ban đầu. \[ CV = e(p_1, u_1) - e(p_1, u_0) = I - e(p_1, u_0) \] Tiêu chí
Kaldor yêu cầu tổng \(\sum CV_i > 0\). Tiêu chí Hicks yêu cầu tổng
\(\sum EV_i > 0\). Bằng cách tích phân đường cong cầu bồi thường
(Hicksian demand) \(\frac{\partial e}{\partial p_k} = h_k(p, u)\), ta có
thể tính chính xác tổn thất vô ích và thặng dư tiêu dùng. Giá trị hiện
tại ròng (NPV) của dự án cộng dồn các luồng EV/CV được chiết khấu theo
tỷ lệ chiết khấu xã hội.

\subsubsection{8.6.2 Điểm mù thường gặp (Cognitive Traps \&
Limitations)}\label{ux111iux1ec3m-muxf9-thux1b0ux1eddng-gux1eb7p-cognitive-traps-limitations-63}

\subsubsection{8.6.2 Điểm mù thường gặp (Cognitive Traps \&
Limitations)}\label{ux111iux1ec3m-muxf9-thux1b0ux1eddng-gux1eb7p-cognitive-traps-limitations-64}

\subsubsection{8.6.2 Điểm mù thường gặp (Cognitive Traps \&
Limitations)}\label{ux111iux1ec3m-muxf9-thux1b0ux1eddng-gux1eb7p-cognitive-traps-limitations-65}

Áp dụng Tiêu chí Kaldor-Hicks ngây thơ. Tiêu chí này cho phép một dự án
được duyệt nếu người được lợi CÓ THỂ đền bù cho người bị thiệt (tiềm
năng). Điểm mù là quên mất rằng nếu sự đền bù đó không thực sự diễn ra
trong thực tế, dự án sẽ làm tăng bất bình đẳng nghiêm trọng.

\subsubsection{8.6.3 Ứng dụng thực tế (Practical
Applications)}\label{ux1ee9ng-dux1ee5ng-thux1ef1c-tux1ebf-practical-applications-31}

Trong việc thẩm định các dự án đường sắt đô thị (Metro) hoặc Cảng hàng
không quốc tế Long Thành, phương pháp CBA được áp dụng nghiêm ngặt. Lợi
ích không chỉ là doanh thu bán vé, mà phải tính toán EV/CV của việc tiết
kiệm thời gian di chuyển, giảm tai nạn giao thông, giảm lượng khí thải
carbon. Việc xác định ``tỷ lệ chiết khấu xã hội'' đóng vai trò then
chốt: một tỷ lệ chiết khấu quá cao sẽ làm các dự án môi trường dài hạn
(trồng rừng, chống biến đổi khí hậu) bị đánh giá thấp và loại bỏ ngay
lập tức.

\subsubsection{8.6.4 Ví dụ minh họa tính toán (Worked Numerical
Examples)}\label{vuxed-dux1ee5-minh-hux1ecda-tuxednh-touxe1n-worked-numerical-examples-30}

Một người tiêu dùng có hàm thỏa dụng \(u(x, y) = x^{0.5} y^{0.5}\). Thu
nhập \(I = 100\). Giá ban đầu: \(p_{x0} = 1, p_{y0} = 1\). Giá sau khi
có thuế/chính sách: \(p_{x1} = 2, p_{y1} = 1\). \textbf{Bước 1: Hàm cầu
Marshall và Thỏa dụng} \(x = \frac{0.5 I}{p_x}, y = \frac{0.5 I}{p_y}\).
Thỏa dụng gián tiếp:
\(V(p_x, p_y, I) = \left(\frac{0.5 I}{p_x}\right)^{0.5} \left(\frac{0.5 I}{p_y}\right)^{0.5} = \frac{0.5 I}{\sqrt{p_x p_y}}\).
Hàm chi tiêu: \(e(p_x, p_y, u) = 2 u \sqrt{p_x p_y}\). Mức thỏa dụng ban
đầu: \(u_0 = \frac{0.5 \times 100}{\sqrt{1 \times 1}} = 50\). Mức thỏa
dụng sau chính sách:
\(u_1 = \frac{0.5 \times 100}{\sqrt{2 \times 1}} = 35.35\). \textbf{Bước
2: Tính Biến thiên Tương đương (EV) và Biến thiên Bù đắp (CV)}
\(EV = e(p_{x0}, p_{y0}, u_1) - e(p_{x0}, p_{y0}, u_0) = e(1, 1, 35.35) - 100\)
\(EV = 2 \times 35.35 \times \sqrt{1 \times 1} - 100 = 70.7 - 100 = -29.3\).
(Người này sẵn sàng trả 29.3 đơn vị tiền để ngăn không cho giá tăng lên
2).
\(CV = e(p_{x1}, p_{y1}, u_1) - e(p_{x1}, p_{y1}, u_0) = 100 - e(2, 1, 50)\)
\(CV = 100 - (2 \times 50 \times \sqrt{2 \times 1}) = 100 - 141.4 = -41.4\).
(Người này cần được bồi thường 41.4 đơn vị tiền sau khi giá đã tăng để
khôi phục mức sống \(u_0=50\)). \textbf{Diễn giải:} Tổn thất tính bằng
tiền khác nhau giữa EV và CV do hiệu ứng thu nhập. Để kiểm tra
Kaldor-Hicks cho dự án, ta sẽ tính tổng EV hoặc CV của mọi cá nhân; nếu
\(\sum CV > 0\), dự án tạo ra lợi ích ròng xã hội.

\begin{center}\rule{0.5\linewidth}{0.5pt}\end{center}

\subsection{8.7 Phân quyền Tài khóa \& Cạnh tranh giữa các Chính quyền
(Fiscal Federalism \&
Tiebout)}\label{phuxe2n-quyux1ec1n-tuxe0i-khuxf3a-cux1ea1nh-tranh-giux1eefa-cuxe1c-chuxednh-quyux1ec1n-fiscal-federalism-tiebout}

\subsubsection{8.7.1 Cốt lõi Lý thuyết (Core
Theory)}\label{cux1ed1t-luxf5i-luxfd-thuyux1ebft-core-theory-31}

Phân quyền tài khóa nghiên cứu sự phân bổ trách nhiệm thu-chi giữa các
cấp chính quyền. \textbf{Định lý Phân quyền của Oates (Decentralization
Theorem)} phát biểu: trong điều kiện không có tính kinh tế nhờ quy mô
(economies of scale) và ngoại tác tràn (spillovers) giữa các địa phương,
việc cung cấp hàng hóa công cục bộ (được điều chỉnh theo sở thích đặc
thù của từng địa phương) luôn mang lại hiệu quả Pareto bằng hoặc cao hơn
so với việc trung ương cung cấp một mức hàng hóa công đồng nhất cho mọi
vùng.

\textbf{Mô hình Tiebout (1956) ``Voting with one's feet'':} Tiebout giải
quyết bài toán free-rider của Samuelson ở cấp độ địa phương bằng cách mô
hình hóa các địa phương như các ``doanh nghiệp'' cạnh tranh cung cấp các
gói ``thuế - hàng hóa công''. Giả sử có một số lượng lớn các cộng đồng
với các mức độ cung cấp hàng hóa công \(G_j\) và thuế \(t_j\). Cư dân có
tính cơ động hoàn hảo (perfect mobility). Cá nhân \(i\) sẽ chọn định cư
tại cộng đồng \(j^*\) nhằm:
\[ \max_{j} u^i(x, G_j) \quad \text{s.t.} \quad x + t_j \leq I^i \] Sự
di cư liên tục sẽ tạo ra sự phân loại (sorting) các cá nhân thành các
cộng đồng đồng nhất về mặt sở thích. Cân bằng Nash không gian này dẫn
đến mức cung cấp hàng hóa công hiệu quả Pareto cục bộ, loại bỏ được vấn
đề giấu diếm thông tin (do người dân thể hiện sự ưu tiên bằng cách ``bỏ
phiếu bằng chân'').

\subsubsection{8.7.2 Điểm mù thường gặp (Cognitive Traps \&
Limitations)}\label{ux111iux1ec3m-muxf9-thux1b0ux1eddng-gux1eb7p-cognitive-traps-limitations-66}

\subsubsection{8.7.2 Điểm mù thường gặp (Cognitive Traps \&
Limitations)}\label{ux111iux1ec3m-muxf9-thux1b0ux1eddng-gux1eb7p-cognitive-traps-limitations-67}

\subsubsection{8.7.2 Điểm mù thường gặp (Cognitive Traps \&
Limitations)}\label{ux111iux1ec3m-muxf9-thux1b0ux1eddng-gux1eb7p-cognitive-traps-limitations-68}

Cạnh tranh giữa các địa phương luôn tốt. Mô hình Tiebout (1956) cho rằng
người dân sẽ ``bầu chọn bằng chân'' (Voting with their feet). Nhưng điều
này giả định tính cơ động hoàn hảo (perfect mobility). Thực tế người
nghèo không dễ di cư, dẫn đến sự phân hóa giàu-nghèo giữa các tiểu bang
(Race to the bottom).

\subsubsection{8.7.3 Ứng dụng thực tế (Practical
Applications)}\label{ux1ee9ng-dux1ee5ng-thux1ef1c-tux1ebf-practical-applications-32}

Tại Việt Nam, Luật Ngân sách Nhà nước phân định rõ nguồn thu ngân sách
trung ương, địa phương, và nguồn thu phân chia. Sự cạnh tranh giữa các
tỉnh thành (thông qua chỉ số PCI - Năng lực cạnh tranh cấp tỉnh) để thu
hút FDI và dân cư tay nghề cao chính là biểu hiện của cơ chế Tiebout.
Tuy nhiên, do chênh lệch phát triển quá lớn, trung ương luôn phải thực
hiện điều tiết thu (ví dụ, điều tiết ngân sách từ TP.HCM, Hà Nội, Bình
Dương về trung ương) để hỗ trợ các tỉnh miền núi khó khăn, thể hiện sự
can thiệp chống lại hệ quả phân hóa tiêu cực của mô hình Tiebout phân
quyền thuần túy.

\subsubsection{8.7.4 Ví dụ minh họa tính toán (Worked Numerical
Examples)}\label{vuxed-dux1ee5-minh-hux1ecda-tuxednh-touxe1n-worked-numerical-examples-31}

Có hai loại người trong xã hội: Người Yêu Giáo Dục (E) muốn mức chi tiêu
giáo dục cao \(G_E = 100\), sẵn sàng trả mức thuế cao \(T_E = 100\).
Người Ít Nhu Cầu (L) muốn \(G_L = 50\), sẵn sàng trả \(T_L = 50\). Có
hai địa phương A và B. Ban đầu, mọi người sống hỗn hợp. \textbf{Trường
hợp 1: Chính quyền Trung ương (Định lý Oates)} Trung ương cung cấp một
gói đồng nhất trung bình: \(G_C = 75\), đánh thuế bình quân
\(T_C = 75\). Người E không hài lòng vì thiếu giáo dục (mất \(25\) đơn
vị thỏa dụng biên tiềm năng). Người L không hài lòng vì bị ép đóng thuế
quá cao cho dịch vụ họ không cần. Hệ thống không tối ưu. \textbf{Trường
hợp 2: Cơ chế Tiebout (Phân quyền cạnh tranh)} Địa phương A công bố gói
\((G_A = 100, T_A = 100)\). Địa phương B công bố gói
\((G_B = 50, T_B = 50)\). Nhóm E lập tức ``bỏ phiếu bằng chân'', chuyển
sang sống tại địa phương A. Nhóm L di cư sang địa phương B. \textbf{Kết
quả cân bằng:} Địa phương A chỉ có người E, địa phương B chỉ có người L
(cân bằng tách biệt - separating equilibrium). Cả hai nhóm đều được thỏa
mãn điểm lý tưởng của mình (tối đa hóa thặng dư), đạt hiệu quả Pareto
theo đúng lý thuyết Phân quyền Tài khóa. Mọi người đều đạt mức thỏa dụng
cao hơn so với việc trung ương cung cấp hàng hóa công đồng nhất.

\begin{center}\rule{0.5\linewidth}{0.5pt}\end{center}

\section{CHƯƠNG 9: KINH TẾ HỌC LAO ĐỘNG \& VỐN NHÂN LỰC (Labor Economics
\& Human
Capital)}\label{chux1b0ux1a1ng-9-kinh-tux1ebf-hux1ecdc-lao-ux111ux1ed9ng-vux1ed1n-nhuxe2n-lux1ef1c-labor-economics-human-capital}

Như đã thiết lập trong các chương thuộc Phần A, thị trường lao động là
một thị trường yếu tố sản xuất đặc biệt, nơi ``hàng hóa'' được giao dịch
gắn liền với con người, hàm chứa những đặc điểm phi chuẩn như nỗ lực bất
đối xứng, ma sát tìm kiếm và bất đồng nhất về kỹ năng. Chương này là cầu
nối quan trọng hướng tới Phần B (Các chủ đề Kinh tế vĩ mô và vi mô nâng
cao), cung cấp khung phân tích về quyết định cung cầu lao động, đầu tư
vốn nhân lực, và các khiếm khuyết thị trường như thông tin bất cân xứng
và phân biệt đối xử.

\begin{center}\rule{0.5\linewidth}{0.5pt}\end{center}

\subsection{9.1 Cung-Cầu Lao động \& Cân bằng Thị trường (Labor Supply
\&
Demand)}\label{cung-cux1ea7u-lao-ux111ux1ed9ng-cuxe2n-bux1eb1ng-thux1ecb-trux1b0ux1eddng-labor-supply-demand}

\subsubsection{9.1.1 Cốt lõi Lý thuyết (Core
Theory)}\label{cux1ed1t-luxf5i-luxfd-thuyux1ebft-core-theory-32}

Phân tích cung lao động cá nhân bắt nguồn từ bài toán tối đa hóa thỏa
dụng giữa tiêu dùng (\(C\)) và thời gian rảnh rỗi (\(L\)). Giả sử cá
nhân có tổng quỹ thời gian \(T\), cung lao động \(h = T - L\). Bài toán
tối ưu được thiết lập như sau: \[ \max_{C, L} U(C, L) \] với điều kiện
ngân sách: \[ pC = w(T - L) + V = wh + V \] Trong đó \(w\) là tiền lương
danh nghĩa, \(p\) là mức giá (chuẩn hóa \(p=1\)), và \(V\) là thu nhập
ngoài lao động.

Điều kiện bậc nhất (FOC) cho tối ưu nội tại (interior solution) đòi hỏi
tỷ lệ thay thế biên (MRS) giữa thời gian rảnh rỗi và tiêu dùng phải bằng
mức lương thực tế:
\[ MRS_{L,C} = \frac{\partial U / \partial L}{\partial U / \partial C} = w \]

\textbf{Phương trình Slutsky} phân rã tác động của một sự thay đổi trong
tiền lương lên cung lao động thành tác động thay thế và tác động thu
nhập:
\[ \frac{\partial h}{\partial w} = \left( \frac{\partial h}{\partial w} \right)_{\bar{U}} + h \frac{\partial h}{\partial V} \]
- \textbf{Tác động thay thế (Substitution Effect - SE)}:
\(\left( \frac{\partial h}{\partial w} \right)_{\bar{U}} > 0\). Khi
lương tăng, chi phí cơ hội của thời gian rảnh rỗi đắt hơn, cá nhân thay
thế \(L\) bằng \(C\), làm tăng \(h\). - \textbf{Tác động thu nhập
(Income Effect - IE)}: \(h \frac{\partial h}{\partial V} < 0\) (giả sử
\(L\) là hàng hóa thông thường). Lương tăng làm tăng thu nhập thực tế,
cá nhân muốn tiêu dùng nhiều \(L\) hơn, làm giảm \(h\).

Nếu \(|IE| > |SE|\) ở mức lương đủ cao, đường cung lao động sẽ uốn cong
về phía sau (backward-bending labor supply curve).

Về phía \textbf{cầu lao động}, hãng cạnh tranh hoàn hảo tối đa hóa lợi
nhuận \(\Pi = pF(E, K) - wE - rK\), với \(E\) là số lao động. FOC cho
\(E\) là: \[ p \frac{\partial F}{\partial E} = w \implies MRPL = w \]
Giá trị sản phẩm cận biên của lao động (\(MRPL\)) bằng mức lương. Cân
bằng thị trường đạt được khi tổng cung bằng tổng cầu.

\subsubsection{9.1.2 Điểm mù thường gặp (Cognitive Traps \&
Limitations)}\label{ux111iux1ec3m-muxf9-thux1b0ux1eddng-gux1eb7p-cognitive-traps-limitations-69}

\textbf{Nhầm lẫn giữa ``thời gian'' và ``nỗ lực''.} Phân tích tĩnh cơ
bản coi thời gian làm việc \(h\) tương đương với nỗ lực. Tuy nhiên, trên
thực tế, người lao động có thể có mặt tại nơi làm việc nhưng không nỗ
lực (shirking), dẫn đến những mô hình phức tạp hơn (xem Phần 9.4). Hơn
nữa, việc giả định cung lao động có thể điều chỉnh liên tục bỏ qua sự
cứng nhắc về giờ làm việc (ví dụ: ca làm việc cố định 40 giờ/tuần) do
chi phí cố định từ phía doanh nghiệp.

\subsubsection{9.1.3 Ứng dụng thực tế (Practical
Applications)}\label{ux1ee9ng-dux1ee5ng-thux1ef1c-tux1ebf-practical-applications-33}

\textbf{Hiện tượng giảm giờ làm khi lương tăng ở giới chuyên gia} Xét
một bác sĩ chuyên khoa có hàm thỏa dụng
\(U(C,L) = C^{\alpha} L^{1-\alpha}\). Giả sử \(T=100\), \(V=0\). Giải
bài toán, cung lao động là \(h = \alpha T = 100\alpha\), không phụ thuộc
vào \(w\) do SE và IE triệt tiêu hoàn toàn. Tuy nhiên, nếu hàm thỏa dụng
là dạng Stone-Geary: \(U = \ln(C - \bar{C}) + \beta \ln L\), với
\(\bar{C}\) là mức tiêu dùng sinh tồn. Giải FOC:
\(\frac{L}{C-\bar{C}} = \frac{\beta}{w}\). Thay vào \(C = w(T-L)\), ta
được \(L = \frac{\beta}{1+\beta} \left( T - \frac{\bar{C}}{w} \right)\).
Cung lao động:
\(h = T - L = \frac{1}{1+\beta} T + \frac{\beta}{1+\beta} \frac{\bar{C}}{w}\).
Rõ ràng \(\frac{\partial h}{\partial w} < 0\). Khi mức lương của bác sĩ
lên rất cao, họ ưu tiên nghỉ ngơi, chơi golf, và giảm số giờ trực, minh
họa đường cung lao động uốn cong về phía sau.

\begin{center}\rule{0.5\linewidth}{0.5pt}\end{center}

\subsection{9.2 Lý thuyết Vốn Nhân lực (Human Capital
Theory)}\label{luxfd-thuyux1ebft-vux1ed1n-nhuxe2n-lux1ef1c-human-capital-theory}

\subsubsection{9.2.1 Cốt lõi Lý thuyết (Core
Theory)}\label{cux1ed1t-luxf5i-luxfd-thuyux1ebft-core-theory-33}

Lý thuyết vốn nhân lực do Gary Becker (1964) tiên phong mô hình hóa giáo
dục và đào tạo như một khoản đầu tư: chấp nhận chi phí hiện tại (học
phí, thu nhập bị mất) để đổi lấy dòng thu nhập cao hơn trong tương lai.

Giả sử một cá nhân quyết định có nên học thêm \(s\) năm học (đi học) so
với đi làm ngay (không đi học). Giá trị hiện tại (Present Value - PV)
của dòng thu nhập nếu đi làm ngay là:
\[ PV_0 = \int_0^{65} Y_0 e^{-rt} dt \] Nếu đi học \(s\) năm, với chi
phí trực tiếp là \(C(t)\) mỗi năm, và đi làm từ năm \(s\) đến 65 với thu
nhập \(Y_s\):
\[ PV_s = \int_0^s -C(t) e^{-rt} dt + \int_s^{65} Y_s e^{-rt} dt \] Cá
nhân sẽ chọn đi học nếu \(PV_s \geq PV_0\). \textbf{Tỷ suất sinh lời nội
bộ (Internal Rate of Return - IRR)} là mức chiết khấu \(r^*\) làm cho
\(PV_s = PV_0\).

Dạng thực nghiệm phổ biến nhất là \textbf{Phương trình Mincer} (Jacob
Mincer, 1974):
\[ \ln(w) = \beta_0 + \rho s + \beta_1 x + \beta_2 x^2 + \varepsilon \]
Trong đó \(w\) là tiền lương, \(s\) là số năm đi học, \(x\) là kinh
nghiệm làm việc (\(x = tuổi - s - 6\)). Hệ số \(\rho\) ước lượng tỷ suất
sinh lời cận biên của một năm đi học bổ sung.

Đào tạo tại chỗ (On-the-job training - OJT) chia làm 2 loại: 1.
\textbf{Đào tạo chung (General training):} Nâng cao năng suất ở mọi công
ty (ví dụ: kỹ năng Excel). Người lao động phải chịu toàn bộ chi phí
thông qua mức lương thấp hơn trong kỳ đào tạo vì họ có thể mang kỹ năng
này đi nơi khác. 2. \textbf{Đào tạo đặc thù (Firm-specific training):}
Chỉ nâng cao năng suất tại công ty hiện tại (ví dụ: phần mềm nội bộ).
Chi phí và lợi ích thường được chia sẻ giữa công ty và người lao động,
tạo ra cơ chế gắn kết (lock-in effect).

\subsubsection{9.2.2 Điểm mù thường gặp (Cognitive Traps \&
Limitations)}\label{ux111iux1ec3m-muxf9-thux1b0ux1eddng-gux1eb7p-cognitive-traps-limitations-70}

\textbf{Tranh luận Phát tín hiệu (Signaling) vs.~Vốn nhân lực.} Liệu đại
học thực sự dạy kỹ năng (Becker), hay chỉ là một bộ lọc khắc nghiệt để
những người vốn đã thông minh chứng minh năng lực của họ với nhà tuyển
dụng (Michael Spence, 1973)? Nếu giáo dục chỉ là phát tín hiệu, thì tỷ
suất sinh lời cá nhân vẫn cao (người có bằng cấp lương cao hơn), nhưng
tỷ suất sinh lời xã hội có thể bằng 0 (chỉ tốn kém tài nguyên để phân
loại mà không tạo thêm năng suất). Sự nhầm lẫn giữa tương quan (học cao,
lương cao) và nhân quả (học đại học làm tăng năng suất) là lỗi phổ biến
nhất.

\subsubsection{9.2.3 Ứng dụng thực tế (Practical
Applications)}\label{ux1ee9ng-dux1ee5ng-thux1ef1c-tux1ebf-practical-applications-34}

\textbf{Thiết kế chính sách trợ cấp giáo dục} Giả sử có 2 loại lao động:
Năng lực cao (H) chiếm 20\% và Năng lực thấp (L) chiếm 80\%. Năng suất
của H là \(w_H = 100\), của L là \(w_L = 50\). Chi phí nỗ lực đi học
\(s\) năm đối với H là \(C_H(s) = 5s\), đối với L là \(C_L(s) = 15s\).
Nhà tuyển dụng không quan sát được năng lực, nhưng quan sát được \(s\).
Họ đặt ngưỡng: nếu \(s \geq s^*\), trả lương \(w_H\); nếu \(s < s^*\),
trả \(w_L\). Để tín hiệu tách biệt (Separating Equilibrium) xảy ra, L
không muốn giả danh H:
\(100 - 15s^* < 50 \implies 15s^* > 50 \implies s^* > 3.33\) Và H muốn
phát tín hiệu: \(100 - 5s^* > 50 \implies 5s^* < 50 \implies s^* < 10\)
Cân bằng hiệu quả nhất là \(s^* = 3.34\). Nếu chính phủ trợ cấp học phí
làm giảm chi phí cho mọi người (ví dụ giảm \(5s\) cho tất cả), thì
\(C_L\) thành \(10s\). Lúc này L sẽ học lên mức \(s\) cao hơn để giả H,
buộc H phải học lên \(s^*\) cao hơn (ví dụ \(s^* > 5\)) để phân biệt.
Trợ cấp giáo dục cào bằng, trong bối cảnh signaling, có thể gây ra lãng
phí xã hội lớn do ``chạy đua vũ trang'' bằng cấp.

\begin{center}\rule{0.5\linewidth}{0.5pt}\end{center}

\subsection{9.3 Mô hình Tìm kiếm \& Ghép cặp (Search \& Matching
Models)}\label{muxf4-huxecnh-tuxecm-kiux1ebfm-ghuxe9p-cux1eb7p-search-matching-models}

\subsubsection{9.3.1 Cốt lõi Lý thuyết (Core
Theory)}\label{cux1ed1t-luxf5i-luxfd-thuyux1ebft-core-theory-34}

Thị trường lao động không khớp ngay lập tức. Người lao động thất nghiệp
mất thời gian tìm việc, công ty mất thời gian tuyển dụng.

\textbf{Mô hình Tìm kiếm Tuần tự của McCall (1970):} Một người thất
nghiệp (U) nhận được trợ cấp \(b\) mỗi kỳ. Mỗi kỳ, họ bốc ngẫu nhiên một
lời đề nghị tiền lương \(w \sim F(w)\). Nếu nhận việc, họ đi làm mãi mãi
với lương \(w\). Nếu từ chối, họ tiếp tục thất nghiệp. Giá trị hiện tại
(Bellman equation) của trạng thái thất nghiệp (\(V_U\)) và làm việc
(\(V_E(w)\)) là: \[ rV_E(w) = w \]
\[ rV_U = b + \lambda \int_0^\infty \max[V_E(w) - V_U, 0] dF(w) \] Người
lao động sẽ chấp nhận công việc nếu \(V_E(w) \geq V_U\), dẫn đến quy tắc
\textbf{Mức lương bảo lưu (Reservation Wage, \(w_R\))}: \[ w_R = rV_U \]
Thay vào phương trình \(V_U\):
\[ w_R = b + \frac{\lambda}{r} \int_{w_R}^\infty (w - w_R) dF(w) \] Tăng
trợ cấp \(b\) làm tăng \(w_R\), dẫn đến người lao động kén chọn hơn và
thời gian thất nghiệp kéo dài hơn.

\textbf{Mô hình Ghép cặp Mortensen-Pissarides (DMP):} Kinh tế vĩ mô mô
hình hóa cả hai phía thông qua hàm ghép cặp (Matching Function):
\(M(U, V)\), với \(U\) là số người thất nghiệp, \(V\) là số vị trí trống
(vacancies). \(M\) có lợi tức không đổi theo quy mô (CRS). Định nghĩa độ
thắt chặt thị trường (market tightness) \(\theta = V / U\). Xác suất
người thất nghiệp tìm được việc:
\(p(\theta) = \frac{M(U, V)}{U} = M(1, \theta)\). Rõ ràng
\(p'(\theta) > 0\). Tiền lương trong mô hình này được quyết định thông
qua \textbf{Đàm phán Nash (Nash Bargaining)}, chia sẻ thặng dư của việc
ghép cặp thành công giữa công ty và người lao động:
\[ w = \beta p + (1-\beta) rV_U \] với \(\beta\) là sức mạnh đàm phán
của người lao động, \(p\) là năng suất biên. Đường cong Beveridge mô tả
mối quan hệ nghịch biến giữa \(U\) và \(V\) trong trạng thái dừng.

\subsubsection{9.3.2 Điểm mù thường gặp (Cognitive Traps \&
Limitations)}\label{ux111iux1ec3m-muxf9-thux1b0ux1eddng-gux1eb7p-cognitive-traps-limitations-71}

\textbf{Thất nghiệp tự nguyện hay không tự nguyện?} Trong mô hình McCall
cổ điển, thất nghiệp thực chất là ``tự nguyện'' hoặc là một ``hoạt động
đầu tư'' - cá nhân từ chối công việc lương thấp để chờ lương cao. Điều
này mâu thuẫn với nhận thức thông thường rằng người thất nghiệp thực sự
không thể tìm được việc làm. Mô hình DMP khắc phục một phần bằng cách
đưa vào ma sát ở cấp độ vĩ mô, nhưng vẫn ngầm định sự cân bằng tối ưu.

\subsubsection{9.3.3 Ứng dụng thực tế (Practical
Applications)}\label{ux1ee9ng-dux1ee5ng-thux1ef1c-tux1ebf-practical-applications-35}

\textbf{Thiết kế Bảo hiểm Thất nghiệp (UI)} Giả sử \(\lambda = 0.1\),
\(r=0.05\). Phân phối tiền lương đều \(F(w)\) trên \([0, 100]\). Phương
trình mức lương bảo lưu:
\(w_R = b + \frac{0.1}{0.05} \int_{w_R}^{100} (w - w_R) \frac{1}{100} dw = b + \frac{2}{100} \left[ \frac{w^2}{2} - w_R w \right]_{w_R}^{100}\)
\(w_R = b + 0.01 (10000 - 200w_R + w_R^2) = b + 100 - 2w_R + 0.01 w_R^2\)
Xét \(b=0\), giải phương trình
\(3w_R - 0.01 w_R^2 = 100 \implies w_R \approx 38.2\). Nếu UI tăng lên
\(b=20\): \(3w_R - 0.01 w_R^2 = 120 \implies w_R \approx 47.2\). Việc
tăng UI (\(b\)) đẩy mức lương tối thiểu mà người lao động chấp nhận lên
cao, làm giảm tỷ lệ chấp nhận việc làm \(\alpha = 1 - F(w_R)\), kéo dài
chu kỳ thất nghiệp. Tuy nhiên, nó cũng giúp người lao động tìm được việc
làm phù hợp với năng suất hơn.

\begin{center}\rule{0.5\linewidth}{0.5pt}\end{center}

\subsection{9.4 Tiền lương Hiệu quả (Efficiency Wage
Theory)}\label{tiux1ec1n-lux1b0ux1a1ng-hiux1ec7u-quux1ea3-efficiency-wage-theory}

\subsubsection{9.4.1 Cốt lõi Lý thuyết (Core
Theory)}\label{cux1ed1t-luxf5i-luxfd-thuyux1ebft-core-theory-35}

Lý thuyết tiền lương hiệu quả giải thích tại sao các hãng sẵn sàng trả
cao hơn mức lương cân bằng thị trường cạnh tranh. Có hai nhánh chính: 1.
\textbf{Mô hình Trốn việc (Shirking Model) của Shapiro \& Stiglitz
(1984):} Người lao động có thể trốn việc vì nỗ lực \(e\) gây mất thỏa
dụng. Doanh nghiệp giám sát với xác suất \(q\). 2. \textbf{Mô hình Trao
đổi quà tặng (Gift Exchange) của Akerlof (1982):} Trả lương cao tạo cảm
giác công bằng, người lao động đáp lại bằng nỗ lực cao (norms of
fairness).

\textbf{Thiết lập mô hình Shapiro-Stiglitz:} Giả sử cá nhân tối đa hóa
\(U = w - e\), với \(e>0\) nếu làm việc và \(e=0\) nếu trốn việc. Tỷ lệ
mất việc ngẫu nhiên là \(b\). Nếu trốn việc, khả năng bị sa thải là
\(b + q\). Giá trị của việc đi làm và nỗ lực (\(V_N\)):
\[ rV_N = (w - e) + b(V_U - V_N) \] Giá trị của việc trốn việc
(\(V_S\)): \[ rV_S = w + (b + q)(V_U - V_S) \] Để người lao động không
trốn việc, cần \textbf{Điều kiện không trốn việc (No-Shirking Condition
- NSC)}: \(V_N \geq V_S\). Giải ra, ta được:
\[ w \geq rV_U + e + \frac{(r+b+q)e}{q} \] Mức lương phải lớn hơn giá
trị bên ngoài (\(rV_U\)) cộng với chi phí nỗ lực (\(e\)) và một ``phần
bù trốn việc'' tỷ lệ nghịch với xác suất bị bắt (\(q\)). Nếu \(q\) thấp
(khó giám sát), tiền lương phải rất cao. Để \(w\) này duy trì như một
mối đe dọa, thất nghiệp \(U\) phải lớn hơn 0 (vì nếu \(U=0\), tìm việc
mới ngay tức khắc, \(V_U = V_E\), lời đe dọa sa thải vô nghĩa). Do đó,
thất nghiệp là tất yếu trong cân bằng.

\subsubsection{9.4.2 Điểm mù thường gặp (Cognitive Traps \&
Limitations)}\label{ux111iux1ec3m-muxf9-thux1b0ux1eddng-gux1eb7p-cognitive-traps-limitations-72}

\textbf{Thất nghiệp theo cơ cấu vs.~Thất nghiệp hiệu quả.} Tiền lương
hiệu quả ngụ ý thất nghiệp là ``không tự nguyện'' (involuntary) - có
những người sẵn sàng làm việc với mức lương hiện hành nhưng hãng không
thuê vì thuê họ sẽ làm giảm mức lương trung bình, vi phạm NSC, dẫn đến
trốn việc lan tràn. Điều này khác biệt sâu sắc với thất nghiệp tự nguyện
do ma sát trong phần 9.3.

\subsubsection{9.4.3 Ứng dụng thực tế (Practical
Applications)}\label{ux1ee9ng-dux1ee5ng-thux1ef1c-tux1ebf-practical-applications-36}

\textbf{Chương trình 5 đô la/ngày của Henry Ford (1914)} Năm 1914, tỷ lệ
luân chuyển nhân công (turnover) tại Ford Motor là 370\%, và tình trạng
vắng mặt (absenteeism) rất cao. Ford đã tăng lương gấp đôi lên \$5/ngày
(rất cao so với thị trường). Theo mô hình shirking, việc này làm tăng
\((w - rV_U)\), khiến giá trị của việc giữ được công việc ở Ford tăng
vọt. Kết quả, nỗ lực dây chuyền được đảm bảo, tỷ lệ nghỉ việc giảm xuống
16\%, và lợi nhuận tăng. Tương tự, ngày nay ở Thung lũng Silicon, các kỹ
sư phần mềm (nơi \(q\) giám sát rất khó do tính chất công việc sáng tạo)
thường được trả cao hơn nhiều so với mặt bằng chung để đảm bảo ``nỗ lực
nội tại'' (intrinsic motivation) và giảm rò rỉ chất xám.

\begin{center}\rule{0.5\linewidth}{0.5pt}\end{center}

\subsection{9.5 Phân biệt Đối xử trên Thị trường Lao động (Labor Market
Discrimination)}\label{phuxe2n-biux1ec7t-ux111ux1ed1i-xux1eed-truxean-thux1ecb-trux1b0ux1eddng-lao-ux111ux1ed9ng-labor-market-discrimination}

\subsubsection{9.5.1 Cốt lõi Lý thuyết (Core
Theory)}\label{cux1ed1t-luxf5i-luxfd-thuyux1ebft-core-theory-36}

Phân biệt đối xử kinh tế xảy ra khi những cá nhân có cùng năng suất lại
được trả lương khác nhau dựa trên các đặc điểm nhân khẩu học (giới tính,
chủng tộc).

\textbf{1. Phân biệt đối xử theo sở thích (Taste-based Discrimination) -
Gary Becker (1957):} Chủ doanh nghiệp có thành kiến (prejudice) không
thích nhóm thiểu số (B) so với nhóm đa số (W). Hàm thỏa dụng của nhà
tuyển dụng khi thuê nhóm B với lương \(w_B\) là:
\[ U = \Pi = pF(E_W + E_B) - w_W E_W - w_B E_B - d \cdot E_B \] Trong đó
\(d > 0\) là \textbf{Hệ số phân biệt đối xử (Discrimination
coefficient)}. Chi phí thực tế để thuê nhóm B trong mắt nhà tuyển dụng
là \(w_B + d\). Họ sẽ chỉ thuê nhóm B nếu \(w_B + d \leq w_W\), tức là
nhóm B phải chấp nhận mức lương thấp hơn (\(w_B = w_W - d\)) để bù đắp
cho ``sự khó chịu'' của ông chủ. Theo Becker, thị trường cạnh tranh dài
hạn sẽ trừng phạt sự phân biệt đối xử, vì các công ty không phân biệt
(d=0) sẽ thuê B với lương rẻ, có chi phí thấp hơn và sẽ đánh bật các
công ty có d\textgreater0 ra khỏi thị trường.

\textbf{2. Phân biệt đối xử thống kê (Statistical Discrimination) -
Arrow (1973), Phelps (1972):} Không có ác cảm cá nhân. Xảy ra do thông
tin bất cân xứng. Nhà tuyển dụng không quan sát được năng suất thực
\(q_i\) của cá nhân, chỉ đo lường được tín hiệu nhiễu
\(\theta_i = q_i + \epsilon_i\), và biết nhóm của họ (W hoặc B). Nếu kỳ
vọng năng suất nhóm trung bình là \(\bar{q}_W > \bar{q}_B\), hoặc phương
sai tín hiệu của nhóm B lớn hơn, hãng sẽ sử dụng lý thuyết Bayes để ước
lượng năng suất:
\[ E[q_i | \theta_i, \text{Group}] = \alpha \theta_i + (1-\alpha) \bar{q}_{\text{Group}} \]
Một cá nhân nhóm B dù có tín hiệu tốt cũng bị ``kéo xuống'' bởi mức
trung bình nhóm \(\bar{q}_B\) thấp hơn.

\subsubsection{9.5.2 Điểm mù thường gặp (Cognitive Traps \&
Limitations)}\label{ux111iux1ec3m-muxf9-thux1b0ux1eddng-gux1eb7p-cognitive-traps-limitations-73}

\textbf{Phân biệt thống kê không mang ý nghĩa biện minh.} Một sai lầm
kinh tế học phổ biến là cho rằng ``phân biệt thống kê'' là hợp lý về mặt
kinh tế (rational) đối với hãng, nên nó ``không sao''. Thực tế, sự cân
bằng phân biệt thống kê có thể tự hoàn thành (self-fulfilling prophecy):
nếu nhóm B biết họ luôn bị đánh giá thấp qua điểm trung bình nhóm, họ sẽ
không đầu tư vào vốn nhân lực, khiến \(\bar{q}_B\) thực sự thấp đi ở thế
hệ sau, tạo ra bẫy nghèo đói dai dẳng.

\subsubsection{9.5.3 Ứng dụng thực tế (Practical
Applications)}\label{ux1ee9ng-dux1ee5ng-thux1ef1c-tux1ebf-practical-applications-37}

\textbf{Phân rã Oaxaca-Blinder cho khoảng cách tiền lương giới} Giả sử
ta hồi quy phương trình Mincer cho nam (M) và nữ (F):
\(\ln w_M = \beta_M X_M + \varepsilon_M\)
\(\ln w_F = \beta_F X_F + \varepsilon_F\) Khoảng cách lương trung bình
\(\overline{\ln w_M} - \overline{\ln w_F}\) được phân rã thành:
\[ \overline{\ln w_M} - \overline{\ln w_F} = (\bar{X}_M - \bar{X}_F)\hat{\beta}_M + \bar{X}_F(\hat{\beta}_M - \hat{\beta}_F) \]
- Phần 1: \textbf{Tác động được giải thích (Endowment effect)}: Khác
biệt do đặc điểm (kinh nghiệm, trình độ). - Phần 2: \textbf{Tác động
không giải thích được (Discrimination/Unexplained effect)}: Khác biệt do
phần lợi suất (giới tính). Nếu phụ nữ bị trả thấp hơn cho cùng một đặc
điểm (ví dụ \(\beta_{kinh\_nghiệm}\) của nữ thấp hơn nam), phần này đo
lường sự phân biệt đối xử.

\begin{center}\rule{0.5\linewidth}{0.5pt}\end{center}

\subsection{9.6 Kinh tế học Di cư (Economics of
Migration)}\label{kinh-tux1ebf-hux1ecdc-di-cux1b0-economics-of-migration}

\subsubsection{9.6.1 Cốt lõi Lý thuyết (Core
Theory)}\label{cux1ed1t-luxf5i-luxfd-thuyux1ebft-core-theory-37}

Quyết định di cư là một dạng đầu tư vốn nhân lực đặc biệt.

\textbf{Mô hình Tự lựa chọn của Roy (Roy Model - 1951), áp dụng bởi
Borjas (1987):} Một công nhân có phân phối kỹ năng
\(s \sim N(\mu_s, \sigma_s^2)\). Thu nhập ở nước gốc (0) và nước đích
(1): \(\ln w_0 = \mu_0 + r_0 s\) \(\ln w_1 = \mu_1 + r_1 s - C\) Với
\(C\) là chi phí di cư. Cá nhân di cư nếu
\(\ln w_1 > \ln w_0 \implies (\mu_1 - \mu_0 - C) + (r_1 - r_0)s > 0\). -
\textbf{Lựa chọn dương (Positive Selection):} Nếu \(r_1 > r_0\) (nước
đích bất bình đẳng hơn, trả lợi suất cao cho kỹ năng), những người di cư
sẽ ở đầu trên của phân phối kỹ năng. - \textbf{Lựa chọn âm (Negative
Selection):} Nếu \(r_1 < r_0\) (nước đích có mạng lưới an sinh xã hội
tốt, phân phối thu nhập nén), những người có kỹ năng thấp lại có động
lực di cư lớn nhất.

\textbf{Thặng dư Nhập cư (Immigration Surplus):} Theo phân tích cung cầu
tĩnh cổ điển, dòng người nhập cư làm tăng tổng cung lao động. Điều này
làm giảm lương nội địa, nhưng làm tăng lợi nhuận của chủ sở hữu vốn. Lợi
ích ròng cho nền kinh tế bản địa (tam giác thặng dư) lớn hơn 0.
\[ \text{Surplus} = \frac{1}{2} \cdot \Delta w \cdot \Delta E \]

\subsubsection{9.6.2 Điểm mù thường gặp (Cognitive Traps \&
Limitations)}\label{ux111iux1ec3m-muxf9-thux1b0ux1eddng-gux1eb7p-cognitive-traps-limitations-74}

\textbf{Ảo tưởng ``Cướp việc làm'' (Lump of Labor Fallacy).} Quan điểm
dân túy cho rằng số lượng việc làm trong nền kinh tế là cố định, một
người nhập cư vào sẽ đẩy một người bản địa ra. Về mặt kinh tế học, nhập
cư không chỉ tăng Cung lao động mà còn tăng Cầu hàng hóa (và do đó Cầu
lao động phái sinh). Tác động thực lên tiền lương thường rất nhỏ hoặc
dương đối với các kỹ năng bổ cả (complementary skills).

\subsubsection{9.6.3 Ứng dụng thực tế (Practical
Applications)}\label{ux1ee9ng-dux1ee5ng-thux1ef1c-tux1ebf-practical-applications-38}

\textbf{Xuất khẩu lao động Việt Nam và Hiện tượng ``Chảy máu chất xám''}
Việt Nam xuất khẩu mạnh lao động phổ thông sang Nhật, Hàn, Đài Loan.
Theo mô hình Roy, do khoảng cách tiền lương cơ bản lớn
(\(\mu_1 \gg \mu_0\)) và chi phí \(C\) có thể tiếp cận được qua vay
mượn, dòng di cư này mang lại lượng kiều hối khổng lồ. Mặt khác, đối với
lao động trình độ cao (visa H-1B sang Mỹ), Mỹ có \(r_1\) (lợi suất cho
kỹ năng IT/STEM) rất cao so với Việt Nam. Điều này tạo ra \emph{Positive
Selection} rõ nét (chảy máu chất xám). Tuy nhiên, một số nghiên cứu gần
đây về ``Brain Gain'' cho rằng triển vọng được đi nước ngoài làm tăng
động lực đầu tư vào giáo dục đại học của toàn bộ dân số, và vì chỉ một
phần đi được, phần chất lượng cao ở lại giúp tăng năng suất trong nước.

\begin{center}\rule{0.5\linewidth}{0.5pt}\end{center}

\subsection{9.7 Bất bình đẳng Tiền lương \& Thay đổi Công nghệ (Wage
Inequality \& Skill-Biased Technical
Change)}\label{bux1ea5t-buxecnh-ux111ux1eb3ng-tiux1ec1n-lux1b0ux1a1ng-thay-ux111ux1ed5i-cuxf4ng-nghux1ec7-wage-inequality-skill-biased-technical-change}

\subsubsection{9.7.1 Cốt lõi Lý thuyết (Core
Theory)}\label{cux1ed1t-luxf5i-luxfd-thuyux1ebft-core-theory-38}

Từ thập niên 1980, phần bù lương đại học (college premium) tăng mạnh ở
Mỹ và toàn cầu, bất chấp nguồn cung sinh viên đại học tăng.

\textbf{Giả thuyết Thay đổi Công nghệ Thiên lệch Kỹ năng (SBTC - Katz \&
Murphy, 1992):} Hàm sản xuất CES với lao động trình độ cao (\(H\)) và
trình độ thấp (\(L\)):
\[ Y = \left[ (\theta_L L)^\rho + (\theta_H H)^\rho \right]^{1/\rho} \]
với \(\rho = \frac{\sigma - 1}{\sigma}\), \(\sigma\) là độ co giãn thay
thế. Tiền lương cân bằng tương đối:
\[ \ln\left(\frac{w_H}{w_L}\right) = \frac{\sigma - 1}{\sigma} \ln\left(\frac{\theta_H}{\theta_L}\right) - \frac{1}{\sigma} \ln\left(\frac{H}{L}\right) \]
Nếu \(\sigma > 1\) (H và L thay thế cho nhau một phần), thì khi
\(\theta_H / \theta_L\) tăng (công nghệ nâng cao năng suất của H nhiều
hơn L, như máy tính), cầu tương đối đối với H sẽ dịch chuyển ra ngoài.
Cuộc ``chạy đua giữa giáo dục và công nghệ'' quyết định phần bù lương:
nếu công nghệ (số hạng đầu) tăng nhanh hơn nguồn cung \(H/L\) (số hạng
sau), bất bình đẳng \(\frac{w_H}{w_L}\) sẽ tăng.

\textbf{Cách tiếp cận dựa trên nhiệm vụ (Task-based Approach - Autor,
Levy, Murnane, 2003):} Máy tính không hoàn toàn thay thế L, mà thay thế
các \textbf{nhiệm vụ thủ tục (routine tasks)} (công việc văn thư, lắp
ráp), thường nằm ở giữa phân phối thu nhập. Ngược lại, máy tính bổ sung
cho các \textbf{nhiệm vụ nhận thức phức tạp (abstract/cognitive tasks)}
(quản lý, lập trình) ở đỉnh, và không thể thay thế các \textbf{nhiệm vụ
thủ công phi cấu trúc (non-routine manual)} (cắt tóc, dọn dẹp) ở đáy.
Kết quả là sự \textbf{Phân cực thị trường lao động (Job Polarization)},
làm ``rỗng ruột'' tầng lớp trung lưu.

\subsubsection{9.7.2 Điểm mù thường gặp (Cognitive Traps \&
Limitations)}\label{ux111iux1ec3m-muxf9-thux1b0ux1eddng-gux1eb7p-cognitive-traps-limitations-75}

\textbf{Đồng nhất ``Công nghệ'' với ``Mất việc làm'' (Technological
Unemployment).} Nhiều người e sợ AI và tự động hóa sẽ tạo ra thất nghiệp
diện rộng. Tuy nhiên, theo Autor (2015), tự động hóa một tác vụ làm giảm
chi phí sản xuất, tăng cầu cho sản phẩm đó, từ đó có thể tăng cầu cho
các tác vụ \emph{không thể tự động hóa} còn lại trong cùng quy trình
(hiệu ứng O-Ring). Bất bình đẳng thu nhập là hệ quả đáng lo ngại hơn là
thất nghiệp tổng thể.

\subsubsection{9.7.3 Ứng dụng thực tế (Practical
Applications)}\label{ux1ee9ng-dux1ee5ng-thux1ef1c-tux1ebf-practical-applications-39}

\textbf{Tác động của AI tạo sinh (Generative AI) lên tương lai việc làm}
Lịch sử cho thấy SBTC và phân cực (1980-2010) đánh vào giới trung lưu.
Tuy nhiên, Generative AI (LLMs) hiện nay có đặc tính khác: chúng tự động
hóa các tác vụ nhận thức (viết mã, soạn thảo văn bản, phân tích y khoa).
Các nghiên cứu thực nghiệm ban đầu (ví dụ Brynjolfsson et al., 2023 trên
nhân viên chăm sóc khách hàng) cho thấy AI đóng vai trò làm \emph{nâng
cao} năng suất của những nhân viên yếu kém/ít kinh nghiệm nhiều hơn là
nhóm top-tier, ngụ ý AI có thể làm \textbf{giảm} bất bình đẳng nội bộ
thay vì làm tăng nó như các làn sóng IT trước đây. Khung phân tích
Autor-Levy-Murnane cần được cập nhật khi định nghĩa về ``routine'' giờ
đây bao hàm cả những suy luận nhận thức ở mức độ cao.

\begin{center}\rule{0.5\linewidth}{0.5pt}\end{center}

\section{CHƯƠNG 10: KINH TẾ HỌC TIỀN TỆ \& TÀI CHÍNH (Monetary \&
Financial
Economics)}\label{chux1b0ux1a1ng-10-kinh-tux1ebf-hux1ecdc-tiux1ec1n-tux1ec7-tuxe0i-chuxednh-monetary-financial-economics}

\subsection{10.1 Lý thuyết Danh mục \& CAPM (Portfolio Theory \& Capital
Asset Pricing
Model)}\label{luxfd-thuyux1ebft-danh-mux1ee5c-capm-portfolio-theory-capital-asset-pricing-model}

\subsubsection{10.1.1 Lý thuyết Danh mục \& CAPM - Cốt lõi Lý
thuyết}\label{luxfd-thuyux1ebft-danh-mux1ee5c-capm---cux1ed1t-luxf5i-luxfd-thuyux1ebft}

Lý thuyết Danh mục Markowitz (1952) giải quyết bài toán phân bổ tài sản
thông qua việc tối ưu hóa kỳ vọng lợi nhuận và phương sai (mean-variance
optimization). Gọi \(w = (w_1, w_2, \dots, w_n)^T\) là vector tỷ trọng
đầu tư vào \(n\) tài sản rủi ro,
\(\mu = (\mu_1, \mu_2, \dots, \mu_n)^T\) là vector lợi nhuận kỳ vọng, và
\(\Sigma\) là ma trận hiệp phương sai \((n \times n)\). Bài toán của nhà
đầu tư là tìm danh mục cực tiểu hóa rủi ro cho một mức lợi nhuận kỳ vọng
\(\bar{\mu}\) cho trước: \[ \min_{w} \frac{1}{2} w^T \Sigma w \] với
điều kiện \(w^T \mu = \bar{\mu}\) và \(w^T \mathbf{1} = 1\). Sử dụng hàm
Lagrangian:
\[ \mathcal{L}(w, \lambda, \gamma) = \frac{1}{2} w^T \Sigma w - \lambda (w^T \mu - \bar{\mu}) - \gamma (w^T \mathbf{1} - 1) \]
Điều kiện bậc nhất (FOC) cho ta:
\[ \Sigma w - \lambda \mu - \gamma \mathbf{1} = 0 \implies w^* = \Sigma^{-1} (\lambda \mu + \gamma \mathbf{1}) \]
Đường cong tập hợp các danh mục tối ưu này tạo thành đường biên hiệu quả
(efficient frontier). Khi đưa thêm tài sản phi rủi ro với lãi suất
\(r_f\), đường biên hiệu quả trở thành đường thẳng Capital Market Line
(CML) tiếp xúc với đường biên của các tài sản rủi ro tại Danh mục Thị
trường (Market Portfolio) \(M\).

Mô hình Định giá Tài sản Vốn (CAPM) do Sharpe (1964) và Lintner (1965)
phát triển từ lý thuyết này. Trong trạng thái cân bằng, lợi nhuận kỳ
vọng của bất kỳ tài sản \(i\) nào phải thỏa mãn:
\[ E[R_i] = r_f + \beta_i (E[R_M] - r_f) \] trong đó phần bù rủi ro thị
trường là \((E[R_M] - r_f)\) và hệ số rủi ro hệ thống \(\beta_i\) được
định nghĩa là: \[ \beta_i = \frac{Cov(R_i, R_M)}{Var(R_M)} \]

\subsubsection{10.1.2 Lý thuyết Danh mục \& CAPM - Điểm mù thường
gặp}\label{luxfd-thuyux1ebft-danh-mux1ee5c-capm---ux111iux1ec3m-muxf9-thux1b0ux1eddng-gux1eb7p}

\subsubsection{10.1.2 Điểm mù thường gặp (Cognitive Traps \&
Limitations)}\label{ux111iux1ec3m-muxf9-thux1b0ux1eddng-gux1eb7p-cognitive-traps-limitations-76}

\subsubsection{10.1.2 Điểm mù thường gặp (Cognitive Traps \&
Limitations)}\label{ux111iux1ec3m-muxf9-thux1b0ux1eddng-gux1eb7p-cognitive-traps-limitations-77}

\subsubsection{10.1.3 Lý thuyết Danh mục \& CAPM - Ứng dụng thực
tế}\label{luxfd-thuyux1ebft-danh-mux1ee5c-capm---ux1ee9ng-dux1ee5ng-thux1ef1c-tux1ebf}

Trong thị trường chứng khoán Việt Nam, việc ước lượng \(\beta\) cho các
cổ phiếu thuộc VN-Index thường được thực hiện thông qua hồi quy chuỗi
thời gian lợi suất. Chẳng hạn, một quỹ đầu tư định giá cổ phiếu VNM
(Vinamilk) bằng cách hồi quy lợi suất của VNM theo lợi suất của
VN-Index. Nếu \(\beta_{VNM} < 1\), VNM được xem là cổ phiếu phòng thủ.
CAPM được sử dụng rộng rãi bởi các công ty chứng khoán để tính toán chi
phí vốn chủ sở hữu (\(K_e\)) khi thực hiện định giá doanh nghiệp bằng
phương pháp dòng tiền chiết khấu (DCF).

\subsubsection{10.1.4 Lý thuyết Danh mục \& CAPM - Ví dụ minh họa tính
toán}\label{luxfd-thuyux1ebft-danh-mux1ee5c-capm---vuxed-dux1ee5-minh-hux1ecda-tuxednh-touxe1n}

\textbf{Bài toán:} Giả sử lãi suất trái phiếu chính phủ kỳ hạn 10 năm
của Việt Nam là \(r_f = 4.0\%\). Lợi nhuận kỳ vọng của VN-Index là
\(E[R_M] = 11.5\%\). Cổ phiếu FPT có hiệp phương sai lợi suất với
VN-Index là \(0.024\) và phương sai lợi suất của VN-Index là \(0.020\).
\textbf{Bước 1:} Tính hệ số Beta của FPT:
\[ \beta_{FPT} = \frac{Cov(R_{FPT}, R_M)}{Var(R_M)} = \frac{0.024}{0.020} = 1.2 \]
\textbf{Bước 2:} Áp dụng công thức CAPM để tính chi phí vốn chủ sở hữu
(lợi nhuận kỳ vọng) cho FPT:
\[ E[R_{FPT}] = r_f + \beta_{FPT} (E[R_M] - r_f) = 4.0\% + 1.2 \times (11.5\% - 4.0\%) \]
\[ E[R_{FPT}] = 4.0\% + 1.2 \times 7.5\% = 4.0\% + 9.0\% = 13.0\% \]
\textbf{Nhận xét:} Cổ phiếu FPT có mức độ rủi ro hệ thống cao hơn thị
trường 20\% (\(\beta = 1.2\)), do đó nhà đầu tư yêu cầu một mức tỷ suất
sinh lợi là 13.0\%, cao hơn so với mức sinh lợi kỳ vọng của thị trường
là 11.5\%.

\begin{center}\rule{0.5\linewidth}{0.5pt}\end{center}

\subsection{10.2 Giả thuyết Thị trường Hiệu quả \& Anomalies (Efficient
Market
Hypothesis)}\label{giux1ea3-thuyux1ebft-thux1ecb-trux1b0ux1eddng-hiux1ec7u-quux1ea3-anomalies-efficient-market-hypothesis}

\subsubsection{10.2.1 Giả thuyết Thị trường Hiệu quả \& Anomalies - Cốt
lõi Lý
thuyết}\label{giux1ea3-thuyux1ebft-thux1ecb-trux1b0ux1eddng-hiux1ec7u-quux1ea3-anomalies---cux1ed1t-luxf5i-luxfd-thuyux1ebft}

Giả thuyết Thị trường Hiệu quả (EMH) của Fama (1970) cho rằng giá cả tài
sản phản ánh đầy đủ mọi thông tin hiện có. Về mặt toán học, giá tài sản
\(P_t\) đi theo bước đi ngẫu nhiên (random walk) kết hợp với một xu
hướng kỳ vọng (drift). Sự thay đổi của giá chỉ phụ thuộc vào thông tin
mới (news), vốn không thể dự đoán được.
\[ E[P_{t+1} | \Omega_t] = P_t (1 + \mu) \] trong đó \(\Omega_t\) là tập
hợp thông tin tại thời điểm \(t\).

Để giải thích các dị thường (anomalies) mà CAPM không giải thích được
(như hiệu ứng quy mô, hiệu ứng giá trị), Fama và French (1993) đã phát
triển mô hình ba yếu tố (Three-Factor Model):
\[ E[R_i] - r_f = \beta_i (E[R_M] - r_f) + s_i E[SMB] + h_i E[HML] \]
Trong đó, \(SMB\) (Small Minus Big) là phần bù rủi ro quy mô, và \(HML\)
(High Minus Low) là phần bù rủi ro giá trị. Sau này, Carhart (1997) thêm
yếu tố động lượng (Momentum - WML). Những mô hình này lập luận rằng các
``dị thường'' thực chất là những yếu tố rủi ro có hệ thống (systematic
risk factors) chưa được mô hình hóa trong CAPM.

\subsubsection{10.2.2 Giả thuyết Thị trường Hiệu quả \& Anomalies - Điểm
mù thường
gặp}\label{giux1ea3-thuyux1ebft-thux1ecb-trux1b0ux1eddng-hiux1ec7u-quux1ea3-anomalies---ux111iux1ec3m-muxf9-thux1b0ux1eddng-gux1eb7p}

\subsubsection{10.2.2 Điểm mù thường gặp (Cognitive Traps \&
Limitations)}\label{ux111iux1ec3m-muxf9-thux1b0ux1eddng-gux1eb7p-cognitive-traps-limitations-78}

\subsubsection{10.2.2 Điểm mù thường gặp (Cognitive Traps \&
Limitations)}\label{ux111iux1ec3m-muxf9-thux1b0ux1eddng-gux1eb7p-cognitive-traps-limitations-79}

\subsubsection{10.2.3 Giả thuyết Thị trường Hiệu quả \& Anomalies - Ứng
dụng thực
tế}\label{giux1ea3-thuyux1ebft-thux1ecb-trux1b0ux1eddng-hiux1ec7u-quux1ea3-anomalies---ux1ee9ng-dux1ee5ng-thux1ef1c-tux1ebf}

Các quỹ đầu tư định lượng (Quant Funds) áp dụng các mô hình đa yếu tố để
xây dựng chiến lược ``Smart Beta''. Thay vì đầu tư vào quỹ chỉ số theo
tỷ trọng vốn hóa, họ xây dựng các danh mục nghiêng về các cổ phiếu có hệ
số sổ sách trên giá thị trường (B/M) cao hoặc động lượng giá mạnh. Tuy
nhiên, ở các thị trường mới nổi như Việt Nam, giới hạn của bán khống và
thanh khoản kém làm cho các yếu tố như Momentum khó thực thi một cách
hoàn hảo.

\subsubsection{10.2.4 Giả thuyết Thị trường Hiệu quả \& Anomalies - Ví
dụ minh họa tính
toán}\label{giux1ea3-thuyux1ebft-thux1ecb-trux1b0ux1eddng-hiux1ec7u-quux1ea3-anomalies---vuxed-dux1ee5-minh-hux1ecda-tuxednh-touxe1n}

\textbf{Bài toán:} Khảo sát lợi suất kỳ vọng của cổ phiếu Ngân hàng ACB
bằng mô hình Fama-French 3 yếu tố. Biết: lãi suất phi rủi ro
\(r_f = 3\%\), phần bù rủi ro thị trường \(E[R_M] - r_f = 8\%\), phần bù
quy mô \(E[SMB] = 3.5\%\), phần bù giá trị \(E[HML] = 4.8\%\). Các hệ số
hồi quy của ACB là: \(\beta = 1.1\), \(s = -0.2\) (vốn hóa lớn),
\(h = 0.5\) (B/M cao). \textbf{Bước 1:} Viết phương trình kỳ vọng:
\[ E[R_{ACB}] = r_f + \beta (E[R_M] - r_f) + s \times E[SMB] + h \times E[HML] \]
\textbf{Bước 2:} Thay số tính toán:
\[ E[R_{ACB}] = 3\% + 1.1 \times 8\% + (-0.2) \times 3.5\% + 0.5 \times 4.8\% \]
\[ E[R_{ACB}] = 3\% + 8.8\% - 0.7\% + 2.4\% = 13.5\% \] \textbf{Nhận
xét:} So với CAPM đơn thuần (cho kết quả
\(3\% + 1.1 \times 8\% = 11.8\%\)), mô hình Fama-French yêu cầu mức lợi
suất 13.5\%. Sự gia tăng này chủ yếu do ACB là cổ phiếu ``giá trị'' mang
rủi ro HML.

\begin{center}\rule{0.5\linewidth}{0.5pt}\end{center}

\subsection{10.3 Cấu trúc Kỳ hạn Lãi suất (Term Structure of Interest
Rates)}\label{cux1ea5u-truxfac-kux1ef3-hux1ea1n-luxe3i-suux1ea5t-term-structure-of-interest-rates}

\subsubsection{10.3.1 Cấu trúc Kỳ hạn Lãi suất - Cốt lõi Lý
thuyết}\label{cux1ea5u-truxfac-kux1ef3-hux1ea1n-luxe3i-suux1ea5t---cux1ed1t-luxf5i-luxfd-thuyux1ebft}

Cấu trúc kỳ hạn mô tả mối quan hệ giữa lãi suất của trái phiếu và thời
gian đáo hạn. Lý thuyết Kỳ vọng (Expectations Hypothesis) khẳng định
rằng lãi suất dài hạn là trung bình hình học của các mức lãi suất ngắn
hạn kỳ vọng trong tương lai. Lãi suất trái phiếu zero-coupon kỳ hạn
\(n\) năm, \(y_n\):
\[ (1 + y_n)^n = (1 + y_1)(1 + E_t[y_{1, t+1}]) \dots (1 + E_t[y_{1, t+n-1}]) \]
Lý thuyết Phần bù Thanh khoản (Liquidity Premium Theory) mở rộng mô hình
này bằng cách thừa nhận nhà đầu tư yêu cầu một phần bù rủi ro
\(L_n > 0\) khi nắm giữ trái phiếu dài hạn:
\[ y_n \approx \frac{1}{n} \sum_{i=0}^{n-1} E_t[y_{1, t+i}] + L_n \]
Trong định giá tài chính, Lãi suất Kỳ hạn (Forward Rate) \(f_{t, t+1}\)
từ năm \(t\) đến \(t+1\) được chiết xuất từ đường cong lợi suất (yield
curve): \[ (1 + f_{t, t+1}) = \frac{(1 + y_{t+1})^{t+1}}{(1 + y_t)^t} \]

\subsubsection{10.3.2 Cấu trúc Kỳ hạn Lãi suất - Điểm mù thường
gặp}\label{cux1ea5u-truxfac-kux1ef3-hux1ea1n-luxe3i-suux1ea5t---ux111iux1ec3m-muxf9-thux1b0ux1eddng-gux1eb7p}

\subsubsection{10.3.2 Điểm mù thường gặp (Cognitive Traps \&
Limitations)}\label{ux111iux1ec3m-muxf9-thux1b0ux1eddng-gux1eb7p-cognitive-traps-limitations-80}

\subsubsection{10.3.2 Điểm mù thường gặp (Cognitive Traps \&
Limitations)}\label{ux111iux1ec3m-muxf9-thux1b0ux1eddng-gux1eb7p-cognitive-traps-limitations-81}

\subsubsection{10.3.3 Cấu trúc Kỳ hạn Lãi suất - Ứng dụng thực
tế}\label{cux1ea5u-truxfac-kux1ef3-hux1ea1n-luxe3i-suux1ea5t---ux1ee9ng-dux1ee5ng-thux1ef1c-tux1ebf}

Đường cong lợi suất đảo ngược (Inverted Yield Curve) - khi lãi suất ngắn
hạn cao hơn dài hạn - là chỉ báo dự báo suy thoái kinh tế. Ở Việt Nam,
Kho bạc Nhà nước sử dụng đường cong lợi suất trái phiếu chính phủ để
định giá các đợt phát hành mới, và các NHTM dùng nó làm tham chiếu để
thiết lập lãi suất cho vay dài hạn.

\subsubsection{10.3.4 Cấu trúc Kỳ hạn Lãi suất - Ví dụ minh họa tính
toán}\label{cux1ea5u-truxfac-kux1ef3-hux1ea1n-luxe3i-suux1ea5t---vuxed-dux1ee5-minh-hux1ecda-tuxednh-touxe1n}

\textbf{Bài toán:} Lãi suất kỳ hạn 1 năm là \(y_1 = 5.0\%\), và lãi suất
kỳ hạn 2 năm là \(y_2 = 5.8\%\). \textbf{Bước 1:} Tính lãi suất kỳ hạn
(forward rate) cho năm thứ 2, ký hiệu \(f_{1,2}\):
\[ (1 + y_2)^2 = (1 + y_1)(1 + f_{1,2}) \]
\[ (1.058)^2 = (1.050) \times (1 + f_{1,2}) \implies f_{1,2} = \frac{1.119364}{1.050} - 1 \approx 6.61\% \]
\textbf{Bước 2:} Phân tích với phần bù thanh khoản. Giả sử
\(L_2 = 0.4\%\).
\[ f_{1,2} \approx E[y_{1, t+1}] + L_2 \implies 6.61\% = E[y_{1, t+1}] + 0.4\% \implies E[y_{1, t+1}] = 6.21\% \]
\textbf{Nhận xét:} Mặc dù lãi suất forward ngụ ý mức 6.61\%, kỳ vọng
thực tế của thị trường về lãi suất giao ngay 1 năm tới chỉ là 6.21\%.

\begin{center}\rule{0.5\linewidth}{0.5pt}\end{center}

\subsection{10.4 Trung gian Tài chính \& Lý thuyết Ngân hàng (Financial
Intermediation)}\label{trung-gian-tuxe0i-chuxednh-luxfd-thuyux1ebft-nguxe2n-huxe0ng-financial-intermediation}

\subsubsection{10.4.1 Trung gian Tài chính - Cốt lõi Lý
thuyết}\label{trung-gian-tuxe0i-chuxednh---cux1ed1t-luxf5i-luxfd-thuyux1ebft}

Mô hình Diamond-Dybvig (1983) giải thích lý do tồn tại của ngân hàng như
nhà cung cấp thanh khoản. Người tiêu dùng có xác suất \(t\) cần tiền ở
\(T=1\) và \((1-t)\) cần tiền ở \(T=2\). Hàm lợi ích:
\[ U = t u(c_1) + (1-t) u(c_2) \] Đầu tư dài hạn sinh lời \(R > 1\) tại
\(T=2\), nhưng chỉ thu về 1 nếu thanh lý ở \(T=1\). Ngân hàng cung cấp
hợp đồng tiền gửi \((c_1^*, c_2^*)\) với \(1 < c_1^* < c_2^* < R\). Điều
này cung cấp bảo hiểm thanh khoản. Tuy nhiên, mô hình có nhiều trạng
thái cân bằng. Nếu người kiên nhẫn tin rằng mọi người sẽ rút ở \(T=1\),
ngân hàng sẽ cạn kiệt tài sản. Phản ứng tối ưu là chạy đi rút tiền, dẫn
đến Bank Run. Diamond (1984) mở rộng lý thuyết ngân hàng với vai trò
giám sát được ủy quyền (delegated monitoring), giảm thiểu chi phí giám
sát dự án so với việc hàng ngàn nhà đầu tư tự đi giám sát độc lập.

\subsubsection{10.4.2 Trung gian Tài chính - Điểm mù thường
gặp}\label{trung-gian-tuxe0i-chuxednh---ux111iux1ec3m-muxf9-thux1b0ux1eddng-gux1eb7p}

\subsubsection{10.4.2 Điểm mù thường gặp (Cognitive Traps \&
Limitations)}\label{ux111iux1ec3m-muxf9-thux1b0ux1eddng-gux1eb7p-cognitive-traps-limitations-82}

\subsubsection{10.4.2 Điểm mù thường gặp (Cognitive Traps \&
Limitations)}\label{ux111iux1ec3m-muxf9-thux1b0ux1eddng-gux1eb7p-cognitive-traps-limitations-83}

\subsubsection{10.4.3 Trung gian Tài chính - Ứng dụng thực
tế}\label{trung-gian-tuxe0i-chuxednh---ux1ee9ng-dux1ee5ng-thux1ef1c-tux1ebf}

Sự sụp đổ của Silicon Valley Bank (SVB) năm 2023. SVB nhận tiền gửi ngắn
hạn và đầu tư vào trái phiếu MBS dài hạn. Khi lãi suất tăng, giá trị MBS
giảm, dẫn đến thanh lý lỗ nặng ở T=1 khi có đột biến rút tiền.

\subsubsection{10.4.4 Trung gian Tài chính - Ví dụ minh họa tính
toán}\label{trung-gian-tuxe0i-chuxednh---vuxed-dux1ee5-minh-hux1ecda-tuxednh-touxe1n}

\textbf{Bài toán:} Có 100 người gửi \$1. Tỷ lệ cần tiền T=1 là
\(t = 25\%\). \(R = 1.4\) ở T=2. Ngân hàng đề xuất hợp đồng
\((c_1 = 1.1, c_2)\). \textbf{Bước 1:} Trạng thái bình thường: Tại
\(T=1\): Trả \(25 \times 1.1 = \$27.5\). Còn lại
\(100 - 27.5 = \$72.5\). Tại \(T=2\): Sinh lời
\(72.5 \times 1.4 = 101.5\) USD. Trả 75 người,
\(c_2 = \frac{101.5}{75} = 1.353\) USD. (\(1.353 > 1.1\)) \textbf{Bước
2:} Trạng thái hoảng loạn (Bank Run): 100 người rút ở T=1. Ngân hàng
thanh lý toàn bộ \$100 tài sản thu về \$100. Chỉ trả đủ \(1.1\) USD cho
\(\frac{100}{1.1} = 90.9\) người. 9 người cuối sẽ mất trắng. Nên việc
chen lấn rút tiền là quyết định kinh tế duy lý đối với cá nhân.

\begin{center}\rule{0.5\linewidth}{0.5pt}\end{center}

\subsection{10.5 Cấu trúc Vốn \& Định lý Modigliani-Miller (Capital
Structure)}\label{cux1ea5u-truxfac-vux1ed1n-ux111ux1ecbnh-luxfd-modigliani-miller-capital-structure}

\subsubsection{10.5.1 Cấu trúc Vốn - Cốt lõi Lý
thuyết}\label{cux1ea5u-truxfac-vux1ed1n---cux1ed1t-luxf5i-luxfd-thuyux1ebft}

Định lý MM (1958) chỉ ra trong thị trường hoàn hảo \(V_L = V_U\). MM
1963 mở rộng với thuế thu nhập doanh nghiệp \(\tau_c\):
\[ V_L = V_U + \tau_c D \] Định đề MM II:
\[ K_e = K_u + (K_u - K_d)(1 - \tau_c)\frac{D}{E} \] Thuyết Đánh đổi
(Trade-off Theory) chỉ ra nợ cao làm tăng Chi phí Kiệt quệ Tài chính:
\[ V_L = V_U + PV(\text{Tax Shield}) - PV(\text{Financial Distress Costs}) \]
WACC được cực tiểu hóa tại mức nợ tối ưu:
\[ WACC = K_e \left(\frac{E}{V_L}\right) + K_d (1 - \tau_c) \left(\frac{D}{V_L}\right) \]

\subsubsection{10.5.2 Cấu trúc Vốn - Điểm mù thường
gặp}\label{cux1ea5u-truxfac-vux1ed1n---ux111iux1ec3m-muxf9-thux1b0ux1eddng-gux1eb7p}

\subsubsection{10.5.2 Điểm mù thường gặp (Cognitive Traps \&
Limitations)}\label{ux111iux1ec3m-muxf9-thux1b0ux1eddng-gux1eb7p-cognitive-traps-limitations-84}

\subsubsection{10.5.2 Điểm mù thường gặp (Cognitive Traps \&
Limitations)}\label{ux111iux1ec3m-muxf9-thux1b0ux1eddng-gux1eb7p-cognitive-traps-limitations-85}

\subsubsection{10.5.3 Cấu trúc Vốn - Ứng dụng thực
tế}\label{cux1ea5u-truxfac-vux1ed1n---ux1ee9ng-dux1ee5ng-thux1ef1c-tux1ebf}

Các doanh nghiệp bất động sản Việt Nam dùng đòn bẩy quá cao. Khi NHNN
siết tín dụng, PV của Chi phí Kiệt quệ tăng vọt. Sự cân bằng bị phá vỡ,
giá trị doanh nghiệp \(V_L\) lao dốc.

\subsubsection{10.5.4 Cấu trúc Vốn - Ví dụ minh họa tính
toán}\label{cux1ea5u-truxfac-vux1ed1n---vuxed-dux1ee5-minh-hux1ecda-tuxednh-touxe1n}

\textbf{Bài toán:} Công ty có \(K_u = 12\%\), \(EBIT = 100\),
\(\tau_c = 20\%\). Phát hành \$200 nợ với \(K_d = 6\%\). \textbf{Bước
1:} Không vay nợ: \(V_U = \frac{100 \times (1 - 0.2)}{0.12} = \$666.67\)
\textbf{Bước 2:} Có vay nợ: \(V_L = 666.67 + 0.2 \times 200 = \$706.67\)
\textbf{Bước 3:} Tính \(K_e\): \(E = 706.67 - 200 = \$506.67\)
\(K_e = 12\% + (12\% - 6\%)(1 - 0.2)\frac{200}{506.67} = 13.89\%\)
\textbf{Bước 4:} Tính WACC:
\(WACC = 13.89\% \times \left(\frac{506.67}{706.67}\right) + 6\% \times 0.8 \times \left(\frac{200}{706.67}\right) = 11.32\%\)
\textbf{Nhận xét:} WACC giảm nhờ tấm chắn thuế, làm tăng giá trị công
ty, nhưng cổ đông phải yêu cầu sinh lời cao hơn.

\begin{center}\rule{0.5\linewidth}{0.5pt}\end{center}

\subsection{10.6 Định giá Phái sinh \& Mô hình Black-Scholes
(Derivatives
Pricing)}\label{ux111ux1ecbnh-giuxe1-phuxe1i-sinh-muxf4-huxecnh-black-scholes-derivatives-pricing}

\subsubsection{10.6.1 Định giá Phái sinh - Cốt lõi Lý
thuyết}\label{ux111ux1ecbnh-giuxe1-phuxe1i-sinh---cux1ed1t-luxf5i-luxfd-thuyux1ebft}

Mô hình Black-Scholes (1973) định giá quyền chọn bằng danh mục mô phỏng
(no-arbitrage). Giá tài sản \(S_t\) theo Brownian hình học
\(dS_t = \mu S_t dt + \sigma S_t dW_t\). PDE định giá:
\[ \frac{\partial V}{\partial t} + \frac{1}{2}\sigma^2 S^2 \frac{\partial^2 V}{\partial S^2} + r S \frac{\partial V}{\partial S} - rV = 0 \]
Giá Call Option kiểu Âu: \[ C = S_0 N(d_1) - K e^{-rT} N(d_2) \] với:
\[ d_1 = \frac{\ln(S_0/K) + (r + \sigma^2/2)T}{\sigma \sqrt{T}}, \quad d_2 = d_1 - \sigma \sqrt{T} \]
Greeks: \(\Delta = N(d_1)\),
\(\Gamma = \frac{N'(d_1)}{S \sigma \sqrt{T}}\),
\(\mathcal{V} = S \sqrt{T} N'(d_1)\).

\subsubsection{10.6.2 Định giá Phái sinh - Điểm mù thường
gặp}\label{ux111ux1ecbnh-giuxe1-phuxe1i-sinh---ux111iux1ec3m-muxf9-thux1b0ux1eddng-gux1eb7p}

\subsubsection{10.6.2 Điểm mù thường gặp (Cognitive Traps \&
Limitations)}\label{ux111iux1ec3m-muxf9-thux1b0ux1eddng-gux1eb7p-cognitive-traps-limitations-86}

\subsubsection{10.6.2 Điểm mù thường gặp (Cognitive Traps \&
Limitations)}\label{ux111iux1ec3m-muxf9-thux1b0ux1eddng-gux1eb7p-cognitive-traps-limitations-87}

\subsubsection{10.6.3 Định giá Phái sinh - Ứng dụng thực
tế}\label{ux111ux1ecbnh-giuxe1-phuxe1i-sinh---ux1ee9ng-dux1ee5ng-thux1ef1c-tux1ebf}

Phòng ngừa rủi ro (hedging) Covered Warrants trên TTCK Việt Nam. Các
CTCK phải Delta hedge và giám sát Vega chặt chẽ để tránh lỗ nặng khi
implied volatility biến động.

\subsubsection{10.6.4 Định giá Phái sinh - Ví dụ minh họa tính
toán}\label{ux111ux1ecbnh-giuxe1-phuxe1i-sinh---vuxed-dux1ee5-minh-hux1ecda-tuxednh-touxe1n}

\textbf{Bài toán:} HPG giá \(S_0 = 30\), \(K = 30\), \(T = 0.25\) năm,
\(r = 5\%\), \(\sigma = 0.3\). \textbf{Bước 1:} Tính \(d_1, d_2\):
\[ d_1 = \frac{\ln(1) + (0.05 + 0.3^2/2) \times 0.25}{0.3 \times \sqrt{0.25}} = 0.1583 \]
\[ d_2 = 0.1583 - 0.3 \times 0.5 = 0.0083 \] \textbf{Bước 2:} Tra bảng
phân phối \(N(x)\): \(N(d_1) \approx 0.5629\),
\(N(d_2) \approx 0.5033\). \textbf{Bước 3:} Tính Call Price:
\[ C = 30 \times 0.5629 - 30 \times e^{-0.0125} \times 0.5033 = 1.975 \]
\textbf{Nhận xét:} Giá chứng quyền là 1.975. Delta \(\Delta = 0.5629\).

\begin{center}\rule{0.5\linewidth}{0.5pt}\end{center}

\subsection{10.7 Cung Tiền \& Cơ chế Truyền dẫn Tiền tệ (Monetary
Transmission)}\label{cung-tiux1ec1n-cux1a1-chux1ebf-truyux1ec1n-dux1eabn-tiux1ec1n-tux1ec7-monetary-transmission}

\subsubsection{10.7.1 Cung Tiền \& Cơ chế Truyền dẫn - Cốt lõi Lý
thuyết}\label{cung-tiux1ec1n-cux1a1-chux1ebf-truyux1ec1n-dux1eabn---cux1ed1t-luxf5i-luxfd-thuyux1ebft}

Cơ chế New Keynesian. Đường IS động:
\[ \tilde{y}_t = E_t[\tilde{y}_{t+1}] - \frac{1}{\sigma} (i_t - E_t[\pi_{t+1}] - r_t^n) \]
Đường cong Phillips NKPC:
\[ \pi_t = \beta E_t[\pi_{t+1}] + \kappa \tilde{y}_t \] Quy tắc Taylor
(Taylor Rule):
\[ i_t = r_t^n + \pi_t + \phi_{\pi}(\pi_t - \pi^*) + \phi_y \tilde{y}_t \]
Nguyên lý Taylor (\(\phi_{\pi} > 1\)) đảm bảo khi lạm phát tăng 1\%, lãi
suất danh nghĩa tăng \(> 1\%\) để lãi suất thực tăng, thắt chặt nền kinh
tế.

\subsubsection{10.7.2 Cung Tiền \& Cơ chế Truyền dẫn - Điểm mù thường
gặp}\label{cung-tiux1ec1n-cux1a1-chux1ebf-truyux1ec1n-dux1eabn---ux111iux1ec3m-muxf9-thux1b0ux1eddng-gux1eb7p}

\subsubsection{10.7.2 Điểm mù thường gặp (Cognitive Traps \&
Limitations)}\label{ux111iux1ec3m-muxf9-thux1b0ux1eddng-gux1eb7p-cognitive-traps-limitations-88}

\subsubsection{10.7.2 Điểm mù thường gặp (Cognitive Traps \&
Limitations)}\label{ux111iux1ec3m-muxf9-thux1b0ux1eddng-gux1eb7p-cognitive-traps-limitations-89}

\subsubsection{10.7.3 Cung Tiền \& Cơ chế Truyền dẫn - Ứng dụng thực
tế}\label{cung-tiux1ec1n-cux1a1-chux1ebf-truyux1ec1n-dux1eabn---ux1ee9ng-dux1ee5ng-thux1ef1c-tux1ebf}

Ngân hàng Nhà nước Việt Nam (SBV) dùng OMO phát hành tín phiếu hút thanh
khoản để nâng lãi suất liên ngân hàng, tạo chênh lệch lãi suất hấp dẫn
giữ vốn ngoại nhằm bảo vệ tỷ giá.

\subsubsection{10.7.4 Cung Tiền \& Cơ chế Truyền dẫn - Ví dụ minh họa
tính
toán}\label{cung-tiux1ec1n-cux1a1-chux1ebf-truyux1ec1n-dux1eabn---vuxed-dux1ee5-minh-hux1ecda-tuxednh-touxe1n}

\textbf{Bài toán:} \(r^n = 2\%\), \(\pi^* = 4\%\), \(\phi_{\pi} = 1.5\),
\(\phi_y = 0.5\). Nền kinh tế có lạm phát \(\pi_t = 7\%\), khoảng trống
sản lượng \(\tilde{y}_t = -3\%\). \textbf{Bước 1:} Áp dụng Taylor Rule:
\[ i_t = 2\% + 7\% + 1.5 \times (7\% - 4\%) + 0.5 \times (-3\%) = 9\% + 4.5\% - 1.5\% = 12.0\% \]
\textbf{Nhận xét:} Để dập tắt lạm phát, NHTW phải nâng lãi suất danh
nghĩa lên 12\%, làm lãi suất thực dương \(5\% > r^n = 2\%\). Điều này sẽ
làm giảm tổng cầu theo đường IS, hi sinh tăng trưởng để hạ nhiệt lạm
phát.

\begin{center}\rule{0.5\linewidth}{0.5pt}\end{center}

\begin{center}\rule{0.5\linewidth}{0.5pt}\end{center}

\section{CHƯƠNG 11: LÝ THUYẾT TRÒ CHƠI NÂNG CAO \& THIẾT KẾ CƠ
CHẾ}\label{chux1b0ux1a1ng-11-luxfd-thuyux1ebft-truxf2-chux1a1i-nuxe2ng-cao-thiux1ebft-kux1ebf-cux1a1-chux1ebf}

\begin{center}\rule{0.5\linewidth}{0.5pt}\end{center}

\subsection{11.1 Tinh chỉnh cân bằng: SPNE \& Trembling-Hand
Perfection}\label{tinh-chux1ec9nh-cuxe2n-bux1eb1ng-spne-trembling-hand-perfection}

\subsubsection{1. Bối cảnh lý thuyết \& Động lực trực
giác}\label{bux1ed1i-cux1ea3nh-luxfd-thuyux1ebft-ux111ux1ed9ng-lux1ef1c-trux1ef1c-giuxe1c}

Trong lý thuyết trò chơi dạng chuẩn (normal-form games), \textbf{Cân
bằng Nash (Nash Equilibrium - NE)} đòi hỏi mỗi người chơi chọn chiến
lược phản ứng tốt nhất (best response) trước chiến lược của đối phương.
Tuy nhiên, trong các trò chơi động (extensive-form games / sequential
games), khái niệm NE không đủ mạnh vì nó cho phép tồn tại những
\textbf{đe dọa không đáng tin (incredible threats)}---những hành động mà
một người chơi đe dọa sẽ thực hiện nếu một nhánh ngoài đường đi cân bằng
(off-equilibrium path) xảy ra, nhưng thực tế khi rơi vào nhánh đó, việc
thực hiện lời đe dọa lại là phi duy lý.

Để loại bỏ các NE mang tính đe dọa rỗng này, Selten (1965, 1975) đã đưa
ra các khái niệm \textbf{tinh chỉnh cân bằng (equilibrium refinements)},
yêu cầu tính hợp lý phải được duy trì không chỉ trên đường đi cân bằng
mà còn tại \emph{mọi nút quyết định} (trong trò chơi con) hoặc dưới một
\emph{xác suất nhiễu loạn nhỏ} (bàn tay run rẩy).

\subsubsection{2. Hình thức hóa toán học \& Định nghĩa chính
thức}\label{huxecnh-thux1ee9c-huxf3a-touxe1n-hux1ecdc-ux111ux1ecbnh-nghux129a-chuxednh-thux1ee9c}

Gọi \(\Gamma_E\) là một trò chơi dạng mở rộng với thông tin hoàn hảo
(perfect information) hoặc không hoàn hảo (imperfect information).

\begin{itemize}
\item
  \textbf{Trò chơi con (Subgame)}: Một trò chơi con \(\Gamma_s\) của
  \(\Gamma_E\) là một cây con bắt đầu tại một nút đơn (singleton
  information set), chứa toàn bộ các nút nối tiếp theo nút đó, và không
  cắt ngang qua bất kỳ tập thông tin (information set) nào khác.
\item
  \textbf{Cân bằng Nash hoàn hảo dưới trò chơi con (Subgame Perfect Nash
  Equilibrium - SPNE)}: Một hồ sơ chiến lược (strategy profile)
  \(s^* = (s_1^*, s_2^*, \dots, s_N^*)\) là một SPNE nếu và chỉ nếu, với
  mọi trò chơi con \(\Gamma_s\) của \(\Gamma_E\), phần giới hạn của hồ
  sơ chiến lược \(s^*\) trên \(\Gamma_s\) (ký hiệu là
  \(s^*|_{\Gamma_s}\)) cấu thành một Cân bằng Nash của trò chơi con
  \(\Gamma_s\):
  \[ u_i\left(s_i^*|_{\Gamma_s}, s_{-i}^*|_{\Gamma_s}\right) \ge u_i\left(s_i, s_{-i}^*|_{\Gamma_s}\right) \quad \forall i \in N, \forall s_i \in S_i(\Gamma_s) \]
\item
  \textbf{Cân bằng hoàn hảo bàn tay run rẩy (Trembling-Hand Perfect
  Equilibrium - THPE)}: Xét trò chơi dạng chuẩn
  \(\Gamma_N = \langle N, \{S_i\}, \{u_i\} \rangle\). Một chiến lược hỗn
  hợp \(\sigma_i \in \Sigma_i\) được gọi là \textbf{hoàn toàn dương
  (totally mixed)} nếu \(\sigma_i(s_i) > 0\) với mọi chiến lược thuần
  túy \(s_i \in S_i\). Một hồ sơ chiến lược
  \(\sigma^* = (\sigma_1^*, \dots, \sigma_N^*)\) là một THPE nếu tồn tại
  một dãy các hồ sơ chiến lược hoàn toàn dương
  \(\{\sigma^k\}_{k=1}^{\infty}\) hội tụ về \(\sigma^*\)
  (\(\lim_{k \to \infty} \sigma^k = \sigma^*\)) sao cho với mọi \(k\),
  \(\sigma_i^*\) là phản ứng tốt nhất trước \(\sigma_{-i}^k\):
  \[ u_i(\sigma_i^*, \sigma_{-i}^k) \ge u_i(s_i, \sigma_{-i}^k) \quad \forall s_i \in S_i, \forall i \in N \]
\end{itemize}

\begin{Shaded}
\begin{Highlighting}[]
\NormalTok{graph TD}
\NormalTok{    A["Nút gốc: Doanh nghiệp 1 (Gia nhập / Không)"]}
\NormalTok{    A {-}{-}\textgreater{}|"Không gia nhập"| B["(0, 100) {-} NE (Đe dọa không đáng tin)"]}
\NormalTok{    A {-}{-}\textgreater{}|"Gia nhập"| C["Nút 2: Doanh nghiệp 2 (Chiến tranh / Nhượng bộ)"]}
\NormalTok{    C {-}{-}\textgreater{}|"Chiến tranh giá"| D["({-}50, {-}50)"]}
\NormalTok{    C {-}{-}\textgreater{}|"Nhượng bộ (Hợp tác)"| E["(40, 40) {-} SPNE chuẩn xác"]}
\end{Highlighting}
\end{Shaded}

\subsubsection{3. Phương pháp giải \& Quyết định chứng
minh}\label{phux1b0ux1a1ng-phuxe1p-giux1ea3i-quyux1ebft-ux111ux1ecbnh-chux1ee9ng-minh}

\begin{itemize}
\tightlist
\item
  \textbf{Quy nạp ngược (Backward Induction)}: Với các trò chơi hữu hạn
  thông tin hoàn hảo, SPNE được tìm bằng thuật toán Zermelo-Kuhn: bắt
  đầu từ các nút quyết định kề với nút lá (terminal nodes), tối ưu hóa
  chiến lược tại các nút này, gán giá trị thỏa dụng tối ưu ngược lên nút
  cha, và lặp lại cho đến nút gốc.
\item
  \textbf{Định lý Kuhn (1953)}: Mọi trò chơi hữu hạn dạng mở rộng với
  thông tin hoàn hảo đều tồn tại ít nhất một SPNE với chiến lược thuần
  túy.
\end{itemize}

\subsubsection{4. Giả định then chốt \& Giới hạn lý
thuyết}\label{giux1ea3-ux111ux1ecbnh-then-chux1ed1t-giux1edbi-hux1ea1n-luxfd-thuyux1ebft}

\begin{itemize}
\tightlist
\item
  \textbf{Tri thức chung về tính duy lý (Common Knowledge of Rationality
  - CKR)}: Quy nạp ngược đòi hỏi người chơi \(i\) phải tin rằng người
  chơi \(j\) sẽ hành động duy lý ở bước \(t+1\), và \(j\) tin rằng \(i\)
  tin \(j\) duy lý, vô hạn cấp.
\item
  \textbf{Nghịch lý trò chơi Rết (Centipede Game Paradox - Rosenthal
  1981)}: Trong trò chơi con Rết 100 bước, lý thuyết SPNE dự đoán hai
  bên sẽ dừng trò chơi ngay tại bước 1 (nhận mức thưởng nhỏ nhất). Tuy
  nhiên, các thí nghiệm thực nghiệm (McKelvey \& Palfrey 1992) cho thấy
  người chơi thường hợp tác kéo dài đến gần cuối trò chơi. Điều này vạch
  trần điểm yếu của SPNE: nếu người chơi 1 thấy người chơi 2 không dừng
  ở bước 1 (vi phạm CKR), niềm tin Bayes của người chơi 1 bị đổ vỡ,
  khiến quy nạp ngược mất hiệu lực.
\end{itemize}

\subsubsection{5. Ứng dụng nghiên cứu \& Mở rộng tiền
duyên}\label{ux1ee9ng-dux1ee5ng-nghiuxean-cux1ee9u-mux1edf-rux1ed9ng-tiux1ec1n-duyuxean}

\begin{itemize}
\tightlist
\item
  \textbf{Mô hình Stackelberg (1934)}: Cạnh tranh lãnh đạo sản lượng
  (Quantity Leadership) là ứng dụng kinh điển của SPNE, nơi doanh nghiệp
  lãnh đạo (Leader) tối ưu hóa sản lượng \(q_1\) dựa trên hàm phản ứng
  tốt nhất của doanh nghiệp đi theo (Follower):
  \(q_2^*(q_1) = \frac{a - c - q_1}{2b}\).
\item
  \textbf{Mặc cả tuần tự Rubinstein (1982)}: Mô hình mặc cả luân phiên
  vô hạn kỳ (Alternating-offers bargaining) với hệ số chiết khấu
  \(\delta_1, \delta_2 \in (0,1)\) có một SPNE duy nhất, trong đó người
  đưa ra lời đề nghị đầu tiên nhận phần chia:
  \[ x_1^* = \frac{1 - \delta_2}{1 - \delta_1 \delta_2} \]
\item
  \textbf{Nghiên cứu tiền duyên}: Tinh chỉnh cân bằng trong trò chơi
  động với thông tin không đầy đủ dẫn đến khái niệm \textbf{Cân bằng
  tuần tự (Sequential Equilibrium - Kreps \& Wilson 1982)} và
  \textbf{Cân bằng Markov hoàn hảo (Markov Perfect Equilibrium - Maskin
  \& Tirole 2001)} trong tổ chức ngành động (Dynamic IO).
\end{itemize}

\begin{center}\rule{0.5\linewidth}{0.5pt}\end{center}

\subsection{11.2 Trò chơi lặp lại \& Định lý Dân gian (Repeated Games \&
Folk
Theorem)}\label{truxf2-chux1a1i-lux1eb7p-lux1ea1i-ux111ux1ecbnh-luxfd-duxe2n-gian-repeated-games-folk-theorem}

\subsubsection{1. Bối cảnh lý thuyết \& Động lực trực
giác}\label{bux1ed1i-cux1ea3nh-luxfd-thuyux1ebft-ux111ux1ed9ng-lux1ef1c-trux1ef1c-giuxe1c-1}

Trong các tương tác một lần (one-shot games), Thế lưỡng nan Người tù
(Prisoner's Dilemma) dẫn đến kết quả kém hiệu quả Pareto: cả hai bên đều
chọn phản bội. Tuy nhiên, khi tương tác được \textbf{lặp lại nhiều lần},
bóng tối của tương lai (shadow of the future) cho phép các bên xây dựng
uy tín và sử dụng các chiến lược trừng phạt (punishment strategies) để
duy trì sự hợp tác. \textbf{Định lý Dân gian (Folk Theorem)} chứng minh
rằng nếu các bên đủ kiên nhẫn (hệ số chiết khấu \(\delta\) gần 1),
\emph{bất kỳ} mức thỏa dụng nào vừa túi tiền và đạt mức tối thiểu cá
nhân đều có thể được duy trì như một cân bằng trong trò chơi lặp lại vô
hạn.

\subsubsection{2. Hình thức hóa toán học \& Định nghĩa chính
thức}\label{huxecnh-thux1ee9c-huxf3a-touxe1n-hux1ecdc-ux111ux1ecbnh-nghux129a-chuxednh-thux1ee9c-1}

Xét trò chơi cơ sở (giai đoạn)
\(\text{G} = \langle N, \{S_i\}, \{u_i\} \rangle\). Trò chơi lặp lại vô
hạn \(\text{G}(\delta)\) có tổng thỏa dụng chiết khấu theo thời gian của
người chơi \(i\):
\[ U_i = (1 - \delta) \sum_{t=0}^{\infty} \delta^t u_i(s_{1,t}, s_{2,t}, \dots, s_{N,t}) \]
trong đó \(\delta \in (0,1)\) là hệ số chiết khấu. Nhân tử
\((1 - \delta)\) dùng để chuẩn hóa tổng lợi ích về cùng đơn vị với lợi
ích một giai đoạn.

\begin{itemize}
\item
  \textbf{Payoffs khả thi (Feasible Payoffs)}: Tập hợp các véc-tơ lợi
  ích \(v = (v_1, \dots, v_N)\) nằm trong bao lồi (convex hull) của các
  kết quả lợi ích của trò chơi cơ sở:
  \[ V = \text{co}\left\{ v \in \mathbb{R}^N \mid \exists s \in S, u(s) = v \right\} \]
\item
  \textbf{Payoffs hợp lý cá nhân nghiêm ngặt (Strictly Individually
  Rational Payoffs)}: Gọi \(v_i^{\text{minmax}}\) là giá trị minmax (chi
  phí tổn thất tối thiểu mà đối thủ có thể ép người chơi \(i\) phải
  chịu):
  \[ v_i^{\text{minmax}} = \min_{s_{-i} \in S_{-i}} \max_{s_i \in S_i} u_i(s_i, s_{-i}) \]
  Một véc-tơ lợi ích \(v\) là hợp lý cá nhân nghiêm ngặt nếu
  \(v_i > v_i^{\text{minmax}}\) với mọi \(i \in N\). Tập hợp này ký hiệu
  là \(V^* = \{v \in V \mid v_i > v_i^{\text{minmax}}, \forall i\}\).
\item
  \textbf{Định lý Dân gian Fudenberg-Maskin (1986)}: Giả sử tập lợi ích
  khả thi \(V^*\) có số chiều đầy đủ (full dimension, tức
  \(\text{dim}(V^*) = N\)). Với mọi véc-tơ lợi ích \(v^* \in V^*\), tồn
  tại một ngưỡng chiết khấu \(\bar{\delta} < 1\) sao cho với mọi
  \(\delta \in (\bar{\delta}, 1)\), tồn tại một Cân bằng Nash Hoàn hảo
  dưới Trò chơi con (SPNE) của \(\text{G}(\delta)\) đem lại lợi ích
  trung bình đúng bằng \(v^*\).
\end{itemize}

\subsubsection{3. Cơ chế trừng phạt \& Chiến
lược}\label{cux1a1-chux1ebf-trux1eebng-phux1ea1t-chiux1ebfn-lux1b0ux1ee3c}

\begin{itemize}
\tightlist
\item
  \textbf{Chiến lược Kích hoàn (Grim Trigger Strategy)}: Hợp tác cho đến
  khi có bất kỳ ai phản bội; sau đó chuyển sang phản bội vĩnh viễn
  (trừng phạt vĩnh viễn).

  \begin{itemize}
  \tightlist
  \item
    Bất đẳng thức duy trì cân bằng trong Thế lưỡng nan Người tù:
    \[ \frac{u_i(C, C)}{1 - \delta} \ge u_i(D, C) + \frac{\delta \cdot u_i(D, D)}{1 - \delta} \iff \delta \ge \frac{u_i(D, C) - u_i(C, C)}{u_i(D, C) - u_i(D, D)} \]
  \end{itemize}
\item
  \textbf{Chiến lược Trừng phạt có thời hạn (Penance / Abreu's Simple
  Penalties)}: Trừng phạt trong \(T\) kỳ bằng cách hạ thỏa dụng xuống
  mức minmax, sau đó quay lại hợp tác nếu kẻ vi phạm chấp nhận hình
  phạt.
\end{itemize}

\begin{center}\rule{0.5\linewidth}{0.5pt}\end{center}

\subsection{11.3 Trò chơi thông tin không hoàn hảo \& Cân bằng Bayesian
hoàn hảo
(PBE)}\label{truxf2-chux1a1i-thuxf4ng-tin-khuxf4ng-houxe0n-hux1ea3o-cuxe2n-bux1eb1ng-bayesian-houxe0n-hux1ea3o-pbe}

\subsubsection{1. Bối cảnh lý thuyết \& Động lực trực
giác}\label{bux1ed1i-cux1ea3nh-luxfd-thuyux1ebft-ux111ux1ed9ng-lux1ef1c-trux1ef1c-giuxe1c-2}

Khi trò chơi diễn ra tuần tự nhưng một số người chơi \textbf{không quan
sát được đặc tính/loại (type)} của người chơi khác (thông tin bất cân
xứng / bất hoàn hảo), khái niệm SPNE không đủ để phân tích vì cây trò
chơi chứa các tập thông tin không phải nút đơn (non-singleton
information sets). Harsanyi (1967) đưa ra cách tiếp cận biến trò chơi
thông tin không đầy đủ thành trò chơi thông tin không hoàn hảo bằng cách
thêm người chơi ``Tự nhiên'' (Nature) chọn loại của người chơi ở bước 0.

Để giải các trò chơi này, Harsanyi, Selten, Kreps \& Wilson, và
Fudenberg \& Tirole (1991) phát triển khái niệm \textbf{Cân bằng
Bayesian Hoàn hảo (Perfect Bayesian Equilibrium - PBE)}, kết hợp hành
động tối ưu với hệ thống \textbf{niềm tin (beliefs)} được cập nhật theo
định lý Bayes dọc theo đường đi cân bằng.

\subsubsection{2. Hình thức hóa toán học \& Định nghĩa chính
thức}\label{huxecnh-thux1ee9c-huxf3a-touxe1n-hux1ecdc-ux111ux1ecbnh-nghux129a-chuxednh-thux1ee9c-2}

Một PBE trong trò chơi dạng mở rộng bao gồm một \textbf{hồ sơ chiến
lược} \(\sigma\) và một \textbf{hệ thống niềm tin} \(\mu\) (gán xác suất
cho mỗi nút trong tập thông tin):

\begin{enumerate}
\def\labelenumi{\arabic{enumi}.}
\item
  \textbf{Tính hợp lý theo trình tự (Sequential Rationality)}: Tại mỗi
  tập thông tin \(I_i\) của người chơi \(i\), chiến lược \(\sigma_i\)
  phải tối đa hóa thỏa dụng kỳ vọng của \(i\), dựa trên hệ thống niềm
  tin \(\mu(I_i)\) và chiến lược của các người chơi khác
  \(\sigma_{-i}\):
  \[ \mathbb{E}[u_i(\sigma_i, \sigma_{-i} \mid I_i, \mu)] \ge \mathbb{E}[u_i(\sigma_i', \sigma_{-i} \mid I_i, \mu)] \quad \forall \sigma_i' \in \Sigma_i \]
\item
  \textbf{Cập nhật niềm tin nhất quán theo Bayes (Bayes' Rule on the
  Equilibrium Path)}: Với mọi tập thông tin \(I\) đạt được với xác suất
  dương dưới \(\sigma\) (\(P(I \mid \sigma) > 0\)), niềm tin \(\mu\)
  phải được cập nhật chính xác theo công thức Bayes:
  \[ \mu(h \mid I) = \frac{P(h \mid \sigma)}{\sum_{h' \in I} P(h' \mid \sigma)} \quad \forall h \in I \]
\item
  \textbf{Cập nhật niềm tin ngoài đường đi cân bằng (Off-Equilibrium
  Beliefs)}: Tại các tập thông tin có xác suất bằng 0 dưới \(\sigma\)
  (\(P(I \mid \sigma) = 0\)), quy tắc Bayes không áp dụng được. Niềm tin
  \(\mu(I)\) có thể gán tùy ý, nhưng phải tuân theo các nguyên tắc tinh
  chế nâng cao (như \textbf{Tiêu chuẩn Trực giác Cho-Kreps 1987}).
\end{enumerate}

\subsubsection{3. Phân loại Cân bằng trong Trò chơi Phát tín hiệu
(Signaling
Games)}\label{phuxe2n-loux1ea1i-cuxe2n-bux1eb1ng-trong-truxf2-chux1a1i-phuxe1t-tuxedn-hiux1ec7u-signaling-games}

Trong trò chơi phát tín hiệu giữa Người phát (Sender - loại \(t \in T\))
và Người nhận (Receiver): * \textbf{Cân bằng Phân tách (Separating
Equilibrium)}: Mỗi loại \(t\) chọn một tín hiệu khác nhau
(\(m(t_1) \neq m(t_2)\)). Người nhận nhận diện hoàn hảo loại của Sender
sau khi thấy tín hiệu: \(\mu(t_1 \mid m(t_1)) = 1\). * \textbf{Cân bằng
Gộp (Pooling Equilibrium)}: Tất cả các loại \(t\) chọn cùng một tín hiệu
(\(m(t_1) = m(t_2) = m^*\)). Tín hiệu không mang thông tin mới; người
nhận giữ nguyên niềm tin tiên nghiệm: \(\mu(t \mid m^*) = p(t)\). *
\textbf{Cân bằng Bán phân tách / Hỗn hợp (Semi-separating / Hybrid
Equilibrium)}: Ít nhất một loại chọn chiến lược hỗn hợp, tiết lộ một
phần thông tin.

\begin{center}\rule{0.5\linewidth}{0.5pt}\end{center}

\subsection{11.4 Lý thuyết thiết kế cơ chế \& Định lý Revelation
(Mechanism Design \& Revelation
Principle)}\label{luxfd-thuyux1ebft-thiux1ebft-kux1ebf-cux1a1-chux1ebf-ux111ux1ecbnh-luxfd-revelation-mechanism-design-revelation-principle}

\subsubsection{1. Bối cảnh lý thuyết \& Động lực trực
giác}\label{bux1ed1i-cux1ea3nh-luxfd-thuyux1ebft-ux111ux1ed9ng-lux1ef1c-trux1ef1c-giuxe1c-3}

Nếu Lý thuyết Trò chơi lấy quy tắc trò chơi làm cố định và dự đoán hành
vi của tác nhân, thì \textbf{Lý thuyết Thiết kế Cơ chế (Mechanism
Design)}---còn gọi là ``Lý thuyết trò chơi ngược'' (Reverse Game Theory,
Hurwicz, Maskin, Myerson - Nobel 2007)---đặt ra câu hỏi: Thiết kế quy
tắc trò chơi (luật chơi, hợp đồng, thể chế) như thế nào để các tác nhân
tự nguyện tiết lộ thông tin riêng tư (private information) và hành xử
dẫn đến kết quả mục tiêu xã hội mong muốn?

\subsubsection{2. Hình thức hóa toán học \& Định nghĩa chính
thức}\label{huxecnh-thux1ee9c-huxf3a-touxe1n-hux1ecdc-ux111ux1ecbnh-nghux129a-chuxednh-thux1ee9c-3}

Xét môi trường với \(N\) người chơi. Mỗi người chơi \(i\) có thông tin
riêng (loại) \(\theta_i \in \Theta_i\). Gọi
\(\Theta = \Theta_1 \times \dots \times \Theta_N\). Một \textbf{Hàm lựa
chọn xã hội (Social Choice Function - SCF)} là một ánh xạ
\(f: \Theta \to X\), gán mỗi bộ loại \(\theta\) với một kết quả xã hội
\(x \in X\).

\begin{itemize}
\item
  \textbf{Cơ chế (Mechanism)}: Một cơ chế
  \(\Gamma = (S_1, \dots, S_N, g(\cdot))\) bao gồm tập không gian chiến
  lược \(S_i\) cho mỗi người chơi và một hàm kết quả
  \(g: S_1 \times \dots \times S_N \to X\).
\item
  \textbf{Cơ chế Trực tiếp (Direct Mechanism)}: Là cơ chế trong đó không
  gian chiến lược của mỗi người chơi chính là tập loại của họ:
  \(S_i = \Theta_i\). Hàm kết quả là \(g: \Theta \to X\).
\item
  \textbf{Định lý Tiết lộ (Revelation Principle - Dasgupta, Hammond,
  Maskin 1979; Myerson 1981)}:

  \begin{itemize}
  \tightlist
  \item
    \textbf{Dạng Chiến lược Trội (Dominant Strategy)}: Nếu một cơ chế
    gián tiếp \(\Gamma = (S, g)\) thực thi một SCF \(f(\theta)\) trong
    cân bằng chiến lược trội, thì tồn tại một \textbf{Cơ chế Trực tiếp
    Trung thực (Direct Incentive Compatible Mechanism)} thực thi chính
    SCF \(f(\theta)\) đó trong chiến lược trội, trong đó báo cáo trung
    thực (\(s_i^*(\theta_i) = \theta_i\)) là một chiến lược trội với mọi
    người chơi.
  \item
    \textbf{Dạng Bayesian}: Nếu một cơ chế thực thi \(f(\theta)\) trong
    Cân bằng Bayes-Nash, thì tồn tại một cơ chế trực tiếp thực thi
    \(f(\theta)\) trong Cân bằng Bayes-Nash trung thực.
  \end{itemize}
\end{itemize}

\[\text{Cơ chế Gián tiếp phức tạp } (S, g) \xrightarrow{\quad \text{Định lý Revelation} \quad} \text{Cơ chế Trực tiếp trung thực } (\Theta, f)\]

Ý nghĩa phương pháp luận: Định lý Revelation cho phép nhà thiết kế cơ
chế giới hạn tìm kiếm trong tập các cơ chế trực tiếp trung thực mà không
lo bỏ sót bất kỳ kết quả xã hội nào có thể đạt được bằng các cơ chế phức
tạp khác.

\subsubsection{3. Điều kiện Tương thích Khuyến khích \& Tự nguyện Tham
gia}\label{ux111iux1ec1u-kiux1ec7n-tux1b0ux1a1ng-thuxedch-khuyux1ebfn-khuxedch-tux1ef1-nguyux1ec7n-tham-gia}

Trong một cơ chế trực tiếp, để đảm bảo tính trung thực: * \textbf{Tương
thích Khuyến khích Bayesian (Bayesian Incentive Compatibility - BIC)}:
\[ \mathbb{E}_{\theta_{-i}}\left[ u_i(f(\theta_i, \theta_{-i}), \theta_i) \mid \theta_i \right] \ge \mathbb{E}_{\theta_{-i}}\left[ u_i(f(\hat{\theta}_i, \theta_{-i}), \theta_i) \mid \theta_i \right] \quad \forall \theta_i, \hat{\theta}_i \in \Theta_i \]
* \textbf{Cá nhân Hợp lý / Tự nguyện Tham gia (Individual Rationality -
IR)}: - \emph{Ex-ante IR}: Trước khi biết \(\theta_i\), lợi ích kỳ vọng
\(\ge\) lợi ích bên ngoài \(\bar{u}_i\). - \emph{Interim IR}: Sau khi
biết \(\theta_i\) nhưng trước khi biết \(\theta_{-i}\):
\(\mathbb{E}_{\theta_{-i}}[u_i(f(\theta), \theta_i) \mid \theta_i] \ge \bar{u}_i(\theta_i)\).
- \emph{Ex-post IR}: Sau khi biết toàn bộ \(\theta\):
\(u_i(f(\theta), \theta_i) \ge \bar{u}_i(\theta_i)\).

\begin{itemize}
\tightlist
\item
  \textbf{Định lý Bất khả thi Myerson-Satterthwaite (1983)}: Trong giao
  dịch song phương giữa một người bán và một người mua có giá trị cá
  nhân độc lập và riêng tư (\(\theta_S, \theta_B\)) với các khoảng giá
  trị giao nhau, \textbf{không tồn tại} một cơ chế nào đồng thời thỏa
  mãn: (1) Hiệu quả Pareto ex-post; (2) Tương thích khuyến khích
  Bayesian (BIC); (3) Tự nguyện tham gia interim (IR); và (4) Ngân sách
  tự cân bằng (No deficit / Budget balance).
\end{itemize}

\begin{center}\rule{0.5\linewidth}{0.5pt}\end{center}

\subsection{11.5 Cơ chế Vickrey-Clarke-Groves (VCG) \& Thiết kế đấu giá
(Auction
Theory)}\label{cux1a1-chux1ebf-vickrey-clarke-groves-vcg-thiux1ebft-kux1ebf-ux111ux1ea5u-giuxe1-auction-theory}

\subsubsection{1. Bối cảnh lý thuyết \& Động lực trực
giác}\label{bux1ed1i-cux1ea3nh-luxfd-thuyux1ebft-ux111ux1ed9ng-lux1ef1c-trux1ef1c-giuxe1c-4}

Làm thế nào để phân bổ hàng hóa công hoặc tài sản công (như tần số vô
tuyến, quyền khai thác khoáng sản) cho người định giá cao nhất mà vẫn
buộc họ tiết lộ đúng định giá thực, đồng thời thu tối đa doanh thu?
\textbf{Cơ chế VCG (Vickrey 1961, Clarke 1971, Groves 1973)} là cơ chế
tổng quát nhất đạt được tính hiệu quả Pareto và tương thích khuyến khích
trong \textbf{chiến lược trội}.

\subsubsection{2. Hình thức hóa toán học Cơ chế
VCG}\label{huxecnh-thux1ee9c-huxf3a-touxe1n-hux1ecdc-cux1a1-chux1ebf-vcg}

Xét bài toán phân bổ hàng hóa công hoặc các dự án công. Mỗi người chơi
\(i\) có hàm thỏa dụng dạng quasilinear:
\(u_i(x, t_i) = v_i(x, \theta_i) + t_i\), trong đó \(x \in X\) là quyết
định dự án và \(t_i\) là khoản chuyển giao tiền tệ.

\begin{enumerate}
\def\labelenumi{\arabic{enumi}.}
\item
  \textbf{Quyết định Hiệu quả Pareto}: Chọn dự án \(x^*(\theta)\) tối đa
  hóa tổng giá trị báo cáo:
  \[ x^*(\theta) = \arg\max_{x \in X} \sum_{i=1}^{N} v_i(x, \theta_i) \]
\item
  \textbf{Quy tắc Thanh toán VCG (Groves Payment Rule)}: Khoản thanh
  toán \(t_i(\theta)\) người chơi \(i\) nhận được là:
  \[ t_i(\theta) = \sum_{j \neq i} v_j(x^*(\theta), \theta_j) + h_i(\theta_{-i}) \]
  Nếu chọn
  \(h_i(\theta_{-i}) = -\max_{x \in X} \sum_{j \neq i} v_j(x_{-i}^*(\theta_{-i}), \theta_j)\)
  (Chi phí Clarke / Internalization of Externality), người chơi \(i\)
  phải trả một khoản thuế bằng \textbf{chi phí ngoại tác} mà sự hiện
  diện của \(i\) gây ra cho các người chơi còn lại:
  \[ p_i(\theta) = -t_i(\theta) = \underbrace{\max_{x \in X} \sum_{j \neq i} v_j(x_{-i}^*(\theta_{-i}), \theta_j)}_{\text{Tổng giá trị người khác nhận được nếu KHÔNG có } i} - \underbrace{\sum_{j \neq i} v_j(x^*(\theta), \theta_j)}_{\text{Tổng giá trị người khác thực nhận khi CÓ } i} \]
\end{enumerate}

\begin{itemize}
\tightlist
\item
  \textbf{Định lý}: Trong cơ chế VCG, báo cáo trung thực
  \(\hat{\theta}_i = \theta_i\) là một \textbf{chiến lược trội} với mọi
  người chơi \(i\).
\end{itemize}

\subsubsection{3. Lý thuyết Đấu giá \& Định lý Doanh thu Tương đương
(Revenue Equivalence
Theorem)}\label{luxfd-thuyux1ebft-ux111ux1ea5u-giuxe1-ux111ux1ecbnh-luxfd-doanh-thu-tux1b0ux1a1ng-ux111ux1b0ux1a1ng-revenue-equivalence-theorem}

Xét đấu giá 1 tài sản duy nhất với \(N\) người mua có giá trị độc lập
riêng tư (IPV) được rút từ cùng phân phối liên tục \(F(v)\) trên
\([\underline{v}, \bar{v}]\).

\begin{itemize}
\tightlist
\item
  \textbf{Bốn hình thức đấu giá tiêu chuẩn}:

  \begin{enumerate}
  \def\labelenumi{\arabic{enumi}.}
  \tightlist
  \item
    \emph{Đấu giá giá tăng dần (English Auction)}: Giá tăng liên tục,
    người rút lui cuối cùng thắng và trả mức giá của người rút lui kế
    cuối.
  \item
    \emph{Đấu giá giá giảm dần (Dutch Auction)}: Giá giảm liên tục,
    người đầu tiên chấp nhận giá thắng và trả đúng giá đó.
  \item
    \emph{Đấu giá kín giá thứ nhất (First-Price Sealed-Bid - FPSB)}: Đặt
    giá kín, người bỏ giá cao nhất thắng và trả đúng giá mình bỏ.
  \item
    \emph{Đấu giá kín giá thứ hai (Second-Price Sealed-Bid / Vickrey
    Auction)}: Đặt giá kín, người bỏ giá cao nhất thắng và trả mức giá
    của người bỏ giá cao thứ hai.
  \end{enumerate}
\item
  \textbf{Định lý Doanh thu Tương đương (Vickrey 1961, Riley-Samuelson
  1981, Myerson 1981)}: Giả sử \(N\) người mua ngại rủi ro trung tính
  (risk-neutral), có giá trị độc lập riêng tư được rút từ cùng một hàm
  phân phối liên tục và nghiêm ngặt \(F(v)\). Khi đó, \textbf{bất kỳ}
  hình thức đấu giá nào thỏa mãn hai điều kiện:

  \begin{enumerate}
  \def\labelenumi{\arabic{enumi}.}
  \tightlist
  \item
    Người có giá trị báo cáo cao nhất luôn thắng tài sản;
  \item
    Người có giá trị thấp nhất \(\underline{v}\) nhận lợi ích kỳ vọng
    bằng 0; đều mang lại \textbf{cùng một doanh thu kỳ vọng} cho người
    bán và cùng một lợi ích kỳ vọng cho mỗi loại người mua.
  \end{enumerate}
\end{itemize}

\[\mathbb{E}[\text{Revenue}_{\text{English}}] = \mathbb{E}[\text{Revenue}_{\text{Dutch}}] = \mathbb{E}[\text{Revenue}_{\text{FPSB}}] = \mathbb{E}[\text{Revenue}_{\text{Vickrey}}] = \mathbb{E}\left[ v_{(2)} \right]\]
trong đó \(v_{(2)}\) là biến ngẫu nhiên thứ tự cao thứ hai
(second-highest order statistic).

\begin{center}\rule{0.5\linewidth}{0.5pt}\end{center}

\subsection{11.6 Ghép đôi thị trường \& Thuật toán Gale-Shapley
(Matching
Theory)}\label{ghuxe9p-ux111uxf4i-thux1ecb-trux1b0ux1eddng-thuux1eadt-touxe1n-gale-shapley-matching-theory}

\subsubsection{1. Bối cảnh lý thuyết \& Động lực trực
giác}\label{bux1ed1i-cux1ea3nh-luxfd-thuyux1ebft-ux111ux1ed9ng-lux1ef1c-trux1ef1c-giuxe1c-5}

Có những thị trường không thể dùng giá cả để điều tiết (thị trường không
giá - bachelorette/bachelor, ghép đôi thận giữa người hiến và bệnh nhân,
phân bổ học sinh vào trường công, phân bổ bác sĩ nội trú vào bệnh viện).
Làm thế nào để phân bổ hai tập hợp tác nhân với nhau một cách \textbf{ổn
định (stable)} sao cho không cặp đôi nào muốn bỏ đối tác hiện tại để đến
với nhau? Roth \& Shapley (Nobel 2012) đã hình thức hóa toán học bài
toán này.

\subsubsection{2. Hình thức hóa toán học Bài toán Hôn nhân Ổn định
(Stable Marriage
Problem)}\label{huxecnh-thux1ee9c-huxf3a-touxe1n-hux1ecdc-buxe0i-touxe1n-huxf4n-nhuxe2n-ux1ed5n-ux111ux1ecbnh-stable-marriage-problem}

Xét hai tập hợp hữu hạn: Tập Nam \(M = \{m_1, \dots, m_n\}\) và Tập Nữ
\(W = \{w_1, \dots, w_p\}\). Mỗi người nam \(m\) có một thứ tự ưa thích
nghiêm ngặt \(P(m)\) trên \(W \cup \{m\}\). Mỗi người nữ \(w\) có một
thứ tự ưa thích nghiêm ngặt \(P(w)\) trên \(M \cup \{w\}\).

\begin{itemize}
\tightlist
\item
  \textbf{Sự ghép đôi (Matching)}: Một ánh xạ song ánh
  \(\mu: M \cup W \to M \cup W\) sao cho:

  \begin{enumerate}
  \def\labelenumi{\arabic{enumi}.}
  \tightlist
  \item
    \(\mu(m) \in W \cup \{m\}\);
  \item
    \(\mu(w) \in M \cup \{w\}\);
  \item
    \(\mu(m) = w \iff \mu(w) = m\).
  \end{enumerate}
\item
  \textbf{Cặp đôi ngăn cản / Cặp đôi mất ổn định (Blocking Pair)}: Một
  cặp \((m, w)\) với \(w \neq \mu(m)\) được gọi là một cặp đôi ngăn cản
  nếu:

  \begin{enumerate}
  \def\labelenumi{\arabic{enumi}.}
  \tightlist
  \item
    \(m\) thích \(w\) hơn người phối ngẫu hiện tại \(\mu(m)\):
    \(w >_m \mu(m)\);
  \item
    \(w\) thích \(m\) hơn người phối ngẫu hiện tại \(\mu(w)\):
    \(m >_w \mu(w)\).
  \end{enumerate}
\item
  \textbf{Sự ghép đôi Ổn định (Stable Matching)}: Một sự ghép đôi
  \(\mu\) được gọi là ổn định nếu nó không bị cản bởi bất kỳ cá nhân độc
  thân nào (individual rationality) và \textbf{không tồn tại bất kỳ cặp
  đôi ngăn cản nào}.
\end{itemize}

\subsubsection{3. Thuật toán Gale-Shapley (Deferred Acceptance Algorithm
-
1962)}\label{thuux1eadt-touxe1n-gale-shapley-deferred-acceptance-algorithm---1962}

Thuật toán Gale-Shapley (Nam cầu hôn, Nữ hoãn chấp nhận):

\begin{itemize}
\item
  \textbf{Vòng 1}: Mỗi người nam \(m\) cầu hôn người nữ xếp đầu tiên
  trong danh sách ưa thích của mình. Mỗi người nữ \(w\) nhận được các
  lời cầu hôn sẽ ``tạm giữ'' người nam mà cô ấy thích nhất và từ chối
  tất cả những người nam còn lại.
\item
  \textbf{Vòng \(k\)}: Những người nam bị từ chối ở vòng trước sẽ cầu
  hôn người nữ tiếp theo trong danh sách ưa thích của mình (người mà họ
  chưa từng cầu hôn). Mỗi người nữ \(w\) so sánh những người cầu hôn mới
  với người nam mà cô ấy đang ``tạm giữ'' từ vòng trước, tiếp tục ``tạm
  giữ'' người tốt nhất theo thứ tự ưa thích của cô ấy và từ chối những
  người còn lại.
\item
  \textbf{Kết thúc}: Thuật toán dừng lại khi không còn người nam nào bị
  từ chối hoặc đã cầu hôn hết danh sách. Tất cả các lời ``tạm giữ'' trở
  thành chính thức.
\item
  \textbf{Các định lý cốt lõi của Gale-Shapley}:

  \begin{enumerate}
  \def\labelenumi{\arabic{enumi}.}
  \tightlist
  \item
    \textbf{Định lý Tồn tại (Gale-Shapley 1962)}: Thuật toán luôn dừng
    lại sau một số hữu hạn bước và cho ra một \textbf{Sự ghép đôi Ổn
    định}.
  \item
    \textbf{Tối ưu hóa bên đề xuất (Proposer-Optimal)}: Thuật toán
    Nam-cầu-hôn tạo ra sự ghép đôi ổn định \textbf{tốt nhất có thể cho
    tất cả người nam} (Nam nhận được người nữ tốt nhất mà họ có thể đạt
    được trong bất kỳ sự ghép đôi ổn định nào). Đồng thời, nó tạo ra sự
    ghép đôi ổn định \textbf{tệ nhất cho tất cả người nữ}
    (Women-pessimal).
  \item
    \textbf{Tính Tương thích Khuyến khích (Strategy-Proofness - Roth
    1982)}: Trong thuật toán Nam-cầu-hôn, báo cáo trung thực thứ tự ưa
    thích là một \textbf{chiến lược trội đối với tất cả người nam}. Tuy
    nhiên, không có cơ chế ghép đôi ổn định nào cho phép báo cáo trung
    thực là chiến lược trội cho \emph{cả hai bên}.
  \end{enumerate}
\end{itemize}

\begin{center}\rule{0.5\linewidth}{0.5pt}\end{center}

\begin{center}\rule{0.5\linewidth}{0.5pt}\end{center}

\section{CHƯƠNG 12: KINH TẾ HỌC THÔNG TIN \& LÝ THUYẾT HỢP ĐỒNG
(Information Economics \& Contract
Theory)}\label{chux1b0ux1a1ng-12-kinh-tux1ebf-hux1ecdc-thuxf4ng-tin-luxfd-thuyux1ebft-hux1ee3p-ux111ux1ed3ng-information-economics-contract-theory}

Kinh tế học thông tin nghiên cứu ảnh hưởng của thông tin bất đối xứng
(asymmetric information) đến các tương tác kinh tế và kết cục thị
trường. Chương này, thuộc Phần B (Chuyên đề nâng cao), mở rộng các nền
tảng của Lý thuyết Trò chơi (đã học ở Chương 11) vào bài toán thiết kế
hợp đồng và cơ chế trong điều kiện thông tin không hoàn hảo.

\begin{center}\rule{0.5\linewidth}{0.5pt}\end{center}

\subsection{12.1 Lựa chọn bất lợi \& Thị trường chanh (Adverse Selection
\& Market for
Lemons)}\label{lux1ef1a-chux1ecdn-bux1ea5t-lux1ee3i-thux1ecb-trux1b0ux1eddng-chanh-adverse-selection-market-for-lemons}

\subsubsection{12.1.1 Cốt lõi: Lý thuyết và Mô hình
Hóa}\label{cux1ed1t-luxf5i-luxfd-thuyux1ebft-vuxe0-muxf4-huxecnh-huxf3a}

Lựa chọn bất lợi (Adverse Selection) xảy ra khi có thông tin bất đối
xứng \emph{trước} khi giao dịch diễn ra (ex-ante hidden information). Mô
hình Akerlof (1970) cho thấy thông tin bất đối xứng có thể dẫn đến sự
sụp đổ của thị trường (market unraveling).

Giả sử chất lượng sản phẩm (ví dụ: xe cũ) được đo bằng biến
\(\theta \in [\underline{\theta}, \overline{\theta}]\), với hàm mật độ
xác suất \(f(\theta)\) và hàm phân phối tích lũy \(F(\theta)\). Người
bán (seller) biết chính xác \(\theta\) của mình, và định giá chiếc xe ở
mức \(V_S(\theta)\). Người mua (buyer) không quan sát được \(\theta\),
nhưng biết phân phối \(F(\theta)\) và định giá chiếc xe ở mức
\(V_B(\theta)\).

Giả định lợi ích từ giao dịch tồn tại với mọi \(\theta\):
\(V_B(\theta) > V_S(\theta)\). Cụ thể, giả sử \(V_B(\theta) = v \theta\)
và \(V_S(\theta) = \theta\), với \(v > 1\). Người mua bàng quan với rủi
ro (risk-neutral), do đó sẵn sàng trả giá \(P\) bằng với giá trị kỳ vọng
của xe trên thị trường: \[ P = E[V_B(\theta) | \text{Xe được bán}] \]

Với mức giá \(P\), người bán chỉ bán xe nếu
\(V_S(\theta) \leq P \iff \theta \leq P\). Do đó, chất lượng kỳ vọng của
xe trên thị trường ở mức giá \(P\) là:
\[ E[\theta | \theta \leq P] = \frac{\int_{\underline{\theta}}^{P} \theta f(\theta) d\theta}{F(P)} \]

\textbf{Định lý 1 (Sự sụp đổ thị trường - Market Unraveling):} Nếu
\(v \cdot E[\theta | \theta \leq P] < P\) với mọi
\(P \in (\underline{\theta}, \overline{\theta}]\), thì không tồn tại cân
bằng giá dương mà tại đó xe chất lượng cao được bán. Thị trường sẽ thu
hẹp về mức giá \(P = \underline{\theta}\), và chỉ có xe chất lượng kém
nhất (chanh) được giao dịch.

\subsubsection{12.1.2 Điểm mù thường
gặp}\label{ux111iux1ec3m-muxf9-thux1b0ux1eddng-gux1eb7p}

\textbf{Lầm tưởng:} Bất kỳ sự chênh lệch thông tin nào cũng làm thị
trường sụp đổ hoàn toàn. \textbf{Thực tế:} Thị trường chỉ sụp đổ hoàn
toàn khi người mua không định giá chất lượng trung bình đủ cao so với
chi phí biên của người bán. Nếu \(v\) đủ lớn, hoặc phân phối
\(F(\theta)\) lệch mạnh về chất lượng cao, thị trường có thể đạt cân
bằng một phần. Ngoài ra, sinh viên thường nhầm lẫn giữa \emph{Lựa chọn
bất lợi} (loại hình ẩn - hidden type) và \emph{Rủi ro đạo đức} (hành
động ẩn - hidden action).

\subsubsection{12.1.3 Ứng dụng thực tế \& Ví dụ số
học}\label{ux1ee9ng-dux1ee5ng-thux1ef1c-tux1ebf-vuxed-dux1ee5-sux1ed1-hux1ecdc}

\subsubsection{12.1.3 Ứng dụng thực tế (Practical
Applications)}\label{ux1ee9ng-dux1ee5ng-thux1ef1c-tux1ebf-practical-applications-40}

\subsubsection{12.1.4 Ví dụ minh họa tính toán (Worked Numerical
Examples)}\label{vuxed-dux1ee5-minh-hux1ecda-tuxednh-touxe1n-worked-numerical-examples-32}

\begin{itemize}
\tightlist
\item
  Với giá \(P\), những người bán có \(\theta \leq P\) sẽ bán.
\item
  Giá trị kỳ vọng của người mua:
  \(E[V_B | \theta \leq P] = \frac{3}{2} \cdot \frac{P}{2} = 0.75P\).
\item
  Điều kiện cân bằng là \(P \leq 0.75P\). Điều này chỉ đúng khi
  \(P = 0\). Kết luận: Dù xe luôn có giá trị với người mua cao hơn người
  bán (\(1.5\theta > \theta\)), thị trường sụp đổ hoàn toàn, không có
  giao dịch nào diễn ra vì thông tin bất đối xứng.
\end{itemize}

\begin{center}\rule{0.5\linewidth}{0.5pt}\end{center}

\subsection{12.2 Rủi ro đạo đức \& Hợp đồng tối ưu (Moral Hazard \&
Optimal
Contracts)}\label{rux1ee7i-ro-ux111ux1ea1o-ux111ux1ee9c-hux1ee3p-ux111ux1ed3ng-tux1ed1i-ux1b0u-moral-hazard-optimal-contracts}

\subsubsection{12.2.1 Cốt lõi: Lý thuyết và Mô hình
Hóa}\label{cux1ed1t-luxf5i-luxfd-thuyux1ebft-vuxe0-muxf4-huxecnh-huxf3a-1}

Rủi ro đạo đức phát sinh do \emph{hành động ẩn} (ex-post hidden action).
Bài toán Principal-Agent thiết lập mô hình nơi Principal (P) thuê Agent
(A) thực hiện nỗ lực \(e \in E\), đem lại kết quả \(x\) với hàm mật độ
\(f(x|e)\). P không quan sát được \(e\), chỉ quan sát được \(x\), và cần
thiết kế hợp đồng trả thưởng \(w(x)\).

\begin{itemize}
\tightlist
\item
  Principal là risk-neutral: \(V(x, w) = x - w(x)\).
\item
  Agent là risk-averse: \(U(w, e) = u(w(x)) - c(e)\), với
  \(u' > 0, u'' < 0, c' > 0, c'' > 0\).
\end{itemize}

\textbf{Bài toán Second-Best:}
\[ \max_{w(x), e} \int [x - w(x)] f(x|e) dx \] Ràng buộc tham gia
(Participation Constraint - PC):
\[ \int u(w(x)) f(x|e) dx - c(e) \geq \bar{u} \] Ràng buộc tương thích
động cơ (Incentive Compatibility - IC):
\[ e \in \arg\max_{e'} \int u(w(x)) f(x|e') dx - c(e') \]

Sử dụng phương pháp First-Order Approach (giả định hàm mục tiêu của A
lõm theo \(e\)), ràng buộc IC trở thành:
\(\int u(w(x)) f_e(x|e) dx - c'(e) = 0\). Hàm Lagrangian là:
\[ \mathcal{L} = \int [x - w(x)] f(x|e) dx + \lambda \left[ \int u(w(x)) f(x|e) dx - c(e) - \bar{u} \right] + \mu \left[ \int u(w(x)) f_e(x|e) dx - c'(e) \right] \]

Tối ưu hóa theo \(w(x)\), ta có điều kiện bậc nhất (FOC):
\[ \frac{1}{u'(w(x))} = \lambda + \mu \frac{f_e(x|e)}{f(x|e)} \] Thuật
ngữ \(\frac{f_e(x|e)}{f(x|e)}\) là \emph{tỷ số khả năng} (likelihood
ratio).

\subsubsection{12.2.2 Điểm mù thường
gặp}\label{ux111iux1ec3m-muxf9-thux1b0ux1eddng-gux1eb7p-1}

\textbf{Lầm tưởng:} Hợp đồng tối ưu luôn phạt Agent khi kết quả thấp và
thưởng khi kết quả cao. \textbf{Thực tế:} Cấu trúc của \(w(x)\) phụ
thuộc vào tỷ số khả năng (Monotone Likelihood Ratio Property - MLRP).
Nếu MLRP không thỏa mãn, kết quả \(x\) cao chưa chắc là tín hiệu tốt cho
nỗ lực \(e\) cao, do đó \(w(x)\) có thể không tăng đơn điệu theo \(x\).
Hơn nữa, nếu Agent là risk-neutral, ta đạt được First-best bằng cách bán
công ty cho Agent (hợp đồng nhượng quyền).

\subsubsection{12.2.3 Ứng dụng thực tế \& Ví dụ số
học}\label{ux1ee9ng-dux1ee5ng-thux1ef1c-tux1ebf-vuxed-dux1ee5-sux1ed1-hux1ecdc-1}

\subsubsection{12.2.3 Ứng dụng thực tế (Practical
Applications)}\label{ux1ee9ng-dux1ee5ng-thux1ef1c-tux1ebf-practical-applications-41}

\subsubsection{12.2.4 Ví dụ minh họa tính toán (Worked Numerical
Examples)}\label{vuxed-dux1ee5-minh-hux1ecda-tuxednh-touxe1n-worked-numerical-examples-33}

Xác suất \(P(x_H|e=1) = 0.8\), \(P(x_H|e=0) = 0.5\). Để kích thích
\(e=1\), Principal giải: (PC):
\(0.8\sqrt{w_H} + 0.2\sqrt{w_L} - 0.1 \geq 1\) (IC):
\(0.8\sqrt{w_H} + 0.2\sqrt{w_L} - 0.1 \geq 0.5\sqrt{w_H} + 0.5\sqrt{w_L}\)
Giải PC và IC như các phương trình đẳng thức:
\(0.3\sqrt{w_H} - 0.3\sqrt{w_L} = 0.1 \Rightarrow \sqrt{w_H} - \sqrt{w_L} = \frac{1}{3}\)
Và \(0.8\sqrt{w_H} + 0.2\sqrt{w_L} = 1.1\). Giải hệ:
\(\sqrt{w_L} = 1.1 - 0.8(\frac{1}{3}) = 1.1 - 0.267 = 0.833 \Rightarrow w_L \approx 0.69\)
\(\sqrt{w_H} = 0.833 + 0.333 = 1.166 \Rightarrow w_H \approx 1.36\). Chi
phí kỳ vọng của Principal là \(0.8(1.36) + 0.2(0.69) \approx 1.23\), lớn
hơn First-best \(w = (1.1)^2 = 1.21\). Sự chênh lệch này là Agency Cost
do rủi ro bù đắp (risk premium).

\begin{center}\rule{0.5\linewidth}{0.5pt}\end{center}

\subsection{12.3 Mô hình Phát tín hiệu Spence (Spence Signaling
Model)}\label{muxf4-huxecnh-phuxe1t-tuxedn-hiux1ec7u-spence-spence-signaling-model}

\subsubsection{12.3.1 Cốt lõi: Lý thuyết và Mô hình
Hóa}\label{cux1ed1t-luxf5i-luxfd-thuyux1ebft-vuxe0-muxf4-huxecnh-huxf3a-2}

Phát tín hiệu (Signaling) là cách phe có thông tin riêng nỗ lực truyền
đạt thông tin cho phe thiếu thông tin nhằm khắc phục Lựa chọn bất lợi.
Spence (1973) chứng minh giáo dục (không nhất thiết tăng năng suất) có
thể làm tín hiệu phân loại năng lực.

Giả sử người lao động có loại \(\theta \in \{\theta_L, \theta_H\}\), với
xác suất \(P(\theta_H) = p\). Năng suất biên đúng bằng \(\theta\). Họ
chọn mức học vấn \(e \ge 0\). Chi phí học vấn là \(c(e, \theta)\), với
giả định \(c_e > 0\), \(c_{ee} > 0\). \textbf{Điều kiện Đơn giao
(Single-Crossing Condition / Spence-Mirrlees):}
\[ \frac{\partial c(e, \theta_H)}{\partial e} < \frac{\partial c(e, \theta_L)}{\partial e} \]
Điều này ngụ ý chi phí biên của việc học thấp hơn đối với người có năng
lực cao. Doanh nghiệp (firm) đưa ra mức lương \(w(e) = E[\theta | e]\).

\textbf{Cân bằng phân tách (Separating Equilibrium):} Các loại chọn
\(e\) khác nhau. - Người năng lực thấp chọn \(e_L = 0\), nhận
\(w_L = \theta_L\). - Người năng lực cao chọn \(e_H = e^* > 0\), nhận
\(w_H = \theta_H\). Điều kiện tự chọn (IC): (IC-L):
\(\theta_L - c(0, \theta_L) \ge \theta_H - c(e^*, \theta_L) \implies c(e^*, \theta_L) \ge \theta_H - \theta_L\)
(IC-H):
\(\theta_H - c(e^*, \theta_H) \ge \theta_L - c(0, \theta_H) \implies c(e^*, \theta_H) \le \theta_H - \theta_L\)
Bất kỳ \(e^*\) nào nằm trong khoảng trên kết hợp với niềm tin (beliefs)
thích hợp ngoài đường dẫn (off-path) sẽ hình thành một cân bằng Cử chỉ
Hoàn hảo (Perfect Bayesian Equilibrium). Cho-Kreps (1987)
\emph{Intuitive Criterion} giúp loại bỏ các cân bằng không hợp lý, hội
tụ về cân bằng Riley (\(e^*\) thấp nhất thỏa mãn).

\subsubsection{12.3.2 Điểm mù thường
gặp}\label{ux111iux1ec3m-muxf9-thux1b0ux1eddng-gux1eb7p-2}

\textbf{Lầm tưởng:} Phát tín hiệu mang lại hiệu quả Pareto. \textbf{Thực
tế:} Cân bằng phân tách trong mô hình Spence thường gây lãng phí nguồn
lực xã hội (Deadweight Loss) do loại \(\theta_H\) phải đầu tư chi phí
\(c(e^*, \theta_H)\) vào một hành động phi năng suất chỉ để chứng minh
năng lực. Một cân bằng gộp (Pooling) có thể vượt trội theo tiêu chuẩn
Pareto nếu sự khác biệt năng lực nhỏ và \(p\) đủ lớn.

\subsubsection{12.3.3 Ứng dụng thực tế \& Ví dụ số
học}\label{ux1ee9ng-dux1ee5ng-thux1ef1c-tux1ebf-vuxed-dux1ee5-sux1ed1-hux1ecdc-2}

\subsubsection{12.3.3 Ứng dụng thực tế (Practical
Applications)}\label{ux1ee9ng-dux1ee5ng-thux1ef1c-tux1ebf-practical-applications-42}

\subsubsection{12.3.4 Ví dụ minh họa tính toán (Worked Numerical
Examples)}\label{vuxed-dux1ee5-minh-hux1ecda-tuxednh-touxe1n-worked-numerical-examples-34}

Tìm cân bằng phân tách Riley: Từ (IC-L):
\(e^* \ge 20 - 10 \implies e^* \ge 10\). Từ (IC-H):
\(e^*/2 \le 20 - 10 \implies e^* \le 20\). Riley Equilibrium (chi phí
thấp nhất cho \(\theta_H\)): \(e^* = 10\). Lương \(w(0) = 10\),
\(w(10) = 20\). Thặng dư của \(\theta_H\) là \(20 - 10/2 = 15\). Nếu
không có giáo dục, ở cân bằng pooling (với \(p=0.5\)), lương chung là
\(15\), thặng dư của \(\theta_H\) cũng là \(15\).

\begin{center}\rule{0.5\linewidth}{0.5pt}\end{center}

\subsection{12.4 Sàng lọc \& Phân biệt giá theo mức độ (Screening \&
Second-degree Price
Discrimination)}\label{suxe0ng-lux1ecdc-phuxe2n-biux1ec7t-giuxe1-theo-mux1ee9c-ux111ux1ed9-screening-second-degree-price-discrimination}

\subsubsection{12.4.1 Cốt lõi: Lý thuyết và Mô hình
Hóa}\label{cux1ed1t-luxf5i-luxfd-thuyux1ebft-vuxe0-muxf4-huxecnh-huxf3a-3}

Sàng lọc (Screening) diễn ra khi bên \emph{không có} thông tin chủ động
thiết kế một tập hợp hợp đồng (menu of contracts) để buộc bên có thông
tin bộc lộ loại của mình thông qua cơ chế tự chọn (self-selection).

Ví dụ: Độc quyền bán thiết kế hợp đồng phân biệt giá cấp hai. Khách hàng
có độ sẵn lòng chi trả \(\theta \in \{\theta_1, \theta_2\}\) với
\(\theta_2 > \theta_1\) (xác suất \(\nu_1, \nu_2\)). Tiện ích của khách
hàng \(U = \theta V(q) - T\), với \(q\) là số lượng, \(T\) là giá phải
trả. Chi phí sản xuất của độc quyền là \(C(q) = c \cdot q\). Độc quyền
giải bài toán:
\[ \max_{(q_1, T_1), (q_2, T_2)} \nu_1(T_1 - c q_1) + \nu_2(T_2 - c q_2) \]
Ràng buộc: (IR1): \(\theta_1 V(q_1) - T_1 \ge 0\) (IR2):
\(\theta_2 V(q_2) - T_2 \ge 0\) (IC1):
\(\theta_1 V(q_1) - T_1 \ge \theta_1 V(q_2) - T_2\) (IC2):
\(\theta_2 V(q_2) - T_2 \ge \theta_2 V(q_1) - T_1\)

Ở mức tối ưu, IR1 và IC2 chặt (binding). Từ đó:
\(T_1 = \theta_1 V(q_1)\) và
\(T_2 = \theta_2 V(q_2) - (\theta_2 - \theta_1)V(q_1)\). Giá trị
\((\theta_2 - \theta_1)V(q_1)\) chính là \textbf{địa tô thông tin}
(information rent) mà loại cao nhận được. Điều kiện FOC cho \(q\): -
\(\theta_2 V'(q_2) = c \implies q_2^{SB} = q_2^{FB}\) (No distortion at
the top). -
\(\left( \theta_1 - \frac{\nu_2}{\nu_1}(\theta_2 - \theta_1) \right) V'(q_1) = c \implies q_1^{SB} < q_1^{FB}\)
(Downward distortion at the bottom).

\subsubsection{12.4.2 Điểm mù thường
gặp}\label{ux111iux1ec3m-muxf9-thux1b0ux1eddng-gux1eb7p-3}

\textbf{Lầm tưởng:} Sàng lọc luôn giúp độc quyền thu được toàn bộ thặng
dư tiêu dùng. \textbf{Thực tế:} Không thể thu toàn bộ thặng dư do phải
trả địa tô thông tin cho loại khách hàng cao để ngăn họ giả vờ làm loại
thấp. Nếu độc quyền cố gắng ép kiệt loại cao, loại cao sẽ chọn gói của
loại thấp, vi phạm ràng buộc IC. Hậu quả là gói của loại thấp bị cắt
giảm chất lượng (hoặc số lượng) so với mức hiệu quả Pareto để làm cho
gói này kém hấp dẫn với loại cao.

\subsubsection{12.4.3 Ứng dụng thực tế \& Ví dụ số
học}\label{ux1ee9ng-dux1ee5ng-thux1ef1c-tux1ebf-vuxed-dux1ee5-sux1ed1-hux1ecdc-3}

\subsubsection{12.4.3 Ứng dụng thực tế (Practical
Applications)}\label{ux1ee9ng-dux1ee5ng-thux1ef1c-tux1ebf-practical-applications-43}

\subsubsection{12.4.4 Ví dụ minh họa tính toán (Worked Numerical
Examples)}\label{vuxed-dux1ee5-minh-hux1ecda-tuxednh-touxe1n-worked-numerical-examples-35}

First best: \(1 / \sqrt{q_1} = 1 \implies q_1^{FB} = 1\).
\(2 / \sqrt{q_2} = 1 \implies q_2^{FB} = 4\). Second best: -
\(q_2^{SB} = q_2^{FB} = 4\). -
\(\left(1 - \frac{0.5}{0.5}(2 - 1)\right) / \sqrt{q_1} = 1 \implies 0 / \sqrt{q_1} = 1\),
nghĩa là độc quyền không phục vụ khách hàng loại 1 (Shutdown). (Thay vào
đó, nếu \(\nu_1 = 0.8, \nu_2 = 0.2\), thì
\(\left(1 - 0.25\right) / \sqrt{q_1} = 1 \implies q_1 = 9/16 < 1\)).

\begin{center}\rule{0.5\linewidth}{0.5pt}\end{center}

\subsection{12.5 Hợp đồng không đầy đủ \& Quyền kiểm soát (Incomplete
Contracts \& Property
Rights)}\label{hux1ee3p-ux111ux1ed3ng-khuxf4ng-ux111ux1ea7y-ux111ux1ee7-quyux1ec1n-kiux1ec3m-souxe1t-incomplete-contracts-property-rights}

\subsubsection{12.5.1 Cốt lõi: Lý thuyết và Mô hình
Hóa}\label{cux1ed1t-luxf5i-luxfd-thuyux1ebft-vuxe0-muxf4-huxecnh-huxf3a-4}

Lý thuyết Grossman-Hart-Moore (GHM) lập luận rằng do chi phí giao dịch,
giới hạn nhận thức, và tính không thể xác minh (non-verifiability) trước
tòa, các hợp đồng luôn không đầy đủ. Mọi tình huống không lường trước
được sẽ được giải quyết bằng thương lượng sau (ex-post bargaining). Khi
hợp đồng không đầy đủ, \textbf{quyền kiểm soát thặng dư} (residual
control rights) hay quyền sở hữu tài sản trở thành công cụ chủ chốt để
định hình vị thế đàm phán (threat points).

Mô hình: - Người mua (B) và Người bán (S). Đầu tư đặc thù \(i\) (do S
thực hiện) và \(e\) (do B thực hiện). - Giá trị tạo ra: \(v(i, e)\),
không thể viết trong hợp đồng do không xác minh được (non-verifiable). -
Sau khi đầu tư, hai bên thương lượng giá chia phần thặng dư theo Nash
Bargaining (chia 50-50 thặng dư vượt mức). Điểm đe dọa (threat point)
phụ thuộc vào ai sở hữu tài sản \(A\). - Nếu B sở hữu (Tích hợp dọc lùi
- Backward integration): S có điểm đe dọa thấp, do đó động cơ đầu tư
\(i\) của S giảm mạnh. - Nếu S sở hữu (Thuê ngoài - Non-integration): S
có điểm đe dọa cao hơn, khuyến khích \(i\), nhưng giảm động cơ \(e\) của
B.

\textbf{Định lý:} Quyền sở hữu tài sản nên được giao cho bên nào có
khoản đầu tư đặc thù quan trọng hơn đối với việc tạo ra giá trị cận
biên, nhằm tối thiểu hóa bài toán Hold-up.

\subsubsection{12.5.2 Điểm mù thường
gặp}\label{ux111iux1ec3m-muxf9-thux1b0ux1eddng-gux1eb7p-4}

\textbf{Lầm tưởng:} Bất khả thi trong việc phác thảo hợp đồng hoàn chỉnh
là lý do chính. \textbf{Thực tế:} Không phải không thể viết, mà là Tòa
án không thể \emph{quan sát và xác minh} các biến số như ``chất lượng nỗ
lực'' hoặc ``đầu tư đặc thù''. Maskin \& Tirole (1999) đã chứng minh nếu
các bên hoàn toàn duy lý, ``hợp đồng không đầy đủ'' không thực sự giới
hạn kết quả nếu ta có thể sử dụng cơ chế báo cáo phức tạp (Mechanism
Design), nhưng trong thực tế, các giới hạn thể chế làm cho GHM trở nên
thiết yếu.

\subsubsection{12.5.3 Ứng dụng thực tế \& Ví dụ số
học}\label{ux1ee9ng-dux1ee5ng-thux1ef1c-tux1ebf-vuxed-dux1ee5-sux1ed1-hux1ecdc-4}

\subsubsection{12.5.3 Ứng dụng thực tế (Practical
Applications)}\label{ux1ee9ng-dux1ee5ng-thux1ef1c-tux1ebf-practical-applications-44}

\subsubsection{12.5.4 Ví dụ minh họa tính toán (Worked Numerical
Examples)}\label{vuxed-dux1ee5-minh-hux1ecda-tuxednh-touxe1n-worked-numerical-examples-36}

Giá trị chung: \(V = 10 + 4\sqrt{i} + 4\sqrt{e}\). Chi phí đầu tư:
\(C_S = i\), \(C_B = e\). First-best (tối đa
\(10 + 4\sqrt{i} + 4\sqrt{e} - i - e\)) \(\implies i^* = 4, e^* = 4\).
Bây giờ giả định Non-integration, nếu thương lượng thất bại, B được
\(0 + \sqrt{e}\), S được \(0 + \sqrt{i}\). Thặng dư đàm phán:
\((10 + 4\sqrt{i} + 4\sqrt{e}) - (\sqrt{e} + \sqrt{i}) = 10 + 3\sqrt{i} + 3\sqrt{e}\).
Theo Nash Bargaining (chia đôi), S nhận được:
\(\text{Payoff}_S = \sqrt{i} + 0.5(10 + 3\sqrt{i} + 3\sqrt{e}) = 5 + 2.5\sqrt{i} + 1.5\sqrt{e}\).
S tối đa hóa
\(5 + 2.5\sqrt{i} + 1.5\sqrt{e} - i \implies \frac{1.25}{\sqrt{i}} = 1 \implies i = 1.5625 < 4\).
Hold-up problem xảy ra, dẫn đến mức đầu tư dưới mức hiệu quả xã hội
(underinvestment).

\begin{center}\rule{0.5\linewidth}{0.5pt}\end{center}

\subsection{12.6 Thiết kế cơ chế ứng dụng: Đấu giá \& Phân bổ nguồn
lực}\label{thiux1ebft-kux1ebf-cux1a1-chux1ebf-ux1ee9ng-dux1ee5ng-ux111ux1ea5u-giuxe1-phuxe2n-bux1ed5-nguux1ed3n-lux1ef1c}

\subsubsection{12.6.1 Cốt lõi: Lý thuyết và Mô hình
Hóa}\label{cux1ed1t-luxf5i-luxfd-thuyux1ebft-vuxe0-muxf4-huxecnh-huxf3a-5}

Thiết kế cơ chế (Mechanism Design) là nhánh ngược của lý thuyết trò
chơi: thiết kế các luật chơi để đạt được mục tiêu mong muốn, giả định
các bên hành xử hợp lý. Đấu giá là ứng dụng nổi bật. Giả sử người bán
muốn bán một vật duy nhất cho \(N\) người mua. Giá trị \(v_i\) của người
mua độc lập và rút từ phân phối \(F(v)\) trên \([0, 1]\).

\textbf{Định lý Cân bằng Doanh thu (Revenue Equivalence Theorem - RET):}
Mọi cơ chế đấu giá thỏa mãn 3 điều kiện: (1) Hàng hóa được bán cho người
định giá cao nhất, (2) Người có định giá bằng 0 có lợi ích bằng 0, (3)
Người chơi hành xử rủi ro trung lập, đều mang lại doanh thu kỳ vọng hoàn
toàn bằng nhau cho người bán (ví dụ: Đấu giá giá thứ nhất và thứ hai là
tương đương).

\textbf{Đấu giá tối ưu (Myerson 1981):} Để tối đa hóa doanh thu, người
bán giải bài toán tối ưu hóa với hàm \textbf{Định giá ảo (Virtual
Valuation):} \[ \psi(v_i) = v_i - \frac{1 - F(v_i)}{f(v_i)} \] Người bán
không cấp hàng cho người có định giá \(v_i\) cao nhất, mà cấp cho người
có \emph{định giá ảo \(\psi(v_i)\) cao nhất}. Hơn nữa, người bán thiết
lập giá khởi điểm (reserve price) \(r^*\) sao cho \(\psi(r^*) = 0\).
\[ r^* - \frac{1 - F(r^*)}{f(r^*)} = 0 \]

\subsubsection{12.6.2 Điểm mù thường
gặp}\label{ux111iux1ec3m-muxf9-thux1b0ux1eddng-gux1eb7p-5}

\textbf{Lầm tưởng:} Đấu giá giá thứ hai (Vickrey) là tối ưu nhất cho
người bán. \textbf{Thực tế:} Vickrey Auction tối ưu hóa \emph{hiệu quả
(efficiency)} nhưng chưa chắc đã tối đa hóa \emph{doanh thu (revenue)}.
Đấu giá Vickrey không có giá khởi điểm sẽ mất một phần doanh thu đáng kể
nếu định giá của những người cao nhất không cạnh tranh sát nhau. Giá
khởi điểm tối ưu độc lập với số lượng người tham gia đấu giá \(N\).

\subsubsection{12.6.3 Ứng dụng thực tế \& Ví dụ số
học}\label{ux1ee9ng-dux1ee5ng-thux1ef1c-tux1ebf-vuxed-dux1ee5-sux1ed1-hux1ecdc-5}

\subsubsection{12.6.3 Ứng dụng thực tế (Practical
Applications)}\label{ux1ee9ng-dux1ee5ng-thux1ef1c-tux1ebf-practical-applications-45}

\subsubsection{12.6.4 Ví dụ minh họa tính toán (Worked Numerical
Examples)}\label{vuxed-dux1ee5-minh-hux1ecda-tuxednh-touxe1n-worked-numerical-examples-37}

Giả sử \(v_i \sim U[0, 1]\). Vậy \(F(v) = v, f(v) = 1\). Định giá ảo:
\(\psi(v) = v - \frac{1-v}{1} = 2v - 1\). Giá khởi điểm tối ưu \(r^*\)
được tìm từ \(\psi(r^*) = 0 \implies 2r^* - 1 = 0 \implies r^* = 0.5\).
Nếu có 2 người chơi (\(N=2\)), doanh thu kỳ vọng từ đấu giá giá thứ 2
không có giá khởi điểm là \(E[v_{(2)}] = \frac{1}{3} \approx 0.33\). Tuy
nhiên, nếu sử dụng \(r^* = 0.5\), doanh thu kỳ vọng tăng lên. Phân tích:
Hàng chỉ được bán nếu ít nhất 1 người \(v > 0.5\). Qua tính toán, doanh
thu kỳ vọng khi có \(r^* = 0.5\) đạt mức \(\frac{5}{12} \approx 0.416\).
Sự gia tăng doanh thu này là do thiết kế cơ chế đấu giá tối ưu
(Myerson).

\begin{center}\rule{0.5\linewidth}{0.5pt}\end{center}

\section{CHƯƠNG 13: KINH TẾ HỌC VĨ MÔ NÂNG CAO (Advanced
Macroeconomics)}\label{chux1b0ux1a1ng-13-kinh-tux1ebf-hux1ecdc-vux129-muxf4-nuxe2ng-cao-advanced-macroeconomics}

Như đã thảo luận trong Chương 7 của Phần A, các mô hình vĩ mô căn bản
như IS-LM, AD-AS hay Solow cung cấp một khung phân tích sơ khởi, nhưng
chúng chứa đựng những hạn chế cơ bản khi đối mặt với thực tiễn. Ở cấp độ
tiến sĩ, kinh tế vĩ mô đòi hỏi sự chính xác và nhất quán từ các nền tảng
vi mô (micro-foundations), xử lý kỳ vọng hợp lý (rational expectations),
và mô hình hóa trạng thái cân bằng tổng thể động ngẫu nhiên (DSGE).
Chương này sẽ cung cấp sự chuyển tiếp sâu sắc vào các biên giới của kinh
tế vĩ mô hiện đại.

\begin{center}\rule{0.5\linewidth}{0.5pt}\end{center}

\subsection{13.1 Mô hình Tăng trưởng Nội sinh (Endogenous Growth
Theory)}\label{muxf4-huxecnh-tux103ng-trux1b0ux1edfng-nux1ed9i-sinh-endogenous-growth-theory}

\subsubsection{13.1.1 Cốt lõi Lý thuyết (Core
Theory)}\label{cux1ed1t-luxf5i-luxfd-thuyux1ebft-core-theory-39}

Khác với mô hình Solow (Chương 7) xem tiến bộ công nghệ (TFP) là một quá
trình ngoại sinh thả từ ``trên trời rơi xuống'', Lý thuyết Tăng trưởng
Nội sinh lập luận rằng công nghệ và tri thức là kết quả của các quyết
định tối đa hóa lợi nhuận có chủ đích.

\textbf{Mô hình Romer (1990) về Nghiên cứu \& Phát triển (R\&D):} Mô
hình Romer nhấn mạnh tính phi cạnh tranh (non-rivalry) của ý tưởng. Một
khi một thiết kế mới được tạo ra, nó có thể được sử dụng bởi vô số hãng
mà không làm giảm đi giá trị sử dụng. Hàm sản xuất vĩ mô có dạng:
\[ Y_t = L_{Y,t}^{1-\alpha} \int_0^{A_t} x_{j,t}^\alpha dj \] Trong đó,
\(L_{Y,t}\) là lao động trong sản xuất hàng hóa cuối cùng, \(A_t\) là số
lượng các loại máy móc trung gian (thước đo tri thức/công nghệ), và
\(x_{j,t}\) là số lượng máy móc loại \(j\).

Động lực cốt lõi nằm ở khu vực R\&D tạo ra thiết kế mới. Giả sử sản
lượng R\&D tuân theo: \[ \dot{A}_t = \theta L_{A,t} A_t^\phi \] Với
\(L_{A,t}\) là lao động nghiên cứu, và \(\theta\) là năng suất. Điểm mấu
t là nếu \(\phi = 1\), ta có:
\[ \frac{\dot{A}_t}{A_t} = g_A = \theta L_{A,t} \] Tốc độ tăng trưởng
công nghệ phụ thuộc trực tiếp vào số lượng nhà nghiên cứu. Dọc theo
Đường Tăng trưởng Cân bằng (Balanced Growth Path - BGP), sản lượng và
tiêu dùng bình quân đầu người tăng trưởng ở mức \(g_A\).

\textbf{Mô hình Lucas (1988) về Vốn nhân lực:} Lucas tiếp cận bằng cách
nhấn mạnh ngoại ứng của vốn nhân lực (human capital externalities). Hàm
sản xuất: \[ Y_t = A K_t^\alpha (u_t h_t L_t)^{1-\alpha} h_a^\gamma \]
Trong đó, \(h_t\) là vốn nhân lực của cá nhân, \(u_t\) là thời gian dành
cho làm việc (phần còn lại \(1-u_t\) dành cho tích lũy vốn nhân lực), và
\(h_a\) là vốn nhân lực trung bình trong nền kinh tế tạo ra hiệu ứng lan
tỏa (\(\gamma > 0\)). Quá trình tích lũy vốn nhân lực:
\[ \dot{h}_t = B(1-u_t)h_t \]

Bên cạnh đó, \textbf{Giải Nobel Kinh tế 2025} đã vinh danh Daron
Acemoglu, Simon Johnson và James Robinson (thể chế, xem Chương 14) cùng
với \textbf{Philippe Aghion, Peter Howitt và Joel Mokyr} cho những đóng
góp đột phá về mô hình \textbf{Hủy diệt Sáng tạo (Creative
Destruction)}. Mô hình Schumpeterian của Aghion-Howitt (1992) giải thích
tăng trưởng không chỉ là sự gia tăng số lượng chủng loại hàng hóa (như
Romer), mà là sự cải tiến chất lượng (quality ladder) nơi công nghệ mới
thay thế và đào thải công nghệ cũ. Mokyr bổ sung góc nhìn lịch sử về
cách các ``cấu trúc vi mô'' (như văn hóa khoa học) nuôi dưỡng đổi mới.

\subsubsection{13.1.2 Điểm mù thường gặp (Cognitive Traps \&
Limitations)}\label{ux111iux1ec3m-muxf9-thux1b0ux1eddng-gux1eb7p-cognitive-traps-limitations-90}

Một điểm mù phổ biến là ``hiệu ứng quy mô'' (scale effects debate). Theo
mô hình Romer nguyên bản, một nền kinh tế có dân số lớn hơn (đồng nghĩa
\(L_A\) lớn hơn) sẽ có tốc độ tăng trưởng BGP vĩnh viễn cao hơn. Thực
nghiệm (như của Jones, 1995) bác bỏ điều này: số lượng nhà khoa học tăng
gấp nhiều lần trong thế kỷ 20 nhưng tốc độ tăng trưởng kinh tế không
tăng tương ứng. Sinh viên thường nhầm lẫn giữa sự gia tăng dân số và sự
gia tăng tốc độ tăng trưởng. Các mô hình tăng trưởng thế hệ thứ hai
(Semi-endogenous growth) giải quyết vấn đề này bằng cách thiết lập
\(\phi < 1\), loại bỏ hiệu ứng quy mô lên tốc độ tăng trưởng.

\subsubsection{13.1.3 Ứng dụng thực tế (Practical
Applications)}\label{ux1ee9ng-dux1ee5ng-thux1ef1c-tux1ebf-practical-applications-46}

Xem xét chính sách trợ cấp R\&D của chính phủ. Giả sử chính phủ đánh
thuế thu nhập \(\tau\) và dùng nó để trợ cấp cho các hãng R\&D. Trong mô
hình Romer, điều này làm tăng tỷ suất sinh lợi biên của lao động trong
khu vực R\&D, khiến lao động dịch chuyển từ \(L_Y\) sang \(L_A\). Ví dụ
số: Nếu dân số là \(L = 100\), \(\theta = 0.05\), ban đầu thị trường
phân bổ \(L_A = 20\). Tốc độ tăng trưởng là
\(g = 0.05 \times 20 = 1\%\). Nếu chính sách trợ cấp tối ưu dịch chuyển
\(L_A = 40\), tăng trưởng tăng lên \(2\%\). Điều này minh họa vai trò
của chính phủ trong việc khắc phục sự thất bại thị trường do tính độc
quyền (monopoly power) cấp bằng sáng chế gây ra phân bổ cận biên dưới
mức tối ưu.

\begin{center}\rule{0.5\linewidth}{0.5pt}\end{center}

\subsection{13.1b Mô hình Ramsey--Cass--Koopmans \& Thế hệ Gối đầu OLG
(Diamond 1965): Nền tảng Vi mô cho Vĩ mô
Động}\label{b-muxf4-huxecnh-ramseycasskoopmans-thux1ebf-hux1ec7-gux1ed1i-ux111ux1ea7u-olg-diamond-1965-nux1ec1n-tux1ea3ng-vi-muxf4-cho-vux129-muxf4-ux111ux1ed9ng}

\paragraph{13.1b.1 Cốt lõi Lý thuyết (Core
Theory)}\label{b.1-cux1ed1t-luxf5i-luxfd-thuyux1ebft-core-theory}

Trong khi mô hình Solow (Chương 7) giả định tỷ lệ tiết kiệm ngoại sinh
\(s\) không đổi, vĩ mô học hiện đại đòi hỏi việc tiết kiệm và tích lũy
vốn phải là kết quả tối ưu hóa hành vi gian kỳ của các tác nhân kinh tế.
Hai khung mô hình nền tảng cho sự chuyển tiếp này là \textbf{Mô hình Hộ
gia đình Đại diện Vô hạn kỳ (Ramsey--Cass--Koopmans 1965)} và \textbf{Mô
hình Thế hệ Gối đầu (Overlapping Generations - OLG, Diamond 1965)}.

\textbf{1. Mô hình Ramsey--Cass--Koopmans (Hộ gia đình Vô hạn kỳ):} Hộ
gia đình đại diện tối đa hóa tổng thỏa dụng chiết khấu gian kỳ từ tiêu
dùng \(c_t\) trong suốt chu kỳ vô hạn với tỷ lệ chiết khấu thời gian
\(\rho > 0\) và độ co giãn thay thế gian kỳ \(\sigma\):
\[\max \int_0^\infty e^{-\rho t} \frac{c_t^{1 - \theta} - 1}{1 - \theta} dt\]
trong đó \(\theta = 1/\sigma\) là hệ số e ngại rủi ro tương đối (CRRA).
Ràng buộc ngân sách động theo đầu người của hộ gia đình phụ thuộc vào
mức vốn \(k_t\), tiền lương \(w_t\), lãi suất \(r_t\) và tốc độ tăng dân
số \(n\): \[\dot{k}_t = w_t + (r_t - n - \delta)k_t - c_t\] Lập hàm
Hamiltonian hiện giá:
\(\mathcal{H} = \frac{c_t^{1-\theta}-1}{1-\theta} + \lambda_t [w_t + (r_t - n - \delta)k_t - c_t]\).
Điều kiện bậc nhất
\(\frac{\partial \mathcal{H}}{\partial c} = 0 \implies c_t^{-\theta} = \lambda_t\)
và phương trình vi phân giá trị biến thiên thái
\(\dot{\lambda}_t = (\rho + n + \delta - r_t)\lambda_t\). Vi phân điều
kiện bậc nhất theo thời gian và thế vào phương trình đồng trạng thái, ta
thu được \textbf{Phương trình Euler cho Tiêu dùng}:
\[\frac{\dot{c}_t}{c_t} = \frac{1}{\theta} (r_t - \rho - \delta)\] Trong
nền kinh tế cạnh tranh hoàn hảo, lãi suất bằng sản phẩm biên của vốn:
\(r_t = f'(k_t)\). Hệ phương trình vi phân động
\((\dot{c}_t, \dot{k}_t)\) tạo ra một không gian pha (phase diagram) với
một đường yên ngựa (saddle-path stability) duy nhất hội tụ về trạng thái
dừng tối ưu modified golden rule: \(f'(k^*) = \rho + \delta + n\).

\textbf{2. Mô hình Thế hệ Gối đầu OLG (Diamond 1965):} Khác với Ramsey,
mô hình OLG nhận thức rằng con người có tuổi thọ hữu hạn. Tại mỗi thời
kỳ \(t\), có hai thế hệ sống đồng thời: Thế hệ trẻ (làm việc, kiếm lương
\(w_t\), tiêu dùng \(c_{1,t}\) và tiết kiệm \(s_t\)) và Thế hệ già (về
hưu, tiêu dùng \(c_{2,t+1}\) từ vốn gốc và lãi tiết kiệm). Bài toán tối
ưu hóa của cá nhân sinh ra ở thời kỳ \(t\):
\[\max U = \ln c_{1,t} + \frac{1}{1+\rho} \ln c_{2,t+1} \quad \text{s.t.} \quad c_{1,t} + \frac{c_{2,t+1}}{1+r_{t+1}} = w_t\]
Giải phương trình Euler gian kỳ
\(\frac{c_{2,t+1}}{c_{1,t}} = \frac{1+r_{t+1}}{1+\rho}\), ta được hàm
tiết kiệm của thế hệ trẻ: \[s_t = \frac{1}{2 + \rho} w_t\] Điều kiện cân
bằng thị trường vốn cho thấy vốn của thời kỳ sau \(K_{t+1}\) chính là
tổng tiết kiệm của thế hệ trẻ thời kỳ này: \((1+n) k_{t+1} = s_t(w_t)\).
Với hàm sản xuất Cobb-Douglas
\(y = k^\alpha \implies w_t = (1-\alpha)k_t^\alpha\), phương trình sai
phân bậc nhất chi phối sự tiến hóa của vốn động là:
\[k_{t+1} = \frac{1-\alpha}{(1+n)(2+\rho)} k_t^\alpha\]

\paragraph{13.1b.2 Điểm mù thường gặp (Cognitive Traps \&
Limitations)}\label{b.2-ux111iux1ec3m-muxf9-thux1b0ux1eddng-gux1eb7p-cognitive-traps-limitations}

\begin{itemize}
\tightlist
\item
  \textbf{Điểm mù thứ nhất (Sự bất khả thi của tính bất biến Ricardian
  trong OLG):} Trong mô hình Ramsey vô hạn kỳ, định lý tương đương
  Ricardo (Ricardian Equivalence) phát biểu rằng việc chính phủ tài trợ
  chi tiêu bằng thuế hôm nay hay phát hành nợ (thuế ngày mai) là hoàn
  toàn tương đương vì hộ gia đình sẽ tăng tiết kiệm để để lại di sản cho
  con cháu trả thuế. Tuy nhiên, điểm mù là áp dụng tính bất biến này vào
  đời sống thực. Trong mô hình OLG, vì các thế hệ mới liên tục sinh ra
  và thế hệ già chết đi mà không có lòng vị tha giữa các thế hệ
  (intergenerational altruism - Barro 1974), việc phát hành nợ công hôm
  nay thực chất là sự chuyển giao tài sản từ thế hệ tương lai sang thế
  hệ hiện tại, làm giảm tiết kiệm quốc gia và lấn át (crowding out) vốn
  đầu tư thực.
\item
  \textbf{Điểm mù thứ hai (Cạm bẫy Tích lũy vốn quá mức - Dynamic
  Inefficiency):} Mô hình Ramsey luôn hội tụ về trạng thái dừng hiệu quả
  Pareto. Nhưng mô hình OLG có một điểm mù lý thuyết vĩ đại: nền kinh tế
  có thể rơi vào trạng thái \textbf{Bất hiệu quả Động (Dynamic
  Inefficiency)}, nơi mức vốn cân bằng \(k^*\) vượt quá mức quy tắc vàng
  \(k_{GR}\), khiến sản phẩm biên của vốn nhỏ hơn tốc độ tăng trưởng
  (\(f'(k^*) < n\)). Khi đó, nền kinh tế đang ``tiết kiệm quá nhiều'',
  và một hệ thống an sinh xã hội chi trả ngay (Pay-As-You-Go / PAYG
  pension) chuyển tiền từ người trẻ sang người già sẽ lập tức cải thiện
  phúc lợi Pareto cho tất cả các thế hệ mà không tốn một đồng chi phí!
\end{itemize}

\paragraph{13.1b.3 Ứng dụng thực tế (Practical
Applications)}\label{b.3-ux1ee9ng-dux1ee5ng-thux1ef1c-tux1ebf-practical-applications}

\begin{itemize}
\tightlist
\item
  \textbf{Khủng hoảng Già hóa Dân số \& Quỹ Hưu trí Việt Nam:} Việt Nam
  đang bước vào giai đoạn già hóa dân số nhanh chóng. Mô hình OLG của
  Diamond là công cụ cốt lõi để các chuyên gia kinh tế thuộc Ngân hàng
  Thế giới (WB) và Bộ LĐ-TB\&XH đánh giá tính bền vững của Quỹ Bảo hiểm
  Xã hội (BHXH). Khi tỷ lệ người phụ thuộc tăng lên (nhóm già về hưu) và
  tốc độ tăng trưởng lực lượng lao động (\(n\)) giảm mạnh, tỷ suất sinh
  lời tự nhiên của hệ thống PAYG (\(r \approx n + g\)) sẽ thấp hơn lãi
  suất thị trường vốn, bắt buộc Việt Nam phải chuyển dịch dần sang hệ
  thống hưu trí đa tầng tích lũy tài khoản cá nhân (Fully Funded OLG
  scheme).
\item
  \textbf{Đánh giá Chính sách Nợ công và Lạm phát:} Phân tích OLG cho
  thấy việc chính phủ liên tục thâm hụt ngân sách để xây dựng hạ tầng
  (như đường cao tốc Bắc - Nam) sẽ tạo ra gánh nặng tài khóa cho thế hệ
  trẻ sau 20-30 năm nữa. Do đó, quy tắc trần nợ công 60\% GDP của Việt
  Nam chính là một ràng buộc thể chế nhằm ngăn chặn sự chuyển giá trị
  bất công gian kỳ giữa các thế hệ.
\end{itemize}

\paragraph{13.1b.4 Ví dụ minh họa tính toán (Worked Numerical
Examples)}\label{b.4-vuxed-dux1ee5-minh-hux1ecda-tuxednh-touxe1n-worked-numerical-examples}

Xét nền kinh tế theo mô hình OLG của Diamond (1965) với hai thời kỳ
sống. Hàm thỏa dụng cá nhân: \(U = \ln c_{1,t} + 0,8 \ln c_{2,t+1}\)
(tức \(\beta = \frac{1}{1+\rho} = 0,8 \implies \rho = 0,25\)). Hàm sản
xuất Cobb-Douglas mỗi lao động: \(y = k^{0,3}\) (hệ số vốn
\(\alpha = 0,3\)). Tốc độ tăng trưởng dân số \(n = 0\%\) (dân số ổn
định, \(1+n = 1\)). 1. \textbf{Xác định mức tiền lương và lãi suất theo
vốn \(k_t\):} \[w_t = (1 - \alpha) k_t^\alpha = 0,7 k_t^{0,3}\]
\[r_t = \alpha k_t^{\alpha - 1} = 0,3 k_t^{-0,7}\] 2. \textbf{Hàm tiết
kiệm tối ưu của thế hệ trẻ:}
\[s_t = \frac{\beta}{1 + \beta} w_t = \frac{0,8}{1 + 0,8} w_t = \frac{4}{9} w_t \approx 0,4444 w_t\]
3. \textbf{Phương trình tiến hóa của vốn cân bằng động
(\(k_{t+1} = s_t\) vì \(n=0\)):}
\[k_{t+1} = \frac{4}{9} \times 0,7 k_t^{0,3} = \frac{2,8}{9} k_t^{0,3} \approx 0,3111 k_t^{0,3}\]
4. \textbf{Tìm trạng thái dừng cân bằng dài hạn
(\(k^* = k_{t+1} = k_t\)):}
\[k^* = 0,3111 (k^*)^{0,3} \implies (k^*)^{0,7} = 0,3111 \implies k^* = (0,3111)^{1/0,7} \approx (0,3111)^{1,4286} \approx 0,1925\]
5. \textbf{Kiểm tra tính Hiệu quả Động (Dynamic Efficiency):} Tính lãi
suất tại trạng thái dừng:
\[r^* = 0,3 \times (0,1925)^{-0,7} \approx 0,3 \times 3,184 \approx 0,955 \quad \text{(95,5\% cho cả kỳ 30 năm, tức ~2,25\%/năm)}\]
Vì \(r^* > n = 0\), nền kinh tế này đang ở trạng thái \textbf{Hiệu quả
Động}, không bị tích lũy vốn dư thừa quá mức quy tắc vàng.

\subsection{13.2 Chu kỳ Kinh doanh Thực (Real Business Cycle -
RBC)}\label{chu-kux1ef3-kinh-doanh-thux1ef1c-real-business-cycle---rbc}

\subsubsection{13.2.1 Cốt lõi Lý thuyết (Core
Theory)}\label{cux1ed1t-luxf5i-luxfd-thuyux1ebft-core-theory-40}

Kydland và Prescott (1982) khởi xướng phương pháp RBC, cho rằng các biến
động kinh tế vĩ mô không phải là kết quả của sự thất bại thị trường, mà
là phản ứng tối ưu của các tác nhân hợp lý trước các cú sốc năng suất
(technology shocks) ngoại sinh.

Xét Bài toán Người lập kế hoạch Xã hội (Social Planner's Problem):
\[ \max_{C_t, L_t} \mathbb{E}_0 \sum_{t=0}^\infty \beta^t U(C_t, 1 - L_t) \]
Với ràng buộc:
\[ C_t + K_{t+1} = Z_t K_t^\alpha L_t^{1-\alpha} + (1-\delta) K_t \] Quá
trình sốc công nghệ \(Z_t\) tuân theo AR(1) log-tuyến tính:
\[ \ln Z_t = \rho \ln Z_{t-1} + \epsilon_t, \quad \epsilon_t \sim \mathcal{N}(0, \sigma^2) \]

Phương trình Bellman (Bellman equation) tương ứng:
\[ V(K_t, Z_t) = \max_{C_t, L_t} \left\{ U(C_t, 1-L_t) + \beta \mathbb{E}_t [V(K_{t+1}, Z_{t+1})] \right\} \]
Giải bài toán này cho ta hai điều kiện bậc nhất (FOCs) quan trọng: 1.
\textbf{Phương trình Euler (Euler equation):}
\[ U_C(t) = \beta \mathbb{E}_t \left[ U_C(t+1) \left( \alpha Z_{t+1} K_{t+1}^{\alpha-1} L_{t+1}^{1-\alpha} + 1 - \delta \right) \right] \]
Định giá trị thời gian của việc tiêu dùng hiện tại so với tương lai. 2.
\textbf{Cân bằng liên khoảng thời gian của Lao động:}
\[ \frac{U_L(t)}{U_C(t)} = (1-\alpha) Z_t K_t^\alpha L_t^{-\alpha} = W_t \]
Trong đó \(W_t\) là tiền lương biên.

Bằng cách tuyến tính hóa theo logarit (log-linearization) quanh trạng
thái dừng (steady state), mô hình có thể giải được thành một hệ động lực
học.

\subsubsection{13.2.2 Điểm mù thường gặp (Cognitive Traps \&
Limitations)}\label{ux111iux1ec3m-muxf9-thux1b0ux1eddng-gux1eb7p-cognitive-traps-limitations-91}

Điểm mù lớn nhất khi tiếp cận RBC là kết luận ``suy thoái là sự phân bổ
tối ưu của thị trường''. Trong RBC nguyên bản, khi có cú sốc công nghệ
âm (\(Z_t\) giảm), năng suất biên của lao động giảm, tiền lương giảm.
Người lao động \emph{tự nguyện} chọn nghỉ ngơi nhiều hơn (giảm giờ làm)
do hiệu ứng thay thế liên thời gian. Do đó, suy thoái theo RBC là một
trạng thái tối ưu Pareto, và bất kỳ can thiệp nào của chính phủ cũng làm
giảm phúc lợi. Điều này hoàn toàn trái ngược với trực giác của kinh tế
học Keynes coi suy thoái là thất bại thị trường có thể được khắc phục
bằng chính sách kích cầu.

\subsubsection{13.2.3 Ứng dụng thực tế (Practical
Applications)}\label{ux1ee9ng-dux1ee5ng-thux1ef1c-tux1ebf-practical-applications-47}

Phương pháp ``hiệu chỉnh'' (Calibration) của RBC. Thay vì hồi quy kinh
tế lượng truyền thống, các thông số được chọn dựa trên dữ liệu vi mô:
\(\alpha = 0.33\) (phần đóng góp của vốn), \(\beta = 0.99\) (tương ứng
lãi suất 4\%/năm theo quý). Sau khi mô phỏng (simulate) mô hình với
chuỗi \(\epsilon_t\) ngẫu nhiên, ta thu được chuỗi thời gian của
\(Y, C, K, L\). Phương sai của sản lượng mô phỏng có thể giải thích tới
70\% biến động của sản lượng thực tế ở Hoa Kỳ. Tuy nhiên, mô hình RBC
truyền thống thất bại trong việc giải thích biến động của số lượng việc
làm, vì giả định độ co giãn cung lao động liên thời gian phải lớn một
cách phi thực tế.

\begin{center}\rule{0.5\linewidth}{0.5pt}\end{center}

\subsection{13.3 New Keynesian DSGE}\label{new-keynesian-dsge}

\subsubsection{13.3.1 Cốt lõi Lý thuyết (Core
Theory)}\label{cux1ed1t-luxf5i-luxfd-thuyux1ebft-core-theory-41}

Các mô hình Cân bằng Tổng thể Động Ngẫu nhiên tân Keynes (New Keynesian
DSGE) hợp nhất khung tối ưu hóa liên thời gian của RBC với giả định về
độ cứng nhắc của giá cả (nominal rigidities).

\textbf{Định giá Calvo (Calvo Pricing, 1983):} Mỗi thời kỳ, chỉ có một
tỷ lệ \((1-\theta)\) các hãng được phép tối ưu hóa giá cả của mình. Tỷ
lệ \(\theta\) các hãng còn lại phải giữ nguyên mức giá cũ. Đây là ``Nàng
tiên Calvo'' (Calvo fairy).

Giải bài toán tối đa hóa lợi nhuận đa kỳ của hãng được đặt giá, và qua
các bước tuyến tính hóa, ta thu được \textbf{Đường cong Phillips Tân
Keynes (New Keynesian Phillips Curve - NKPC):}
\[ \pi_t = \beta \mathbb{E}_t [\pi_{t+1}] + \kappa \tilde{y}_t \] Với
\(\pi_t\) là lạm phát, \(\tilde{y}_t\) là khoảng cách sản lượng (output
gap), và \(\kappa = \frac{(1-\theta)(1-\beta\theta)}{\theta} \xi\) (với
\(\xi\) là độ co giãn). Khác với Phillips Curve truyền thống (dựa vào
lạm phát quá khứ), NKPC nhấn mạnh lạm phát hoàn toàn mang tính
\emph{hướng về tương lai} (forward-looking).

Mô hình New Keynesian 3 phương trình (Three-equation NK model) cơ bản:
1. \textbf{Đường IS Động (Dynamic IS Curve):}
\[ \tilde{y}_t = \mathbb{E}_t [\tilde{y}_{t+1}] - \frac{1}{\sigma} (i_t - \mathbb{E}_t[\pi_{t+1}] - r_t^n) \]
Có được từ phương trình Euler, liên hệ output gap với lãi suất thực. 2.
\textbf{NKPC:} như trên. 3. \textbf{Quy tắc Taylor (Taylor Rule):}
\[ i_t = \rho i_{t-1} + (1-\rho)[r_t^n + \phi_\pi \pi_t + \phi_y \tilde{y}_t] + v_t \]
Để hệ thống có nghiệm duy nhất (determinacy), nó phải thỏa mãn điều kiện
\textbf{Blanchard-Kahn}: Ngân hàng trung ương phải tăng lãi suất danh
nghĩa mạnh hơn mức tăng của lạm phát (Nguyên lý Taylor:
\(\phi_\pi > 1\)).

\subsubsection{13.3.2 Điểm mù thường gặp (Cognitive Traps \&
Limitations)}\label{ux111iux1ec3m-muxf9-thux1b0ux1eddng-gux1eb7p-cognitive-traps-limitations-92}

Người học thường nhầm NKPC với đường cong Phillips nguyên bản. Trong
NKPC, ``kỳ vọng lạm phát'' \(\mathbb{E}_t[\pi_{t+1}]\) quyết định hiện
tại. Do đó, một cam kết đáng tin cậy của NHTW về việc duy trì lạm phát
thấp trong tương lai có thể làm giảm lạm phát ngay lập tức hôm nay mà
không gây ra bất kỳ suy thoái nào (sự sụp đổ của chi phí hy sinh -
sacrifice ratio). Điểm mù là tin rằng điều này luôn đúng trên thực tế;
trong khủng hoảng, giới hạn hợp lý của con người và kỳ vọng thích nghi
(adaptive expectations) thường lấn át, khiến NKPC thất bại cục bộ.

\subsubsection{13.3.3 Ứng dụng thực tế (Practical
Applications)}\label{ux1ee9ng-dux1ee5ng-thux1ef1c-tux1ebf-practical-applications-48}

Phân tán giá (Price Dispersion) do Calvo pricing. Bởi vì các hãng không
được thay đổi giá đồng loạt, sinh ra chênh lệch giá cho các mặt hàng
giống hệt nhau, gây ra sự phân bổ sai lệch (misallocation) nguồn lực.
Hàm tổn thất (loss function) của NHTW được xấp xỉ bằng:
\[ \mathcal{L} = \mathbb{E}_t \sum_{t=0}^\infty \beta^t (\pi_t^2 + \lambda \tilde{y}_t^2) \]
NHTW mục tiêu kép (như Fed) sẽ cân nhắc lựa chọn quỹ đạo lãi suất để tối
thiểu hóa tổn thất này, đặc biệt trong các mô phỏng bằng phần mềm
Dynare.

\begin{center}\rule{0.5\linewidth}{0.5pt}\end{center}

\subsection{13.4 Chính sách Tiền tệ Không Truyền thống (Unconventional
Monetary
Policy)}\label{chuxednh-suxe1ch-tiux1ec1n-tux1ec7-khuxf4ng-truyux1ec1n-thux1ed1ng-unconventional-monetary-policy}

\subsubsection{13.4.1 Cốt lõi Lý thuyết (Core
Theory)}\label{cux1ed1t-luxf5i-luxfd-thuyux1ebft-core-theory-42}

Khi lãi suất điều hành chạm mức không (Zero Lower Bound - ZLB, hoặc
Effective Lower Bound - ELB), quy tắc Taylor truyền thống bị phá vỡ. Lãi
suất danh nghĩa \(i_t \ge 0\). Nếu nền kinh tế bị một cú sốc cầu âm lớn,
lãi suất thực cân bằng tự nhiên \(r_t^n\) rơi xuống mức âm. Lãi suất
thực thực tế
\(r_t = i_t - \mathbb{E}_t[\pi_{t+1}] \ge -\mathbb{E}_t[\pi_{t+1}]\). Do
lạm phát kỳ vọng thường thấp trong suy thoái, \(r_t\) bị kẹt ở mức dương
quá cao, làm hụt cầu tổng thể.

Để phá vỡ bế tắc này, NHTW sử dụng: 1. \textbf{Nới lỏng Định lượng
(Quantitative Easing - QE):} Mua tài sản dài hạn. Hoạt động qua kênh tái
cân bằng danh mục (portfolio balance channel) làm giảm phí rủi ro kỳ
hạn, và kênh tín hiệu (signaling channel) cam kết giữ lãi suất thấp. 2.
\textbf{Định hướng Tương lai (Forward Guidance):} Cam kết đường lối
chính sách tương lai. Trong mô hình NK, phương trình Euler cho thấy:
\[ \tilde{y}_t = \mathbb{E}_t [\tilde{y}_{t+k}] - \frac{1}{\sigma} \sum_{j=0}^{k-1} (i_{t+j} - \mathbb{E}_t[\pi_{t+j+1}] - r_{t+j}^n) \]
Bằng cách hứa giữ \(i_{t+j} = 0\) trong thời gian dài hơn mức quy tắc
Taylor yêu cầu (kênh Odyssean), kỳ vọng lạm phát
\(\mathbb{E}_t[\pi_{t+j+1}]\) tăng lên, làm giảm lãi suất thực hiện tại
và kích thích \(\tilde{y}_t\).

\subsubsection{13.4.2 Điểm mù thường gặp (Cognitive Traps \&
Limitations)}\label{ux111iux1ec3m-muxf9-thux1b0ux1eddng-gux1eb7p-cognitive-traps-limitations-93}

Cái bẫy ``Forward Guidance Puzzle''. Trong các tiêu chuẩn mô hình DSGE,
một cam kết cắt giảm lãi suất ở tương lai rất xa (ví dụ 10 năm nữa) có
sức mạnh vô hạn đối với sản lượng hiện tại. Đây là một lỗ hổng lý thuyết
toán học do giả định độ dài vòng đời vô hạn và thị trường tài chính hoàn
hảo (không có rủi ro vỡ nợ cá nhân). Trên thực tế, người tiêu dùng chiết
khấu mạnh các lời hứa xa xôi (myopia/bounded rationality).

Sự nhầm lẫn giữa \textbf{Odyssean} (cam kết ràng buộc tay mình vào cột
buồm) và \textbf{Delphic} (NHTW đơn thuần dự báo kinh tế mù mịt) cũng là
điểm mù. Nếu NHTW hứa giữ lãi suất thấp vì ``kinh tế sẽ tồi tệ'', nó
(Delphic) có thể gây hoảng loạn thay vì kích thích.

\subsubsection{13.4.3 Ứng dụng thực tế (Practical
Applications)}\label{ux1ee9ng-dux1ee5ng-thux1ef1c-tux1ebf-practical-applications-49}

Trong Đại suy thoái 2008-2020, Fed triển khai QE1 đến QE3, làm bảng cân
đối mở rộng từ \$900 tỷ lên \$4.5 nghìn tỷ. Bằng toán học, sửa quy tắc
Taylor tại ZLB: \[ i_t^* = \dots, \quad i_t = \max(0, i_t^*) \] Các NHTW
như ECB và BOJ thậm chí thử nghiệm Lãi suất Âm (Negative Interest Rates
- NIRP). Sự phi tuyến (non-linearity) tại ranh giới ZLB đòi hỏi các
thuật toán giải số trị phức tạp (như projection methods) thay vì
log-linearization thông thường.

\begin{center}\rule{0.5\linewidth}{0.5pt}\end{center}

\subsection{13.5 Chính sách Tài khóa: Số nhân \& Tương đương Ricardo
(Fiscal
Policy)}\label{chuxednh-suxe1ch-tuxe0i-khuxf3a-sux1ed1-nhuxe2n-tux1b0ux1a1ng-ux111ux1b0ux1a1ng-ricardo-fiscal-policy}

\subsubsection{13.5.1 Cốt lõi Lý thuyết (Core
Theory)}\label{cux1ed1t-luxf5i-luxfd-thuyux1ebft-core-theory-43}

\textbf{Định lý Tương đương Ricardo (Ricardian Equivalence Theorem -
RET)} do Barro (1974) phát biểu: Dưới những điều kiện chặt chẽ, việc
chính phủ tài trợ chi tiêu bằng thuế hay phát hành nợ (trái phiếu) không
làm thay đổi các biến thực (như tiêu dùng, đầu tư, lãi suất).

Ràng buộc ngân sách liên thời gian của Chính phủ:
\[ \sum_{t=0}^\infty \frac{G_t}{(1+r)^t} + D_0 = \sum_{t=0}^\infty \frac{T_t}{(1+r)^t} \]
Ràng buộc ngân sách của Hộ gia đình đại diện:
\[ \sum_{t=0}^\infty \frac{C_t}{(1+r)^t} = \sum_{t=0}^\infty \frac{Y_t - T_t}{(1+r)^t} + (1+r)B_0 \]
Bởi vì tổng tài sản tư nhân ban đầu \(B_0\) bằng nợ chính phủ \(D_0\),
thế ràng buộc của chính phủ vào của hộ gia đình:
\[ \sum_{t=0}^\infty \frac{C_t}{(1+r)^t} = \sum_{t=0}^\infty \frac{Y_t - G_t}{(1+r)^t} \]
Thời điểm thu thuế \(T_t\) biến mất hoàn toàn khỏi phương trình. Một đợt
cắt giảm thuế bằng nợ hôm nay ngụ ý thuế cao hơn trong tương lai với giá
trị hiện tại tương đương; người tiêu dùng hợp lý sẽ tiết kiệm toàn bộ
khoản cắt giảm thuế đó.

\subsubsection{13.5.2 Điểm mù thường gặp (Cognitive Traps \&
Limitations)}\label{ux111iux1ec3m-muxf9-thux1b0ux1eddng-gux1eb7p-cognitive-traps-limitations-94}

Người học thường áp dụng RET một cách giáo điều. RET sẽ \textbf{sụp đổ}
khi: 1. \textbf{Ràng buộc thanh khoản (Liquidity constraints):} Hộ gia
đình không thể vay mượn. Nếu chính phủ giảm thuế, hộ gia đình bị thắt
chặt tín dụng sẽ tiêu dùng ngay lập tức (Số nhân lớn hơn 0). 2.
\textbf{Thuế bóp méo (Distortionary taxes):} Nếu thuế tương lai đánh vào
lao động, nó làm thay đổi giá tương đối, kích hoạt hiệu ứng thay thế. 3.
\textbf{Vòng đời hữu hạn (Finite horizons - OLG models):} Nợ chuyển gánh
nặng thuế cho thế hệ chưa sinh ra.

\subsubsection{13.5.3 Ứng dụng thực tế (Practical
Applications)}\label{ux1ee9ng-dux1ee5ng-thux1ef1c-tux1ebf-practical-applications-50}

\textbf{Số nhân Tài khóa (Fiscal Multipliers).} Trong các mô hình RBC,
số nhân chi tiêu chính phủ thường nhỏ hơn 1 (hiệu ứng lấn át - crowding
out). Trong mô hình New Keynesian có ranh giới ZLB, số nhân tài khóa có
thể lớn hơn 1 đáng kể (thậm chí lên mức 1.5 - 2.5 theo Christiano,
Eichenbaum, and Rebelo, 2011). Lý do (toán học): Tăng \(G\) làm tăng kỳ
vọng lạm phát \(\mathbb{E}_t[\pi_{t+1}]\). Vì \(i_t = 0\) tại ZLB, lãi
suất thực \(r_t\) giảm cực mạnh. Điều này kích thích tiêu dùng tư nhân
thay vì lấn át nó. Các gói kích thích (stimulus packages) như ARRA 2009
ở Mỹ phụ thuộc chặt chẽ vào cơ chế này. Ngược lại, chính sách thắt lưng
buộc bụng (austerity) ở Hy Lạp dẫn đến suy thoái sâu hơn lý thuyết
truyền thống dự báo.

\begin{center}\rule{0.5\linewidth}{0.5pt}\end{center}

\subsection{13.6 Khủng hoảng Tài chính \& Mô hình Thế hệ (Financial
Crises
Models)}\label{khux1ee7ng-houx1ea3ng-tuxe0i-chuxednh-muxf4-huxecnh-thux1ebf-hux1ec7-financial-crises-models}

\subsubsection{13.6.1 Cốt lõi Lý thuyết (Core
Theory)}\label{cux1ed1t-luxf5i-luxfd-thuyux1ebft-core-theory-44}

Các lý thuyết vĩ mô tiền 2008 thường bỏ qua vai trò của các trung gian
tài chính.

\textbf{Mô hình Rút tiền ồ ạt Diamond-Dybvig (1983):} Ngân hàng thực
hiện chức năng chuyển đổi kỳ hạn (maturity transformation) từ tiền gửi
thanh khoản ngắn hạn sang các khoản vay kém thanh khoản dài hạn. Mô hình
có 3 thời kỳ: \(t=0, 1, 2\). Người tiêu dùng nộp 1 đơn vị tài sản. Có tỷ
lệ \(t_1\) người tiêu dùng là ``nhân vật sớm'' (cần tiêu dùng ở \(t=1\))
và \(t_2\) là ``nhân vật muộn''. Công nghệ đầu tư: \(1\) đơn vị đầu tư ở
\(t=0\) mang lại \(1\) ở \(t=1\), hoặc \(R > 1\) ở \(t=2\). Ngân hàng
cung cấp hợp đồng bảo hiểm thanh khoản \((c_1, c_2)\) sao cho
\(1 < c_1 < c_2 < R\).

Vì Ràng buộc Dịch vụ Tuần tự (Sequential Service Constraint), người đến
trước được phục vụ trước. Điều này dẫn đến \textbf{Đa cân bằng (Multiple
Equilibria)}: - Cân bằng tốt: Chỉ những người thực sự cần tiền rút ở
\(t=1\). Ngân hàng tồn tại. - Cân bằng xấu (Bank run): Do tin đồn, ngay
cả nhân vật muộn cũng hoảng loạn chạy đến rút tiền ở \(t=1\). Ngân hàng
thanh lý sớm dự án, vỡ nợ, kinh tế suy sụp.

\textbf{Gia tốc Tài chính Bernanke-Gertler (Financial Accelerator):} Do
thông tin bất xứng (asymmetric information) và chi phí giám sát tốn kém,
các hãng phải chịu Phí rủi ro ngoại sinh (External Finance Premium -
EFP). EFP tỷ lệ nghịch với Giá trị ròng (Net Worth) của hãng. Một cú sốc
vĩ mô nhỏ làm giảm giá trị tài sản thế chấp \(\implies\) EFP tăng vọt
\(\implies\) Đầu tư sụp đổ \(\implies\) Giá tài sản tiếp tục giảm (Vòng
xoáy tử thần).

\subsubsection{13.6.2 Điểm mù thường gặp (Cognitive Traps \&
Limitations)}\label{ux111iux1ec3m-muxf9-thux1b0ux1eddng-gux1eb7p-cognitive-traps-limitations-95}

Giả thuyết Thị trường Hiệu quả (EMH) che khuất \textbf{Giả thuyết Minsky
(Minsky Hypothesis)}: ``Sự ổn định sinh ra bất ổn'' (Stability breeds
instability). Trong giai đoạn hưng thịnh kéo dài, các tác nhân tưởng lầm
rủi ro đã biến mất, tự nguyện tăng đòn bẩy (Ponzi finance). Rất khó mô
hình hóa cú sốc nội sinh này trong cấu trúc DSGE tuyến tính thông
thường. Thêm vào đó, khái niệm ``Suy thoái Bảng cân đối kế toán''
(Balance Sheet Recession - Koo 2008) cho thấy khi nợ quá cao, các hãng
chuyển từ tối đa hóa lợi nhuận sang tối thiểu hóa nợ, làm vô hiệu hóa
hoàn toàn chính sách tiền tệ dẫu lãi suất bằng 0.

\subsubsection{13.6.3 Ứng dụng thực tế (Practical
Applications)}\label{ux1ee9ng-dux1ee5ng-thux1ef1c-tux1ebf-practical-applications-51}

Khủng hoảng Tài chính Toàn cầu (GFC 2008) là minh chứng kinh điển của
Bank run hệ thống trên thị trường Shadow Banking (repo, thương phiếu)
chứ không chỉ ở ngân hàng thương mại bán lẻ. Áp dụng bảo hiểm tiền gửi
(Deposit Insurance) như FDIC là giải pháp tối ưu từ Diamond-Dybvig để
loại bỏ cân bằng Bank run xấu. Khủng hoảng châu Á 1997 thể hiện sự cộng
hưởng của ``currency run'' và ``bank run'' do nợ định giá bằng USD.

\begin{center}\rule{0.5\linewidth}{0.5pt}\end{center}

\subsection{13.7 Lý thuyết Tiền tệ Hiện đại \& Tranh luận MMT (MMT vs
Mainstream)}\label{luxfd-thuyux1ebft-tiux1ec1n-tux1ec7-hiux1ec7n-ux111ux1ea1i-tranh-luux1eadn-mmt-mmt-vs-mainstream}

\subsubsection{13.7.1 Cốt lõi Lý thuyết (Core
Theory)}\label{cux1ed1t-luxf5i-luxfd-thuyux1ebft-core-theory-45}

Thập kỷ qua chứng kiến sự trỗi dậy của MMT (Modern Monetary Theory). Cốt
lõi của MMT cho rằng một quốc gia có chủ quyền phát hành đồng tiền pháp
định của riêng mình, có tỷ giá thả nổi tự do và không vay nợ bằng ngoại
tệ, sẽ \textbf{không bao giờ phá sản} do không thể thiếu tiền. Chi tiêu
chính phủ tạo ra tiền, và thuế không tài trợ cho chi tiêu mà chỉ dùng để
tiêu hủy tiền (kiểm soát lạm phát) và phân phối lại thu nhập.

Ràng buộc ngân sách chính phủ (Government Budget Identity):
\[ \frac{\Delta B_{t}}{P_t} + \frac{\Delta M_{t}}{P_t} = G_t - T_t + r \frac{B_{t-1}}{P_t} \]
Quan điểm Mainstream nhấn mạnh rằng nếu phát hành tiền (\(\Delta M_t\))
quá mức để bù đắp thâm hụt (seigniorage revenue), lạm phát \(P_t\) sẽ
phi mã. Khái niệm cốt lõi của Mainstream để chống lại MMT là \textbf{Lý
thuyết Tài khóa về Mức giá (Fiscal Theory of the Price Level - FTPL)}.
FTPL phát biểu rằng mức giá chung không chỉ do cung tiền quyết định, mà
còn đóng vai trò điều chỉnh sao cho giá trị thực của nợ chính phủ hiện
tại bằng tổng giá trị thặng dư tài khóa kỳ vọng trong tương lai:
\[ \frac{B_{t-1}}{P_t} = \mathbb{E}_t \sum_{j=0}^\infty \frac{T_{t+j} - G_{t+j}}{(1+r)^j} \]
Nếu chính phủ in tiền hoặc phát hành nợ vô tội vạ mà không có kế hoạch
thu thuế trong tương lai, \(P_t\) phải tăng lên (lạm phát) để cân bằng
phương trình này.

\subsubsection{13.7.2 Điểm mù thường gặp (Cognitive Traps \&
Limitations)}\label{ux111iux1ec3m-muxf9-thux1b0ux1eddng-gux1eb7p-cognitive-traps-limitations-96}

Người ủng hộ MMT rơi vào điểm mù khi cho rằng giới hạn duy nhất của chi
tiêu là ``toàn dụng lao động'' (full employment) và ``nguồn lực thực
tế''. Thực tế toán học chỉ ra nguồn thu từ đúc tiền (seigniorage) có
giới hạn do dạng đường Laffer: lạm phát cao làm giảm cầu tiền thực tế
(vận tốc vòng quay tiền tăng), khiến nguồn thu từ in tiền sụp đổ (siêu
lạm phát kiểu Zimbabwe hay Venezuela). Ngược lại, điểm mù của những
người chỉ trích Mainstream là quá ám ảnh về ``chính phủ phá sản''. Thực
tế, Mỹ hay Nhật Bản (nợ \textgreater{} 250\% GDP) chứng minh rằng ở mức
độ nhất định, khi lãi suất thực \(r\) nhỏ hơn tốc độ tăng trưởng \(g\)
(\(r < g\)), nợ công có thể cuộn lại vô thời hạn mà không cần tăng thuế.

\subsubsection{13.7.3 Ứng dụng thực tế (Practical
Applications)}\label{ux1ee9ng-dux1ee5ng-thux1ef1c-tux1ebf-practical-applications-52}

Cuộc tranh luận bùng nổ trong đại dịch COVID-19 (2020-2021). Chính phủ
các nước phát triển bơm lượng cung tiền M2 khổng lồ, chi tiêu tài khóa
tài trợ bởi NHTW (dường như ứng dụng MMT). MMT đã đúng ở đoạn đầu khi
không có sự sụp đổ trái phiếu. Nhưng khi lạm phát toàn cầu bùng nổ vào
2022 do năng lực sản xuất (supply constraints) bị vượt quá và kỳ vọng
lạm phát mất neo, giới hạn thực sự của MMT đã bị phơi bày. Điều này củng
cố quan điểm của FTPL: khi chính phủ không có chiến lược tài khóa bền
vững, lạm phát chính là công cụ tàn nhẫn nhất để xóa bỏ giá trị nợ công.

\begin{center}\rule{0.5\linewidth}{0.5pt}\end{center}

\emph{Phần tiếp theo của tài liệu sẽ thảo luận về các công cụ kinh tế
lượng ứng dụng vào chuỗi thời gian vĩ mô (VAR, SVAR) tại Chương 14.}

\begin{center}\rule{0.5\linewidth}{0.5pt}\end{center}

\section{CHƯƠNG 14: KINH TẾ HỌC PHÁT TRIỂN \& THỂ CHẾ (Development
Economics \& Institutional
Economics)}\label{chux1b0ux1a1ng-14-kinh-tux1ebf-hux1ecdc-phuxe1t-triux1ec3n-thux1ec3-chux1ebf-development-economics-institutional-economics}

Trong chương này, chúng ta sẽ khảo sát hai nhánh nghiên cứu gắn bó chặt
chẽ với nhau: Kinh tế học Phát triển và Kinh tế học Thể chế. Đối với
sinh viên bậc Tiến sĩ, việc nghiên cứu sự phát triển không chỉ dừng lại
ở các thước đo định lượng bề mặt như GDP bình quân đầu người, mà còn
phải đi sâu vào cơ chế cấu trúc, đặc điểm vi mô và nền tảng thể chế định
hình quỹ đạo dài hạn của một quốc gia. Những mô hình được giới thiệu
trong chương này đòi hỏi tư duy hệ thống và phân tích định lượng mạnh mẽ
(như đã thảo luận về kinh tế lượng vĩ mô trong Chương 15).

\begin{center}\rule{0.5\linewidth}{0.5pt}\end{center}

\subsection{14.1 Bẫy Nghèo đói \& Mô hình Big Push (Poverty Traps \& Big
Push
Models)}\label{bux1eaby-nghuxe8o-ux111uxf3i-muxf4-huxecnh-big-push-poverty-traps-big-push-models}

\subsubsection{14.1.1 Cốt lõi: Cú huých Lớn và Bẫy Nghèo
đói}\label{cux1ed1t-luxf5i-cuxfa-huuxfdch-lux1edbn-vuxe0-bux1eaby-nghuxe8o-ux111uxf3i}

Khái niệm \textbf{bẫy nghèo đói (poverty trap)} chỉ trạng thái cân bằng
thấp (low-level equilibrium) mà từ đó một nền kinh tế không thể tự thoát
ra thông qua các cơ chế thị trường thông thường do thất bại phối hợp
(coordination failure) và tính bổ trợ (complementarities). Ý tưởng này,
bắt nguồn từ Rosenstein-Rodan (1943) và được hình thức hóa thành
\textbf{Mô hình Big Push} bởi Murphy, Shleifer và Vishny (MSV, 1989),
cho thấy việc công nghiệp hóa thành công đòi hỏi sự chuyển đổi đồng loạt
trên nhiều lĩnh vực.

\textbf{Thiết lập mô hình hình thức:} Giả sử có một nền kinh tế đóng với
\(N\) ngành sản xuất liên tục được đánh chỉ số \(i \in [0, 1]\). Trong
mỗi ngành có hai công nghệ: công nghệ thủ công (cottage) và công nghệ
công nghiệp hiện đại. Hàm sản xuất công nghệ thủ công có lợi tức không
đổi theo quy mô (CRS): \[ x_{1i} = l_{1i} \] với \(l_{1i}\) là lao động
sử dụng trong khu vực thủ công.

Công nghệ hiện đại có lợi tức tăng dần theo quy mô (IRS), đòi hỏi chi
phí cố định (fixed cost) bằng lao động \(F > 0\) và có năng suất biên
cao hơn (\(\alpha > 1\)): \[ x_{2i} = \max(0, \alpha(l_{2i} - F)) \]

Tổng cung lao động là vô thư giãn (inelastic) và bằng \(L\). Tiền lương
được chuẩn hóa thành 1, và tiền lương ở khu vực hiện đại có thể cao hơn,
giả sử là \(w \geq 1\). Để đơn giản, ta xét \(w=1\). Cầu đối với sản
phẩm của mỗi ngành có dạng Cobb-Douglas đối xứng, tức là người tiêu dùng
chi tiêu tỷ lệ thu nhập \(Y\) đều đặn cho mọi hàng hóa \(i\):
\(p_i x_i = Y\).

\textbf{Phân tích Cân bằng:} Giả sử \(\alpha > 1\) sao cho việc công
nghiệp hóa là hiệu quả. Nếu một hãng hiện đại sản xuất, nó sẽ định giá
bằng với chi phí biên của ngành thủ công, tức là \(p_i = 1\) để chiếm
lĩnh thị trường. Lợi nhuận của một hãng độc quyền hiện đại là:
\[ \pi = Y - \frac{Y}{\alpha} - F = Y\left(1 - \frac{1}{\alpha}\right) - F \]

Nếu hãng này là hãng duy nhất hiện đại hóa, thu nhập \(Y = L\) (vì lợi
nhuận các ngành thủ công bằng 0 và tổng lương bằng \(L\)). Khi đó lợi
nhuận của hãng độc quyền duy nhất là:
\[ \pi_1 = L\left(1 - \frac{1}{\alpha}\right) - F \] Nếu \(\pi_1 < 0\),
không hãng đơn lẻ nào muốn hiện đại hóa.

Tuy nhiên, nếu \emph{tất cả} các ngành cùng hiện đại hóa, thu nhập quốc
dân tăng lên bằng \(Y = \alpha(L - F)\). Khi đó lợi nhuận của mỗi hãng
trong trạng thái công nghiệp hóa toàn diện là:
\[ \pi_2 = \alpha(L - F) \left(1 - \frac{1}{\alpha}\right) - F = (L - F)(\alpha - 1) - F = L(\alpha - 1) - \alpha F \]

\textbf{Định lý 1 (Đa cân bằng - Multiple Equilibria):} Nếu \(F\) thỏa
mãn điều kiện
\(L\left(1 - \frac{1}{\alpha}\right) < F < L\left(1 - \frac{1}{\alpha}\right) \alpha\),
thì tồn tại hai cân bằng: 1. \textbf{Cân bằng bẫy nghèo đói (Low-level
equilibrium):} Không có hãng nào hiện đại hóa (\(Y = L\)). 2.
\textbf{Cân bằng công nghiệp hóa (High-level equilibrium):} Tất cả các
hãng đều hiện đại hóa (\(Y = \alpha(L - F)\)).

Sự hiện diện của đa cân bằng chứng minh tầm quan trọng của chính sách
(Cú huých Lớn - Big Push) để điều phối nền kinh tế từ cân bằng thấp lên
cân bằng cao.

\subsubsection{14.1.2 Điểm mù thường
gặp}\label{ux111iux1ec3m-muxf9-thux1b0ux1eddng-gux1eb7p-6}

Một điểm mù lớn khi tiếp cận mô hình Big Push là hiểu sai nguồn gốc của
tính bổ trợ. Tính bổ trợ ở đây không nhất thiết phải là liên kết công
nghệ (input-output linkages) mà xuất phát từ \textbf{tổng cầu (aggregate
demand externality)}. Việc một hãng trả chi phí cố định sẽ làm tăng thu
nhập, từ đó tăng cầu cho \emph{các hãng khác}. Lỗi tư duy phổ biến là
cho rằng ``tự do hóa thị trường'' sẽ giải quyết mọi vấn đề; trên thực
tế, khi đối mặt với đa cân bằng do ngoại ứng cầu, thị trường tự do hoàn
toàn có thể tự khóa mình vào bẫy nghèo (như phân tích đường cong S trong
hàm sản xuất có điểm ngoặt).

\subsubsection{14.1.3 Ứng dụng thực
tế}\label{ux1ee9ng-dux1ee5ng-thux1ef1c-tux1ebf}

\subsubsection{14.1.4 Ví dụ minh họa tính toán (Worked Numerical
Examples)}\label{vuxed-dux1ee5-minh-hux1ecda-tuxednh-touxe1n-worked-numerical-examples-38}

Giả sử nền kinh tế có \(L = 100\) đơn vị lao động. Năng suất khu vực
hiện đại \(\alpha = 2\). Chi phí cố định \(F = 60\). - Nếu một hãng đơn
lẻ tính toán lợi nhuận:
\(\pi_1 = 100(1 - 1/2) - 60 = 50 - 60 = -10 < 0\). Do đó, không ai đầu
tư tiên phong. - Nếu tất cả cùng đầu tư, tổng thu nhập
\(Y = 2(100 - 60) = 80\). Lợi nhuận mỗi hãng
\(\pi_2 = 80(1 - 1/2) - 60 = -20\). Ở đây \(F\) quá lớn, ngay cả Big
Push cũng không hiệu quả.

Tuy nhiên, nếu \(F = 40\): - \(\pi_1 = 100(0.5) - 40 = 10 > 0\). Cân
bằng duy nhất là công nghiệp hóa.

Thú vị nhất là khi \(F = 55\): - \(\pi_1 = 100(0.5) - 55 = -5 < 0\).
(Bẫy nghèo đói) - Nếu tất cả cùng đầu tư:
\(\pi_2 = 2(100 - 55)(0.5) - 55 = 45 - 55 = -10\) (lỗ nếu chỉ có 1 doanh
nghiệp)

\textbf{Trường hợp Big Push (2 doanh nghiệp công nghiệp hóa đồng thời):}

Khi cả hai ngành công nghiệp hóa, thu nhập tăng lên Y = wL = w·(L - 2F),
tạo cầu chéo cho nhau. Lợi nhuận mỗi doanh nghiệp:

\[\Pi_i = (\alpha - c) \cdot L - F = (1 - 0)(100 - 2 \times 55) - 55 = 1 \times (-10) + \text{cầu bổ sung từ lương tăng}\]

Với thu nhập tổng thể tăng, cầu tăng đủ để: \(\Pi_i > 0\) tại cân bằng
phối hợp.

\textbf{Kết luận:} Hai cân bằng Nash tồn tại --- cân bằng thấp (không ai
công nghiệp hóa) và cân bằng cao (cả hai cùng công nghiệp hóa). Big Push
là chính sách phối hợp để chuyển từ bẫy thấp sang cân bằng cao.

\[ \pi = Y(1 - 1/2) - 55 \] Trong cân bằng Big Push,
\(Y = L + N \cdot \pi\). Nếu \(\pi > 0\), \(Y\) tăng. Thu nhập tổng quát
\(Y = wL + \Pi = 100 + \pi\) (với \(\pi\) là tổng lợi nhuận).
\(\pi = (100 + \pi)(0.5) - 55 \implies 0.5\pi = 50 - 55 = -5 \implies \pi < 0\).

Sự phụ thuộc vào ODA ở Sub-Saharan Africa thường được biện minh bằng lý
thuyết bẫy nghèo đói, cho rằng các quốc gia này không đủ tiết kiệm để
vượt ngưỡng \(F\). Một lượng vốn viện trợ khổng lồ (Sachs' Millennium
Villages) được thiết kế để đẩy vốn qua ngưỡng tới hạn, mặc dù thực tiễn
chứng minh kết quả không luôn đồng nhất do bỏ qua yếu tố thể chế
(Easterly).

\begin{center}\rule{0.5\linewidth}{0.5pt}\end{center}

\subsection{14.2 Thể chế \& Tăng trưởng Dài hạn (Institutions \&
Long-run
Growth)}\label{thux1ec3-chux1ebf-tux103ng-trux1b0ux1edfng-duxe0i-hux1ea1n-institutions-long-run-growth}

\subsubsection{14.2.1 Cốt lõi: Nguồn gốc Thể chế và Khung Phân tích
IV}\label{cux1ed1t-luxf5i-nguux1ed3n-gux1ed1c-thux1ec3-chux1ebf-vuxe0-khung-phuxe2n-tuxedch-iv}

Douglass North (1990) định nghĩa \textbf{thể chế (institutions)} là
``luật chơi trong một xã hội'' (rules of the game). Tuy nhiên, để lượng
hóa tác động của thể chế đối với tăng trưởng, giới học thuật vấp phải
thách thức nội sinh (endogeneity). Acemoglu, Johnson và Robinson (AJR,
2001) đã cách mạng hóa lĩnh vực này bằng phương pháp biến công cụ
(Instrumental Variables - IV).

AJR phân biệt giữa \textbf{Thể chế dung hợp (Inclusive institutions)}
bảo vệ quyền tài sản và khuyến khích đổi mới, với \textbf{Thể chế tước
đoạt (Extractive institutions)} tập trung quyền lực và của cải vào tay
giới tinh hoa.

\textbf{Khung ước lượng 2SLS (Two-Stage Least Squares):} Mô hình cấu
trúc (Structural model) để ước lượng tác động của chất lượng thể chế
\(R_i\) đến thu nhập bình quân \(Y_i\):
\[ \ln Y_i = \beta_0 + \beta_1 R_i + \mathbf{X}_i'\gamma + \epsilon_i \]
Tuy nhiên, \(cov(R_i, \epsilon_i) \neq 0\) do nhân quả ngược (các nước
giàu xây dựng thể chế tốt hơn) và biến bị bỏ sót (văn hóa, địa lý).

AJR sử dụng \textbf{tỷ lệ tử vong của thực dân định cư (settler
mortality)}, \(M_i\), làm biến công cụ. Phương trình giai đoạn một
(First-stage):
\[ R_i = \alpha_0 + \alpha_1 \ln M_i + \mathbf{X}_i'\delta + \nu_i \]

Cơ chế nhận diện: Tỷ lệ tử vong sớm \(\rightarrow\) Quyết định định cư
\(\rightarrow\) Thể chế sơ khai (Early institutions) \(\rightarrow\) Thể
chế hiện tại \(\rightarrow\) Thu nhập hiện tại.

\textbf{Định lý 2 (Điều kiện IV hợp lệ):} Để \(\ln M_i\) là biến công cụ
hợp lệ cho 2SLS, nó phải thỏa mãn: 1. \textbf{Sự liên quan (Relevance):}
\(cov(\ln M_i, R_i) \neq 0\) (Kiểm định qua First-stage F-statistic
\textgreater{} 10). 2. \textbf{Hạn chế loại trừ (Exclusion
Restriction):} \(cov(\ln M_i, \epsilon_i) = 0\), nghĩa là tỷ lệ tử vong
trong quá khứ chỉ tác động đến thu nhập hiện tại \emph{thông qua} kênh
thể chế.

\subsubsection{14.2.2 Điểm mù thường
gặp}\label{ux111iux1ec3m-muxf9-thux1b0ux1eddng-gux1eb7p-7}

Lỗi tư duy cốt lõi trong nghiên cứu thể chế là nhầm lẫn hệ quả thiết kế
chính sách với thay đổi thể chế nền tảng. Khi phân tích mô hình AJR,
sinh viên thường quên rằng ``Hạn chế loại trừ'' là một giả định không
thể kiểm định trực tiếp và rất dễ bị vi phạm. Chẳng hạn, tỷ lệ tử vong
bệnh tật lịch sử có thể tương quan với tỷ lệ gánh nặng bệnh tật (disease
burden) hiện đại (như sốt rét), trực tiếp ảnh hưởng đến vốn con người và
thu nhập ngày nay, làm chệch (bias) ước lượng \(\beta_1\) lên trên.

\subsubsection{14.2.3 Ứng dụng thực
tế}\label{ux1ee9ng-dux1ee5ng-thux1ef1c-tux1ebf-1}

So sánh Bắc và Nam Triều Tiên, hay Nogales (Arizona, Mỹ) và Nogales
(Sonora, Mexico) là những ``thí nghiệm tự nhiên'' kinh điển. Dù địa lý,
tài nguyên, văn hóa ban đầu giống hệt nhau, ranh giới thể chế tạo ra sự
phân hóa GDP bình quân gấp nhiều lần.

\subsubsection{14.2.4 Ví dụ minh họa tính toán (Worked Numerical
Examples)}\label{vuxed-dux1ee5-minh-hux1ecda-tuxednh-touxe1n-worked-numerical-examples-39}

Giả sử hệ số \(\alpha_1 = -0.5\) (tăng 1\% tử vong làm giảm 0.5 điểm
chất lượng thể chế). Trong 2SLS, ta thu được giá trị dự báo
\(\hat{R}_i\). Nếu OLS cho \(\hat{\beta}_{OLS} = 0.4\), nhưng IV ước
lượng cho \(\hat{\beta}_{IV} = 0.9\). Khoảng cách
\(\hat{\beta}_{IV} > \hat{\beta}_{OLS}\) thường ám chỉ sai số đo lường
(attenuation bias) trong \(R_i\) lớn hơn ảnh hưởng của biến bị bỏ sót
(vốn thường làm OLS bị chệch dương). Điều này chứng minh tác động thực
sự của thể chế còn mạnh hơn so với đánh giá thông thường.

\begin{center}\rule{0.5\linewidth}{0.5pt}\end{center}

\subsection{14.3 Vốn Con người \& Giáo dục (Human Capital \&
Education)}\label{vux1ed1n-con-ngux1b0ux1eddi-giuxe1o-dux1ee5c-human-capital-education}

\subsubsection{14.3.1 Cốt lõi: Phương trình Mincer và Quyết định Đầu tư
Vốn Con
người}\label{cux1ed1t-luxf5i-phux1b0ux1a1ng-truxecnh-mincer-vuxe0-quyux1ebft-ux111ux1ecbnh-ux111ux1ea7u-tux1b0-vux1ed1n-con-ngux1b0ux1eddi}

Mô hình nền tảng cho vốn con người là \textbf{Phương trình thu nhập
Mincer (Mincer Earnings Equation, 1974)}.

\textbf{Phát triển hình thức:} Giả sử thu nhập \(Y_s\) của một cá nhân
có \(s\) năm đi học. Bỏ qua chi phí trực tiếp, cá nhân trì hoãn đi làm
\(s\) năm. Nếu tỷ suất sinh lời của một năm đi học là \(r\), thì ở trạng
thái cân bằng, hiện giá của chuỗi thu nhập suốt đời phải làm cho cá nhân
bàng quan giữa học tiếp và đi làm ngay. Vi phân:
\(\frac{d Y_s}{ds} = r Y_s \implies \frac{d \ln Y_s}{ds} = r\). Lấy tích
phân, ta có phương trình Mincer cơ bản:
\[ \ln Y_i = \alpha + r S_i + \beta_1 X_i + \beta_2 X_i^2 + u_i \] Trong
đó \(S_i\) là số năm đi học, \(X_i\) là số năm kinh nghiệm làm việc
(kinh nghiệm thường bị suy giảm lợi tức biên nên có dạng bậc 2). \(r\)
đo lường \textbf{Tỷ suất sinh lời của giáo dục (Returns to Education)}.

\textbf{Sự thiên lệch chọn lọc (Selection Bias):} Lại một lần nữa,
\(\mathbb{E}[u_i | S_i] \neq 0\) vì những người có ``năng lực bẩm sinh''
(ability) cao vừa đi học nhiều hơn, vừa kiếm nhiều tiền hơn bất kể học
vấn. Heckman (1979) cung cấp quy trình hiệu chỉnh hai bước (Heckman
correction) dựa trên tỷ lệ nghịch đảo Mills (Inverse Mills Ratio) để xử
lý thiên lệch chọn lọc mẫu trong hàm cung lao động.

Thay vào đó, Angrist \& Krueger (1991) sử dụng \textbf{Quý sinh (Quarter
of Birth)} làm biến công cụ. Trẻ sinh cuối năm đi học sớm hơn và bị luật
giáo dục bắt buộc giữ lại trường lâu hơn trẻ sinh đầu năm, tạo ra biến
thiên ngoại sinh trong \(S_i\).

\subsubsection{14.3.2 Điểm mù thường
gặp}\label{ux111iux1ec3m-muxf9-thux1b0ux1eddng-gux1eb7p-8}

Tranh luận ``Vốn con người'' vs ``Tín hiệu'' (Human Capital vs Signaling
- Spence, 1973; Caplan, 2018). Điểm mù của chính sách là cho rằng \(r\)
phản ánh thuần túy năng suất tăng thêm. Nếu giáo dục chỉ là cơ chế cấp
chứng chỉ (signaling) để phân loại người thông minh, việc trợ cấp giáo
dục đại trà không làm tăng năng suất tổng thể mà chỉ gây ra lạm phát
bằng cấp. Trọng tâm của nghiên cứu sinh Tiến sĩ là sử dụng các thiết kế
thí nghiệm (regression discontinuity) để phân tách (decompose) giá trị
tín hiệu và giá trị năng suất thực.

\subsubsection{14.3.3 Ứng dụng thực
tế}\label{ux1ee9ng-dux1ee5ng-thux1ef1c-tux1ebf-2}

\subsubsection{14.3.4 Ví dụ minh họa tính toán (Worked Numerical
Examples)}\label{vuxed-dux1ee5-minh-hux1ecda-tuxednh-touxe1n-worked-numerical-examples-40}

Trong dữ liệu cắt ngang OLS: \(\ln Y_i = 10.2 + 0.08 S_i\). (Tỷ suất
sinh lời 8\%/năm). Tuy nhiên, khi dùng luật giáo dục bắt buộc làm IV,
kết quả IV thường dao động từ \(0.10 - 0.12\). Tỷ suất sinh lời IV lớn
hơn OLS ngụ ý rằng: quy định bắt buộc đi học tác động chủ yếu đến nhóm
cận biên (margin) - những trẻ có chi phí chiết khấu rất cao hoặc bị gia
đình cản trở. Đối với nhóm \textbf{LATE (Local Average Treatment
Effect)} này, lợi ích của việc hoàn thành trung học là vô cùng lớn so
với ước lượng trung bình của dân số (ATE).

\begin{center}\rule{0.5\linewidth}{0.5pt}\end{center}

\subsection{14.4 Bất bình đẳng \& Đường cong Kuznets (Inequality \&
Kuznets
Curve)}\label{bux1ea5t-buxecnh-ux111ux1eb3ng-ux111ux1b0ux1eddng-cong-kuznets-inequality-kuznets-curve}

\subsubsection{14.4.1 Cốt lõi: Phân phối Thu nhập và Động học R và
G}\label{cux1ed1t-luxf5i-phuxe2n-phux1ed1i-thu-nhux1eadp-vuxe0-ux111ux1ed9ng-hux1ecdc-r-vuxe0-g}

Simon Kuznets (1955) đặt ra giả thuyết \textbf{Đường cong Kuznets chữ U
ngược (Inverted-U hypothesis)}: Bất bình đẳng tăng trong giai đoạn đầu
của công nghiệp hóa khi người dân chuyển từ nông nghiệp (nghèo, bình
đẳng) sang công nghiệp (giàu, phân hóa), sau đó giảm xuống nhờ thuế lũy
tiến và giáo dục.

Để đo lường, ta dùng \textbf{Đường cong Lorenz} (biểu diễn tỷ lệ tích
lũy của thu nhập so với tỷ lệ tích lũy của dân số) và \textbf{Hệ số
Gini}. Gọi \(L(p)\) là đường cong Lorenz, hệ số Gini \(G\) là tỷ lệ diện
tích giữa đường phân giác 45 độ và đường Lorenz:
\[ G = 1 - 2\int_{0}^{1} L(p) dp \]

Trong thời đại mới, Piketty (2014) lật đổ giả thuyết Kuznets bằng phương
trình động học bất bình đẳng: \textbf{\(r > g\)}. Giả sử \(r\) là tỷ
suất lợi nhuận trung bình từ vốn (capital), và \(g\) là tốc độ tăng
trưởng của thu nhập quốc dân. Phân thức vốn trong nền kinh tế
\(\alpha_K = r \times \frac{K}{Y}\). Khi tỷ lệ tiết kiệm \(s\) và tăng
trưởng \(g\) không đổi, tỷ lệ vốn-thu nhập dài hạn (luật tư bản thứ hai
của Piketty) hội tụ về: \[ \beta = \frac{K}{Y} = \frac{s}{g} \] Nếu
\(r > g\), của cải tích lũy từ quá khứ (kế thừa) sẽ tăng nhanh hơn thu
nhập từ sức lao động, dẫn đến bất bình đẳng tài sản tiến đến mức báo
động như thời kỳ Gilded Age, phá vỡ tính hội tụ của Kuznets.

\subsubsection{14.4.2 Điểm mù thường
gặp}\label{ux111iux1ec3m-muxf9-thux1b0ux1eddng-gux1eb7p-9}

Người học thường nhầm lẫn giữa bất bình đẳng \emph{thu nhập} (income) và
bất bình đẳng \emph{của cải} (wealth). Hệ số Gini về thu nhập thường
khoảng 0.3 - 0.5, nhưng Gini về của cải luôn cao hơn hẳn (0.7 - 0.9).
Điểm mù tiếp theo là bỏ qua giả định độ co giãn thay thế giữa vốn và lao
động (\(\sigma\)). Nếu \(\sigma > 1\), khi tích lũy vốn (\(K/Y\) tăng),
\(r\) không giảm đủ nhanh, làm phần chia của vốn \(\alpha_K\) tăng lên.
Nếu \(\sigma < 1\), sự hội tụ phân phối sẽ diễn ra nhanh chóng.

\subsubsection{14.4.3 Ứng dụng thực
tế}\label{ux1ee9ng-dux1ee5ng-thux1ef1c-tux1ebf-3}

\subsubsection{14.4.4 Ví dụ minh họa tính toán (Worked Numerical
Examples)}\label{vuxed-dux1ee5-minh-hux1ecda-tuxednh-touxe1n-worked-numerical-examples-41}

Một quốc gia phát triển có tăng trưởng dài hạn \(g = 1.5\%\). Tỷ lệ tiết
kiệm \(s = 9\%\). Tỷ lệ vốn-thu nhập hội tụ \(\beta = 0.09 / 0.015 = 6\)
(Tài sản quốc gia gấp 6 lần GDP năm). Nếu tỷ suất sinh lời vốn
\(r = 5\%\). Phần chia của vốn
\(\alpha_K = r \times \beta = 0.05 \times 6 = 30\%\). Nếu \(g\) giảm
xuống 1\%, \(\beta = 9\), \(\alpha_K = 0.05 \times 9 = 45\%\). Điều này
chứng tỏ tăng trưởng chậm lại (secular stagnation) ở phương Tây làm trầm
trọng thêm sự tập trung của cải, đòi hỏi chính sách thuế tài sản toàn
cầu (global wealth tax).

\begin{center}\rule{0.5\linewidth}{0.5pt}\end{center}

\subsection{14.5 Tài chính Vi mô \& RCT trong Kinh tế Phát triển
(Microfinance \&
RCTs)}\label{tuxe0i-chuxednh-vi-muxf4-rct-trong-kinh-tux1ebf-phuxe1t-triux1ec3n-microfinance-rcts}

\subsubsection{14.5.1 Cốt lõi: Lựa chọn Ngẫu nhiên và Phân tích
ATE}\label{cux1ed1t-luxf5i-lux1ef1a-chux1ecdn-ngux1eabu-nhiuxean-vuxe0-phuxe2n-tuxedch-ate}

Mô hình Ngân hàng Grameen của Muhammad Yunus (Nobel 2006) dùng bảo lãnh
chéo (joint-liability) thay cho thế chấp để vượt qua rủi ro đạo đức
(moral hazard) và lựa chọn ngược (adverse selection) trong tín dụng vi
mô. Tuy nhiên, đánh giá hiệu quả thực sự đòi hỏi phương pháp \textbf{Thử
nghiệm ngẫu nhiên có đối chứng (Randomized Controlled Trials - RCTs)} do
Banerjee, Duflo, Kremer (Nobel 2019) tiên phong.

\textbf{Hình thức hóa Lợi ứng Xử lý Trung bình (ATE - Average Treatment
Effect):} Theo Khung hệ quả tiềm năng (Potential Outcomes Framework) của
Rubin: Gọi \(Y_{1i}\) là kết cục nếu cá nhân \(i\) được can thiệp (ví
dụ, nhận vi mô), và \(Y_{0i}\) nếu không được can thiệp. Biến trạng thái
\(D_i \in \{0,1\}\). Kết cục quan sát được:
\(Y_i = D_i Y_{1i} + (1 - D_i) Y_{0i}\).

Ta muốn tính \(ATE = \mathbb{E}[Y_{1i} - Y_{0i}]\). Sự chênh lệch thô
(Naive difference) là:
\[ \mathbb{E}[Y_i | D_i = 1] - \mathbb{E}[Y_i | D_i = 0] = \underbrace{\mathbb{E}[Y_{1i} - Y_{0i} | D_i = 1]}_{\text{ATT}} + \underbrace{\mathbb{E}[Y_{0i} | D_i = 1] - \mathbb{E}[Y_{0i} | D_i = 0]}_{\text{Selection Bias}} \]

Thông qua phân công ngẫu nhiên trong RCT, \(D_i\) độc lập với
\((Y_{1i}, Y_{0i})\), do đó Selection Bias = 0, và khác biệt thô chính
xác bằng ATE. Khi có sự không tuân thủ (non-compliance), thiết kế RCT
chuyển sang sử dụng phân công ban đầu (\(Z_i\)) làm biến công cụ cho sự
nhận can thiệp thực tế (\(D_i\)) để ước lượng \textbf{LATE (Local
Average Treatment Effect)} cho nhóm Compliers.
\[ LATE = \frac{\mathbb{E}[Y_i | Z_i = 1] - \mathbb{E}[Y_i | Z_i = 0]}{\mathbb{P}(D_i = 1 | Z_i = 1) - \mathbb{P}(D_i = 1 | Z_i = 0)} \]
Tử số chính là Intent-to-Treat (ITT) effect.

\subsubsection{14.5.2 Điểm mù thường
gặp}\label{ux111iux1ec3m-muxf9-thux1b0ux1eddng-gux1eb7p-10}

Điểm mù lớn nhất trong RCT là ranh giới giữa \textbf{Giá trị nội suy
(Internal validity)} và \textbf{Giá trị ngoại suy (External validity)}.
RCT tại một vùng nông thôn Kenya có internal validity cực cao nhờ ngẫu
nhiên hóa hoàn hảo. Nhưng liệu ta có thể áp dụng chính sách đó cho Ấn Độ
(external validity)? Hơn nữa, RCT ở quy mô nhỏ có thể bỏ qua \textbf{Cân
bằng tổng thể (General equilibrium effects)}: Cung cấp tín dụng cho vài
hộ giúp họ thoát nghèo; cung cấp tín dụng cho toàn làng có thể làm giảm
giá nông sản do bùng nổ nguồn cung, khiến thu nhập thực không đổi.

\subsubsection{14.5.3 Ứng dụng thực
tế}\label{ux1ee9ng-dux1ee5ng-thux1ef1c-tux1ebf-4}

\subsubsection{14.5.4 Ví dụ minh họa tính toán (Worked Numerical
Examples)}\label{vuxed-dux1ee5-minh-hux1ecda-tuxednh-touxe1n-worked-numerical-examples-42}

Một quỹ vi mô tung ra chiến dịch marketing tín dụng (\(Z_i = 1\)) ở 50
làng và giữ 50 làng đối chứng (\(Z_i = 0\)). Trong làng được can thiệp,
tỷ lệ vay (\(D_i=1\)) là 40\%. Lợi nhuận kinh doanh trung bình
\(\mathbb{E}[Y|Z=1] = 120\). Trong làng đối chứng, tỷ lệ vay tự tìm đến
(\(D_i=1\)) là 10\%. Lợi nhuận trung bình \(\mathbb{E}[Y|Z=0] = 105\).

\begin{itemize}
\tightlist
\item
  Hiệu ứng ITT (Intent-to-treat) = \(120 - 105 = 15\).
\item
  Tỷ lệ tuân thủ (Kênh lan truyền) = \(0.40 - 0.10 = 0.30\).
\item
  LATE = \(15 / 0.30 = 50\). Kết luận: Tín dụng vi mô thực sự tạo ra gia
  tăng lợi nhuận 50 đơn vị cho những hộ \emph{bị thay đổi quyết định}
  nhờ chiến dịch tiếp thị.
\end{itemize}

\begin{center}\rule{0.5\linewidth}{0.5pt}\end{center}

\subsection{14.6 Kinh tế Chính trị của Phát triển (Political Economy of
Development)}\label{kinh-tux1ebf-chuxednh-trux1ecb-cux1ee7a-phuxe1t-triux1ec3n-political-economy-of-development}

\subsubsection{14.6.1 Cốt lõi: Tìm kiếm Đặc lợi và Lời nguyền Tài
nguyên}\label{cux1ed1t-luxf5i-tuxecm-kiux1ebfm-ux111ux1eb7c-lux1ee3i-vuxe0-lux1eddi-nguyux1ec1n-tuxe0i-nguyuxean}

Tại sao các quốc gia thi hành các chính sách kinh tế kém hiệu quả dù đã
có mô hình tối ưu? Mancur Olson (1965) chỉ ra vấn đề hành động tập thể
(collective action): Các nhóm lợi ích nhỏ có cấu kết cao dễ dàng vận
động hành lang tước đoạt lợi ích của số đông phân tán (người tiêu
dùng/free-rider).

Sự bóp méo này thể hiện rõ qua \textbf{Tìm kiếm Đặc lợi (Rent-seeking)}
(Tullock, 1967; Krueger, 1974). Thay vì cạnh tranh sản xuất, nguồn lực
bị ném vào cạnh tranh quyền lực chính trị.

\textbf{Mô hình Tìm kiếm Đặc lợi (Tullock Contest):} Giả sử có 2 nhóm,
mỗi nhóm đầu tư tài nguyên \(x_i\) (\(i \in \{1,2\}\)) để giành được một
khoản đặc quyền (rent) trị giá \(V\). Xác suất chiến thắng của nhóm 1
(Hàm thành công Tullock) là:
\[ p_1(x_1, x_2) = \frac{x_1^\gamma}{x_1^\gamma + x_2^\gamma} \] với
\(\gamma > 0\) là thông số đo lường mức độ quyết định của tiền bạc đối
với quyền lực (discriminatory power).

Hàm mục tiêu của người chơi \(i\): \[ \max_{x_i} \pi_i = p_i V - x_i \]
Điều kiện bậc nhất (FOC) đối với nhóm 1 (giả sử đối xứng
\(x_1 = x_2 = x^*\)):
\[ \frac{\partial p_1}{\partial x_1} V - 1 = 0 \implies \frac{\gamma x_1^{\gamma-1} x_2^\gamma}{(x_1^\gamma + x_2^\gamma)^2} V = 1 \]
Tại cân bằng đối xứng:
\[ \frac{\gamma (x^*)^{2\gamma - 1}}{(2(x^*)^\gamma)^2} V = \frac{\gamma}{4 x^*} V = 1 \implies x^* = \frac{\gamma V}{4} \]
Tổng nguồn lực lãng phí cho xã hội là
\(\Sigma x_i = \frac{\gamma V}{2}\). Nếu \(\gamma = 2\), toàn bộ giá trị
\(V\) bị tiêu hao vào lobby (Full rent dissipation).

Điều này giải thích \textbf{Lời nguyền tài nguyên (Resource Curse)} hay
\textbf{Căn bệnh Hà Lan (Dutch Disease)}: Khi có mỏ dầu, \(V\) khổng lồ.
Vốn và nhân tài chuyển dịch từ khu vực chế tạo (có lợi tức tăng dần và
học hỏi bằng thực hành) sang phe cánh chính trị, kìm hãm phát triển dài
hạn, như lập luận của Acemoglu \& Robinson (2006).

\subsubsection{14.6.2 Điểm mù thường
gặp}\label{ux111iux1ec3m-muxf9-thux1b0ux1eddng-gux1eb7p-11}

Hiểu sai phổ biến là đánh đồng tham nhũng/rent-seeking ở mọi quốc gia có
cùng hệ quả. Shleifer và Vishny phân biệt tham nhũng phi tập trung (độc
lập, nhiều điểm chốt - như Nga thập niên 90) và tham nhũng tập trung
(một cửa - như các quốc gia Đông Á). Tham nhũng phi tập trung gây hại
tàn khốc hơn do mỗi quan chức đối xử người hối lộ như mồi nhện độc quyền
(tragedies of the anticommons).

\subsubsection{14.6.3 Ứng dụng thực
tế}\label{ux1ee9ng-dux1ee5ng-thux1ef1c-tux1ebf-5}

Các quốc gia phụ thuộc xuất khẩu dầu thô như Venezuela là ví dụ điển
hình về Lời nguyền tài nguyên. Biến động giá dầu tạo ra cú sốc ngoại
sinh. Viện trợ ODA (Foreign Aid) có điều kiện (conditionality) thường
thất bại do nước nhận xem đó như một nguồn Rent \(V\). Các nhà lãnh đạo
trích xuất aid để phân phát cho phe cánh duy trì quyền lực, khiến
\(\gamma\) (sự khốc liệt của cạnh tranh chính trị) tăng cao, bào mòn
chính cấu trúc kinh tế nền tảng thay vì thúc đẩy đầu tư.

\begin{center}\rule{0.5\linewidth}{0.5pt}\end{center}

\subsection{14.7 Di cư \& Chuyển tiền (Migration \&
Remittances)}\label{di-cux1b0-chuyux1ec3n-tiux1ec1n-migration-remittances}

\subsubsection{14.7.1 Cốt lõi: Lựa chọn Tự quyết và Cân bằng
Harris-Todaro}\label{cux1ed1t-luxf5i-lux1ef1a-chux1ecdn-tux1ef1-quyux1ebft-vuxe0-cuxe2n-bux1eb1ng-harris-todaro}

Quá trình phát triển đồng hành với sự dịch chuyển lao động. Trong nội
địa, đó là di cư nông thôn - thành thị, được giải thích qua \textbf{Mô
hình Harris-Todaro (1970)}. Trong bối cảnh quốc tế, nó đặt ra bài toán
Chảy máu chất xám (Brain drain) và \textbf{Mô hình Roy (1951)} về tự
chọn lọc (self-selection).

\textbf{Mô hình Harris-Todaro:} Trái với cân bằng thị trường cạnh tranh
hoàn hảo, thị trường lao động đô thị thường bị cứng nhắc (tiền lương tối
thiểu do công đoàn/luật pháp), \(w_U > w_R\) (lương nông thôn). Sự thất
nghiệp ở thành thị là cơ chế điều chỉnh cân bằng. Cân bằng đạt được
không phải khi \(w_U = w_R\), mà khi \textbf{Kỳ vọng tiền lương} ngang
nhau. Xác suất tìm được việc làm ở đô thị \(p = \frac{L_U}{L_U + U}\),
với \(L_U\) là việc làm có sẵn, \(U\) là người thất nghiệp. Điều kiện
cân bằng di cư:
\[ w_R = p \cdot w_U + (1 - p) \cdot 0 \implies w_R = \frac{L_U}{L_U + U} w_U \]
Một nghịch lý chính sách (Todaro paradox): Việc tạo thêm việc làm ở
thành thị (tăng \(L_U\)) có thể làm \emph{tăng} số lượng người thất
nghiệp thành thị do lực lượng nông thôn di cư ồ ạt để tranh giành cơ hội
cao hơn.

\textbf{Mô hình Roy và Di cư Quốc tế:} Một lao động có kỹ năng
\(s \sim N(\mu_s, \sigma_s^2)\). Thu nhập ở quốc gia gốc (Source):
\(\ln w_0 = \mu_0 + r_0 s\). Thu nhập ở quốc gia đích (Destination):
\(\ln w_1 = \mu_1 + r_1 s\). Chi phí di cư: \(C\). Cá nhân di cư nếu
\(\ln w_1 - \ln w_0 > C\). - Nếu \(r_1 > r_0\) (nước đích trọng dụng kỹ
năng hơn, ví dụ Mỹ), di cư mang tính chọn lọc dương (Positive Selection)
- những người giỏi nhất ra đi (Brain Drain). - Nếu \(r_1 < r_0\) (nước
đích có mạng lưới an sinh xã hội cao đánh thuế nặng, hoặc công việc tay
chân), di cư mang tính chọn lọc âm (Negative Selection) - lao động kỹ
năng thấp có động lực di cư lớn hơn.

\subsubsection{14.7.2 Điểm mù thường
gặp}\label{ux111iux1ec3m-muxf9-thux1b0ux1eddng-gux1eb7p-12}

Nhiều nhà hoạch định chính sách mặc định Brain Drain hoàn toàn là tiêu
cực. Quan điểm mới (Stark, 2004) chỉ ra ``Brain Gain'': Chính cơ hội
được xuất khẩu lao động (hoặc du học) với tiền lương tiềm năng khổng lồ
\(\mu_1\) tạo động lực cho \emph{tất cả} mọi người ở quê nhà tích lũy
vốn con người \(s\). Vì không phải ai học xong cũng có thể hoặc muốn đi,
phần kỹ năng thặng dư ở lại trong nước có thể lớn hơn lượng thất thoát,
bù đắp cho Brain Drain.

\subsubsection{14.7.3 Ứng dụng thực
tế}\label{ux1ee9ng-dux1ee5ng-thux1ef1c-tux1ebf-6}

\subsubsection{14.7.4 Ví dụ minh họa tính toán (Worked Numerical
Examples)}\label{vuxed-dux1ee5-minh-hux1ecda-tuxednh-touxe1n-worked-numerical-examples-43}

Lương nông nghiệp \(w_R = 100\) USD. Lương tối thiểu đô thị
\(w_U = 200\) USD. Theo điều kiện cân bằng:
\(100 = p \cdot 200 \implies p = 0.5\). Nếu đô thị có 1000 việc làm, sẽ
có 2000 người chen chúc ở đô thị (1000 có việc, 1000 thất nghiệp). Nếu
chính phủ lập gói kích cầu tạo thêm 200 việc làm ở đô thị
(\(L_U = 1200\)). Đang có \(p = 1200 / 2000 = 0.6 > 0.5\), kỳ vọng lương
đô thị cao hơn nông thôn, dòng di cư tiếp tục nổ ra. Cân bằng mới:
\(p = 1200 / (1200 + U') = 0.5 \implies 1200 + U' = 2400 \implies U' = 1200\).
Vậy, tạo thêm 200 việc làm mới lại làm số người thất nghiệp tăng thêm
200 người!

Đồng thời, sự xuất khẩu lao động và lượng kiều hối (Remittances) khổng
lồ ở các nước như Việt Nam, Philippines đóng vai trò như một quỹ bảo
hiểm vi mô, làm giảm hạn chế tín dụng (credit constraints) cho đầu tư
kinh doanh ở nông thôn, thúc đẩy kinh tế vi mô mạnh mẽ.

\begin{center}\rule{0.5\linewidth}{0.5pt}\end{center}

\emph{Ghi chú: Việc giải các bài toán về bẫy nghèo đói và hành vi người
tìm việc đan xen chặt chẽ với các phương pháp quy hoạch động (xem Mục
15.8 và Phụ lục Bellman)).}

\begin{center}\rule{0.5\linewidth}{0.5pt}\end{center}

\section{CHƯƠNG 15: KINH TẾ LƯỢNG \& SUY LUẬN NHÂN QUẢ (Econometrics \&
Causal
Inference)}\label{chux1b0ux1a1ng-15-kinh-tux1ebf-lux1b0ux1ee3ng-suy-luux1eadn-nhuxe2n-quux1ea3-econometrics-causal-inference}

Kinh tế lượng hiện đại không chỉ dừng lại ở việc thiết lập các mô hình
dự báo, mà trọng tâm đã dịch chuyển sang \textbf{suy luận nhân quả
(causal inference)}. Chương này cung cấp nền tảng phương pháp luận cho
kinh tế học thực nghiệm, đi từ các nguyên lý cơ bản của OLS đến các
phương pháp thiết kế tự nhiên (quasi-experimental designs) tiên tiến
nhất trong ``Cuộc cách mạng Độ tin cậy'' (Credibility Revolution).

\begin{center}\rule{0.5\linewidth}{0.5pt}\end{center}

\subsection{15.1 Hồi quy Tuyến tính \& Giả định OLS (Linear Regression
\& OLS
Assumptions)}\label{hux1ed3i-quy-tuyux1ebfn-tuxednh-giux1ea3-ux111ux1ecbnh-ols-linear-regression-ols-assumptions}

\subsubsection{15.1.1 Cốt lõi Lý thuyết (Core
Theory)}\label{cux1ed1t-luxf5i-luxfd-thuyux1ebft-core-theory-46}

Phương pháp bình phương tối thiểu thông thường (Ordinary Least Squares -
OLS) là nền tảng của kinh tế lượng. Giả sử ta có mô hình hồi quy tổng
thể dưới dạng ma trận: \[ Y = X\beta + u \] Trong đó \(Y\) là vector
\((n \times 1)\) của biến phụ thuộc, \(X\) là ma trận \((n \times k)\)
của các biến độc lập, \(\beta\) là vector \((k \times 1)\) các hệ số, và
\(u\) là vector sai số ngẫu nhiên \((n \times 1)\).

Bài toán tối ưu của OLS là tìm \(\hat{\beta}\) sao cho tổng bình phương
phần dư (SSR) là nhỏ nhất:
\[ \min_{\hat{\beta}} (Y - X\hat{\beta})'(Y - X\hat{\beta}) \] Đạo hàm
bậc nhất theo \(\hat{\beta}\) và cho bằng 0 (First Order Condition), ta
thu được \textbf{hệ số ước lượng OLS}:
\[ \hat{\beta}_{OLS} = (X'X)^{-1}X'Y \] \textbf{Định lý Gauss-Markov}
chỉ ra rằng, dưới 5 giả định cổ điển: 1. Tuyến tính trong tham số
(Linearity in parameters) 2. Lấy mẫu ngẫu nhiên (Random sampling) 3.
Không có cộng tuyến hoàn hảo (No perfect multicollinearity): Điều kiện
\(X'X\) khả nghịch. 4. Trung bình sai số có điều kiện bằng 0 (Zero
conditional mean): \(\mathbb{E}(u|X) = 0\) 5. Phương sai sai số không
đổi và không có tự tương quan (Homoskedasticity and No autocorrelation):
\(Var(u|X) = \sigma^2 I_n\)

Khi đó, \(\hat{\beta}_{OLS}\) là \textbf{BLUE} (Best Linear Unbiased
Estimator - Ước lượng tuyến tính không chệch tốt nhất). Ma trận hiệp
phương sai của \(\hat{\beta}\) được tính bằng:
\[ Var(\hat{\beta}|X) = \sigma^2 (X'X)^{-1} \] Nếu giả định (5) bị vi
phạm (phương sai sai số thay đổi - heteroskedasticity), ước lượng OLS
vẫn không chệch nhưng không còn hiệu quả (không còn là ``Best''), đòi
hỏi sử dụng phương sai chuẩn vững (robust standard errors) của
White/Huber.

\subsubsection{15.1.2 Điểm mù thường gặp (Cognitive Traps \&
Limitations)}\label{ux111iux1ec3m-muxf9-thux1b0ux1eddng-gux1eb7p-cognitive-traps-limitations-97}

Nhiều nhà nghiên cứu quá chú trọng vào việc giá trị \(R^2\) cao mà bỏ
qua giả định \(\mathbb{E}(u|X) = 0\). \(R^2\) đo lường độ phù hợp của dữ
liệu, nhưng một mô hình có \(R^2 = 0.99\) vẫn có thể hoàn toàn vô nghĩa
về mặt nhân quả nếu biến \(X\) nội sinh (endogenous). Đừng nhầm lẫn giữa
sự phù hợp mô hình (goodness-of-fit) và sự nhận dạng nhân quả (causal
identification).

\subsubsection{15.1.3 Ứng dụng thực tế (Practical
Applications)}\label{ux1ee9ng-dux1ee5ng-thux1ef1c-tux1ebf-practical-applications-53}

\subsubsection{15.1.4 Ví dụ minh họa tính toán (Worked Numerical
Examples)}\label{vuxed-dux1ee5-minh-hux1ecda-tuxednh-touxe1n-worked-numerical-examples-44}

Khảo sát 1,000 công nhân, ta hồi quy
\(ln(Wage) = \beta_0 + \beta_1 \times Education + \beta_2 \times Experience + u\).
Giả sử kết quả thu được \(\hat{\beta}_1 = 0.08\). Điều này có nghĩa là
(giả định các yếu tố khác không đổi), thêm một năm học vấn có tương quan
với mức lương tăng thêm xấp xỉ 8\%. Tuy nhiên, liệu đây có phải là hiệu
ứng nhân quả? Nếu năng lực bẩm sinh (Ability) không được đưa vào mô
hình, OLS sẽ bị chệch.

\begin{center}\rule{0.5\linewidth}{0.5pt}\end{center}

\subsection{15.2 Sai lệch Biến thiếu \& Nội sinh (Omitted Variable Bias
\&
Endogeneity)}\label{sai-lux1ec7ch-biux1ebfn-thiux1ebfu-nux1ed9i-sinh-omitted-variable-bias-endogeneity}

\subsubsection{15.2.1 Cốt lõi Lý thuyết (Core
Theory)}\label{cux1ed1t-luxf5i-luxfd-thuyux1ebft-core-theory-47}

Biến nội sinh (Endogeneity) xảy ra khi \(\mathbb{E}(u|X) \neq 0\). Điều
này có thể đến từ ba nguồn chính: Sai lệch biến thiếu (Omitted
Variables), Sai số đo lường (Measurement Error), và Tính đồng thời
(Simultaneity).

\textbf{Sai lệch biến thiếu (OVB):} Giả sử mô hình thực (Population
Regression Function - PRF) là:
\[ Y = \beta_0 + \beta_1 X_1 + \beta_2 X_2 + v \] Nhưng ta lại ước lượng
mô hình thiếu biến \(X_2\): \[ Y = \alpha_0 + \alpha_1 X_1 + u \] Trong
đó \(u = \beta_2 X_2 + v\). Hệ số ước lượng OLS cho \(\alpha_1\) (ký
hiệu là \(\tilde{\beta}_1\)) sẽ hội tụ về (probability limit):
\[ plim(\tilde{\beta}_1) = \beta_1 + \beta_2 \frac{Cov(X_1, X_2)}{Var(X_1)} = \beta_1 + \beta_2 \delta_{1} \]
Với \(\delta_1\) là hệ số từ hồi quy phụ (auxiliary regression) của
\(X_2\) lên \(X_1\). OVB chính là đại lượng \(\beta_2 \delta_{1}\).

\subsubsection{15.2.2 Điểm mù thường gặp (Cognitive Traps \&
Limitations)}\label{ux111iux1ec3m-muxf9-thux1b0ux1eddng-gux1eb7p-cognitive-traps-limitations-98}

Có một quan niệm sai lầm phổ biến là ``càng thêm nhiều biến kiểm soát
càng tốt'' (controlling for everything). Thêm các biến kiểm soát nội
sinh, hoặc tồi tệ hơn, các ``biến kết quả trung gian'' (bad controls /
colliders), không những không giải quyết được vấn đề nội sinh mà còn gây
ra sai lệch (collider bias).

\subsubsection{15.2.3 Ứng dụng thực tế (Practical
Applications)}\label{ux1ee9ng-dux1ee5ng-thux1ef1c-tux1ebf-practical-applications-54}

\textbf{Phân tích chiều sai lệch (Direction of Bias):} Trong ví dụ lợi
tức giáo dục, giả sử \(Y\) là Lương, \(X_1\) là Học vấn (Education),
\(X_2\) là Năng lực (Ability, không quan sát được). Ta biết
\(\beta_2 > 0\) (Năng lực cao dẫn đến lương cao) và \(\delta_1 > 0\)
(Người có năng lực cao thường học nhiều hơn). Vậy
\(OVB = \beta_2 \delta_1 > 0\). Suy ra, ước lượng OLS đơn thuần thường
\emph{phóng đại} (upward bias) lợi ích thực sự của giáo dục.

\begin{center}\rule{0.5\linewidth}{0.5pt}\end{center}

\subsection{15.3 Biến Công cụ \& Ước lượng 2SLS (Instrumental Variables
\&
2SLS)}\label{biux1ebfn-cuxf4ng-cux1ee5-ux1b0ux1edbc-lux1b0ux1ee3ng-2sls-instrumental-variables-2sls}

\subsubsection{15.3.1 Cốt lõi Lý thuyết (Core
Theory)}\label{cux1ed1t-luxf5i-luxfd-thuyux1ebft-core-theory-48}

Khi \(X\) nội sinh, OLS không nhất quán (inconsistent). Phương pháp biến
công cụ (IV) khắc phục bằng cách sử dụng một biến \(Z\) thỏa mãn hai
điều kiện nghiêm ngặt: 1. \textbf{Tính thích hợp (Relevance):}
\(Cov(Z, X) \neq 0\) (Biến công cụ có tương quan mạnh với biến nội
sinh). 2. \textbf{Tính ngoại sinh (Exclusion Restriction):}
\(Cov(Z, u) = 0\) (Biến công cụ không tương quan với sai số; nó chỉ tác
động lên Y thông qua X).

Ước lượng \textbf{Bình phương tối thiểu hai giai đoạn (2SLS)} được tiến
hành như sau: * \textbf{Giai đoạn 1 (First Stage):} Hồi quy biến nội
sinh \(X\) lên \(Z\) (và các biến ngoại sinh khác \(W\)):
\[ X = \pi_0 + \pi_1 Z + \gamma W + v \] Lấy giá trị dự báo \(\hat{X}\).
* \textbf{Giai đoạn 2 (Second Stage):} Hồi quy \(Y\) lên \(\hat{X}\) (và
\(W\)): \[ Y = \beta_0 + \beta_{IV} \hat{X} + \alpha W + u \]

Dưới dạng ma trận, với \(P_Z = Z(Z'Z)^{-1}Z'\) là ma trận chiếu:
\[ \hat{\beta}_{2SLS} = (X' P_Z X)^{-1} X' P_Z Y \] Trong trường hợp một
công cụ \(Z\) và một biến nội sinh \(X\), ước lượng IV chính là
\textbf{Ước lượng Wald}:
\[ \hat{\beta}_{IV} = \frac{Cov(Y, Z)}{Cov(X, Z)} = \frac{\text{Reduced Form}}{\text{First Stage}} \]
Theo khuôn khổ Imbens \& Angrist (1994), nếu tác động không đồng nhất
(heterogeneous treatment effects), IV ước lượng \textbf{LATE} (Local
Average Treatment Effect) - hiệu ứng nhân quả cho nhóm người phản ứng
với biến công cụ (compliers).

\subsubsection{15.3.2 Điểm mù thường gặp (Cognitive Traps \&
Limitations)}\label{ux111iux1ec3m-muxf9-thux1b0ux1eddng-gux1eb7p-cognitive-traps-limitations-99}

\textbf{Vấn đề Công cụ yếu (Weak Instruments):} Nếu \(Z\) và \(X\) có
tương quan yếu (First-stage F-statistic \textless{} 10, theo quy tắc
ngón cái của Stock-Yogo), 2SLS không chỉ mất tính chính xác mà hệ số còn
bị chệch trầm trọng theo chiều hướng của OLS, đặc biệt trong các mẫu hữu
hạn (finite samples).

\subsubsection{15.3.3 Ứng dụng thực tế (Practical
Applications)}\label{ux1ee9ng-dux1ee5ng-thux1ef1c-tux1ebf-practical-applications-55}

\textbf{Acemoglu, Johnson \& Robinson (AJR, 2001):} Để ước lượng tác
động của Thể chế (Institutions - biến nội sinh) lên Thu nhập (Income),
AJR sử dụng ``Tỷ lệ tử vong của người định cư'' (Settler Mortality) làm
biến công cụ. - \emph{First Stage:} Tử vong định cư cao \(\rightarrow\)
Mô hình khai thác, thể chế kém. (\(Cov(Z,X) \neq 0\)) - \emph{Exclusion
Restriction:} Tỷ lệ tử vong ở thế kỷ 19 không trực tiếp ảnh hưởng đến
GDP ngày nay ngoại trừ thông qua kênh thể chế.

\begin{center}\rule{0.5\linewidth}{0.5pt}\end{center}

\subsection{15.4 Khác biệt trong Khác biệt (Difference-in-Differences -
DiD)}\label{khuxe1c-biux1ec7t-trong-khuxe1c-biux1ec7t-difference-in-differences---did}

\subsubsection{15.4.1 Cốt lõi Lý thuyết (Core
Theory)}\label{cux1ed1t-luxf5i-luxfd-thuyux1ebft-core-theory-49}

DiD tận dụng thực nghiệm tự nhiên (Natural Experiments) bằng cách so
sánh sự thay đổi của nhóm xử lý (Treatment Group) trước và sau can thiệp
với sự thay đổi tương ứng của nhóm đối chứng (Control Group). Mô hình
hồi quy kinh điển cho DiD:
\[ Y_{it} = \alpha + \beta D_i + \gamma T_t + \delta_{DiD} (D_i \times T_t) + \epsilon_{it} \]
- \(D_i = 1\) nếu cá nhân \(i\) thuộc nhóm xử lý. - \(T_t = 1\) nếu thời
gian ở giai đoạn sau can thiệp (Post). - Hê số \(\delta_{DiD}\) là hiệu
ứng nhân quả (Treatment Effect).

Ước lượng của DiD:
\[ \hat{\delta}_{DiD} = (\bar{Y}_{Treatment, Post} - \bar{Y}_{Treatment, Pre}) - (\bar{Y}_{Control, Post} - \bar{Y}_{Control, Pre}) \]
\textbf{Giả định Cốt lõi (Parallel Trends Assumption):} Nếu không có sự
kiện can thiệp, xu hướng thời gian của nhóm xử lý và nhóm đối chứng phải
song song với nhau.
\[ \mathbb{E}[Y_1(0) - Y_0(0) | D=1] = \mathbb{E}[Y_1(0) - Y_0(0) | D=0] \]
(Trong đó \(Y_t(0)\) là kết quả tiềm năng khi không nhận can thiệp tại
thời điểm \(t\)).

\subsubsection{15.4.2 Điểm mù thường gặp (Cognitive Traps \&
Limitations)}\label{ux111iux1ec3m-muxf9-thux1b0ux1eddng-gux1eb7p-cognitive-traps-limitations-100}

Giả định Parallel Trends không thể kiểm định trực tiếp cho giai đoạn
Post (do ta không quan sát được counterfactual). Việc kiểm định xu hướng
trong quá khứ (Pre-trends test) là cần thiết nhưng chưa đủ. Gần đây, sự
cố với \textbf{Mô hình tác động cố định hai chiều (TWFE - Two-Way Fixed
Effects)} khi các nhóm nhận xử lý ở các thời điểm khác nhau (Staggered
adoption) đã được chỉ ra: TWFE gán trọng số sai lệch, thậm chí âm, cho
một số nhóm (như Sun-Abraham 2021, Callaway-Sant'Anna 2021 đã chứng
minh).

\subsubsection{15.4.3 Ứng dụng thực tế (Practical
Applications)}\label{ux1ee9ng-dux1ee5ng-thux1ef1c-tux1ebf-practical-applications-56}

\textbf{Card \& Krueger (1994):} Ước lượng hiệu ứng của việc tăng lương
tối thiểu ở New Jersey lên việc làm tại cửa hàng thức ăn nhanh, dùng
Pennsylvania làm nhóm đối chứng. -
\(\Delta Y_{NJ} = 21.03 - 20.44 = 0.59\) (Việc làm tại NJ tăng). -
\(\Delta Y_{PA} = 21.17 - 23.33 = -2.16\) (Việc làm tại PA giảm). -
\(\hat{\delta}_{DiD} = 0.59 - (-2.16) = +2.75\). Kết quả ngược trực giác
vi mô truyền thống: việc tăng lương tối thiểu ở đây không làm giảm việc
làm.

\begin{center}\rule{0.5\linewidth}{0.5pt}\end{center}

\subsection{15.5 Hồi quy Gián đoạn (Regression Discontinuity Design -
RDD)}\label{hux1ed3i-quy-giuxe1n-ux111oux1ea1n-regression-discontinuity-design---rdd}

\subsubsection{15.5.1 Cốt lõi Lý thuyết (Core
Theory)}\label{cux1ed1t-luxf5i-luxfd-thuyux1ebft-core-theory-50}

RDD được sử dụng khi can thiệp \(D_i\) được quyết định hoàn toàn bởi một
biến chạy (running variable) \(X_i\) vượt qua một ngưỡng (cutoff) \(c\).
- \textbf{Sharp RDD (Gián đoạn sắc nét):} Xác suất nhận can thiệp chuyển
từ 0 sang 1 ngay tại \(c\). \(P(D_i=1|X_i=c) = 1\). - \textbf{Fuzzy RDD
(Gián đoạn mờ):} Tồn tại một cú nhảy (jump) trong xác suất nhận can
thiệp tại \(c\), nhưng nhỏ hơn 1.

Hiệu ứng LATE tại điểm ngưỡng (Sharp RD) được tính bằng giới hạn trái và
phải:
\[ \tau_{SRD} = \lim_{x \downarrow c} \mathbb{E}[Y_i | X_i = x] - \lim_{x \uparrow c} \mathbb{E}[Y_i | X_i = x] \]

Ước lượng thường dùng Hồi quy đa thức cục bộ (Local Polynomial
Regression) quanh vùng cửa sổ băng thông (bandwidth) \(h\):
\([c-h, c+h]\). Lựa chọn băng thông tối ưu \(h^*\) (ví dụ
Imbens-Kalyanaraman) đối mặt với sự đánh đổi giữa phương sai (variance)
và độ chệch (bias).

\subsubsection{15.5.2 Điểm mù thường gặp (Cognitive Traps \&
Limitations)}\label{ux111iux1ec3m-muxf9-thux1b0ux1eddng-gux1eb7p-cognitive-traps-limitations-101}

Sự nhận dạng nhân quả dựa vào \textbf{giả định liên tục (Continuity
assumption)}: Các đặc điểm của quan sát ngay trước và ngay sau điểm
ngưỡng \(c\) là tương đồng nhau. Nếu các tác nhân có khả năng thao túng
vị trí của mình (sorting/manipulation), RDD sẽ vô hiệu. Cần áp dụng
\textbf{Kiểm định mật độ McCrary (McCrary Density Test)} để kiểm tra sự
đột biến bất thường của phân phối biến chạy \(X_i\) tại \(c\).

\subsubsection{15.5.3 Ứng dụng thực tế (Practical
Applications)}\label{ux1ee9ng-dux1ee5ng-thux1ef1c-tux1ebf-practical-applications-57}

\textbf{Lee (2008) - Lợi thế đương nhiệm trong bầu cử:} Biến chạy là tỷ
lệ phiếu bầu chênh lệch (Margin of Victory). Điểm ngưỡng \(c=0\). Những
ứng viên thắng sát nút (với biên độ +0.01\%) rất giống những ứng viên
thua sát nút (-0.01\%). So sánh xác suất đắc cử ở kỳ tiếp theo của hai
nhóm này cho ra hệ số nhân quả chính xác về ``Lợi thế đương nhiệm''.

\begin{center}\rule{0.5\linewidth}{0.5pt}\end{center}

\subsection{15.6 Nền tảng Dữ liệu Bảng (Panel Data): Fixed vs Random
Effects \& Sai số chuẩn
Cụm}\label{nux1ec1n-tux1ea3ng-dux1eef-liux1ec7u-bux1ea3ng-panel-data-fixed-vs-random-effects-sai-sux1ed1-chuux1ea9n-cux1ee5m}

\paragraph{15.6.1 Cốt lõi Lý thuyết (Core
Theory)}\label{cux1ed1t-luxf5i-luxfd-thuyux1ebft-core-theory-51}

Trước khi tiến tới các phương pháp suy luận nhân quả tiên tiến như Khác
biệt trong Khác biệt (DiD) hay kiểm định phê bình TWFE (Two-Way Fixed
Effects), nhà nghiên cứu phải nắm vững cấu trúc nền tảng của \textbf{Dữ
liệu Bảng (Panel Data / Longitudinal Data)} --- bộ dữ liệu theo dõi
\(N\) thực thể (cá nhân, doanh nghiệp, quốc gia) qua \(T\) thời kỳ.

Mô hình tuyến tính dữ liệu bảng tổng quát có dạng:
\[y_{it} = \beta_0 + \mathbf{x}_{it}'\boldsymbol{\beta} + c_i + u_{it} \quad (i = 1, \dots, N; t = 1, \dots, T)\]
trong đó \(c_i\) là hiệu ứng riêng không quan sát được theo thực thể và
bất biến theo thời gian (time-invariant unobserved individual
heterogeneity như trí thông minh, văn hóa doanh nghiệp, vị trí địa lý
gốc), và \(u_{it}\) là sai số ngẫu nhiên đặc thù theo thời kỳ
(idiosyncratic error).

\textbf{1. Mô hình Tác động Cố định (Fixed Effects - FE / Within
Estimator):} Mô hình FE cho phép hiệu ứng riêng \(c_i\) có tương quan
tùy ý với các biến giải thích quan sát được:
\(\text{Cov}(\mathbf{x}_{it}, c_i) \neq 0\). Để loại bỏ sai lệch biến
thiếu do \(c_i\), phương pháp FE thực hiện biến đổi trừ trung bình theo
thời gian (time-demeaning / within transformation):
\[\bar{y}_i = \frac{1}{T} \sum_{t=1}^T y_{it} = \beta_0 + \bar{\mathbf{x}}_i'\boldsymbol{\beta} + c_i + \bar{u}_i\]
Trừ phương trình trung bình khỏi phương trình gốc, ta tiêu diệt hoàn
toàn \(c_i\):
\[(y_{it} - \bar{y}_i) = (\mathbf{x}_{it} - \bar{\mathbf{x}}_i)'\boldsymbol{\beta} + (u_{it} - \bar{u}_i) \implies \tilde{y}_{it} = \tilde{\mathbf{x}}_{it}'\boldsymbol{\beta} + \tilde{u}_{it}\]
Ước lượng OLS trên phương trình đã biến đổi cho ra bộ ước lượng không
chệch \(\hat{\boldsymbol{\beta}}_{FE}\).

\textbf{2. Mô hình Tác động Ngẫu nhiên (Random Effects - RE / Feasible
GLS):} Nếu giả định cực kỳ mạnh được thỏa mãn là hiệu ứng riêng \(c_i\)
hoàn toàn độc lập và không tương quan với mọi biến giải thích:
\(\text{Cov}(\mathbf{x}_{it}, c_i) = 0\), mô hình RE coi
\(v_{it} = c_i + u_{it}\) là một sai số phức hợp (composite error). Khi
đó, phương pháp Bình phương Tối thiểu Tổng quát Khả thi (FGLS) thực hiện
biến đổi quasi-demeaning:
\[y_{it} - \theta \bar{y}_i = (1-\theta)\beta_0 + (\mathbf{x}_{it} - \theta \bar{\mathbf{x}}_i)'\boldsymbol{\beta} + (v_{it} - \theta \bar{v}_i)\]
với trọng số
\(\theta = 1 - \sqrt{\frac{\sigma_u^2}{\sigma_u^2 + T\sigma_c^2}} \in [0, 1]\).
Ước lượng RE hiệu quả hơn FE (độ lệch chuẩn nhỏ hơn) và cho phép ước
lượng hệ số của các biến bất biến theo thời gian (như giới tính, khoảng
cách).

\textbf{3. Kiểm định Hausman (1978):} Để lựa chọn giữa FE và RE, kiểm
định Hausman đánh giá giả thuyết không
\(H_0: \text{Cov}(\mathbf{x}_{it}, c_i) = 0\). Dưới \(H_0\), cả RE và FE
đều nhất quán nhưng RE hiệu quả hơn; dưới \(H_1\), RE bị chệch và không
nhất quán trong khi FE vẫn nhất quán. Thống kê kiểm định Wald có dạng:
\[H = (\hat{\boldsymbol{\beta}}_{FE} - \hat{\boldsymbol{\beta}}_{RE})' [\text{Var}(\hat{\boldsymbol{\beta}}_{FE}) - \text{Var}(\hat{\boldsymbol{\beta}}_{RE})]^{-1} (\hat{\boldsymbol{\beta}}_{FE} - \hat{\boldsymbol{\beta}}_{RE}) \sim \chi^2(k)\]
Nếu \(H > \chi_{crit}^2\) (p-value \textless{} 0,05), bác bỏ \(H_0\),
bắt buộc phải sử dụng Fixed Effects.

\textbf{4. Sai số chuẩn Cụm (Clustered Standard Errors - Cameron \&
Miller 2015):} Trong dữ liệu bảng, sai số \(u_{it}\) thường có sự tương
quan chuỗi (serial correlation) nặng nề bên trong cùng một thực thể qua
các năm: \(\text{Cov}(u_{is}, u_{it}) \neq 0\). Nếu sử dụng sai số chuẩn
OLS thông thường hoặc Robust White, độ lệch chuẩn sẽ bị đánh giá thấp
nghiêm trọng (underestimated standard errors), dẫn đến thống kê \(t\) bị
thổi phồng và tỷ lệ bác bỏ sai \(H_0\) tăng vọt lên 40-50\% thay vì 5\%
(Bertrand, Duflo, \& Mullainathan 2004). Giải pháp là bắt buộc phải hiệu
chỉnh ma trận hiệp phương sai theo cụm thực thể (Clustered by ID /
Entity):
\[\widehat{\text{Var}}_{cluster}(\hat{\boldsymbol{\beta}}) = (\mathbf{X}'\mathbf{X})^{-1} \left( \sum_{i=1}^N \mathbf{X}_i' \hat{\mathbf{u}}_i \hat{\mathbf{u}}_i' \mathbf{X}_i \right) (\mathbf{X}'\mathbf{X})^{-1}\]

\paragraph{15.6.2 Điểm mù thường gặp (Cognitive Traps \&
Limitations)}\label{ux111iux1ec3m-muxf9-thux1b0ux1eddng-gux1eb7p-cognitive-traps-limitations-102}

\begin{itemize}
\tightlist
\item
  \textbf{Điểm mù thứ nhất (Sự bất lực của FE trước các biến bất biến
  theo thời gian và tác động ngược động):} Khi chạy mô hình FE, tất cả
  các biến giải thích không thay đổi theo thời gian (như khoảng cách địa
  lý đến cảng biển, giới tính, chủng tộc, thể chế pháp lý gốc) bị tiêu
  diệt hoàn toàn trong phép trừ trung bình. Nhiều nhà nghiên cứu rơi vào
  điểm mù khi cố gắng đưa các biến này vào mô hình FE và bị phần mềm
  thống kê (Stata/R) tự động loại bỏ (dropped due to collinearity).
  Furthermore, FE đòi hỏi giả định ngoại sinh nghiêm ngặt theo mọi thời
  kỳ (Strict Exogeneity:
  \(E[u_{it} | \mathbf{x}_{i1}, \dots, \mathbf{x}_{iT}, c_i] = 0\)). Nếu
  biến phụ thuộc trong quá khứ ảnh hưởng đến biến giải thích hiện tại
  (như cú sốc doanh thu năm trước làm thay đổi chi tiêu quảng cáo năm
  nay), FE sẽ bị chệch nghiêm trọng (Nickell bias 1981), buộc phải
  chuyển sang mô hình GMM phương pháp phản xạ (Arellano-Bond GMM).
\item
  \textbf{Điểm mù thứ hai (Lạm dụng Kiểm định Hausman mà quên kiểm tra
  giả định phân phối sai số):} Kiểm định Hausman truyền thống phụ thuộc
  cực mạnh vào giả định sai số \(u_{it}\) có phương sai đồng nhất
  (homoskedasticity) và không tương quan chuỗi. Trong thực tế dữ liệu
  tài chính hoặc vi mô, sai số gần như luôn dị phương sai và tự tương
  quan. Khi đó, ma trận hiệu
  \([\text{Var}(\hat{\boldsymbol{\beta}}_{FE}) - \text{Var}(\hat{\boldsymbol{\beta}}_{RE})]\)
  không còn bán xác định dương, khiến kiểm định Hausman ra giá trị âm
  hoặc sai lệch hoàn toàn. Nhà làm lượng hóa chuyên nghiệp phải luôn sử
  dụng \textbf{Kiểm định Hausman bền vững (Robust / Wooldridge
  Regression-based Hausman test)} với sai số chuẩn cụm để ra quyết định
  chính xác.
\end{itemize}

\paragraph{15.6.3 Ứng dụng thực tế (Practical
Applications)}\label{ux1ee9ng-dux1ee5ng-thux1ef1c-tux1ebf-practical-applications-58}

\begin{itemize}
\tightlist
\item
  \textbf{Đánh giá Tác động của Quản trị Doanh nghiệp đến Hiệu quả Tài
  chính tại HOSE:} Khi nghiên cứu ảnh hưởng của tỷ lệ thành viên Hội
  đồng Quản trị độc lập đến chỉ số Tobin's Q của các công ty niêm yết
  trên Sở Giao dịch Chứng khoán TP.HCM (HOSE) giai đoạn 2015--2023, nếu
  dùng OLS gộp (Pooled OLS), kết quả thường bị chệch do biến thiếu
  \(c_i\) (như uy tín thương hiệu lâu đời hoặc tài năng của người sáng
  lập). Áp dụng mô hình Fixed Effects giúp kiểm soát toàn bộ các đặc
  tính bất biến này, chỉ đo lường sự thay đổi thực sự trong nội bộ công
  ty qua các năm.
\item
  \textbf{Nghiên cứu Kinh tế Địa phương \& Chi tiêu Công tại Việt Nam:}
  Phân tích dữ liệu bảng 63 tỉnh/thành phố Việt Nam về tác động của đầu
  tư công đến tăng trưởng GRDP. Vì 63 tỉnh có vị trí địa lý, điều kiện
  tự nhiên khác nhau hoàn toàn, kiểm định Hausman gần như luôn bác bỏ
  \(H_0\) với \(p < 0,001\), khẳng định việc dùng Random Effects là sai
  lầm nghiêm trọng. Đồng thời, do chu kỳ kinh tế tỉnh có tính bền vững
  qua nhiều năm, việc cụm sai số chuẩn theo cấp Tỉnh (với \(N=63\) cụm)
  là tiêu chuẩn bắt buộc trong các báo cáo tư vấn chính sách cho Quốc
  hội.
\end{itemize}

\paragraph{15.6.4 Ví dụ minh họa tính toán (Worked Numerical
Examples)}\label{vuxed-dux1ee5-minh-hux1ecda-tuxednh-touxe1n-worked-numerical-examples-45}

Xét bộ dữ liệu bảng nhỏ gồm \(N = 2\) doanh nghiệp (ID 1 và ID 2) trong
\(T = 2\) năm (2022 và 2023) để đo lường tác động của Chi phí đào tạo
(\(X\), triệu đồng) đến Năng suất lao động (\(Y\), triệu đồng/người). Dữ
liệu quan sát: * Doanh nghiệp 1 (Công nghệ, gốc cao \(c_1 = 10\)): Năm
2022 (\(X_{1,1}=2, Y_{1,1}=14\)), Năm 2023 (\(X_{1,2}=4, Y_{1,2}=18\)).
* Doanh nghiệp 2 (Dệt may, gốc thấp \(c_2 = 2\)): Năm 2022
(\(X_{2,1}=1, Y_{2,1}=5\)), Năm 2023 (\(X_{2,2}=3, Y_{2,2}=9\)).

\textbf{Bước 1: Tính trung bình theo thời gian cho từng doanh nghiệp
(\(\bar{X}_i, \bar{Y}_i\)):} * Doanh nghiệp 1:
\(\bar{X}_1 = \frac{2 + 4}{2} = 3\);
\(\bar{Y}_1 = \frac{14 + 18}{2} = 16\). * Doanh nghiệp 2:
\(\bar{X}_2 = \frac{1 + 3}{2} = 2\);
\(\bar{Y}_2 = \frac{5 + 9}{2} = 7\).

\textbf{Bước 2: Biến đổi trừ trung bình (Within Transformation /
Time-demeaning):} * \(\tilde{X}_{1,1} = 2 - 3 = -1\);
\(\tilde{Y}_{1,1} = 14 - 16 = -2\). * \(\tilde{X}_{1,2} = 4 - 3 = 1\);
\(\tilde{Y}_{1,2} = 18 - 16 = 2\). * \(\tilde{X}_{2,1} = 1 - 2 = -1\);
\(\tilde{Y}_{2,1} = 5 - 7 = -2\). * \(\tilde{X}_{2,2} = 3 - 2 = 1\);
\(\tilde{Y}_{2,2} = 9 - 7 = 2\).

\textbf{Bước 3: Ước lượng hệ số hồi quy Fixed Effects
(\(\hat{\beta}_{FE}\)):} Công thức OLS không hệ số chặn trên dữ liệu đã
trừ trung bình:
\[\hat{\beta}_{FE} = \frac{\sum_{i=1}^2 \sum_{t=1}^2 \tilde{X}_{it} \tilde{Y}_{it}}{\sum_{i=1}^2 \sum_{t=1}^2 \tilde{X}_{it}^2}\]
Tính tử số:
\((-1 \times -2) + (1 \times 2) + (-1 \times -2) + (1 \times 2) = 2 + 2 + 2 + 2 = 8\).
Tính mẫu số: \((-1)^2 + (1)^2 + (-1)^2 + (1)^2 = 1 + 1 + 1 + 1 = 4\).
\[\hat{\beta}_{FE} = \frac{8}{4} = 2,0\] \textbf{Biện luận:} Tác động
thực sự nội bộ (within effect) của chi phí đào tạo là: cứ tăng 1 triệu
đồng đào tạo thì năng suất lao động tăng thêm chính xác \textbf{2,0
triệu đồng/người}, hoàn toàn không bị vấy bẩn bởi sự chênh lệch nền tảng
công nghệ \(c_1, c_2\) giữa hai doanh nghiệp.

\subsection{15.7 Phương pháp Ghép cặp \& Điểm xu hướng (Matching \&
Propensity Score
Methods)}\label{phux1b0ux1a1ng-phuxe1p-ghuxe9p-cux1eb7p-ux111iux1ec3m-xu-hux1b0ux1edbng-matching-propensity-score-methods}

\begin{figure}
\centering
\pandocbounded{\includegraphics[keepaspectratio,alt={Minh họa VAR Impulse Response}]{C:/Users/Admin/.gemini/antigravity/brain/e5791226-a607-45b8-984e-554eb13f4da1/var_impulse_response_1785203333263.jpg}}
\caption{Minh họa VAR Impulse Response}
\end{figure}

\subsubsection{15.7.1 Cốt lõi Lý thuyết (Core
Theory)}\label{cux1ed1t-luxf5i-luxfd-thuyux1ebft-core-theory-52}

Khi việc chọn mẫu hoàn toàn dựa trên các đặc điểm có thể quan sát
(Selection on Observables), ta có \textbf{Giả định Độc lập Có điều kiện
(Conditional Independence Assumption - CIA)}:
\[ (Y(0), Y(1)) \perp \!\!\! \perp D_i | X_i \] Tuy nhiên, nếu số chiều
của \(X\) quá lớn, việc ghép cặp (Matching) sẽ gặp ``lời nguyền đa
chiều'' (Curse of Dimensionality). Định lý Rosenbaum và Rubin (1983)
chứng minh rằng: Nếu CIA đúng với \(X\), nó cũng đúng với Điểm xu hướng
\(e(X)\): \[ (Y(0), Y(1)) \perp \!\!\! \perp D_i | e(X_i) \] Trong đó,
\textbf{Điểm xu hướng (Propensity Score)} được định nghĩa là
\(e(X) = \mathbb{P}(D=1 | X)\), thường được ước lượng bằng Logit hoặc
Probit. Các phương pháp phổ biến: K-Nearest Neighbor, Kernel Matching,
hay Inverse Probability Weighting (IPW).

\subsubsection{15.7.2 Điểm mù thường gặp (Cognitive Traps \&
Limitations)}\label{ux111iux1ec3m-muxf9-thux1b0ux1eddng-gux1eb7p-cognitive-traps-limitations-103}

Điểm mù nguy hiểm nhất của PSM là các nhà nghiên cứu tin rằng phương
pháp này giải quyết được vấn đề nội sinh tổng quát. \textbf{Sự thật là
PSM KHÔNG giải quyết được Selection on Unobservables} (các biến không
quan sát được). Khác với IV hay DiD, giả định CIA là giả định không thể
kiểm định (untestable assumption) và thường quá mạnh trong khoa học xã
hội.

\subsubsection{15.7.3 Ứng dụng thực tế (Practical
Applications)}\label{ux1ee9ng-dux1ee5ng-thux1ef1c-tux1ebf-practical-applications-59}

\textbf{Đánh giá chương trình đào tạo việc làm (LaLonde, 1986):} Tính
điểm xu hướng cho một người nhận tham gia khóa học nghề dựa trên tuổi,
bằng cấp, sắc tộc. Sau đó ``ghép'' (match) một người trong nhóm có học
với một người trong nhóm không học có cùng điểm xu hướng \(e(X)\) để
tính chênh lệch thu nhập. Đòi hỏi vùng trùng khớp chung (Common Support)
đủ lớn giữa hai nhóm.

\begin{center}\rule{0.5\linewidth}{0.5pt}\end{center}

\subsection{15.8 Phương pháp Tổng hợp (Synthetic Control
Method)}\label{phux1b0ux1a1ng-phuxe1p-tux1ed5ng-hux1ee3p-synthetic-control-method}

\subsubsection{15.8.1 Cốt lõi Lý thuyết (Core
Theory)}\label{cux1ed1t-luxf5i-luxfd-thuyux1ebft-core-theory-53}

Abadie, Diamond, và Hainmueller (2010) đề xuất Phương pháp Đối chứng
Tổng hợp (SCM) cho các nghiên cứu vĩ mô/chính sách áp dụng trên một đơn
vị duy nhất (vd: một quốc gia, một bang). Thay vì chọn một bang duy nhất
làm đối chứng như DiD, SCM tạo ra một ``phiên bản tổng hợp'' bằng cách
kết hợp lồi (convex combination) của các bang không bị tác động (Donor
pool). Vector trọng số \(W^* = (w_1, w_2, ..., w_J)\) (với
\(w_j \ge 0, \sum w_j = 1\)) được chọn để tối thiểu hóa khoảng cách giữa
đặc điểm của đơn vị xử lý (\(X_1\)) và tổng hợp của nhóm đối chứng
(\(X_0 W\)):
\[ \min_{W} ||X_1 - X_0 W||_V = \sqrt{(X_1 - X_0 W)' V (X_1 - X_0 W)} \]
Trong đó \(V\) là ma trận trọng số thể hiện tầm quan trọng của các biến
dự báo.

Suy luận thống kê (Inference) không dùng sai số chuẩn cổ điển, mà sử
dụng \textbf{Thực nghiệm giả dược (Placebo Tests/Permutation)}: Lần lượt
áp dụng phương pháp SCM cho các đơn vị trong Donor pool để tạo ra một
phân phối ``ảo'' của các tác động, từ đó tính Root Mean Square
Prediction Error (RMSPE) và giá trị p-value giả.

\subsubsection{15.8.2 Điểm mù thường gặp (Cognitive Traps \&
Limitations)}\label{ux111iux1ec3m-muxf9-thux1b0ux1eddng-gux1eb7p-cognitive-traps-limitations-104}

Phương pháp SCM có rủi ro cao bị quá khớp (overfitting) trong giai đoạn
Pre-treatment. Nếu giai đoạn trước can thiệp quá ngắn hoặc thị trường
biến động mạnh, vector \(W\) sẽ chỉ bám vào các nhiễu ngẫu nhiên (noise)
thay vì các đặc tính cấu trúc.

\subsubsection{15.8.3 Ứng dụng thực tế (Practical
Applications)}\label{ux1ee9ng-dux1ee5ng-thux1ef1c-tux1ebf-practical-applications-60}

\textbf{Dự luật 99 của California (1988):} Đánh giá việc tăng thuế thuốc
lá ở California. Không có một tiểu bang đơn lẻ nào giống hoàn toàn với
California. SCM tạo ra ``Synthetic California'' bằng cách gộp Colorado
(16\%), Connecticut (69\%), Montana (20\%),\ldots{} Kết quả cho thấy tỷ
lệ hút thuốc ở California giảm mạnh hơn hẳn so với ``phiên bản tổng
hợp'' sau 1988, xác lập một hiệu ứng nhân quả rõ ràng.

\begin{center}\rule{0.5\linewidth}{0.5pt}\end{center}

\subsection{15.9 Suy luận Nhân quả Hiện đại \& Credibility
Revolution}\label{suy-luux1eadn-nhuxe2n-quux1ea3-hiux1ec7n-ux111ux1ea1i-credibility-revolution}

\subsubsection{15.9.1 Cốt lõi Lý thuyết (Core
Theory)}\label{cux1ed1t-luxf5i-luxfd-thuyux1ebft-core-theory-54}

``Cuộc cách mạng độ tin cậy'' (Credibility Revolution), do Angrist và
Pischke đại diện với tinh thần ``Mostly Harmless Econometrics'', đặt
thiết kế nghiên cứu (research design) lên trên việc đặc tả phức tạp các
mô hình thống kê.

Nền tảng của cuộc cách mạng này là \textbf{Mô hình Nhân quả Rubin (Rubin
Causal Model / Potential Outcomes Framework)}. Ký hiệu hai kết quả tiềm
năng của cá nhân \(i\): - \(Y_i(1)\): Kết quả nếu nhận xử lý. -
\(Y_i(0)\): Kết quả nếu không nhận xử lý. Hiệu ứng nhân quả cá nhân là
\(\tau_i = Y_i(1) - Y_i(0)\). Tuy nhiên, ta đối mặt với \textbf{Vấn đề
Cốt lõi của Suy luận Nhân quả}: ta chỉ quan sát được
\(Y_i = D_i Y_i(1) + (1-D_i)Y_i(0)\), chứ không bao giờ thấy đồng thời
cả hai.

Các đại lượng ước lượng tập trung vào trung bình: - \textbf{ATE (Average
Treatment Effect):} \(\mathbb{E}[Y_i(1) - Y_i(0)]\) - \textbf{ATT
(Average Treatment Effect on the Treated):}
\(\mathbb{E}[Y_i(1) - Y_i(0) | D_i = 1]\) - \textbf{LATE (Local Average
Treatment Effect):} \(\mathbb{E}[Y_i(1) - Y_i(0) | \text{Compliers}]\)

Hệ thống này đòi hỏi \textbf{SUTVA (Stable Unit Treatment Value
Assumption)}: Không có sự lan tỏa (no spillover effects) - việc can
thiệp vào cá nhân này không ảnh hưởng đến kết quả của cá nhân khác. Bên
cạnh khuôn khổ Rubin, Đồ thị Có hướng Không chu trình (DAGs - Directed
Acyclic Graphs) của Judea Pearl cung cấp một hệ thống toán học hình học
trực quan để nhận diện các nút thắt nội sinh, biến nhiễu (confounders)
và biến chặn (colliders).

\subsubsection{15.9.2 Điểm mù thường gặp (Cognitive Traps \&
Limitations)}\label{ux111iux1ec3m-muxf9-thux1b0ux1eddng-gux1eb7p-cognitive-traps-limitations-105}

Sự phân mảnh giữa nhóm dùng kinh tế lượng ứng dụng (Rubin framework) và
nhóm khoa học máy tính (Pearl's DAGs). Nhiều nhà kinh tế học e ngại DAGs
vì chúng không tham số hóa (non-parametric), nhưng DAGs lại tối ưu hơn
để tư duy về cơ chế truyền dẫn cấu trúc trước khi viết phương trình.

\subsubsection{15.9.3 Ứng dụng thực tế (Practical
Applications)}\label{ux1ee9ng-dux1ee5ng-thux1ef1c-tux1ebf-practical-applications-61}

Trong ứng dụng thực tế (thí nghiệm thực địa RCTs hay thiết kế tự nhiên),
nhà nghiên cứu PhD cần ưu tiên: (1) Trả lời câu hỏi chính sách rõ ràng,
(2) Có cơ chế nhận dạng nguồn biến thiên ngoại sinh (exogenous shock),
hơn là (3) Kỹ thuật ước lượng hàn lâm phức tạp nhưng dữ liệu rác
(Garbage in, garbage out). Ưu tiên RDD và thiết kế thực nghiệm ngẫu
nhiên trước các phương pháp phụ thuộc quá nhiều vào quan sát như
Propensity Score Matching.

\begin{center}\rule{0.5\linewidth}{0.5pt}\end{center}

\begin{center}\rule{0.5\linewidth}{0.5pt}\end{center}

\subsection{15.10 Kinh tế lượng Chuỗi thời gian (Time-Series
Econometrics)}\label{kinh-tux1ebf-lux1b0ux1ee3ng-chuux1ed7i-thux1eddi-gian-time-series-econometrics}

Như đã thảo luận ở Chương 13 (Kinh tế vĩ mô nâng cao), phần này sẽ bổ
sung các nền tảng toán học và công cụ phân tích dữ liệu chuỗi thời gian,
hoàn thiện cho các mô hình kinh tế lượng vi mô/nhân quả đã được trình
bày ở các phần trước của Chương 15.

\subsubsection{15.10.1 Tính dừng \& Kiểm định Nghiệm đơn vị
(Stationarity \& Unit Root
Tests)}\label{tuxednh-dux1eebng-kiux1ec3m-ux111ux1ecbnh-nghiux1ec7m-ux111ux1a1n-vux1ecb-stationarity-unit-root-tests}

\subsubsection{15.10.1 Cốt lõi Lý thuyết (Core
Theory)}\label{cux1ed1t-luxf5i-luxfd-thuyux1ebft-core-theory-55}

Một quá trình ngẫu nhiên \(\{Y_t\}\) được gọi là \textbf{dừng yếu}
(weakly stationary) hay dừng hiệp phương sai nếu nó thỏa mãn ba điều
kiện: 1. Kỳ vọng hữu hạn và không đổi:
\(\mathbb{E}[Y_t] = \mu < \infty, \quad \forall t\) 2. Phương sai hữu
hạn và không đổi:
\(\text{Var}(Y_t) = \mathbb{E}[(Y_t - \mu)^2] = \gamma_0 < \infty, \quad \forall t\)
3. Tự hiệp phương sai chỉ phụ thuộc vào khoảng cách độ trễ \(k\), không
phụ thuộc vào thời điểm \(t\):
\(\text{Cov}(Y_t, Y_{t-k}) = \mathbb{E}[(Y_t - \mu)(Y_{t-k} - \mu)] = \gamma_k, \quad \forall t, k\)

Ngược lại, một quá trình \textbf{bước đi ngẫu nhiên} (Random Walk) là ví
dụ kinh điển của chuỗi không dừng: \[ Y_t = Y_{t-1} + \varepsilon_t \]
trong đó \(\varepsilon_t \sim \text{iid}(0, \sigma^2)\) là nhiễu trắng
(white noise). Bằng cách lặp lùi về \(Y_0\), ta có
\(Y_t = Y_0 + \sum_{i=1}^t \varepsilon_i\). Lấy phương sai hai vế, ta
thấy \(\text{Var}(Y_t) = t\sigma^2\), tăng dần tới vô cực theo thời gian
\(t\), vi phạm điều kiện 2. Quá trình này được gọi là tích hợp bậc 1, ký
hiệu \(I(1)\), và chứa một ``nghiệm đơn vị'' (unit root).

Để kiểm tra tính dừng, ta sử dụng \textbf{Kiểm định Dickey-Fuller (DF)}.
Mô hình kiểm định cơ bản:
\[ \Delta Y_t = \delta Y_{t-1} + \varepsilon_t \] Với giả thuyết không
(Null hypothesis) \(H_0: \delta = 0\) (tương đương với \(\rho = 1\),
chuỗi có nghiệm đơn vị) và đối thiết \(H_1: \delta < 0\) (chuỗi dừng).
Nếu nhiễu \(\varepsilon_t\) có tự tương quan, ta sử dụng kiểm định
\textbf{Augmented Dickey-Fuller (ADF)} bằng cách thêm các biến trễ:
\[ \Delta Y_t = \alpha + \beta t + \delta Y_{t-1} + \sum_{i=1}^{p} \gamma_i \Delta Y_{t-i} + \varepsilon_t \]
Thống kê kiểm định
\(t_{\delta} = \frac{\hat{\delta}}{\text{SE}(\hat{\delta})}\) không tuân
theo phân phối \(t\)-Student chuẩn (vì mẫu không có phương sai hữu hạn
dưới \(H_0\)) mà tuân theo phân phối Dickey-Fuller phi chuẩn.

\textbf{Hệ quả: Hồi quy giả mạo (Spurious Regression):} Granger và
Newbold (1974) chứng minh rằng nếu ta hồi quy hai chuỗi không dừng, độc
lập hoàn toàn với nhau \(Y_t \sim I(1)\) và \(X_t \sim I(1)\) bằng OLS:
\(Y_t = \alpha + \beta X_t + u_t\), thống kê \(t\) của \(\hat{\beta}\)
sẽ không hội tụ về phân phối chuẩn mà sẽ tiến tới vô cực khi
\(T \to \infty\). Điều này dẫn đến xác suất bác bỏ \(H_0: \beta=0\) tiến
tới 1, tạo ra mối quan hệ giả mạo mặc dù \(R^2\) rất cao.

\subsubsection{15.10.2 Điểm mù thường gặp (Cognitive Traps \&
Limitations)}\label{ux111iux1ec3m-muxf9-thux1b0ux1eddng-gux1eb7p-cognitive-traps-limitations-106}

\begin{itemize}
\tightlist
\item
  \textbf{Nghiệm đơn vị vs.~Biến đổi cấu trúc (Structural Break):} Bài
  báo nổi tiếng của Perron (1989) chỉ ra rằng một chuỗi dừng có biến đổi
  cấu trúc trong xu hướng (ví dụ khủng hoảng 1929 hoặc cú sốc dầu mỏ
  1973) sẽ làm cho kiểm định ADF thất bại trong việc bác bỏ \(H_0\).
  Nhiều nhà nghiên cứu nhầm lẫn kết luận chuỗi có nghiệm đơn vị trong
  khi thực tế chỉ là có structural break.
\item
  \textbf{Sức mạnh (Power) của kiểm định:} ADF có sức mạnh thống kê rất
  thấp đối với các quá trình gần với nghiệm đơn vị (near unit root, ví
  dụ \(\rho = 0.98\)).
\end{itemize}

\subsubsection{15.10.3 Ứng dụng thực tế (Practical
Applications)}\label{ux1ee9ng-dux1ee5ng-thux1ef1c-tux1ebf-practical-applications-62}

Kiểm tra tính dừng của lạm phát \(\pi_t\). Giả sử dữ liệu quý trong 25
năm (\(T=100\)). Ước lượng phương trình ADF cho kết quả
\(\hat{\delta} = -0.04\) với sai số chuẩn
\(\text{SE}(\hat{\delta}) = 0.02\). Thống kê tính toán:
\(t = -0.04 / 0.02 = -2.00\). Tra bảng giá trị tới hạn của MacKinnon cho
kiểm định Dickey-Fuller (có hằng số, không có xu hướng) ở mức ý nghĩa
5\% là \(-2.89\). Do \(-2.00 > -2.89\), ta \textbf{không thể bác bỏ} giả
thuyết \(H_0\). Kết luận: chuỗi lạm phát chứa nghiệm đơn vị \(I(1)\).

\begin{center}\rule{0.5\linewidth}{0.5pt}\end{center}

\subsubsection{15.10.2 Đồng liên kết \& Mô hình Hiệu chỉnh Sai số
(Cointegration \&
VECM)}\label{ux111ux1ed3ng-liuxean-kux1ebft-muxf4-huxecnh-hiux1ec7u-chux1ec9nh-sai-sux1ed1-cointegration-vecm}

\subsubsection{15.10.1 Cốt lõi Lý thuyết (Core
Theory)}\label{cux1ed1t-luxf5i-luxfd-thuyux1ebft-core-theory-56}

Định nghĩa của Engle \& Granger (1987): Hai chuỗi \(Y_t \sim I(1)\) và
\(X_t \sim I(1)\) được gọi là \textbf{đồng liên kết} (cointegrated) nếu
tồn tại một tổ hợp tuyến tính của chúng là chuỗi dừng, tức là
\(u_t = Y_t - \beta X_t \sim I(0)\). Vector \((1, -\beta)'\) được gọi là
vector đồng liên kết, đặc trưng cho trạng thái cân bằng dài hạn.

Quy trình Engle-Granger gồm 2 bước: 1. Hồi quy phương trình đồng liên
kết dài hạn bằng OLS: \(Y_t = \alpha + \beta X_t + u_t\), lấy phần dư
\(\hat{u}_t\). 2. Kiểm định nghiệm đơn vị trên \(\hat{u}_t\). Nếu
\(\hat{u}_t \sim I(0)\), hai chuỗi đồng liên kết. (Lưu ý: phải sử dụng
bảng giá trị tới hạn riêng của Engle-Granger).

Tuy nhiên, với \(n > 2\) biến, quy trình Engle-Granger mất tác dụng. Ta
phải sử dụng \textbf{phương pháp Johansen}, dựa trên mô hình VECM
(Vector Error Correction Model):
\[ \Delta \mathbf{Y}_t = \Pi \mathbf{Y}_{t-1} + \sum_{i=1}^{p-1} \Gamma_i \Delta \mathbf{Y}_{t-i} + \mathbf{\varepsilon}_t \]
Với \(\mathbf{Y}_t\) là vector \(n \times 1\). Hạng (rank \(r\)) của ma
trận \(\Pi\) (\(n \times n\)) cho biết số lượng vector đồng liên kết. -
Nếu \(r = 0\): không có đồng liên kết. Mô hình trở về VAR dạng sai phân.
- Nếu \(0 < r < n\): \(\Pi\) có thể phân rã thành
\(\Pi = \mathbf{\alpha} \mathbf{\beta}'\), trong đó \(\mathbf{\beta}\)
(\(n \times r\)) chứa các vector đồng liên kết và \(\mathbf{\alpha}\)
(\(n \times r\)) là ma trận tốc độ điều chỉnh. Phương pháp Johansen dựa
trên cấu trúc trị riêng để tính hai thống kê: Trace statistic
(\(\lambda_{trace}\)) và Maximum Eigenvalue (\(\lambda_{max}\)).

\subsubsection{15.10.2 Điểm mù thường gặp (Cognitive Traps \&
Limitations)}\label{ux111iux1ec3m-muxf9-thux1b0ux1eddng-gux1eb7p-cognitive-traps-limitations-107}

\begin{itemize}
\tightlist
\item
  \textbf{Siêu hội tụ (Super-consistency):} Sinh viên thường lo lắng về
  hiện tượng nội sinh (endogeneity) trong bước 1 của Engle-Granger. Tuy
  nhiên, nếu đồng liên kết tồn tại, ước lượng OLS cho \(\hat{\beta}\) là
  siêu hội tụ ở tốc độ \(O(T)\) thay vì \(O(\sqrt{T})\) như thông
  thường. Độ lệch nội sinh hội tụ về 0 cực nhanh ở mẫu lớn.
\item
  \textbf{Tốc độ điều chỉnh:} Ma trận \(\mathbf{\alpha}\) phải có dấu
  phù hợp để đảm bảo hệ thống hội tụ (stable). Nếu \(\alpha\) dương
  trong phương trình của biến đang ở mức cao hơn cân bằng, hệ thống sẽ
  bùng nổ.
\end{itemize}

\subsubsection{15.10.3 Ứng dụng thực tế (Practical
Applications)}\label{ux1ee9ng-dux1ee5ng-thux1ef1c-tux1ebf-practical-applications-63}

Kiểm định Thuyết sức mua tương đương (PPP): \(e_t\) (log tỷ giá),
\(p_t\) (log giá trong nước), \(p^*_t\) (log giá nước ngoài), đều là
\(I(1)\). Hồi quy cân bằng:
\(e_t = \beta_0 + \beta_1 (p_t - p^*_t) + u_t\). Giả sử VECM cho kết quả
ước lượng \(\hat{\alpha} = -0.15\) cho phương trình của \(\Delta e_t\).
Hệ số \(\alpha\) biểu thị tốc độ điều chỉnh: khi tỷ giá thực chệch khỏi
cân bằng PPP (overvalued/undervalued), nó sẽ điều chỉnh tự động thu hẹp
lại 15\% phần sai lệch đó trong mỗi kỳ (ví dụ mỗi quý).

\begin{center}\rule{0.5\linewidth}{0.5pt}\end{center}

\subsubsection{15.10.3 Mô hình VAR \& SVAR (Vector Autoregression \&
Structural
VAR)}\label{muxf4-huxecnh-var-svar-vector-autoregression-structural-var}

\subsubsection{15.10.1 Cốt lõi Lý thuyết (Core
Theory)}\label{cux1ed1t-luxf5i-luxfd-thuyux1ebft-core-theory-57}

Được Christopher Sims (Nobel 2011) tiên phong, mô hình VAR không đòi hỏi
phân loại tiên nghiệm biến nào nội sinh/ngoại sinh. Một VAR(\(p\)) dạng
rút gọn (reduced-form):
\[ \mathbf{Y}_t = \mathbf{A}_1 \mathbf{Y}_{t-1} + \dots + \mathbf{A}_p \mathbf{Y}_{t-p} + \mathbf{u}_t \]
Trong đó \(\mathbf{Y}_t\) là vector \(n \times 1\),
\(\mathbf{u}_t \sim (\mathbf{0}, \Sigma)\) là vector nhiễu rút gọn.

Dùng toán tử trễ \(L\), \(\mathbf{A}(L)\mathbf{Y}_t = \mathbf{u}_t\).
Nếu hệ thống dừng, ta có thể nghịch đảo \(\mathbf{A}(L)\) để thu được
biểu diễn trung bình trượt dạng vector (VMA):
\[ \mathbf{Y}_t = \mathbf{A}(L)^{-1} \mathbf{u}_t = \mathbf{u}_t + \mathbf{\Phi}_1 \mathbf{u}_{t-1} + \mathbf{\Phi}_2 \mathbf{u}_{t-2} + \dots \]
Biểu diễn VMA là cốt lõi để tính \textbf{Hàm Phản ứng Xung lực} (IRF),
truy xuất chuỗi thời gian của \(\mathbf{Y}\) thay đổi thế nào sau 1 cú
sốc vào \(\mathbf{u}\).

Tuy nhiên, các thành phần của \(\mathbf{u}_t\) tương quan với nhau (ma
trận hiệp phương sai \(\Sigma\) không chéo). Để diễn dịch ý nghĩa kinh
tế, ta cần nhận dạng \textbf{SVAR} với các cú sốc cấu trúc (structural
shocks) \(\mathbf{\varepsilon}_t \sim (\mathbf{0}, \Omega)\) với
\(\Omega\) là ma trận đường chéo:
\[ \mathbf{A}_0 \mathbf{Y}_t = \mathbf{A}_1^* \mathbf{Y}_{t-1} + \dots + \mathbf{\varepsilon}_t \]
Liên hệ: \(\mathbf{u}_t = \mathbf{A}_0^{-1} \mathbf{\varepsilon}_t\) và
\(\Sigma = \mathbf{A}_0^{-1} \Omega (\mathbf{A}_0^{-1})'\). Bởi ma trận
\(\Sigma\) (đối xứng) chỉ cho \(n(n+1)/2\) tham số, ta cần áp đặt ít
nhất \(n(n-1)/2\) ràng buộc lên \(\mathbf{A}_0\). Phương pháp Cholesky
áp đặt \(\mathbf{A}_0\) là ma trận tam giác dưới, ngụ ý một cấu trúc đệ
quy (recursive identification).

\textbf{Phân rã phương sai (Variance Decomposition):} Tính toán phần
trăm biến động của biến \(i\) tại bước \(h\) được giải thích bởi cú sốc
cấu trúc của biến \(j\).

\subsubsection{15.10.2 Điểm mù thường gặp (Cognitive Traps \&
Limitations)}\label{ux111iux1ec3m-muxf9-thux1b0ux1eddng-gux1eb7p-cognitive-traps-limitations-108}

\begin{itemize}
\tightlist
\item
  \textbf{Nhân quả Granger (Granger Causality):} \(X\) là ``nguyên nhân
  Granger'' của \(Y\) nếu các độ trễ của \(X\) trong phương trình của
  \(Y\) có ý nghĩa thống kê. Điểm mù ở đây là nhầm lẫn với ``nhân quả
  thực sự'' (structural causality). Granger causality thuần túy chỉ sự
  cải thiện trong năng lực dự báo.
\item
  \textbf{Thứ tự Cholesky:} Kết quả IRF phụ thuộc mạnh mẽ vào cách xếp
  thứ tự biến trong \(\mathbf{Y}_t\). Cần phải có lý thuyết kinh tế vĩ
  mô cực kỳ vững để biện minh (ví dụ tại sao biến này phản ứng chậm,
  biến kia phản ứng nhanh với cú sốc).
\end{itemize}

\subsubsection{15.10.3 Ứng dụng thực tế (Practical
Applications)}\label{ux1ee9ng-dux1ee5ng-thux1ef1c-tux1ebf-practical-applications-64}

Phân tích truyền dẫn chính sách tiền tệ.
\(\mathbf{Y}_t = [y_t, \pi_t, r_t]'\) (Sản lượng, Lạm phát, Lãi suất Fed
Funds). Xếp lãi suất \(r_t\) cuối cùng ngụ ý: cú sốc chính sách tiền tệ
(lãi suất) có độ trễ, không ảnh hưởng tức thời đến sản lượng và giá cả
trong cùng 1 tháng. Ngược lại, Fed quan sát cả \(y_t\) và \(\pi_t\) ngay
lập tức để ra quyết định lãi suất. Khi vẽ đồ thị IRF, cú sốc dương 1\%
vào \(r_t\) thường dẫn đến ``Price Puzzle'' -- lạm phát tăng nhẹ trong
1-2 quý đầu, trước khi giảm đúng như lý thuyết vĩ mô.

\begin{center}\rule{0.5\linewidth}{0.5pt}\end{center}

\subsubsection{15.10.4 Mô hình ARIMA \& Dự báo (ARIMA \&
Forecasting)}\label{muxf4-huxecnh-arima-dux1ef1-buxe1o-arima-forecasting}

\subsubsection{15.10.1 Cốt lõi Lý thuyết (Core
Theory)}\label{cux1ed1t-luxf5i-luxfd-thuyux1ebft-core-theory-58}

Phương pháp Box-Jenkins là trụ cột của dự báo kinh tế đơn biến, dựa trên
cấu trúc ARIMA(\(p, d, q\)):
\[ \Phi(L) (1-L)^d Y_t = \Theta(L) \varepsilon_t \] Trong đó: - Phần tự
hồi quy AR(\(p\)): \(\Phi(L) = 1 - \phi_1 L - \dots - \phi_p L^p\). -
Bậc lấy sai phân \(d\): làm chuỗi dừng, áp dụng bộ lọc \((1-L)^d\). -
Phần trung bình trượt MA(\(q\)):
\(\Theta(L) = 1 + \theta_1 L + \dots + \theta_q L^q\).

Chu trình gồm 3 bước: 1. \textbf{Định dạng (Identification):} Vẽ biểu đồ
Hàm tự tương quan (ACF) và Tự tương quan riêng phần (PACF) của chuỗi đã
dừng. - Quá trình AR(\(p\)): PACF cắt đứt sau lag \(p\), ACF giảm dần. -
Quá trình MA(\(q\)): ACF cắt đứt sau lag \(q\), PACF giảm dần. 2.
\textbf{Ước lượng \& Lựa chọn (Estimation):} Sử dụng MLE. Lựa chọn độ
trễ dựa trên tiêu chuẩn thông tin Akaike (AIC) và Bayesian (BIC/SBC):
\[ AIC = -2 \ln(\mathcal{L}) + 2k, \quad BIC = -2 \ln(\mathcal{L}) + k \ln(T) \]
(Với \(\mathcal{L}\) là hàm hợp lý cực đại, \(k=p+q\)). BIC phạt nặng
tham số thừa hơn AIC, giúp tránh overfitting. 3. \textbf{Chẩn đoán
(Diagnostics):} Kiểm tra phần dư bằng kiểm định Ljung-Box Q-statistic để
đảm bảo nhiễu trắng.

\subsubsection{15.10.2 Điểm mù thường gặp (Cognitive Traps \&
Limitations)}\label{ux111iux1ec3m-muxf9-thux1b0ux1eddng-gux1eb7p-cognitive-traps-limitations-109}

\begin{itemize}
\tightlist
\item
  \textbf{Rễ của đa thức đặc trưng (Roots of characteristic
  polynomials):} Mô hình chỉ hợp lệ nếu tất cả nghiệm của phương trình
  \(\Phi(z) = 0\) và \(\Theta(z) = 0\) đều nằm hoàn toàn bên ngoài vòng
  tròn đơn vị (\(|z| > 1\)) trong mặt phẳng phức. Bỏ qua kiểm tra này
  dẫn đến mô hình không đảo ngược được hoặc không dừng.
\item
  \textbf{Over-differencing:} Việc lấy sai phân thêm một lần đối với
  chuỗi vốn dĩ đã dừng (\(d\) vượt quá mức cần thiết) sẽ tự tạo ra một
  cấu trúc MA nhân tạo (non-invertible MA process), gây khó khăn trầm
  trọng cho ước lượng.
\end{itemize}

\subsubsection{15.10.3 Ứng dụng thực tế (Practical
Applications)}\label{ux1ee9ng-dux1ee5ng-thux1ef1c-tux1ebf-practical-applications-65}

Dự báo tăng trưởng GDP \(Y_t\). Sau khi lấy sai phân bậc 1,
\(\Delta Y_t\) dừng. ACF có 1 gai (spike) mạnh ở lag 1 rồi biến mất,
PACF giảm dần. Ta chỉ định mô hình ARIMA(0,1,1):
\[ \Delta Y_t = \mu + \varepsilon_t + \theta_1 \varepsilon_{t-1} \] Giả
sử ước lượng ra \(\mu = 0.5, \theta_1 = 0.4\). Để dự báo một bước tới
tương lai (\(t+1\)):
\[ \mathbb{E}_t[\Delta Y_{t+1}] = \mu + \mathbb{E}_t[\varepsilon_{t+1}] + \theta_1 \mathbb{E}_t[\varepsilon_t] = 0.5 + 0 + 0.4 \hat{\varepsilon}_t \]
Từ đó, \(\hat{Y}_{t+1} = Y_t + 0.5 + 0.4 \hat{\varepsilon}_t\).

\begin{center}\rule{0.5\linewidth}{0.5pt}\end{center}

\subsubsection{15.10.5 Mô hình ARCH/GARCH \& Biến động Tài chính
(Volatility
Modeling)}\label{muxf4-huxecnh-archgarch-biux1ebfn-ux111ux1ed9ng-tuxe0i-chuxednh-volatility-modeling}

\subsubsection{15.10.1 Cốt lõi Lý thuyết (Core
Theory)}\label{cux1ed1t-luxf5i-luxfd-thuyux1ebft-core-theory-59}

Trong kinh tế học tài chính, dữ liệu lợi suất thường có ``volatility
clustering'' (nhóm biến động): các giai đoạn thị trường hoảng loạn có
biên độ dao động lớn nối tiếp nhau. Robert Engle (Nobel 2003) giới thiệu
ARCH để mô hình hóa phương sai có điều kiện \(\sigma_t^2\) phụ thuộc
trực tiếp vào bình phương cú sốc quá khứ. Tim Bollerslev (1986) mở rộng
thành mô hình GARCH(\(p, q\)): - Phương trình trung bình:
\(Y_t = \mu + \varepsilon_t \quad\) với
\(\quad \varepsilon_t = \sigma_t z_t, \quad z_t \sim \mathcal{N}(0,1)\)
- Phương trình phương sai:
\(\sigma_t^2 = \omega + \sum_{i=1}^q \alpha_i \varepsilon_{t-i}^2 + \sum_{j=1}^p \beta_j \sigma_{t-j}^2\)

Đối với GARCH(1,1) tiêu chuẩn phổ biến nhất:
\[ \sigma_t^2 = \omega + \alpha \varepsilon_{t-1}^2 + \beta \sigma_{t-1}^2 \]
Ràng buộc: \(\omega > 0, \alpha \ge 0, \beta \ge 0\) để phương sai
dương; và \(\alpha + \beta < 1\) để mô hình có tính dừng, với phương sai
vô điều kiện dài hạn là
\(\bar{\sigma}^2 = \frac{\omega}{1 - \alpha - \beta}\).

\subsubsection{15.10.2 Điểm mù thường gặp (Cognitive Traps \&
Limitations)}\label{ux111iux1ec3m-muxf9-thux1b0ux1eddng-gux1eb7p-cognitive-traps-limitations-110}

\begin{itemize}
\tightlist
\item
  \textbf{Đuôi dày (Fat tails):} Một điểm mù tinh tế là dù nhiễu có điều
  kiện \(z_t\) được giả định tuân theo phân phối chuẩn, mô hình GARCH
  vẫn tự sinh ra một phân phối vô điều kiện cho \(Y_t\) với độ nhọn
  (kurtosis) \(> 3\) (đuôi dày). Tuy nhiên, đôi khi độ nhọn nội sinh này
  chưa đủ khớp với dữ liệu thật, buộc nhà nghiên cứu phải dùng \(z_t\)
  tuân theo phân phối t-Student.
\item
  \textbf{Tác động phi đối xứng (Leverage Effect):} GARCH(1,1) bình
  phương nhiễu \(\varepsilon_{t-1}^2\), nghĩa là tin tốt và tin xấu ảnh
  hưởng như nhau đến biến động. Thực tế ``thị trường con gấu'' (bear
  market) tạo biến động lớn hơn. Cần dùng các mô hình mở rộng như TGARCH
  hay EGARCH.
\end{itemize}

\subsubsection{15.10.3 Ứng dụng thực tế (Practical
Applications)}\label{ux1ee9ng-dux1ee5ng-thux1ef1c-tux1ebf-practical-applications-66}

Quản trị rủi ro Value-at-Risk (VaR) cho danh mục S\&P 500. Ước lượng mô
hình GARCH(1,1) trên lợi suất hàng ngày cho kết quả:
\(\omega = 0.000001, \alpha = 0.08, \beta = 0.90\). Phương sai dài hạn
của thị trường là
\(\bar{\sigma}^2 = \frac{0.000001}{1 - 0.08 - 0.90} = 0.00005\) (độ lệch
chuẩn \(\approx 0.7\%\)/ngày). Giả sử ngày hôm qua (\(t-1\)) thị trường
giảm rất mạnh, cú sốc \(\varepsilon_{t-1} = -0.05\) (giảm 5\%), với mức
phương sai dự báo ngày hôm qua là \(\sigma_{t-1}^2 = 0.0001\). Dự báo
biến động cho ngày hôm nay (\(t\)):
\[ \sigma_t^2 = 0.000001 + 0.08(-0.05)^2 + 0.90(0.0001) = 0.000001 + 0.0002 + 0.00009 = 0.000291 \]
Độ lệch chuẩn hôm nay tăng mạnh lên
\(\sigma_t = \sqrt{0.000291} \approx 1.7\%\). Với một danh mục 10 triệu
USD, VaR 1 ngày ở mức tin cậy 95\% (tương đương \(-1.645\sigma\)) sẽ là
\(10 \times 1.645 \times 0.017 \approx 0.28\) triệu USD. Cú sốc ngày hôm
qua đã làm rủi ro của danh mục tăng lên đáng kể.

\begin{center}\rule{0.5\linewidth}{0.5pt}\end{center}

\section{CHƯƠNG 16: KINH TẾ HỌC HÀNH VI NÂNG CAO \& KINH TẾ HỌC THẦN
KINH}\label{chux1b0ux1a1ng-16-kinh-tux1ebf-hux1ecdc-huxe0nh-vi-nuxe2ng-cao-kinh-tux1ebf-hux1ecdc-thux1ea7n-kinh}

Như đã hứa hẹn ở Chương 6, chương này sẽ hình thức hóa các khái niệm
trực quan về hành vi con người bằng ngôn ngữ toán học chặt chẽ của kinh
tế học vi mô ở bậc sau đại học. Nếu như lý thuyết hữu dụng kỳ vọng
(Expected Utility Theory - EUT) và chiết khấu hàm mũ (Exponential
Discounting) cung cấp một khung phân tích tiêu chuẩn về tính duy lý, thì
thực nghiệm lại chỉ ra những sự sai lệch hệ thống. Chương này không chỉ
trình bày các mô hình thay thế mà còn đi sâu vào cơ chế thần kinh học
của các quyết định, thiết kế chính sách hành vi, và đánh giá một cách có
phê phán giới hạn của các lý thuyết này.

\begin{center}\rule{0.5\linewidth}{0.5pt}\end{center}

\subsection{16.1 Hình thức hóa Lý thuyết Triển vọng (Formal Prospect
Theory)}\label{huxecnh-thux1ee9c-huxf3a-luxfd-thuyux1ebft-triux1ec3n-vux1ecdng-formal-prospect-theory}

\subsubsection{16.1.1 Cốt lõi Lý thuyết (Core
Theory)}\label{cux1ed1t-luxf5i-luxfd-thuyux1ebft-core-theory-60}

Lý thuyết Triển vọng Tích lũy (Cumulative Prospect Theory - CPT) do
Tversky và Kahneman (1992) phát triển giải quyết những hạn chế của EUT
bằng cách thay đổi cả hàm giá trị (value function) và hàm trọng số xác
suất (probability weighting function).

Trong CPT, tính hữu dụng không dựa trên mức tài sản cuối cùng mà dựa
trên \textbf{sự thay đổi tài sản} so với một điểm tham chiếu (reference
point). Giả sử một cá nhân đối mặt với một tập hợp các kết cục
\(X = \{x_{-m}, \dots, x_{-1}, x_0, x_1, \dots, x_n\}\) với
\(x_i < x_j\) khi \(i < j\), trong đó \(x_0 = 0\) là điểm tham chiếu.
Xác suất xảy ra kết cục \(x_i\) là \(p_i\).

Hàm giá trị \(v(x)\) được định nghĩa là: \[
 v(x) = \begin{cases} 
\] x\^{}\alpha \& \text{nếu } x \geq 0 \textbackslash{} \[
\] -\lambda(-x)\^{}\beta \& \text{nếu } x \textless{} 0 \[
 \end{cases} 
\] với \(\alpha, \beta \in (0,1)\) phản ánh độ lõm ở miền lời (gains) và
độ lồi ở miền lỗ (losses), biểu thị e ngại rủi ro khi có lời và thích
rủi ro khi bị lỗ. Tham số \(\lambda > 1\) là hệ số \textbf{ác cảm mất
mát} (loss aversion), cho thấy một đơn vị lỗ gây ra sự suy giảm hữu dụng
lớn hơn so với sự gia tăng hữu dụng từ một đơn vị lời.

Hàm trọng số xác suất tích lũy (rank-dependent weighting) cho kết cục
\(x_i\) được tính bằng: Đối với miền lời (\(i > 0\)): \[
\] \pi\_i\^{}+ = w\textsuperscript{+\left(\sum\_\{j=i\}}n p\_j\right) -
w\textsuperscript{+\left(\sum\_\{j=i+1\}}n p\_j\right) \[
\] Đối với miền lỗ (\(i < 0\)): \[
\] \pi\_i\^{}- = w\textsuperscript{-\left(\sum\_\{j=-m\}}i p\_j\right) -
w\textsuperscript{-\left(\sum\_\{j=-m\}}\{i-1\} p\_j\right) \[
\] Trong đó, hàm trọng số \(w(p)\) thường có dạng ngược chữ S (inverse
S-shaped), ví dụ như dạng Prelec (1998):
\(w(p) = \exp(-(-\ln p)^\gamma)\) với \(\gamma \in (0,1)\). Điều này cho
thấy cá nhân đánh giá cao quá mức các xác suất nhỏ (overweight small
probabilities) và đánh giá thấp các xác suất lớn (underweight large
probabilities).

Hữu dụng tổng thể của triển vọng \(f\) là: \[
\] V(f) = \sum\_\{i=-m\}\^{}\{-1\} \pi\emph{i\^{}- v(x\_i) +
\sum}\{i=1\}\^{}n \pi\_i\^{}+ v(x\_i) \[
\]

\subsubsection{16.1.2 Điểm mù thường gặp (Cognitive Traps \&
Limitations)}\label{ux111iux1ec3m-muxf9-thux1b0ux1eddng-gux1eb7p-cognitive-traps-limitations-111}

Nhiều nhà nghiên cứu lầm tưởng rằng CPT chỉ là EUT với một hàm hữu dụng
phi tuyến tính chứa điểm gãy (kink). Tuy nhiên, sự khác biệt căn bản nằm
ở việc biến đổi xác suất không tuyến tính (non-linear probability
weighting) thông qua phân phối tích lũy (rank-dependence) nhằm tránh vi
phạm nguyên lý chi phối bậc nhất (first-order stochastic dominance). Nếu
chỉ dùng \(w(p_i)\) một cách độc lập cho từng kết cục như Prospect
Theory (1979) ban đầu, mô hình sẽ dẫn đến việc lựa chọn các triển vọng
bị chi phối (dominated prospects).

\subsubsection{16.1.3 Ứng dụng thực tế (Practical
Applications)}\label{ux1ee9ng-dux1ee5ng-thux1ef1c-tux1ebf-practical-applications-67}

Trong tài chính, nhà đầu tư thường có xu hướng bán sớm các cổ phiếu đang
sinh lời và giữ lại các cổ phiếu đang thua lỗ. Giả sử nhà đầu tư mua cổ
phiếu ở giá tham chiếu \(P_0 = 100\). Hiện tại giá là \(P_1 = 110\). Họ
đối mặt với việc chốt lời (nhận chắc chắn \(+10\)) hoặc giữ lại với xác
suất \(50\%\) giá lên \(120\) (\(+20\)) và \(50\%\) giá về \(100\)
(\(0\)). Với \(\alpha = 0.88\), giá trị chốt lời:
\(v(10) = 10^{0.88} \approx 7.58\). Giữ lại:
\(\pi^+ v(20) + (1-\pi^+)v(0)\). Với \(w(0.5) \approx 0.42\), giá trị
giữ lại là \(0.42 \times 20^{0.88} \approx 0.42 \times 13.93 = 5.85\).
Vì \(7.58 > 5.85\), nhà đầu tư chọn chốt lời. Ngược lại, nếu giá hiện
tại là \(P_1 = 90\) (lỗ 10), họ có xu hướng đánh cược để gỡ hòa do hàm
lồi ở miền lỗ, dẫn đến việc giữ lại cổ phiếu thua lỗ.

\begin{center}\rule{0.5\linewidth}{0.5pt}\end{center}

\subsection{16.2 Mô hình Chiết khấu Phi Chuẩn (Non-Standard Discounting
Models)}\label{muxf4-huxecnh-chiux1ebft-khux1ea5u-phi-chuux1ea9n-non-standard-discounting-models}

\subsubsection{16.2.1 Cốt lõi Lý thuyết (Core
Theory)}\label{cux1ed1t-luxf5i-luxfd-thuyux1ebft-core-theory-61}

Chiết khấu chuẩn (exponential discounting) giả định tỷ lệ chiết khấu
giữa hai thời kỳ bất kỳ cách nhau \(t\) đơn vị thời gian là không đổi
(time-consistent). Mô hình chiết khấu Hyperbolic, đặc biệt là dạng
\(\beta-\delta\) (Quasi-hyperbolic discounting) do Laibson (1997) đề
xuất, nắm bắt hiện tượng \textbf{thiên kiến hiện tại} (present bias),
khi cá nhân chiết khấu cực kỳ nặng nề giữa ``ngay bây giờ'' và ``ngày
mai'', nhưng chiết khấu đều đặn giữa các ngày trong tương lai.

Hữu dụng liên thời gian tại thời điểm \(t=0\) của chuỗi tiêu dùng
\(\{c_t\}\) là: \[
\] U\_0 = u(c\_0) + \beta \sum\_\{t=1\}\^{}\infty \delta\^{}t u(c\_t) \[
\] Trong đó: - \(\delta \in (0,1)\): hệ số chiết khấu chuẩn (long-run
discount factor). - \(\beta \in (0,1]\): tham số thiên kiến hiện tại
(short-run discount factor). Khi \(\beta = 1\), ta trở lại mô hình
chuẩn.

Sự mâu thuẫn thời gian (time-inconsistency) xuất hiện. Tại \(t=0\), cá
nhân sẵn sàng đánh đổi giữa \(t=1\) và \(t=2\) với tỷ lệ \(\delta\).
Nhưng khi \(t=1\) đến, tỷ lệ đánh đổi lại trở thành \(\beta\delta\).

Định nghĩa cân bằng phụ thuộc vào nhận thức của cá nhân về sự thiên kiến
của mình: 1. \textbf{Naive agents} (Ngây thơ): Không nhận thức được mình
có \(\beta < 1\) trong tương lai, luôn tin rằng tương lai mình sẽ hành
xử với \(\beta = 1\). 2. \textbf{Sophisticated agents} (Khôn ngoan):
Nhận thức rõ mình có \(\beta < 1\) và có thể sử dụng các công cụ cam kết
(commitment devices). Cân bằng của họ là một Trò chơi Hoàn hảo Cân bằng
(Subgame Perfect Equilibrium) giữa các ``bản ngã'' (selves) tại các thời
điểm khác nhau.

\subsubsection{16.2.2 Điểm mù thường gặp (Cognitive Traps \&
Limitations)}\label{ux111iux1ec3m-muxf9-thux1b0ux1eddng-gux1eb7p-cognitive-traps-limitations-112}

Sự nhầm lẫn phổ biến là đánh đồng việc thiếu tự chủ (self-control
problem) với tỷ lệ chiết khấu cao (\(\delta\) nhỏ). Một người có
\(\delta\) thấp đơn giản là người thiếu kiên nhẫn một cách nhất quán.
Một người có \(\beta < 1\) mới là người trì hoãn: hôm nay họ lập kế
hoạch ngày mai sẽ tiết kiệm, nhưng khi ngày mai đến, họ lại tiêu dùng
hết.

\subsubsection{16.2.3 Ứng dụng thực tế (Practical
Applications)}\label{ux1ee9ng-dux1ee5ng-thux1ef1c-tux1ebf-practical-applications-68}

Giả sử chi phí của việc làm bài tập là \(c=10\), lợi ích trong tương lai
quy về hiện tại là \(v=15\). Nếu làm ngay hôm nay (t=0), hữu dụng là
\(-10 + \beta \delta \times 15\). Nếu để ngày mai (t=1), hữu dụng đánh
giá tại t=0 là \(\beta \delta (-10 + \delta \times 15)\). Với
\(\delta=1, \beta=0.5\): - Làm ngay: \(-10 + 0.5 \times 15 = -2.5\) -
Làm ngày mai: \(0.5 \times (-10 + 15) = 2.5\) Cá nhân quyết định: ``Để
ngày mai làm!''. Nhưng khi ngày mai đến, bài toán quay lại y hệt, và họ
tiếp tục trì hoãn.

\begin{center}\rule{0.5\linewidth}{0.5pt}\end{center}

\subsection{16.3 Kinh tế học Thần kinh
(Neuroeconomics)}\label{kinh-tux1ebf-hux1ecdc-thux1ea7n-kinh-neuroeconomics}

\subsubsection{16.3.1 Cốt lõi Lý thuyết (Core
Theory)}\label{cux1ed1t-luxf5i-luxfd-thuyux1ebft-core-theory-62}

Kinh tế học thần kinh cung cấp nền tảng sinh học cho các mô hình hành vi
thông qua Lý thuyết Tiến trình Kép (Dual-process theory). Kahneman định
nghĩa Hệ thống 1 (System 1) là tự động, cảm xúc và nhanh chóng, trong
khi Hệ thống 2 (System 2) là cân nhắc, logic và tốn năng lượng.

Mô hình Fudenberg-Levine (2006) hình thức hóa hệ thống này dưới dạng
``bản ngã ngắn hạn'' (short-run self) bị chi phối bởi cảm xúc và ``bản
ngã dài hạn'' (long-run self) có khả năng tự chủ nhưng phải chịu chi phí
nhận thức (cognitive cost). Bản ngã dài hạn tối đa hóa: \[
\] U = \sum\_\{t=0\}\^{}\infty \delta\^{}t {[}u(a\_t, r\_t) -
\Gamma(a\_t, r\_t){]} \[
\] Trong đó \(u(a_t, r_t)\) là hữu dụng từ hành động \(a_t\) khi đối mặt
với kích thích \(r_t\) (reward cue). \(\Gamma\) là chi phí tự kiểm soát
(cost of self-control) khi bản ngã dài hạn muốn ép bản ngã ngắn hạn làm
trái với sự thôi thúc (urge).

Về mặt thần kinh, sai số dự đoán phần thưởng (Reward Prediction Error -
RPE) do tín hiệu dopamine điều khiển được hình thức hóa tương tự như quá
trình cập nhật Bayes trong học tập tăng cường (Reinforcement Learning):
\[
\] \Delta V(s) = \alpha [R(s) + \gamma V(s') - V(s)] \[
\] với \([R(s) + \gamma V(s') - V(s)]\) chính là RPE.

\subsubsection{16.3.2 Điểm mù thường gặp (Cognitive Traps \&
Limitations)}\label{ux111iux1ec3m-muxf9-thux1b0ux1eddng-gux1eb7p-cognitive-traps-limitations-113}

Nhiều nhà kinh tế truyền thống bác bỏ neuroeconomics, cho rằng ``không
cần thiết phải nhìn vào bên trong hộp đen của bộ não'' (as-if
assumption). Tuy nhiên, điểm mù ở đây là giả định as-if không thể dự báo
chính xác hành vi khi có các yếu tố ngoại sinh làm suy giảm nhận thức
(ví dụ: stress làm tăng chi phí \(\Gamma\), dẫn đến giảm tự chủ), điều
mà mô hình Fudenberg-Levine giải thích được một cách tự nhiên.

\subsubsection{16.3.3 Ứng dụng thực tế (Practical
Applications)}\label{ux1ee9ng-dux1ee5ng-thux1ef1c-tux1ebf-practical-applications-69}

Trong mô hình của Bernheim và Rangel (2004) về nghiện (addiction), não
bộ có hai trạng thái: ``nóng'' (hot) và ``lạnh'' (cold). Ở trạng thái
lạnh, cá nhân tối đa hóa hữu dụng duy lý. Tuy nhiên, nếu gặp môi trường
kích hoạt (cue), não chuyển sang trạng thái nóng, nơi việc tiêu dùng
chất gây nghiện được Hệ thống 1 ra lệnh một cách phản xạ (visceral
override) bất chấp Hệ thống 2. Marketing khai thác điều này bằng cách
đặt các mặt hàng kích thích (đồ ngọt) ở quầy thu ngân để kích hoạt xung
động khi tài nguyên nhận thức của khách hàng đã cạn kiệt.

\begin{center}\rule{0.5\linewidth}{0.5pt}\end{center}

\subsection{16.4 Kinh tế học Xã hội \& Sở thích về Công bằng (Social
Preferences \&
Fairness)}\label{kinh-tux1ebf-hux1ecdc-xuxe3-hux1ed9i-sux1edf-thuxedch-vux1ec1-cuxf4ng-bux1eb1ng-social-preferences-fairness}

\subsubsection{16.4.1 Cốt lõi Lý thuyết (Core
Theory)}\label{cux1ed1t-luxf5i-luxfd-thuyux1ebft-core-theory-63}

Con người không chỉ quan tâm đến lợi ích bản thân mà còn quan tâm đến
công bằng. Mô hình ác cảm với sự bất bình đẳng (Inequality Aversion) của
Fehr và Schmidt (1999) hình thức hóa điều này. Xét nhóm \(N\) người,
phân bổ tiền là \(x = (x_1, \dots, x_N)\). Hữu dụng của cá nhân \(i\)
là:

\[
\] U\_i(x) = x\_i - \frac{\alpha_i}{N-1} \sum\emph{\{j \neq i\}
\max(x\_j - x\_i, 0) - \frac{\beta_i}{N-1} \sum}\{j \neq i\} \max(x\_i -
x\_j, 0) \[
\]

Trong đó: - \(\alpha_i \geq 0\): Mức độ khó chịu khi bị thiệt thòi
(disadvantageous inequality). - \(0 \leq \beta_i < 1\) và
\(\beta_i \leq \alpha_i\): Mức độ khó chịu khi được hưởng lợi nhiều hơn
người khác (advantageous inequality). Ràng buộc \(\beta_i < 1\) đảm bảo
người \(i\) thà đốt tiền của mình đi còn hơn là cho không nếu việc đó
làm giảm sự chênh lệch có lợi cho họ, nhưng họ không sẵn sàng hi sinh 1
đồng để giảm bất bình đẳng có lợi lớn hơn 1 đồng. Ràng buộc
\(\beta_i \leq \alpha_i\) phản ánh việc con người ghét bị thiệt thòi hơn
là ghét việc người khác bị thiệt.

Ngoài ra, mô hình của Rabin (1993) đưa vào yếu tố \textbf{có qua có lại}
(reciprocity), trong đó hữu dụng phụ thuộc vào niềm tin bậc hai
(second-order beliefs) về ý định của đối phương.

\subsubsection{16.4.2 Điểm mù thường gặp (Cognitive Traps \&
Limitations)}\label{ux111iux1ec3m-muxf9-thux1b0ux1eddng-gux1eb7p-cognitive-traps-limitations-114}

Một điểm mù lớn là cho rằng kết quả trong phòng thí nghiệm (ví dụ Trò
chơi Tối hậu thư - Ultimatum Game) sẽ phản ánh chính xác thị trường thực
tế. Trên thực tế, thiết chế thị trường cạnh tranh có thể làm mờ đi sở
thích công bằng. Fehr và Schmidt chứng minh bằng toán học rằng trong một
thị trường có thiết kế kiểu đấu giá (competitive market), chỉ cần một tỷ
lệ nhỏ những người hoàn toàn ích kỷ (\(\alpha_i = \beta_i = 0\)), họ có
thể chi phối toàn bộ kết quả, đẩy giá về mức cân bằng Walrasian cạnh
tranh khốc liệt.

\subsubsection{16.4.3 Ứng dụng thực tế (Practical
Applications)}\label{ux1ee9ng-dux1ee5ng-thux1ef1c-tux1ebf-practical-applications-70}

Giả sử một công ty muốn trả lương cho nhân viên. Theo lý thuyết chuẩn,
công ty trả mức lương tối thiểu thỏa mãn điều kiện tham gia
(Participation Constraint). Tuy nhiên, theo lý thuyết quà tặng trao đổi
(Gift Exchange), công ty trả mức \(w > w_{min}\), và nhân viên có
\(\beta > 0\) (cảm thấy bất công nếu nhận lương cao mà không cống hiến)
sẽ đáp trả bằng nỗ lực \(e > e_{min}\). Sự hiểu biết về sở thích xã hội
giúp giải thích hiện tượng tiền lương cứng nhắc (wage rigidity) và tại
sao các công ty không giảm lương trong suy thoái kinh tế.

\begin{center}\rule{0.5\linewidth}{0.5pt}\end{center}

\subsection{16.5 Thiết kế Chính sách Hành vi (Behavioral Policy
Design)}\label{thiux1ebft-kux1ebf-chuxednh-suxe1ch-huxe0nh-vi-behavioral-policy-design}

\subsubsection{16.5.1 Cốt lõi Lý thuyết (Core
Theory)}\label{cux1ed1t-luxf5i-luxfd-thuyux1ebft-core-theory-64}

Dựa trên những hạn chế về sự duy lý, Thaler và Sunstein (2008) đề xuất
khái niệm \textbf{Nudge} (Cú hích), được định nghĩa là bất kỳ khía cạnh
nào của ``cấu trúc lựa chọn'' (choice architecture) thay đổi hành vi của
con người một cách dự đoán được mà không cấm đoán bất kỳ lựa chọn nào
hoặc không làm thay đổi đáng kể động cơ kinh tế. Triết lý đằng sau là
Chủ nghĩa bảo trợ tự do (Libertarian Paternalism).

Hình thức hóa hiệu ứng mặc định (Default Effects): Giả sử cá nhân chọn
giữa hai trạng thái \(x \in \{0, 1\}\). Hữu dụng nội tại là \(u(x)\).
Tuy nhiên, việc thay đổi khỏi trạng thái mặc định \(x_D\) đòi hỏi một
chi phí ma sát (friction cost) \(k > 0\), bao gồm cả chi phí nhận thức
(cognitive cost) và chi phí điền đơn (sludge). Bài toán tối đa hóa: \[
\] \max\_\{x \in \{0, 1\}\} U(x) = u(x) - k \cdot \mathbb{I}(x
\neq x\_D) \[
\] Nếu \(x_D = 0\) (opt-in policy), cá nhân sẽ chỉ chọn \(1\) nếu
\(u(1) - u(0) > k\). Nếu chính phủ đổi sang \(x_D = 1\) (opt-out
policy), cá nhân sẽ chỉ đổi lại \(0\) nếu \(u(0) - u(1) > k\). Đối với
những cá nhân có \(|u(1) - u(0)| < k\), chính sách mặc định hoàn toàn
định hình sự lựa chọn của họ.

\subsubsection{16.5.2 Điểm mù thường gặp (Cognitive Traps \&
Limitations)}\label{ux111iux1ec3m-muxf9-thux1b0ux1eddng-gux1eb7p-cognitive-traps-limitations-115}

Nhiều nhà hoạch định chính sách lạm dụng Nudge, dẫn đến \textbf{Sludge}
- những rào cản hành vi được dựng lên có chủ ý nhằm ngăn cản người dân
tiếp cận quyền lợi (ví dụ: quy trình hủy đăng ký dịch vụ quá phức tạp).
Điểm mù là tin rằng mọi Cú hích đều nhằm tối đa hóa phúc lợi xã hội
(social welfare maximization), trong khi thực tế có thể bị các nhóm lợi
ích thao túng.

\subsubsection{16.5.3 Ứng dụng thực tế (Practical
Applications)}\label{ux1ee9ng-dux1ee5ng-thux1ef1c-tux1ebf-practical-applications-71}

Chính phủ Anh, thông qua Behavioural Insights Team (BIT), đã áp dụng
khuôn khổ EAST (Easy, Attractive, Social, Timely) để tăng cường tiết
kiệm hưu trí. Thay vì yêu cầu người lao động điền đơn tham gia (Opt-in),
họ tự động đăng ký cho người lao động vào quỹ hưu trí (Automatic
Enrolment - Opt-out). Mô hình \(k\) ma sát đã thành công rực rỡ, đưa tỷ
lệ tham gia từ dưới 50\% lên trên 90\%.

\begin{center}\rule{0.5\linewidth}{0.5pt}\end{center}

\subsection{16.6 Phê bình \& Giới hạn của Kinh tế học Hành vi (Critiques
\&
Limits)}\label{phuxea-buxecnh-giux1edbi-hux1ea1n-cux1ee7a-kinh-tux1ebf-hux1ecdc-huxe0nh-vi-critiques-limits}

\subsubsection{16.6.1 Cốt lõi Lý thuyết (Core
Theory)}\label{cux1ed1t-luxf5i-luxfd-thuyux1ebft-core-theory-65}

Kinh tế học hành vi đối mặt với nhiều phê phán nghiêm khắc từ kinh tế
học dòng chính (Neoclassical Economics), đặc biệt về độ giá trị ngoại
tại (external validity) và sự chọn lọc thị trường (market selection).

List (2003) chứng minh rằng hiệu ứng sở hữu (Endowment Effect) rất mạnh
ở những người nghiệp dư nhưng gần như biến mất ở những nhà buôn chuyên
nghiệp. Quá trình học tập (learning) và giải trừ thiên kiến (de-biasing)
khiến các tác nhân hội tụ về sự duy lý.

Hơn nữa, Levine (2009) lập luận dựa trên lý thuyết tiến hóa của thị
trường (evolutionary market selection): Trong một thị trường cạnh tranh
có rủi ro, những tác nhân hành động phi duy lý (ví dụ: giao dịch quá mức
- overconfidence, hay vi phạm quy tắc Bayes) sẽ dần đánh mất tài sản vào
tay những tác nhân duy lý. Giả sử tác nhân duy lý có hàm cầu tài sản
\(D_R(p)\) và phi duy lý có hàm cầu \(D_I(p)\). Hàm động lực học tích
lũy tài sản (wealth dynamics): \[
\] W\_I(t+1) = W\_I(t) \cdot (1 - \rho) \[
\] với \(\rho > 0\) là tỷ lệ hao hụt do các quyết định sai lầm. Khi
\(t \to \infty\), \(W_I(t) \to 0\). Trọng lượng của họ trong việc xác
định giá cả thị trường sẽ triệt tiêu. Do đó, các dị thường hành vi
(behavioral anomalies) không thể tồn tại lâu dài ở cấp độ vĩ mô.

\subsubsection{16.6.2 Điểm mù thường gặp (Cognitive Traps \&
Limitations)}\label{ux111iux1ec3m-muxf9-thux1b0ux1eddng-gux1eb7p-cognitive-traps-limitations-116}

Một số nhà kinh tế học hành vi bỏ qua \textbf{cuộc khủng hoảng lặp lại}
(Replication Crisis) trong tâm lý học. Rất nhiều hiện tượng như ``social
priming'' không thể tái tạo được (failed to replicate) trong các mẫu
lớn, dẫn đến việc xây dựng mô hình toán học trên một nền tảng thực
nghiệm thiếu vững chắc.

\subsubsection{16.6.3 Ứng dụng thực tế (Practical
Applications)}\label{ux1ee9ng-dux1ee5ng-thux1ef1c-tux1ebf-practical-applications-72}

Sự hiểu biết sâu sắc nhất dành cho nghiên cứu sinh tiến sĩ là sự kết
hợp: Kinh tế học chuẩn là công cụ tuyệt vời để dự báo cân bằng dài hạn
của các thị trường tài chính thanh khoản cao và các quyết định lặp lại
nhiều lần. Kinh tế học hành vi lại không thể thay thế trong việc thiết
kế vi mô (mechanism design), bảo vệ người tiêu dùng bối rối trước các
quyết định một lần (one-shot decisions) như mua bảo hiểm y tế, vay thế
chấp, hay thiết kế chính sách hưu trí. Cả hai khung lý thuyết đều là
công cụ (toolkit) chứ không phải tôn giáo, và nhà nghiên cứu cần biết áp
dụng đúng công cụ cho đúng bài toán.

\begin{center}\rule{0.5\linewidth}{0.5pt}\end{center}

\emph{(Xem tiếp Chương 17: Kinh tế học Tổ chức Ngành Nâng cao)}

\section{CHƯƠNG 17: KINH TẾ HỌC TỔ CHỨC NGÀNH NÂNG
CAO}\label{chux1b0ux1a1ng-17-kinh-tux1ebf-hux1ecdc-tux1ed5-chux1ee9c-nguxe0nh-nuxe2ng-cao}

\subsection{17.1 Phân biệt Giá Hoàn hảo \& Bậc 1/2/3 (Price
Discrimination \& Pigovian
Taxonomy)}\label{phuxe2n-biux1ec7t-giuxe1-houxe0n-hux1ea3o-bux1eadc-123-price-discrimination-pigovian-taxonomy}

\subsubsection{17.1.1 Cốt lõi Lý
thuyết}\label{cux1ed1t-luxf5i-luxfd-thuyux1ebft}

Phân biệt giá (Price Discrimination) là chiến lược tối đa hóa lợi nhuận
của doanh nghiệp độc quyền thông qua việc tính các mức giá khác nhau cho
các nhóm khách hàng khác nhau đối với cùng một sản phẩm, dựa trên sự
khác biệt về độ sẵn sàng chi trả (Willingness to Pay - WTP). Theo phân
loại kinh điển của A.C. Pigou (1920), phân biệt giá được chia thành ba
bậc:

\textbf{1. Phân biệt giá bậc 1 (Hoàn hảo - First-degree/Perfect Price
Discrimination):} Doanh nghiệp có thông tin hoàn hảo về WTP của từng cá
nhân và thu toàn bộ thặng dư tiêu dùng (Consumer Surplus - CS). Nếu hàm
cầu nghịch là \(P(q)\), doanh nghiệp sẽ đặt giá \(p = P(q)\) cho từng
đơn vị. Lợi nhuận đạt mức tối đa:
\[ \pi = \int_{0}^{q^*} [P(q) - MC(q)] dq \] trong đó \(q^*\) là mức sản
lượng tại \(P(q^*) = MC(q^*)\). Doanh nghiệp cũng có thể sử dụng biểu
giá hai phần (Two-part tariff) \(T(q) = A + pq\), với phí cố định
\(A = CS\) và giá biên \(p = MC\). Mặc dù thặng dư tiêu dùng bằng 0,
tổng thặng dư (Total Surplus) đạt mức tối đa (Pareto hiệu quả).

\textbf{2. Phân biệt giá bậc 2 (Phi tuyến - Second-degree/Nonlinear
Screening):} Xảy ra khi thông tin bất cân xứng: doanh nghiệp biết phân
phối các loại khách hàng nhưng không nhận diện được từng cá nhân. Doanh
nghiệp thiết kế một danh mục các gói hợp đồng \((q_i, T_i)\) để khách
hàng tự lựa chọn (menu pricing). Giả sử có hai loại khách hàng có tham
số định giá \(\theta \in \{\theta_L, \theta_H\}\) với
\(\theta_H > \theta_L\) (điều kiện Spence-Mirrlees). Hàm hữu dụng:
\(U(q, T, \theta) = \theta V(q) - T\). Bài toán tối ưu của doanh nghiệp
(Screening Problem) là tối đa hóa lợi nhuận kỳ vọng, tuân theo các ràng
buộc: - Ràng buộc tham gia (Individual Rationality - IR):
\(\theta_i V(q_i) - T_i \ge 0 \quad \forall i \in \{L, H\}\) - Ràng buộc
tương thích khuyến khích (Incentive Compatibility - IC):
\(\theta_i V(q_i) - T_i \ge \theta_i V(q_j) - T_j \quad \forall i, j\)
Tại điểm cân bằng (Separating Equilibrium), ràng buộc \(IR_L\) và
\(IC_H\) thắt chặt (bind). Người tiêu dùng loại cao (\(\theta_H\)) nhận
được mức thông tin tô (Information Rent) \(\Delta \theta V(q_L) > 0\) và
tiêu dùng lượng \(q_H\) tối ưu Pareto (\(\theta_H V'(q_H) = MC\)), trong
khi người loại thấp (\(\theta_L\)) không có thông tin tô nhưng bị giới
hạn sản lượng \(q_L\) dưới mức tối ưu (Downward Distortion ở đáy để giảm
thông tin tô cho loại cao).

\textbf{3. Phân biệt giá bậc 3 (Phân khúc thị trường - Third-degree):}
Doanh nghiệp có thể quan sát một tín hiệu ngoại sinh (ví dụ: thẻ sinh
viên, độ tuổi) để chia thị trường thành các phân khúc
\(i \in \{1, 2, ..., N\}\) với hàm cầu \(q_i(p_i)\). Bài toán tối ưu:
\[ \max_{p_1, ..., p_N} \sum_{i=1}^N [p_i q_i(p_i) - C(q_i)] \] Điều
kiện bậc nhất (FOC) cho từng phân khúc \(i\):
\[ MR_i(q_i) = MC \implies p_i \left( 1 - \frac{1}{\epsilon_i} \right) = MC \]
trong đó
\(\epsilon_i = -\frac{\partial q_i}{\partial p_i} \frac{p_i}{q_i}\) là
độ co giãn của cầu theo giá. Phân khúc có độ co giãn thấp (ít nhạy cảm
với giá) sẽ chịu mức giá cao hơn:
\(\epsilon_1 < \epsilon_2 \implies p_1 > p_2\).

\subsubsection{17.1.2 Điểm mù thường
gặp}\label{ux111iux1ec3m-muxf9-thux1b0ux1eddng-gux1eb7p-13}

\subsubsection{17.1.2 Điểm mù thường gặp (Cognitive Traps \&
Limitations)}\label{ux111iux1ec3m-muxf9-thux1b0ux1eddng-gux1eb7p-cognitive-traps-limitations-117}

Một sai lầm phổ biến là cho rằng phân biệt giá luôn gây tổn hại cho nền
kinh tế do chuyển giao thặng dư từ người tiêu dùng sang doanh nghiệp.
Trong thực tế, phân biệt giá bậc 3 có thể làm tăng tổng phúc lợi (Total
Welfare) nếu và chỉ nếu nó làm tăng tổng sản lượng bán ra so với định
giá đồng nhất (uniform pricing). Nếu sản lượng không đổi, phân biệt giá
bậc 3 làm giảm phúc lợi do sự phân bổ sai lệch (misallocation) hàng hóa
giữa các người tiêu dùng. Ngược lại, phân biệt giá bậc 1 tối đa hóa hoàn
toàn tổng phúc lợi (loại bỏ mất mát vô ích - Deadweight Loss).

\subsubsection{17.1.2 Điểm mù thường gặp (Cognitive Traps \&
Limitations)}\label{ux111iux1ec3m-muxf9-thux1b0ux1eddng-gux1eb7p-cognitive-traps-limitations-118}

Sinh viên thường cố gắng giải hệ phương trình với tất cả các ràng buộc
IC và IR được cho là có dấu bằng. Tuy nhiên, theo lý thuyết cơ chế
(Mechanism Design), chỉ có ràng buộc tham gia của nhóm có định giá thấp
(\(IR_L\)) và ràng buộc tương thích khuyến khích của nhóm có định giá
cao (\(IC_H\)) mới thực sự ràng buộc (bind) tại điểm tối ưu (the binding
constraints). Ràng buộc \(IR_H\) và \(IC_L\) thường được thỏa mãn một
cách tự động (slack).

\subsubsection{17.1.3 Ứng dụng thực
tế}\label{ux1ee9ng-dux1ee5ng-thux1ef1c-tux1ebf-7}

Trong ngành hàng không Việt Nam (như Vietnam Airlines, Vietjet Air),
phân biệt giá được áp dụng triệt để. - \textbf{Bậc 3:} Giảm giá cho sinh
viên, người cao tuổi, hoặc định giá khác nhau giữa hành khách mua vé sớm
(thường là du lịch, \(\epsilon\) cao) và hành khách mua sát giờ bay
(thương gia, \(\epsilon\) thấp). - \textbf{Bậc 2:} Hệ thống vé được chia
thành nhiều hạng (Eco, Deluxe, Business) với các quyền lợi (hành lý, đổi
vé, phòng chờ) và mức giá khác nhau. Các hãng hàng không sử dụng chất
lượng dịch vụ (quality degradation cho hạng rẻ) như một cơ chế sàng lọc
(screening) để phân biệt các loại hành khách theo mức sẵn sàng chi trả,
ép những người sẵn sàng trả cao chọn hạng Business thay vì Eco. Điều này
giải thích tại sao đôi khi các hãng hàng không cố tình làm giảm trải
nghiệm hạng phổ thông (ví dụ: không cho chọn chỗ) để duy trì ràng buộc
tương thích khuyến khích (\(IC\)).

\subsubsection{17.1.4 Ví dụ minh họa tính
toán}\label{vuxed-dux1ee5-minh-hux1ecda-tuxednh-touxe1n}

\textbf{Bài toán (Phân biệt giá bậc 3):} Một hãng độc quyền đối diện với
hai thị trường phân biệt: - Thị trường A (nội địa): \(P_A = 100 - Q_A\)
- Thị trường B (quốc tế): \(P_B = 80 - 2Q_B\) Chi phí biên không đổi
\(MC = 20\). Tìm giá, sản lượng và lợi nhuận khi có phân biệt giá, so
với khi phải áp dụng giá đồng nhất (uniform pricing).

\textbf{Giải:} 1. Khi có phân biệt giá bậc 3: Thị trường A:
\(MR_A = 100 - 2Q_A = MC = 20 \implies Q_A^* = 40, P_A^* = 60\). Thị
trường B:
\(MR_B = 80 - 4Q_B = MC = 20 \implies Q_B^* = 15, P_B^* = 50\). Lợi
nhuận:
\(\pi_D = (P_A - MC)Q_A + (P_B - MC)Q_B = (60-20)40 + (50-20)15 = 1600 + 450 = 2050\).

\begin{enumerate}
\def\labelenumi{\arabic{enumi}.}
\setcounter{enumi}{1}
\tightlist
\item
  Khi định giá đồng nhất (không phân biệt giá): Hàm cầu tổng:
\end{enumerate}

\begin{itemize}
\tightlist
\item
  Nếu \(P > 80\): chỉ có A mua, \(Q = 100 - P\).
\item
  Nếu \(P \le 80\): cả hai mua,
  \(Q = Q_A + Q_B = (100 - P) + (40 - 0.5P) = 140 - 1.5P\). Hàm cầu
  nghịch tổng hợp (khi \(P \le 80\)):
  \(P = \frac{140 - Q}{1.5} \approx 93.33 - 0.67Q\). Doanh thu biên:
  \(MR = 93.33 - 1.33Q\). Đặt
  \(MR = MC \implies 93.33 - 1.33Q = 20 \implies 1.33Q = 73.33 \implies Q^* \approx 55\).
  Giá đồng nhất: \(P^* = 93.33 - 0.67(55) = 56.67\). Lợi nhuận:
  \(\pi_U = (56.67 - 20)55 = 2016.85\). So sánh:
  \(\pi_D = 2050 > \pi_U = 2016.85\). Phân biệt giá mang lại lợi nhuận
  cao hơn. (Lưu ý: tại \(P^* = 56.67 \le 80\), giả định cả hai thị
  trường đều mua là hợp lệ).
\end{itemize}

\begin{center}\rule{0.5\linewidth}{0.5pt}\end{center}

\subsection{17.2 Khác biệt hóa Sản phẩm \& Mô hình Địa tính Hotelling
(Product Differentiation \& Spatial
Competition)}\label{khuxe1c-biux1ec7t-huxf3a-sux1ea3n-phux1ea9m-muxf4-huxecnh-ux111ux1ecba-tuxednh-hotelling-product-differentiation-spatial-competition}

\subsubsection{17.2.1 Cốt lõi Lý
thuyết}\label{cux1ed1t-luxf5i-luxfd-thuyux1ebft-1}

Mô hình Thành phố tuyến tính (Linear City Model) của Hotelling (1929) là
khuôn khổ nền tảng cho lý thuyết khác biệt hóa sản phẩm theo chiều ngang
(horizontal differentiation). Khác biệt hóa ngang phản ánh sự đa dạng
trong sở thích của người tiêu dùng (ví dụ: màu sắc, vị trí địa lý),
ngược với khác biệt hóa dọc (chất lượng).

Giả sử người tiêu dùng phân bố đều trên một đoạn thẳng \([0, 1]\) với
mật độ 1. Có hai doanh nghiệp, Firm 1 tại vị trí \(a \ge 0\) và Firm 2
tại vị trí \(1 - b\) (với \(a \le 1 - b\)). Khách hàng tại vị trí
\(x \in [0, 1]\) phải chịu chi phí vận chuyển để tiêu dùng.

\textbf{Mô hình chi phí vận chuyển tuyến tính (Hotelling gốc 1929):} Chi
phí vận chuyển là \(\tau \cdot d\). Hữu dụng ròng của khách hàng \(x\):
- Mua từ Firm 1: \(U_1 = V - p_1 - \tau |x - a|\) - Mua từ Firm 2:
\(U_2 = V - p_2 - \tau |1 - b - x|\) (Giả định \(V\) đủ lớn để tất cả
mọi người đều mua - thị trường được bao phủ hoàn toàn). Khách hàng bàng
quan \(\bar{x}\) nằm giữa \(a\) và \(1-b\):
\[ p_1 + \tau (\bar{x} - a) = p_2 + \tau (1 - b - \bar{x}) \implies \bar{x} = \frac{p_2 - p_1}{2\tau} + \frac{1 - a - b}{2} + a \]
Tuy nhiên, d'Aspremont, Gabszewicz \& Thisse (1979) đã chứng minh rằng
nếu \(a\) và \(b\) quá gần nhau, cân bằng Nash theo giá thuần túy không
tồn tại do lợi ích từ việc hạ giá một chút để chiếm toàn bộ thị trường
vượt quá lợi ích chia sẻ thị trường.

\textbf{Mô hình chi phí vận chuyển bậc hai (d'Aspremont et al., 1979):}
Chi phí vận chuyển là \(t \cdot d^2\). Hữu dụng ròng: -
\(U_1 = V - p_1 - t(x - a)^2\) - \(U_2 = V - p_2 - t(1 - b - x)^2\)
Khách hàng bàng quan \(\bar{x}\):
\[ p_1 + t(x - a)^2 = p_2 + t(1 - b - x)^2 \implies \bar{x} = \frac{p_2 - p_1}{2t(1 - a - b)} + \frac{a + 1 - b}{2} \]
Sản lượng của hai doanh nghiệp: \(D_1 = \bar{x}\) và
\(D_2 = 1 - \bar{x}\). Lợi nhuận (với \(MC=c\)):
\[ \pi_1 = (p_1 - c)\bar{x} \quad ; \quad \pi_2 = (p_2 - c)(1 - \bar{x}) \]
Giải điều kiện bậc nhất FOC \(\frac{\partial \pi_i}{\partial p_i} = 0\),
ta được hàm phản ứng giá. Cân bằng Nash trong giai đoạn định giá:
\[ p_1^* = c + t(1 - a - b)\left(1 + \frac{a - b}{3}\right) \]
\[ p_2^* = c + t(1 - a - b)\left(1 + \frac{b - a}{3}\right) \] Nguyên lý
Phân hóa Cực đại (Principle of Maximum Differentiation): Khi các doanh
nghiệp chọn vị trí \(a, b\) trước khi định giá, họ đối diện với hai lực
ngược chiều: - Hiệu ứng thị phần (Market-share effect): Kéo họ lại gần
trung tâm để tăng lượng khách hàng. - Hiệu ứng cạnh tranh giá
(Price-competition effect): Đẩy họ ra xa nhau để làm dịu sự khốc liệt
của cạnh tranh giá (tăng sức mạnh độc quyền cục bộ). Với chi phí bậc
hai, hiệu ứng cạnh tranh giá chiếm ưu thế, dẫn đến cân bằng tại các cực:
\(a^* = 0, b^* = 0\) (Firm 1 ở 0, Firm 2 ở 1).

Trong khi đó, Mô hình Thành phố Tròn (Circular City) của Salop (1979) mở
rộng Hotelling bằng cách cho phép sự xâm nhập thị trường nội sinh
(endogenous entry) trên một vòng tròn chu vi 1. Ở cân bằng, khoảng cách
giữa các doanh nghiệp là \(\frac{1}{n}\). Lợi nhuận của mỗi doanh nghiệp
\(\pi = \frac{t}{n^2} - F\) (với F là chi phí cố định). Số lượng doanh
nghiệp tự do gia nhập cân bằng là \(n^c = \sqrt{\frac{t}{F}}\), lớn hơn
mức tối ưu xã hội \(n^* = \frac{1}{2}\sqrt{\frac{t}{F}}\), cho thấy xu
hướng khác biệt hóa quá mức (excessive differentiation) trong cạnh tranh
độc quyền.

\subsubsection{17.2.2 Điểm mù thường
gặp}\label{ux111iux1ec3m-muxf9-thux1b0ux1eddng-gux1eb7p-14}

\subsubsection{17.2.2 Điểm mù thường gặp (Cognitive Traps \&
Limitations)}\label{ux111iux1ec3m-muxf9-thux1b0ux1eddng-gux1eb7p-cognitive-traps-limitations-119}

Nguyên lý phân hóa tối thiểu (Hai quán kem đứng cạnh nhau ở giữa bãi
biển) chỉ đúng khi giá bị cố định (không có cạnh tranh giá, chỉ cạnh
tranh không gian) hoặc chi phí vận chuyển là tuyến tính (nhưng khi đó
không có cân bằng giá). Khi giá là nội sinh và chi phí vận chuyển lồi
(bậc hai), các doanh nghiệp có động lực phân hóa tối đa (đứng ở 2 đầu
bãi biển) để tránh một cuộc chiến tranh giá tàn khốc, theo đúng mô hình
của d'Aspremont et al.~(1979).

\subsubsection{17.2.2 Điểm mù thường gặp (Cognitive Traps \&
Limitations)}\label{ux111iux1ec3m-muxf9-thux1b0ux1eddng-gux1eb7p-cognitive-traps-limitations-120}

Trong kinh tế học tổ chức ngành, \(t\) đại diện cho mức độ trung thành
với thương hiệu hoặc chi phí chuyển đổi (switching cost). Giá trị \(t\)
càng lớn nghĩa là khách hàng càng coi trọng sự phù hợp chính xác với sở
thích của họ và ngần ngại thay đổi sang sản phẩm khác. Do đó, \(t\)
chính là thước đo sức mạnh thị trường (market power); khi \(t \to 0\),
sản phẩm trở nên đồng nhất, mô hình quay về Bertrand Paradox (giá bằng
chi phí biên). Sinh viên thường quên rằng các nỗ lực marketing chính là
cố gắng làm tăng giá trị \(t\) trong tâm trí người tiêu dùng.

\subsubsection{17.2.3 Ứng dụng thực
tế}\label{ux1ee9ng-dux1ee5ng-thux1ef1c-tux1ebf-8}

Khác biệt hóa sản phẩm ngang giải thích rất rõ chiến lược định vị trên
thị trường Việt Nam. - \textbf{Ứng dụng công nghệ xe công nghệ
(Ride-hailing):} Các hãng như Grab, Be, và Gojek đối mặt với chi phí vận
chuyển không gian \(t\) tương đối thấp (ứng dụng dễ cài đặt, chi phí
chuyển đổi thấp). Vì \(t\) thấp, cạnh tranh giá rất khốc liệt. Để đối
phó, Grab cố gắng khác biệt hóa (tăng \(t\)) thông qua hệ sinh thái
(Giao đồ ăn, thanh toán Moca/ZaloPay, tích điểm GrabRewards), dịch
chuyển từ cạnh tranh thuần túy về giá sang cạnh tranh hệ sinh thái. -
\textbf{Ngành bán lẻ / Siêu thị tiện lợi:} Việc lựa chọn vị trí của
Circle K, FamilyMart, hay Winmart+ tại các ngã tư hoặc khu dân cư đông
đúc mô phỏng hoàn hảo mô hình Thành phố Tròn của Salop. Sự gia nhập thị
trường ồ ạt cho đến khi lợi nhuận biên tiệm cận chi phí cố định (\(F\)),
giải thích tại sao có rất nhiều cửa hàng tiện lợi nằm sát nhau trong một
đoạn đường ngắn ở TP.HCM hay Hà Nội (excessive entry).

\subsubsection{17.2.4 Ví dụ minh họa tính
toán}\label{vuxed-dux1ee5-minh-hux1ecda-tuxednh-touxe1n-1}

\textbf{Bài toán (Mô hình Hotelling chi phí bậc hai):} Hai doanh nghiệp
1 và 2 đặt tại hai đầu mút của một thành phố tuyến tính độ dài 1 (Firm 1
ở \(x=0\), Firm 2 ở \(x=1\)). Người tiêu dùng phân bố đều với mật độ 1.
Chi phí vận chuyển bậc hai \(t = 10\). Chi phí biên của cả hai doanh
nghiệp là \(MC = 2\). a) Tìm hàm phản ứng giá của hai doanh nghiệp. b)
Tìm giá cân bằng, sản lượng và lợi nhuận của mỗi doanh nghiệp.

\textbf{Giải:} a) Với Firm 1 ở \(a=0\) và Firm 2 ở \(b=0\) (tức là
\(1-b=1\)). Khách hàng bàng quan \(\bar{x}\) nằm giữa 0 và 1:
\[ p_1 + 10x^2 = p_2 + 10(1-x)^2 \]
\[ p_1 + 10x^2 = p_2 + 10(1 - 2x + x^2) \]
\[ p_1 = p_2 + 10 - 20x \implies x = \bar{x} = \frac{p_2 - p_1 + 10}{20} \]
Lượng cầu: \(D_1 = \bar{x} = \frac{1}{2} + \frac{p_2 - p_1}{20}\);
\(D_2 = 1 - \bar{x} = \frac{1}{2} + \frac{p_1 - p_2}{20}\). Lợi nhuận
Firm 1:
\(\pi_1 = (p_1 - 2)\left(\frac{1}{2} + \frac{p_2 - p_1}{20}\right)\).
FOC cho Firm 1 (đạo hàm theo \(p_1\) bằng 0):
\[ \frac{\partial \pi_1}{\partial p_1} = \frac{1}{2} + \frac{p_2 - p_1}{20} - \frac{1}{20}(p_1 - 2) = 0 \]
\[ 10 + p_2 - 2p_1 + 2 = 0 \implies 2p_1 = p_2 + 12 \implies p_1 = 6 + 0.5p_2 \]
Do tính đối xứng, hàm phản ứng của Firm 2 là: \(p_2 = 6 + 0.5p_1\).

\begin{enumerate}
\def\labelenumi{\alph{enumi})}
\setcounter{enumi}{1}
\tightlist
\item
  Tại cân bằng Nash, hai hàm phản ứng giao nhau. Vì đối xứng,
  \(p_1 = p_2 = p^*\).
  \[ p^* = 6 + 0.5p^* \implies 0.5p^* = 6 \implies p^* = 12 \] Giá cân
  bằng: \(p_1^* = p_2^* = 12\). (Biên lợi nhuận là
  \(p - MC = 12 - 2 = 10 = t\)). Sản lượng: \(D_1 = D_2 = 1/2\). Lợi
  nhuận: \(\pi_1 = \pi_2 = (12 - 2) \times 0.5 = 5\).
\end{enumerate}

\begin{center}\rule{0.5\linewidth}{0.5pt}\end{center}

\subsection{17.3 Cấu kết Động \& Cartel: Mô hình Green-Porter \&
Rotemberg-Saloner (Dynamic Collusion \& Secret Price
Cuts)}\label{cux1ea5u-kux1ebft-ux111ux1ed9ng-cartel-muxf4-huxecnh-green-porter-rotemberg-saloner-dynamic-collusion-secret-price-cuts}

\subsubsection{17.3.1 Cốt lõi Lý
thuyết}\label{cux1ed1t-luxf5i-luxfd-thuyux1ebft-2}

Để duy trì cấu kết (collusion) hay thiết lập một Cartel thành công, các
doanh nghiệp phải giải quyết được bài toán động (dynamic game) với sự
cám dỗ ``lật lọng'' (temptation to cheat) từ các thành viên. Trong trò
chơi lặp lại vô hạn (infinitely repeated games), sự cấu kết được duy trì
thông qua chiến lược đe dọa, nổi bật là chiến lược kích hoạt khắc nghiệt
(Grim Trigger Strategy).

Cho \(\pi^M\) là lợi nhuận cấu kết (Cartel), \(\pi^D\) là lợi nhuận khi
một bên lật lọng (Defect) bằng cách hạ giá, và \(\pi^N\) là lợi nhuận
cân bằng Cournot/Bertrand một lần (Nash). Điều kiện để duy trì cấu kết
(Incentive Compatibility Constraint) phụ thuộc vào hệ số chiết khấu
\(\delta \in (0,1)\): Giá trị hiện tại của cấu kết mãi mãi \(\ge\) Lợi
ích ngắn hạn từ việc lật lọng + Lợi nhuận trừng phạt mãi mãi
\[ \frac{\pi^M}{1 - \delta} \ge \pi^D + \frac{\delta \pi^N}{1 - \delta} \]
Rút gọn lại, cấu kết ổn định nếu các doanh nghiệp đủ kiên nhẫn:
\[ \delta \ge \bar{\delta} = \frac{\pi^D - \pi^M}{\pi^D - \pi^N} \]

Tuy nhiên, trong thực tế, các doanh nghiệp không thể quan sát giá của
nhau mà chỉ quan sát được mức cầu của riêng mình. \textbf{Mô hình
Green-Porter (1984):} Nghiên cứu trạng thái bất cân xứng thông tin
(imperfect monitoring). Khi một doanh nghiệp quan sát thấy sản lượng bán
ra của mình sụt giảm, họ không biết do đối thủ đã bí mật giảm giá
(secret price cuts) hay do cú sốc cầu tiêu cực ngoại sinh (unobservable
negative demand shocks). Để ngăn chặn lật lọng, các doanh nghiệp phải
thiết lập một ``vùng trừng phạt'' (reversionary episode). Nếu giá thị
trường rơi xuống dưới một ngưỡng \(\bar{p}\), tất cả sẽ tự động bước vào
``cuộc chiến tranh giá'' (price war - sản xuất ở mức Cournot) trong
\(T\) giai đoạn, bất kể nguyên nhân thực sự là gì, sau đó mới quay lại
cấu kết. Cân bằng của Green-Porter chứng minh rằng \textbf{chiến tranh
giá xảy ra định kỳ ngay cả trong cân bằng} như một chi phí tất yếu để
duy trì kỷ luật Cartel dưới điều kiện thông tin không hoàn hảo.

\textbf{Mô hình Rotemberg-Saloner (1986):} Nghiên cứu hành vi cấu kết
theo chu kỳ kinh tế khi cú sốc cầu là quan sát được. Trong giai đoạn
bùng nổ (Boom), nhu cầu hiện tại rất cao, do đó lợi nhuận ngắn hạn từ
việc lật lọng \(\pi^D_{boom}\) trở nên cực kỳ hấp dẫn. Ngược lại, hình
phạt trong tương lai không thay đổi nhiều nếu tính trung bình qua các
chu kỳ. Để ngăn chặn sự phá vỡ Cartel trong thời kỳ bùng nổ, Cartel phải
chủ động hạ mức giá cấu kết xuống. Kết luận ngược trực giác: \textbf{Cấu
kết khó duy trì nhất trong thời kỳ bùng nổ kinh tế, dẫn đến định giá
ngược chu kỳ (counter-cyclical pricing)} - giá có xu hướng giảm hoặc
không tăng tương xứng khi cầu cao.

\subsubsection{17.3.2 Điểm mù thường
gặp}\label{ux111iux1ec3m-muxf9-thux1b0ux1eddng-gux1eb7p-15}

\subsubsection{17.3.2 Điểm mù thường gặp (Cognitive Traps \&
Limitations)}\label{ux111iux1ec3m-muxf9-thux1b0ux1eddng-gux1eb7p-cognitive-traps-limitations-121}

Quan sát bề ngoài, một cuộc chiến tranh giá tàn khốc dường như là dấu
hiệu Cartel đã sụp đổ. Tuy nhiên, theo logic của Green-Porter, chiến
tranh giá chính là bằng chứng cho thấy cơ chế thực thi của Cartel đang
hoạt động bình thường! Nó là rủi ro được tính toán trước để duy trì uy
tín của lời đe dọa (credible threat). Không có các cuộc chiến tranh giá
(price wars in equilibrium), không ai có động lực duy trì mức giá cao do
sự cám dỗ lật lọng ngầm.

\subsubsection{17.3.2 Điểm mù thường gặp (Cognitive Traps \&
Limitations)}\label{ux111iux1ec3m-muxf9-thux1b0ux1eddng-gux1eb7p-cognitive-traps-limitations-122}

Sinh viên thường dùng nguyên lý cung cầu cơ bản (cầu tăng thì giá tăng)
và áp dụng cho Cartel. Rotemberg-Saloner chứng minh điều ngược lại. Khi
cầu bùng nổ, phần thưởng của việc ``ăn cắp'' toàn bộ thị trường trong
một kỳ là vô cùng lớn, phá vỡ ràng buộc IC. Để giữ chân các thành viên
khỏi việc phản bội, Cartel buộc phải ``bớt tham'', giảm giá xuống mức
thấp hơn giá độc quyền tối ưu để làm giảm sức hấp dẫn của sự lật lọng.

\subsubsection{17.3.3 Ứng dụng thực
tế}\label{ux1ee9ng-dux1ee5ng-thux1ef1c-tux1ebf-9}

\begin{itemize}
\tightlist
\item
  \textbf{Thị trường Dầu mỏ và OPEC+:} OPEC+ là ví dụ kinh điển nhất về
  một Cartel. Mô hình Green-Porter giải thích các giai đoạn sụt giảm giá
  dầu mạnh mẻ (như cuộc chiến giá giữa Saudi Arabia và Nga đầu năm 2020)
  không hẳn là dấu chấm hết cho OPEC+, mà là các giai đoạn ``trừng
  phạt'' (reversionary phases) được tính toán trước nhằm ép các quốc gia
  vi phạm hạn ngạch (sản xuất quá mức) phải quay lại tuân thủ.
\item
  ** 담합 (Đàm phán ngầm) trong ngành Thép/Xi măng:** Tại Việt Nam và
  các quốc gia đang phát triển, các nhà sản xuất vật liệu xây dựng
  thường tạo ra các nhóm ``câu lạc bộ'' để bình ổn giá. Mô hình
  Rotemberg-Saloner cảnh báo rằng các thỏa thuận ngầm này dễ bị phá vỡ
  nhất vào mùa xây dựng cao điểm (mùa khô), khi các nhà máy bị cám dỗ
  giảm giá ngầm (chiết khấu sâu cho đại lý lớn) để tối đa hóa doanh số
  ngắn hạn, dẫn đến các đợt cạnh tranh khốc liệt nhất lại rơi vào lúc
  thị trường bùng nổ.
\end{itemize}

\subsubsection{17.3.4 Ví dụ minh họa tính
toán}\label{vuxed-dux1ee5-minh-hux1ecda-tuxednh-touxe1n-2}

\textbf{Bài toán (Cân bằng Grim Trigger Strategy):} Hai doanh nghiệp
tham gia cạnh tranh Bertrand với hàng hóa đồng nhất, \(MC = 10\). Hàm
cầu thị trường \(Q = 100 - P\). - Nếu cấu kết (định giá độc quyền
\(P^M\)), họ chia đôi lợi nhuận. - Nếu một bên lật lọng, họ định giá
\(P = P^M - \epsilon\), chiếm toàn bộ thị trường. - Nếu chiến tranh giá
nổ ra, họ quay về cân bằng Bertrand \(P = MC = 10\). Tìm giá trị chiết
khấu \(\delta\) tối thiểu để duy trì sự cấu kết.

\textbf{Giải:} 1. Lợi nhuận độc quyền (Cartel):
\(\max \pi = (P - 10)(100 - P)\).
\(FOC \implies 100 - 2P + 10 = 0 \implies P^M = 55\). Lượng độc quyền:
\(Q^M = 100 - 55 = 45\). Tổng lợi nhuận:
\(\pi^{Total} = (55 - 10)45 = 2025\). Mỗi doanh nghiệp nhận (cấu kết):
\(\pi^M = 2025 / 2 = 1012.5\).

\begin{enumerate}
\def\labelenumi{\arabic{enumi}.}
\setcounter{enumi}{1}
\item
  Lợi nhuận lật lọng (Defect): Bằng cách giảm giá một lượng siêu nhỏ
  \(\epsilon \to 0\) (đặt \(P \approx 55\)), doanh nghiệp lật lọng chiếm
  toàn bộ cầu \(Q = 45\). Lợi nhuận lật lọng: \(\pi^D = 2025\).
\item
  Lợi nhuận trừng phạt (Punishment/Nash): Theo cạnh tranh Bertrand,
  \(P = MC = 10 \implies Q = 90\). Lợi nhuận: \(\pi^N = 0\).
\item
  Điều kiện duy trì cấu kết:
  \[ \delta \ge \frac{\pi^D - \pi^M}{\pi^D - \pi^N} = \frac{2025 - 1012.5}{2025 - 0} = \frac{1012.5}{2025} = 0.5 \]
  Vậy các doanh nghiệp phải có hệ số chiết khấu \(\delta \ge 0.5\) (đủ
  quan tâm đến tương lai) để duy trì cấu kết Cartel thành công.
\end{enumerate}

\begin{center}\rule{0.5\linewidth}{0.5pt}\end{center}

\subsection{17.4 Kinh tế học Nền tảng \& Thị trường Hai mặt (Platform
Economics \& Two-Sided
Markets)}\label{kinh-tux1ebf-hux1ecdc-nux1ec1n-tux1ea3ng-thux1ecb-trux1b0ux1eddng-hai-mux1eb7t-platform-economics-two-sided-markets}

\subsubsection{17.4.1 Cốt lõi Lý
thuyết}\label{cux1ed1t-luxf5i-luxfd-thuyux1ebft-3}

Lý thuyết thị trường hai mặt (Two-sided markets), được định hình bởi
Rochet \& Tirole (2003, 2006) và Armstrong (2006), phân tích các nền
tảng (Platforms) phục vụ hai (hoặc nhiều) nhóm khách hàng khác biệt,
trong đó tồn tại \textbf{hiệu ứng mạng chéo (cross-side network
externalities)}. Lợi ích của nhóm A phụ thuộc vào số lượng hoặc sự tham
gia của nhóm B, và ngược lại (ví dụ: người mua và người bán trên Shopee,
hành khách và tài xế trên Uber).

Một thị trường được gọi là ``hai mặt'' theo Rochet \& Tirole (2006) khi
nền tảng không chỉ quan tâm đến tổng mức giá (Price level) \(p_A + p_B\)
mà cơ cấu giá (Price structure) ảnh hưởng trực tiếp đến tổng khối lượng
giao dịch. Nói cách khác, việc phân bổ lại phí từ mặt A sang mặt B trong
khi giữ nguyên tổng giá không có tính trung lập (non-neutrality).

\textbf{Mô hình Rochet \& Tirole (2003):} Giả sử nền tảng tối đa hóa lợi
nhuận từ việc định giá giao dịch cho hai nhóm A và B. Hàm cầu tham gia
của nhóm A là \(D_A(p_A)\) và nhóm B là \(D_B(p_B)\). Tổng giao dịch là
\(D_A(p_A) \times D_B(p_B)\). Gọi \(\alpha_A\) là ngoại ứng mạng chéo mà
nhóm B mang lại cho mỗi thành viên nhóm A, và \(\alpha_B\) là ngoại ứng
mà nhóm A mang lại cho nhóm B. Bài toán của nền tảng:
\[ \max_{p_A, p_B} \pi = (p_A + p_B - c) D_A(p_A) D_B(p_B) \] trong đó
\(c\) là chi phí biên để hoàn thành một giao dịch (marginal cost per
interaction). Giải điều kiện FOC, ta thu được công thức định giá
Rochet-Tirole mở rộng:
\[ p_A - (c - p_B) = \frac{p_A}{\epsilon_A} \quad \text{hoặc} \quad \frac{p_A - (c - \alpha_B)}{p_A} = \frac{1}{\epsilon_A} \]
Viết lại một cách trực quan hơn đối với mức giá mặt A (với chi phí biên
\(MC_A\) thực tế phân bổ cho A):
\[ p_A - MC_A = \frac{p_A}{\epsilon_A} - \alpha_B \] Công thức này là sự
tổng quát hóa của nguyên lý Lerner. Mức giá (markup) trên mặt A không
chỉ phụ thuộc vào độ co giãn của cầu \(\epsilon_A\) mà còn bị trừ đi
\(\alpha_B\) (giá trị mà một khách hàng A tạo ra cho mặt B).

\textbf{Mô hình Armstrong (2006) - Single-homing vs Multi-homing:}
Armstrong tập trung vào phí thành viên (membership fees) thay vì phí
giao dịch. Nếu một mặt chỉ sử dụng duy nhất một nền tảng (Single-homing
- ví dụ người dùng iOS) và mặt kia tham gia nhiều nền tảng (Multi-homing
- ví dụ nhà phát triển ứng dụng), nền tảng sẽ cạnh tranh khốc liệt để
giành giật mặt Single-homing (trợ giá mạnh) nhằm biến họ thành ``điểm
nghẽn'' độc quyền (competitive bottleneck) để thu lợi từ mặt
Multi-homing.

\subsubsection{17.4.2 Điểm mù thường
gặp}\label{ux111iux1ec3m-muxf9-thux1b0ux1eddng-gux1eb7p-16}

\subsubsection{17.4.2 Điểm mù thường gặp (Cognitive Traps \&
Limitations)}\label{ux111iux1ec3m-muxf9-thux1b0ux1eddng-gux1eb7p-cognitive-traps-limitations-123}

Trong thị trường truyền thống, định giá dưới chi phí biên (\(P < MC\))
được coi là hành vi bán phá giá tiêu diệt (Predatory Pricing) mang tính
chống độc quyền. Tuy nhiên, trong thị trường hai mặt, việc \(p_A < 0\)
(trợ cấp cho nhóm A) là hoàn toàn bình thường và tối ưu nếu \(\alpha_B\)
(giá trị của A đối với B) đủ lớn. Các cơ quan quản lý (Antitrust) thường
mắc sai lầm khi áp dụng khung phân tích một mặt cho thị trường hai mặt,
dẫn đến việc trừng phạt sai các nền tảng khi họ miễn phí hoặc trợ cấp
cho người dùng.

\subsubsection{17.4.2 Điểm mù thường gặp (Cognitive Traps \&
Limitations)}\label{ux111iux1ec3m-muxf9-thux1b0ux1eddng-gux1eb7p-cognitive-traps-limitations-124}

Sinh viên thường giả định thị trường hai mặt sẽ tồn tại cạnh tranh dài
hạn. Thực tế, do hiệu ứng mạng chéo dương (positive feedback loop), thị
trường hai mặt có xu hướng tự nhiên tiến về ``Winner-takes-all'' (Kẻ
chiến thắng lấy tất cả). Nền tảng nào đạt được quy mô tới hạn (Critical
Mass) trước sẽ hút hết người dùng từ nền tảng nhỏ hơn. Điểm mù là không
nhận ra sự cân bằng trong thị trường hai mặt thường là trạng thái không
ổn định (unstable equilibria) trừ khi có sự khác biệt hóa rất lớn.

\subsubsection{17.4.3 Ứng dụng thực
tế}\label{ux1ee9ng-dux1ee5ng-thux1ef1c-tux1ebf-10}

\begin{itemize}
\tightlist
\item
  \textbf{Thương mại điện tử (Shopee, TikTok Shop):} Shopee Việt Nam ban
  đầu áp dụng trợ cấp sâu cho cả người mua (Freeship, mã giảm giá) và
  người bán (miễn phí hoa hồng) để xây dựng hiệu ứng mạng. Khi đạt vị
  thế thống lĩnh, phí bắt đầu chuyển dịch mạnh: họ tăng phí hoa hồng,
  phí thanh toán, và phí cố định (Mức \(p_B\) tăng vọt) đối với người
  bán (mặt Multi-homing, ít nhạy cảm với việc rời bỏ hơn do phụ thuộc
  vào lượng traffic khổng lồ của Shopee), trong khi vẫn giữ trợ giá cho
  người mua (mặt Single-homing, độ nhạy giá \(\epsilon_A\) cao).
\item
  \textbf{Bất đối xứng định giá thẻ tín dụng / Ứng dụng Hẹn hò:} Tinder
  miễn phí đối với nữ giới (\(p_A = 0\) do \(\alpha_B\) - giá trị của nữ
  đối với nam - rất lớn), trong khi thu phí đắt đỏ đối với nam giới qua
  Tinder Plus/Gold. Hệ thống thanh toán Visa/Mastercard thu phí cao đối
  với Merchant (điểm bán hàng) và dùng số tiền đó để hoàn tiền
  (Cashback) hoặc tích điểm cho chủ thẻ (Cardholders) nhằm kích thích họ
  quẹt thẻ.
\end{itemize}

\subsubsection{17.4.4 Ví dụ minh họa tính
toán}\label{vuxed-dux1ee5-minh-hux1ecda-tuxednh-touxe1n-3}

\textbf{Bài toán (Định giá nền tảng Rochet-Tirole):} Một nền tảng kết
nối Người mua (A) và Người bán (B). Nền tảng ấn định phí tham gia
\(p_A\) và \(p_B\). Lượng cầu của người mua: \(N_A = 100 - p_A + 2N_B\)
(Hiệu ứng chéo: thêm 1 người bán làm tăng WTP của người mua thêm 2).
Lượng cầu của người bán: \(N_B = 100 - p_B + 2N_A\) Chi phí biên cung
cấp dịch vụ cho mỗi người tham gia là \(c_A = 10, c_B = 10\). Giả sử nền
tảng là độc quyền. Để đơn giản hóa trong bài toán tối ưu bậc nhất của
mạng lưới đối xứng (giả định \(p_A = p_B = p, N_A = N_B = N\)), hãy tìm
mức giá \(p\) tối ưu.

\textbf{Giải:} Do tính đối xứng \(N_A = N_B = N\), thay vào phương trình
cầu: \(N = 100 - p + 2N \implies -N = 100 - p \implies p = 100 + N\) (Vô
lý nếu xét hàm cầu nghịch thông thường, nhưng hệ số ngoại ứng 2 là quá
lớn gây bùng nổ, ta thiết lập lại hàm cầu để phản ánh ngoại ứng biên
\(\alpha < 1\)). \emph{Điều chỉnh hàm cầu chuẩn mực (với ngoại ứng
\(\alpha = 0.5\)):} \(N_A = 100 - p_A + 0.5 N_B\)
\(N_B = 100 - p_B + 0.5 N_A\) Thay \(N_A = N_B = N\) và
\(p_A = p_B = p\):
\(N = 100 - p + 0.5N \implies 0.5N = 100 - p \implies N = 200 - 2p\).

Lợi nhuận của nền tảng đối xứng:
\(\pi = (p_A - c_A)N_A + (p_B - c_B)N_B = 2(p - 10)N = 2(p - 10)(200 - 2p)\)
\(\pi = 4(p - 10)(100 - p)\). Tối đa hóa lợi nhuận
\(\frac{\partial \pi}{\partial p} = 0\):
\(\frac{\partial}{\partial p} [100p - p^2 - 1000 + 10p] = 110 - 2p = 0 \implies p^* = 55\).
Lượng người tham gia: \(N^* = 200 - 2(55) = 90\). Lợi nhuận:
\(\pi^* = 4(55 - 10)(100 - 55) = 4 \times 45 \times 45 = 8100\).

So sánh với trường hợp KHÔNG có ngoại ứng mạng chéo (coi
\(\alpha = 0\)): \(N = 100 - p\). \(\pi = 2(p - 10)(100 - p)\). Tối đa
hóa
\(\pi \implies FOC: 100 - 2p + 10 = 0 \implies p_{no-network} = 55\).
\(N_{no-network} = 45\). Nhờ có hiệu ứng mạng, ở cùng mức giá \(p=55\),
nền tảng bùng nổ lượng người dùng từ 45 lên 90 ở mỗi mặt, làm tăng gấp
đôi lợi nhuận!

\section{PHỤ LỤC A: NỀN TẢNG TOÁN HỌC \& LÝ THUYẾT ĐỐI NGẪU TRONG VI
MÔ/VĨ MÔ
PHD}\label{phux1ee5-lux1ee5c-a-nux1ec1n-tux1ea3ng-touxe1n-hux1ecdc-luxfd-thuyux1ebft-ux111ux1ed1i-ngux1eabu-trong-vi-muxf4vux129-muxf4-phd}

\subsection{A.1 Cân bằng Tổng quát Arrow--Debreu \& Định lý Tồn
tại}\label{a.1-cuxe2n-bux1eb1ng-tux1ed5ng-quuxe1t-arrowdebreu-ux111ux1ecbnh-luxfd-tux1ed3n-tux1ea1i}

\subsubsection{A.1.1 Cốt lõi Lý
thuyết}\label{a.1.1-cux1ed1t-luxf5i-luxfd-thuyux1ebft}

Mô hình Cân bằng Tổng quát (General Equilibrium - GE) Arrow-Debreu là
nền tảng của kinh tế học tân cổ điển, mô tả sự tương tác đồng thời trên
tất cả các thị trường.

\textbf{Thiết lập mô hình:} Giả sử nền kinh tế có \(L\) loại hàng hóa
(với không gian hàng hóa \(\mathbb{R}^L\)), \(I\) người tiêu dùng (index
\(i \in \{1, \dots, I\}\)), và \(J\) hãng sản xuất (index
\(j \in \{1, \dots, J\}\)). - Mỗi người tiêu dùng \(i\) có tập tiêu dùng
\(X_i \subset \mathbb{R}^L_+\), hàm thỏa dụng
\(u_i: X_i \to \mathbb{R}\), và nguồn lực ban đầu
\(\omega_i \in \mathbb{R}^L_+\). Tổng nguồn lực ban đầu là
\(\bar{\omega} = \sum_{i=1}^I \omega_i\). Người tiêu dùng sở hữu cổ phần
\(\theta_{ij} \geq 0\) tại hãng \(j\), với
\(\sum_{i=1}^I \theta_{ij} = 1, \forall j\). - Mỗi hãng \(j\) có tập sản
xuất (công nghệ) \(Y_j \subset \mathbb{R}^L\). Một véc-tơ sản xuất
\(y_j \in Y_j\) có thành phần âm cho yếu tố đầu vào và dương cho sản
phẩm đầu ra. - Hệ thống giá là véc-tơ
\(p \in \mathbb{R}^L_+ \setminus \{0\}\).

\textbf{Định nghĩa Cân bằng Walras (Walrasian Equilibrium):} Một phân bổ
\((x^*, y^*)\) cùng với véc-tơ giá \(p^*\) cấu thành một Cân bằng Walras
nếu: 1. \emph{Tối đa hóa lợi nhuận:} Với mọi hãng \(j\),
\(y_j^* \in Y_j\) tối đa hóa \(p^* \cdot y_j\), tức là
\(p^* \cdot y_j^* \geq p^* \cdot y_j\) cho mọi \(y_j \in Y_j\). 2.
\emph{Tối đa hóa thỏa dụng:} Với mọi người tiêu dùng \(i\),
\(x_i^* \in X_i\) tối đa hóa \(u_i(x_i)\) với điều kiện ràng buộc ngân
sách
\(p^* \cdot x_i \leq p^* \cdot \omega_i + \sum_{j=1}^J \theta_{ij} (p^* \cdot y_j^*)\).
3. \emph{Thanh khoản thị trường (Market Clearing):}
\(\sum_{i=1}^I x_i^* = \bar{\omega} + \sum_{j=1}^J y_j^*\).

\textbf{Hàm Cầu vượt (Excess Demand Function) \& Định luật Walras:} Cầu
vượt của nền kinh tế tại mức giá \(p\) được định nghĩa là
\(z(p) = \sum_{i=1}^I x_i(p) - \bar{\omega} - \sum_{j=1}^J y_j(p)\).
\emph{Định luật Walras:} Dưới giả định về tính không bão hòa cục bộ
(local nonsatiation) của sở thích, ràng buộc ngân sách giữ dấu bằng, dẫn
đến: \[ p \cdot z(p) = 0 \quad \text{với mọi } p \] Điều này ngụ ý rằng
tổng giá trị của cầu vượt luôn bằng không; nếu \(L-1\) thị trường thanh
khoản (\(z_l(p)=0\)), thị trường thứ \(L\) cũng tự động thanh khoản.

\textbf{Sự tồn tại của Cân bằng Tổng quát (Arrow-Debreu 1954):} Chứng
minh sự tồn tại của GE dựa trên Định lý Điểm bất động. Ta chuẩn hóa
không gian giá thành một đơn hình đơn vị
\(\Delta = \{ p \in \mathbb{R}^L_+ \mid \sum_{l=1}^L p_l = 1 \}\). Theo
Brouwer (hoặc Kakutani đối với ánh xạ đa trị), định lý đảm bảo tồn tại
một \(p^* \in \Delta\) sao cho \(z(p^*) \leq 0\). \emph{Cơ chế phác
thảo:} Định nghĩa một hàm điều chỉnh giá kiểu đấu giá viên Walras
(Walrasian auctioneer):
\[ g_l(p) = \frac{p_l + \max\{0, z_l(p)\}}{1 + \sum_{k=1}^L \max\{0, z_k(p)\}} \]
Hàm \(g: \Delta \to \Delta\) là liên tục. Theo Định lý Điểm bất động
Brouwer, tồn tại \(p^*\) sao cho \(g(p^*) = p^*\). Khi kết hợp với Định
luật Walras \(p^* \cdot z(p^*) = 0\), ta suy ra \(z(p^*) \leq 0\), và do
đó \(p^*\) là một mức giá cân bằng.

\subsubsection{A.1.2 Điểm mù thường
gặp}\label{a.1.2-ux111iux1ec3m-muxf9-thux1b0ux1eddng-gux1eb7p}

\subsubsection{A.1.2 Điểm mù thường gặp (Cognitive Traps \&
Limitations)}\label{a.1.2-ux111iux1ec3m-muxf9-thux1b0ux1eddng-gux1eb7p-cognitive-traps-limitations}

Nhiều nghiên cứu sinh tin rằng việc chứng minh sự tồn tại của cân bằng
Arrow-Debreu đồng nghĩa với việc nền kinh tế hội tụ về điểm cân bằng đó
(stable) và cân bằng là duy nhất. Thực tế, Định lý
Sonnenschein--Mantel--Debreu (SMD) chứng minh rằng hàm cầu vượt tổng
cộng \(z(p)\) có thể nhận bất kỳ hình dạng nào miễn là liên tục, thuần
nhất bậc 0 và tuân thủ Định luật Walras. Do đó, cân bằng có thể đa trị
(multiple equilibria) và quá trình điều chỉnh giá (tâtonnement) có thể
không bao giờ hội tụ, tạo ra chu kỳ hoặc hỗn loạn.

\subsubsection{A.1.2 Điểm mù thường gặp (Cognitive Traps \&
Limitations)}\label{a.1.2-ux111iux1ec3m-muxf9-thux1b0ux1eddng-gux1eb7p-cognitive-traps-limitations-1}

Mô hình Arrow-Debreu ngầm giả định các thị trường là đầy đủ cho mọi
trạng thái của tự nhiên trong tương lai (contingent commodities). Trong
thực tế, các thị trường tài chính và bảo hiểm rất không đầy đủ
(incomplete markets). Khi thị trường thiếu vắng, định lý tồn tại có thể
sụp đổ hoặc phân bổ cân bằng không còn đạt Hiệu quả Pareto (First
Welfare Theorem không còn đúng). Việc áp dụng bừa bãi mô hình
Arrow-Debreu vào thực tiễn mà quên mất giới hạn này dẫn đến các chính
sách kinh tế vĩ mô sai lầm.

\subsubsection{A.1.3 Ứng dụng thực
tế}\label{a.1.3-ux1ee9ng-dux1ee5ng-thux1ef1c-tux1ebf}

Lý thuyết GE được ứng dụng rộng rãi trong các mô hình Cân bằng Tổng quát
Có thể Tính toán (Computable General Equilibrium - CGE) và các mô hình
Cân bằng Tổng quát Động Ngẫu nhiên (DSGE) trong kinh tế vĩ mô hiện đại.
Tại Việt Nam, các mô hình CGE được Bộ Kế hoạch \& Đầu tư sử dụng để đánh
giá tác động của các hiệp định thương mại tự do (FTA) như CPTPP hay
EVFTA. Các cú sốc thuế quan làm thay đổi giá tương đối, kích hoạt sự
dịch chuyển vốn và lao động giữa các ngành công nghiệp. Mô hình CGE sử
dụng thuật toán tìm điểm bất động để tính toán mức giá cân bằng mới và
sự thay đổi trong phúc lợi xã hội thông qua các hàm tiện ích tương
đương.

\subsubsection{A.1.4 Ví dụ minh họa tính
toán}\label{a.1.4-vuxed-dux1ee5-minh-hux1ecda-tuxednh-touxe1n}

Giả sử có nền kinh tế trao đổi (Edgeworth box) với 2 hàng hóa \(X, Y\)
và 2 người tiêu dùng \(A, B\). Hàm thỏa dụng:
\(U_A(X_A, Y_A) = X_A Y_A\) và \(U_B(X_B, Y_B) = X_B Y_B\). Tài nguyên
ban đầu: \(\omega_A = (8, 2)\) và \(\omega_B = (2, 8)\). Tổng cung:
\(\bar{X} = 10, \bar{Y} = 10\). Gọi giá của \(X\) là \(p_X = p\) và giá
của \(Y\) là \(p_Y = 1\) (chuẩn hóa). Thu nhập: \(I_A = 8p + 2\) và
\(I_B = 2p + 8\). Từ bài toán tối đa hóa thỏa dụng Cobb-Douglas, cầu
Marshallian là:
\(X_A(p) = \frac{I_A}{2p} = \frac{8p + 2}{2p} = 4 + \frac{1}{p}\)
\(X_B(p) = \frac{I_B}{2p} = \frac{2p + 8}{2p} = 1 + \frac{4}{p}\) Điều
kiện thanh khoản thị trường cho hàng hóa \(X\):
\(X_A(p) + X_B(p) = 10 \implies (4 + \frac{1}{p}) + (1 + \frac{4}{p}) = 10 \implies 5 + \frac{5}{p} = 10 \implies p^* = 1\).
Vậy mức giá cân bằng là \(p_X/p_Y = 1\). Phân bổ cân bằng:
\(X_A^* = 5, Y_A^* = 5\) và \(X_B^* = 5, Y_B^* = 5\). Cầu vượt bằng 0.

\begin{center}\rule{0.5\linewidth}{0.5pt}\end{center}

\subsection{A.2 Lý thuyết Đối ngẫu Tiêu dùng \& Sản
xuất}\label{a.2-luxfd-thuyux1ebft-ux111ux1ed1i-ngux1eabu-tiuxeau-duxf9ng-sux1ea3n-xuux1ea5t}

\subsubsection{A.2.1 Cốt lõi Lý
thuyết}\label{a.2.1-cux1ed1t-luxf5i-luxfd-thuyux1ebft}

Lý thuyết Đối ngẫu (Duality Theory) thiết lập mối liên hệ toán học chặt
chẽ giữa hàm mục tiêu được tối ưu hóa (Primal) và bài toán phản chiếu
của nó (Dual).

\textbf{Tối đa hóa Thỏa dụng (Primal) và Tối thiểu hóa Chi tiêu (Dual):}
- \emph{Bài toán Primal (UMP):} \(\max u(x)\) s.t. \(p \cdot x \leq w\).
Lời giải là Cầu Marshallian \(x(p, w)\). Thay lời giải vào hàm mục tiêu,
ta có Hàm Thỏa dụng Gián tiếp (Indirect Utility Function):
\(v(p, w) = u(x(p, w))\). - \emph{Bài toán Dual (EMP):}
\(\min p \cdot x\) s.t. \(u(x) \geq u\). Lời giải là Cầu Hicksian (Cầu
bồi thường) \(h(p, u)\). Thay lời giải vào, ta có Hàm Chi tiêu
(Expenditure Function): \(e(p, u) = p \cdot h(p, u)\).

\textbf{Bổ đề Shephard (Shephard's Lemma):} Theo định lý bao (Envelope
Theorem), đạo hàm của hàm chi tiêu theo giá chính là cầu Hicksian:
\[ \frac{\partial e(p, u)}{\partial p_i} = h_i(p, u) \]

\textbf{Đồng nhất thức Roy (Roy's Identity):} Cho phép tìm hàm cầu
Marshallian từ hàm thỏa dụng gián tiếp mà không cần giải bài toán UMP
trực tiếp:
\[ x_i(p, w) = - \frac{\partial v(p, w) / \partial p_i}{\partial v(p, w) / \partial w} \]

\textbf{Phương trình Slutsky:} Liên kết sự thay đổi của cầu theo giá qua
hai kênh: hiệu ứng thay thế (substitution effect) và hiệu ứng thu nhập
(income effect). Ghi nhận rằng \(h_i(p, u) = x_i(p, e(p, u))\). Đạo hàm
hai vế theo \(p_j\):
\[ \frac{\partial h_i(p, u)}{\partial p_j} = \frac{\partial x_i(p, w)}{\partial p_j} + \frac{\partial x_i(p, w)}{\partial w} \frac{\partial e(p, u)}{\partial p_j} \]
Áp dụng Bổ đề Shephard
\(\frac{\partial e}{\partial p_j} = h_j(p, u) = x_j(p, w)\) và sắp xếp
lại:
\[ \frac{\partial x_i(p, w)}{\partial p_j} = \frac{\partial h_i(p, u)}{\partial p_j} - x_j(p, w) \frac{\partial x_i(p, w)}{\partial w} \]

\textbf{Đối ngẫu trong Sản xuất:} Tương tự, bài toán tối đa hóa lợi
nhuận tạo ra Hàm Lợi nhuận \(\pi(p, w)\). \emph{Bổ đề Hotelling} phát
biểu rằng \(\frac{\partial \pi(p, w)}{\partial p} = y(p, w)\) (Cung sản
phẩm) và \(\frac{\partial \pi(p, w)}{\partial w_i} = -x_i(p, w)\) (Cầu
yếu tố sản xuất).

\subsubsection{A.2.2 Điểm mù thường
gặp}\label{a.2.2-ux111iux1ec3m-muxf9-thux1b0ux1eddng-gux1eb7p}

\subsubsection{A.2.2 Điểm mù thường gặp (Cognitive Traps \&
Limitations)}\label{a.2.2-ux111iux1ec3m-muxf9-thux1b0ux1eddng-gux1eb7p-cognitive-traps-limitations}

Sinh viên thường đạo hàm tự động bất kỳ hàm \(e(p, u)\) nào bằng Bổ đề
Shephard để tìm cầu Hicksian. Điểm mù là hàm \(e(p, u)\) về bản chất
phải là hàm lõm (concave) theo giá \(p\). Nếu hàm chi tiêu giả định
không thỏa mãn tính chất lõm này (ví dụ do ước lượng kinh tế lượng tuyến
tính hóa sai lầm), ma trận Slutsky thu được sẽ không xác định âm
(negative semi-definite), vi phạm Quy luật Cầu bồi thường (Luật Cầu
yếu).

\subsubsection{A.2.2 Điểm mù thường gặp (Cognitive Traps \&
Limitations)}\label{a.2.2-ux111iux1ec3m-muxf9-thux1b0ux1eddng-gux1eb7p-cognitive-traps-limitations-1}

Trong các phân tích chính sách công, nhiều nhà kinh tế đo lường tổn thất
vô ích (Deadweight Loss) sử dụng thặng dư tiêu dùng Marshallian
(Consumer Surplus) - tích phân đường cầu Marshallian. Tuy nhiên, Cầu
Marshallian chứa đựng hiệu ứng thu nhập, khiến CS phụ thuộc vào lộ trình
thay đổi giá (path-dependent). Đối ngẫu lý thuyết chứng minh rằng chỉ có
Biến thiên Tương đương (Equivalent Variation - EV) và Biến thiên Bù đắp
(Compensating Variation - CV), đo lường qua tích phân cầu Hicksian, mới
là các thước đo phúc lợi chính xác.

\subsubsection{A.2.3 Ứng dụng thực
tế}\label{a.2.3-ux1ee9ng-dux1ee5ng-thux1ef1c-tux1ebf}

Trong kinh tế học phát triển và phân tích chính sách thuế, lý thuyết đối
ngẫu được sử dụng để ước lượng Hệ thống Cầu Lý tưởng Gần đúng (Almost
Ideal Demand System - AIDS) của Deaton và Muellbauer (1980). Tại Việt
Nam, hệ thống AIDS được ứng dụng để phân tích mô hình tiêu dùng hộ gia
đình (thông qua dữ liệu VHLSS), từ đó tính toán độ co giãn của cầu theo
giá và thu nhập đối với thực phẩm, năng lượng, giáo dục. Roy's Identity
cho phép thiết lập hệ phương trình ước lượng kinh tế lượng một cách tao
nhã trực tiếp từ hàm chi tiêu (expenditure function) log-tuyến tính, bảo
toàn các ràng buộc lý thuyết vi mô (tính thuần nhất, đối xứng Slutsky).

\subsubsection{A.2.4 Ví dụ minh họa tính
toán}\label{a.2.4-vuxed-dux1ee5-minh-hux1ecda-tuxednh-touxe1n}

Cho hàm chi tiêu của người tiêu dùng:
\(e(p_1, p_2, u) = 2 u (p_1 p_2)^{0.5}\). Yêu cầu: Tìm cầu Hicksian, Cầu
Marshallian, Hàm thỏa dụng gián tiếp.

\emph{Bước 1: Tìm cầu Hicksian bằng Bổ đề Shephard}
\(h_1(p, u) = \frac{\partial e}{\partial p_1} = 2 u (p_2)^{0.5} \cdot 0.5 (p_1)^{-0.5} = u \left( \frac{p_2}{p_1} \right)^{0.5}\)
\(h_2(p, u) = \frac{\partial e}{\partial p_2} = u \left( \frac{p_1}{p_2} \right)^{0.5}\)

\emph{Bước 2: Tìm Hàm thỏa dụng gián tiếp} Do
\(w = e(p_1, p_2, u) = 2 u (p_1 p_2)^{0.5}\), giải \(u\) theo \(w, p\):
\(v(p_1, p_2, w) = \frac{w}{2 (p_1 p_2)^{0.5}}\)

\emph{Bước 3: Tìm cầu Marshallian thông qua Đồng nhất thức Roy} Tử số
cho \(x_1\):
\(\frac{\partial v}{\partial p_1} = \frac{w}{2 (p_2)^{0.5}} (-0.5) p_1^{-1.5} = - \frac{w}{4 p_1 (p_1 p_2)^{0.5}}\)
Mẫu số: \(\frac{\partial v}{\partial w} = \frac{1}{2 (p_1 p_2)^{0.5}}\)
\(x_1(p, w) = - \frac{- \frac{w}{4 p_1 (p_1 p_2)^{0.5}}}{\frac{1}{2 (p_1 p_2)^{0.5}}} = \frac{w}{2 p_1}\).
Tương tự, \(x_2(p, w) = \frac{w}{2 p_2}\). Điều này tương ứng với hàm
thỏa dụng Cobb-Douglas \(U(x_1, x_2) = x_1^{0.5} x_2^{0.5}\) với chi
tiêu ngân sách chia đều (50-50).

\begin{center}\rule{0.5\linewidth}{0.5pt}\end{center}

\subsection{A.3 Lựa chọn trong Bất định \& Thỏa dụng Kỳ vọng von
Neumann--Morgenstern}\label{a.3-lux1ef1a-chux1ecdn-trong-bux1ea5t-ux111ux1ecbnh-thux1ecfa-dux1ee5ng-kux1ef3-vux1ecdng-von-neumannmorgenstern}

\subsubsection{A.3.1 Cốt lõi Lý
thuyết}\label{a.3.1-cux1ed1t-luxf5i-luxfd-thuyux1ebft}

Sự ra quyết định dưới điều kiện bất định được mô hình hóa qua không gian
các xổ số (Lotteries). Gọi \(\mathcal{L}\) là không gian của các phân
phối xác suất trên tập các kết cục. \textbf{Các tiên đề von
Neumann-Morgenstern (vNM):} Sở thích \(\succsim\) trên \(\mathcal{L}\)
phải thỏa mãn: 1. \emph{Đầy đủ (Completeness) \& Bắc cầu
(Transitivity).} 2. \emph{Liên tục (Continuity):} Với mọi
\(L_1 \succ L_2 \succ L_3\), tồn tại \(\alpha \in (0, 1)\) sao cho
\(\alpha L_1 + (1-\alpha) L_3 \sim L_2\). 3. \emph{Độc lập
(Independence):} Với mọi \(L_1 \succ L_2\) và mọi \(L_3\), cùng với
\(\alpha \in (0,1]\), ta có
\(\alpha L_1 + (1-\alpha) L_3 \succ \alpha L_2 + (1-\alpha) L_3\).

\textbf{Định lý Thỏa dụng Kỳ vọng vNM:} Sở thích \(\succsim\) thỏa mãn 4
tiên đề trên khi và chỉ khi tồn tại một hàm thỏa dụng Bernoulli \(u(x)\)
sao cho \(L_1 \succsim L_2\) tương đương với:
\[ U(L_1) = \sum p_i u(x_i) \geq U(L_2) = \sum q_i u(x_i) \]

\textbf{Phân tích Rủi ro và Hàm Arrow-Pratt:} Sự e ngại rủi ro (Risk
Aversion) tương đương với việc hàm \(u(x)\) là lõm (\(u''(x) < 0\)). -
\emph{Độ e ngại rủi ro tuyệt đối (Absolute Risk Aversion - ARA):}
\(ARA(x) = - \frac{u''(x)}{u'(x)}\). - \emph{Độ e ngại rủi ro tương đối
(Relative Risk Aversion - RRA):} \(RRA(x) = - x \frac{u''(x)}{u'(x)}\).
Các dạng hàm điển hình: - \emph{CARA (Constant ARA):}
\(u(x) = -e^{-\alpha x}\) với \(ARA(x) = \alpha\). - \emph{CRRA
(Constant RRA):} \(u(x) = \frac{x^{1-\gamma}}{1-\gamma}\) với
\(RRA(x) = \gamma\).

\textbf{Tương đương Chắc chắn (Certainty Equivalent - \(CE\)) và Phần bù
Rủi ro (Risk Premium - \(\pi_R\)):} - \(CE\) là số tiền chắc chắn mang
lại cùng mức thỏa dụng với xổ số: \(u(CE) = E[u(x)]\). - Phần bù rủi ro
\(\pi_R\) là chi phí người đó sẵn sàng trả để loại bỏ rủi ro:
\(\pi_R = E[x] - CE\). Khai triển Taylor bậc 2 quanh giá trị kỳ vọng
\(\mu = E[x]\):
\(u(x) \approx u(\mu) + u'(\mu)(x - \mu) + \frac{1}{2}u''(\mu)(x - \mu)^2\).
Lấy kỳ vọng hai vế:
\(E[u(x)] \approx u(\mu) + \frac{1}{2}u''(\mu) \sigma^2\). Tương tự,
khai triển \(u(\mu - \pi_R) \approx u(\mu) - \pi_R u'(\mu)\). Đồng nhất
hai vế, ta có:
\(\pi_R \approx \frac{1}{2} \sigma^2 \left( - \frac{u''(\mu)}{u'(\mu)} \right) = \frac{1}{2} \sigma^2 ARA(\mu)\).

\subsubsection{A.3.2 Điểm mù thường
gặp}\label{a.3.2-ux111iux1ec3m-muxf9-thux1b0ux1eddng-gux1eb7p}

\subsubsection{A.3.2 Điểm mù thường gặp (Cognitive Traps \&
Limitations)}\label{a.3.2-ux111iux1ec3m-muxf9-thux1b0ux1eddng-gux1eb7p-cognitive-traps-limitations}

Trong mô hình vNM, một cá nhân e ngại rủi ro là do tính lõm của hàm
\(u(x)\), dẫn đến lợi ích biên của của cải giảm dần (Diminishing
Marginal Utility of Wealth). Điểm mù nằm ở chỗ vNM không thể giải thích
được hiện tượng ``E ngại rủi ro vi mô'' (Small-stakes Risk Aversion).
Theo Rabin (2000), nếu \(u(x)\) đủ lõm để từ chối một ván cược nhỏ (ví
dụ cược 50/50 mất 100\$ hoặc được
110\() ở bất kỳ mức tài sản nào, toán học ngụ ý rằng họ sẽ phải từ chối một ván cược khổng lồ phi lý (50/50 mất 1000\)
hoặc được tài sản vô hạn). Điều này đòi hỏi Lý thuyết Viễn cảnh
(Prospect Theory) với tính phụ thuộc vào điểm tham chiếu
(reference-dependence), phá vỡ mô hình vNM cổ điển.

\subsubsection{A.3.2 Điểm mù thường gặp (Cognitive Traps \&
Limitations)}\label{a.3.2-ux111iux1ec3m-muxf9-thux1b0ux1eddng-gux1eb7p-cognitive-traps-limitations-1}

Nhiều nghiên cứu sinh tin rằng, giống như trong lý thuyết tiêu dùng chắc
chắn (Ordinally Utility), mọi phép biến đổi đơn điệu (monotonic
transformation) của hàm thỏa dụng vNM \(U(P)\) đều bảo toàn sở thích. Sự
thật là hàm \(U(P)\) chỉ bất biến (invariant) dưới \textbf{các phép biến
đổi tuyến tính dương (Positive Affine Transformations)}:
\(v(x) = a \cdot u(x) + b\) với \(a>0\). Nếu thực hiện phép biến đổi phi
tuyến (ví dụ \(v(x) = u(x)^3\)), mức độ e ngại rủi ro \(ARA(x)\) bị thay
đổi đột ngột, và cá thể trong mô hình sẽ đưa ra các lựa chọn rủi ro hoàn
toàn khác biệt.

\subsubsection{A.3.3 Ứng dụng thực
tế}\label{a.3.3-ux1ee9ng-dux1ee5ng-thux1ef1c-tux1ebf}

Lý thuyết thỏa dụng kỳ vọng là lõi của toán tài chính, lý thuyết hợp
đồng (contract theory) và định giá tài sản rủi ro (CAPM, CCAPM). Trong
nghiên cứu vĩ mô, các mô hình Chu kỳ Kinh doanh Thực (RBC) và DSGE
thường sử dụng hàm thỏa dụng CRRA cho các hộ gia đình đại diện để đảm
bảo rằng độ e ngại rủi ro tương đối không thay đổi theo xu hướng tăng
trưởng kinh tế dài hạn, giúp duy trì tính ổn định của lãi suất cân bằng.
Tại Việt Nam, lý thuyết này được các ngân hàng thương mại ứng dụng trong
việc đo lường khẩu vị rủi ro và thiết lập các ma trận tổn thất dự kiến
(Expected Loss) để tuân thủ chuẩn mực Basel II/III.

\subsubsection{A.3.4 Ví dụ minh họa tính
toán}\label{a.3.4-vuxed-dux1ee5-minh-hux1ecda-tuxednh-touxe1n}

Một nhà đầu tư có hàm thỏa dụng CRRA: \(u(x) = \ln(x)\) (trường hợp
\(\gamma = 1\)). Tài sản hiện tại của người này là \(W_0 = 100\). Người
này đang cân nhắc một dự án đầu tư rủi ro (Lottery \(L\)): - Xác suất
50\%: Thắng 21\$ (Tổng tài sản \(x = 121\)). - Xác suất 50\%: Thua 19\$
(Tổng tài sản \(x = 81\)). Yêu cầu: Tính Kỳ vọng tài sản, Tương đương
Chắc chắn (\(CE\)), và Phần bù Rủi ro (\(\pi_R\)).

\emph{Bước 1: Tính Thỏa dụng kỳ vọng}
\(E[u(x)] = 0.5 \ln(121) + 0.5 \ln(81) = 0.5 (4.7958) + 0.5 (4.3944) = 4.5951\).
\emph{(Lưu ý toán học: \(\ln(121) = 2\ln(11)\), \(\ln(81) = 2\ln(9)\),
\(E[u(x)] = \ln(11) + \ln(9) = \ln(99) \approx 4.5951\)).}

\emph{Bước 2: Tính Tương đương Chắc chắn (CE)}
\(u(CE) = E[u(x)] \implies \ln(CE) = \ln(99) \implies CE = 99\). Điều
này có nghĩa là, một tài sản chắc chắn ở mức 99 mang lại cùng mức thỏa
dụng với việc tham gia dự án rủi ro.

\emph{Bước 3: Tính Kỳ vọng tài sản và Phần bù Rủi ro}
\(E[x] = 0.5(121) + 0.5(81) = 101\). Phần bù rủi ro:
\(\pi_R = E[x] - CE = 101 - 99 = 2\). Nhà đầu tư sẵn sàng trả (hy sinh)
tới 2 đơn vị tài sản kỳ vọng để tránh sự bất định. Dù kỳ vọng của dự án
mang lại lợi nhuận +1 (\(E[x] - W_0 = 101 - 100 = 1\)), nhưng do
\(CE = 99 < W_0 = 100\), nhà đầu tư e ngại rủi ro này sẽ \textbf{từ
chối} tham gia dự án. Cùng kiểm tra bằng xấp xỉ Taylor: \(\mu = 101\).
Độ lệch chuẩn \(\sigma^2 = 0.5(20)^2 + 0.5(-20)^2 = 400\).
\(ARA(\mu) = -u''(\mu)/u'(\mu) = 1/\mu = 1/101\).
\(\pi_R \approx \frac{1}{2} (400) (1/101) = 200/101 \approx 1.98\), rất
sát với kết quả chính xác là 2.

\section{PHỤ LỤC B: QUY HOẠCH ĐỘNG \& PHƯƠNG TRÌNH BELLMAN (Dynamic
Programming \&
Bellman)}\label{phux1ee5-lux1ee5c-b-quy-houx1ea1ch-ux111ux1ed9ng-phux1b0ux1a1ng-truxecnh-bellman-dynamic-programming-bellman}

\begin{figure}
\centering
\pandocbounded{\includegraphics[keepaspectratio,alt={Minh họa Quy hoạch động Bellman}]{C:/Users/Admin/.gemini/antigravity/brain/e5791226-a607-45b8-984e-554eb13f4da1/bellman_value_iteration_1785203348626.jpg}}
\caption{Minh họa Quy hoạch động Bellman}
\end{figure}

\subsubsection{B.1.1 Cốt lõi Lý thuyết (Core
Theory)}\label{b.1.1-cux1ed1t-luxf5i-luxfd-thuyux1ebft-core-theory}

Quy hoạch động (Dynamic Programming) là phương pháp tối ưu hóa toán học
do Richard Bellman phát triển để giải quyết các bài toán ra quyết định
theo thời gian. Nguyên lý Tối ưu Bellman phát biểu rằng: ``Một chính
sách tối ưu có tính chất là bất kể trạng thái ban đầu và quyết định ban
đầu là gì, các quyết định còn lại phải tạo thành một chính sách tối ưu
liên quan đến trạng thái là kết quả của quyết định ban đầu.''

\textbf{Định dạng cơ bản (Discrete Time):} Một bài toán quy hoạch động
được cấu thành bởi: - \textbf{Biến trạng thái (State variable - \$ x\_t
\$):} Tóm tắt toàn bộ lịch sử hệ thống cần thiết để ra quyết định hiện
tại. - \textbf{Biến kiểm soát (Control variable - \$ u\_t \$):} Quyết
định được đưa ra tại thời kỳ \$ t \$. - \textbf{Hàm phần thưởng
(Reward/Return function - \$ r(x\_t, u\_t) \$):} Lợi ích tức thời. -
\textbf{Quy luật chuyển động (Transition law):} \$ x\_\{t+1\} = h(x\_t,
u\_t) \$. - \textbf{Hệ số chiết khấu (Discount factor - \$
\beta \in (0,1) \$).}

Bài toán tối ưu chuỗi thời gian vô hạn được viết dưới dạng Hàm Giá trị
(Value Function - \$ V(x) \$), đại diện cho tổng hữu dụng chiết khấu tối
đa có thể đạt được khi bắt đầu tại trạng thái \$ x \$. Thay vì tìm một
chuỗi \(\{u_t\}_{t=0}^\infty\), ta giải \textbf{Phương trình Bellman}:
\[ V(x) = \max_{u} \left\{ r(x,u) + \beta \mathbb{E}[V(x')] \right\} \]
với \$ x' = h(x,u) \$ là trạng thái ở kỳ tiếp theo. Lời giải của phương
trình này là một \textbf{Hàm chính sách (Policy function - \$ u = g(x)
\$)}, chỉ định hành động tối ưu cho mỗi trạng thái.

\textbf{Cơ sở Toán học:} Phương trình Bellman có thể được xem như một
toán tử \$ T \$ tác động lên không gian các hàm liên tục:
\[ (TV)(x) = \max_u \{ r(x,u) + \beta V(h(x,u)) \} \] Theo \textbf{Định
lý Ánh xạ Co (Contraction Mapping Theorem)} (Định lý Banach), vì \$
\beta \textless{} 1 \$, toán tử \$ T \$ là một ánh xạ co. Do đó, nó có
một điểm bất động (fixed point) duy nhất \$ V = TV \$, và ta có thể tìm
\$ V \$ bằng phương pháp Lặp Hàm Giá Trị (Value Function Iteration -
VFI): bắt đầu từ một hàm dự đoán \$ V\_0 \$, tính \$ V\_\{n+1\} = TV\_n
\$ cho đến khi hội tụ.

\subsubsection{B.1.2 Điểm mù thường gặp (Cognitive Traps \&
Limitations)}\label{b.1.2-ux111iux1ec3m-muxf9-thux1b0ux1eddng-gux1eb7p-cognitive-traps-limitations}

\begin{enumerate}
\def\labelenumi{\arabic{enumi}.}
\tightlist
\item
  \textbf{Quên tính chất Markov (Markov Property):} Quy hoạch động chỉ
  áp dụng được nếu phương trình chuyển động \$ x\_\{t+1\} \$ chỉ phụ
  thuộc vào \$ x\_t \$ và \$ u\_t \$. Sinh viên thường cố gắng đưa các
  biến bị trễ (ví dụ \$ x\_\{t-2\} \$) vào thẳng phương trình mà không
  mở rộng không gian trạng thái. Để giải quyết, mọi biến trong quá khứ
  có ảnh hưởng tới tương lai đều phải được tái định nghĩa và đưa vào
  véc-tơ trạng thái \$ \mathbf{x}\_t \$.
\item
  \textbf{Hiểu lầm về V(x'):} Một số học giả nhầm tưởng \$ V(x') \$ là
  giá trị \emph{hiện tại}. Thực tế, \$ V(x') \$ là giá trị của bài toán
  nếu bài toán \emph{mới bắt đầu} vào ngày mai tại trạng thái \$ x' \$.
  Chính hệ số \$ \beta \$ mới là thứ kéo giá trị tương lai đó về hiện
  tại.
\end{enumerate}

\subsubsection{B.1.3 Ứng dụng thực tế (Practical
Applications)}\label{b.1.3-ux1ee9ng-dux1ee5ng-thux1ef1c-tux1ebf-practical-applications}

\subsubsection{B.1.4 Ví dụ minh họa tính toán (Worked Numerical
Examples)}\label{b.1.4-vuxed-dux1ee5-minh-hux1ecda-tuxednh-touxe1n-worked-numerical-examples}

Một người tiêu dùng có tài sản \$ W \$, chọn tiêu dùng \$ c \$ và tiết
kiệm \$ W' \$. Phương trình Bellman:
\[ V(W) = \max_{0 \le c \le W} \left\{ \ln(c) + \beta V(W') \right\} \]
với ràng buộc \$ W' = (W - c)(1+r) \$. Khi giải, ta thu được hàm chính
sách \$ c(W) = (1-\beta) W \$, một kết quả kinh điển giải thích động
thái làm phẳng tiêu dùng.

\subsubsection{B.1.4 Ví dụ minh họa tính toán (Worked Numerical
Examples)}\label{b.1.4-vuxed-dux1ee5-minh-hux1ecda-tuxednh-touxe1n-worked-numerical-examples-1}

Một người thất nghiệp nhận được một đề nghị mức lương \$ w \$ mỗi kỳ từ
phân phối \$ F(w) \$. Người này có thể chọn chấp nhận (A) hoặc từ chối
để tiếp tục tìm kiếm (R). Biến trạng thái là mức lương đề nghị \$ w \$.
Nếu chấp nhận, làm việc mãi mãi với mức lương đó:
\[ V^A(w) = \frac{w}{1-\beta} \] Nếu từ chối, nhận trợ cấp thất nghiệp
\$ c \$ và bốc thăm lại vào ngày mai:
\[ V^R = c + \beta \int V(w') dF(w') \] Phương trình Bellman tổng quát:
\[ V(w) = \max \{ V^A(w), V^R \} \]

\textbf{Bài toán số học:} Giả sử hệ số chiết khấu \$ \beta = 0.9 \$, trợ
cấp thất nghiệp \$ c = 20 \$. Phân phối lương có 2 khả năng với xác suất
50-50: \$ w\_1 = 10 \$, \$ w\_2 = 50 \$. Tìm giá trị bảo lưu
(reservation wage) \$ w\^{}* \$ để quyết định. Nếu nhận \$ w=10 \$: \$
V\^{}A(10) = \frac{10}{1 - 0.9} = 100 \$. Nếu nhận \$ w=50 \$: \$
V\^{}A(50) = \frac{50}{1 - 0.9} = 500 \$. Người lao động chắc chắn từ
chối \$ w\_1 = 10 \$ vì tiền trợ cấp \$ c=20 \$ còn cao hơn. Giá trị của
việc chờ đợi \$ V\^{}R \$ (giả định từ chối \$ w\_1 \$, nhận \$ w\_2
\$):
\[ V^R = 20 + 0.9 \left[ 0.5 \times V^R + 0.5 \times V^A(50) \right] \]
\[ V^R = 20 + 0.45 V^R + 0.45(500) \implies 0.55 V^R = 20 + 225 = 245 \implies V^R \approx 445.45 \]
Vì \$ V\^{}A(50) = 500 \textgreater{} V\^{}R \$, người lao động sẽ nhận
công việc có mức lương 50. Ngưỡng lương bảo lưu \$ w\^{}* = V\^{}R
(1-\beta) = 445.45 \times 0.1 \approx 44.55 \$. Điều này chứng minh bằng
toán học rằng người lao động sẽ chỉ chấp nhận các lời đề nghị công việc
có mức lương cao hơn 44.55, giải thích tỷ lệ thất nghiệp tự nhiên trong
nền kinh tế.

\emph{Tài liệu này hoàn chỉnh hệ thống kiến thức kinh tế học từ các
nguyên lý cơ bản đến các khung lý thuyết nâng cao trình độ Tiến sĩ.}

\section{PHỤ LỤC C: NGÂN HÀNG CÂU HỎI TỰ KIỂM TRA \& ÔN LUYỆN TỪ ĐIỂM MÙ
NHẬN THỨC (Self-Assessment
Bank)}\label{phux1ee5-lux1ee5c-c-nguxe2n-huxe0ng-cuxe2u-hux1ecfi-tux1ef1-kiux1ec3m-tra-uxf4n-luyux1ec7n-tux1eeb-ux111iux1ec3m-muxf9-nhux1eadn-thux1ee9c-self-assessment-bank}

Phụ lục này tổng hợp các câu hỏi phân tích biện luận chuyên sâu được rút
ra trực tiếp từ \textbf{các Điểm mù Nhận thức (Cognitive Traps)} trong
17 Chương của giáo trình. Đây là công cụ tự kiểm tra lý tưởng cho nghiên
cứu sinh tiến sĩ (PhD) và các chuyên gia hoạch định chính sách kinh tế.

\begin{enumerate}
\def\labelenumi{\arabic{enumi}.}
\tightlist
\item
  \textbf{(Chương 1)} Tại sao giá trần nhằm bảo vệ người tiêu dùng nghèo
  thường dẫn đến tổn thất vô ích (DWL) và sự phân phối bất công qua thị
  trường chợ đen hoặc xếp hàng?
\item
  \textbf{(Chương 2 \& 8)} Tại sao Định lý Coase (1960) thường thất bại
  trong việc giải quyết ngoại tác biến đổi khí hậu toàn cầu dù quyền sở
  hữu tài sản đã được xác định?
\item
  \textbf{(Chương 3 \& 6)} Khi nào một cá nhân vi phạm tiên đề Độc lập
  (Independence Axiom) của von Neumann-Morgenstern? Phân tích hiện tượng
  Nghịch lý Allais và cách Lý thuyết Triển vọng (Prospect Theory) giải
  quyết nghịch lý này.
\item
  \textbf{(Chương 4)} Theo mô hình Melitz (2003), tại sao việc tự do hóa
  thương mại lại khiến năng suất trung bình của ngành tăng lên dù công
  nghệ sản xuất của từng doanh nghiệp không đổi?
\item
  \textbf{(Chương 5 \& 11)} Trong trò chơi lặp vô hạn, tại sao sự tồn
  tại của đa cân bằng theo Định lý Dân gian (Folk Theorem) lại làm giảm
  sức mạnh dự báo của khái niệm Cân bằng Nash hoàn hảo cục bộ (SPNE)?
\item
  \textbf{(Chương 7 \& 13)} Phân tích sự sụp đổ của Định lý tương đương
  Ricardo (Ricardian Equivalence) khi chuyển từ mô hình Ramsey hộ gia
  đình vô hạn kỳ sang mô hình thế hệ gối đầu OLG của Diamond (1965).
\item
  \textbf{(Chương 8)} Theo công thức Mirrlees-Saez, tại sao việc áp đặt
  thuế suất biên thu nhập 90\% lên giới siêu giàu có thể làm giảm nguồn
  thu ngân sách để tái phân phối cho người nghèo nếu độ co giãn
  trốn/tránh thuế (ETI) cao?
\item
  \textbf{(Chương 9)} Trong mô hình tìm kiếm việc làm
  Mortensen-Pissarides, tại sao trợ cấp thất nghiệp quá cao có thể làm
  dịch chuyển đường cong Beveridge ra ngoài và làm tăng tỷ lệ thất
  nghiệp tự nhiên?
\item
  \textbf{(Chương 10)} Tại sao cân bằng Nash trong trò chơi rút tiền
  ngân hàng theo mô hình Diamond-Dybvig (1983) lại có thể dẫn đến sự sụp
  đổ của một ngân hàng hoàn toàn dung mạo và có khả năng thanh toán
  (solvent)?
\item
  \textbf{(Chương 12)} Phân tích điều kiện Đơn điệu Spence-Mirrlees
  trong thiết kế hợp đồng sàng lọc (Screening). Tại sao hợp đồng tối ưu
  luôn gây ra sự méo mó (distortion) đối với nhóm có năng suất thấp để
  ngăn cản nhóm năng suất cao hưởng địa tô thông tin?
\item
  \textbf{(Chương 14)} Tại sao sự thịnh vượng kinh tế không thể duy trì
  dài hạn bằng việc đơn thuần tăng vốn đầu tư (theo Solow) nếu thiếu các
  thể chế bao trùm (Inclusive Institutions theo Acemoglu \& Robinson)?
\item
  \textbf{(Chương 15)} Phân tích sai lầm của phép hồi quy Fixed Effects
  (FE) khi biến phụ thuộc trong quá khứ có tương quan động với biến giải
  thích (Nickell bias 1981). Khi nào cần dùng kiểm định Hausman bền vững
  thay cho kiểm định Hausman truyền thống?
\item
  \textbf{(Chương 15)} Tại sao việc sử dụng mô hình Two-Way Fixed
  Effects (TWFE) truyền thống lại tạo ra trọng số âm (negative weights)
  và ước lượng sai tác động nhân quả khi các nhóm đối tượng được can
  thiệp ở các thời điểm khác nhau (Staggered DiD theo Callaway \&
  Sant'Anna 2021)?
\item
  \textbf{(Chương 17)} Theo Rochet \& Tirole (2003), tại sao một nền
  tảng thị trường hai mặt (Two-Sided Platform) có thể đặt mức giá dưới
  chi phí biên (hoặc bằng 0) cho một bên người dùng mà vẫn tối đa hóa
  lợi nhuận chung mà không mang tính bán phá giá (predatory pricing)?
\item
  \textbf{(Phụ lục A)} Tại sao Định lý Bất khả thi Arrow (1951) chứng
  minh rằng không tồn tại bất kỳ một hàm phúc lợi xã hội nào có thể
  chuyển hóa sở thích cá nhân thành sở thích tập thể mà thỏa mãn đồng
  thời 4 tiên đề tự do hợp lý?
\end{enumerate}

\begin{center}\rule{0.5\linewidth}{0.5pt}\end{center}

\begin{enumerate}
\def\labelenumi{\arabic{enumi}.}
\setcounter{enumi}{15}
\tightlist
\item
  \textbf{(Chương 16)} Dưới lăng kính của Kinh tế học Thần kinh
  (Neuroeconomics), tại sao con người lại có xu hướng vi phạm các tiên
  đề hợp lý truyền thống khi đối mặt với rủi ro và phần thưởng? Vai trò
  của hệ thống Dopamine trong các mô hình học hỏi củng cố (Reinforcement
  Learning) là gì?
\item
  \textbf{(Phụ lục B)} Trong quy hoạch động, tại sao nguyên lý Bellman
  lại cho phép chia một bài toán tối ưu chuỗi thời gian vô hạn kỳ thành
  một chuỗi các quyết định hai thời kỳ (hiện tại và tương lai)? Khi nào
  thì toán tử Bellman không phải là một ánh xạ co (contraction mapping)?
\end{enumerate}

\section{DANH MỤC TÀI LIỆU THAM KHẢO TỔNG HỢP (COMPREHENSIVE
BIBLIOGRAPHY)}\label{danh-mux1ee5c-tuxe0i-liux1ec7u-tham-khux1ea3o-tux1ed5ng-hux1ee3p-comprehensive-bibliography}

\begin{enumerate}
\def\labelenumi{\arabic{enumi}.}
\tightlist
\item
  \textbf{Acemoglu, D., \& Robinson, J. A. (2012).} \emph{Why Nations
  Fail: The Origins of Power, Prosperity, and Poverty.} Crown Business.
\item
  \textbf{Akerlof, G. A. (1970).} The Market for ``Lemons'': Quality
  Uncertainty and the Market Mechanism. \emph{The Quarterly Journal of
  Economics}, 84(3), 488-500.
\item
  \textbf{Angrist, J. D., \& Pischke, J.-S. (2009).} \emph{Mostly
  Harmless Econometrics: An Empiricist's Companion.} Princeton
  University Press.
\item
  \textbf{Arellano, M., \& Bond, S. (1991).} Some Tests of Specification
  for Panel Data: Monte Carlo Evidence and an Application to Employment
  Equations. \emph{The Review of Economic Studies}, 58(2), 277-297.
\item
  \textbf{Arrow, K. J. (1951).} \emph{Social Choice and Individual
  Values.} Yale University Press.
\item
  \textbf{Arrow, K. J., \& Debreu, G. (1954).} Existence of an
  Equilibrium for a Competitive Economy. \emph{Econometrica}, 22(3),
  265-290.
\item
  \textbf{Becker, G. S. (1964).} \emph{Human Capital: A Theoretical and
  Empirical Analysis, with Special Reference to Education.} National
  Bureau of Economic Research.
\item
  \textbf{Bellman, R. (1957).} \emph{Dynamic Programming.} Princeton
  University Press.
\item
  \textbf{Bertrand, M., Duflo, E., \& Mullainathan, S. (2004).} How Much
  Should We Trust Differences-in-Differences Estimates? \emph{The
  Quarterly Journal of Economics}, 119(1), 249-275.
\item
  \textbf{Black, D. (1948).} On the Rationale of Group Decision-making.
  \emph{Journal of Political Economy}, 56(1), 23-34.
\item
  \textbf{Callaway, B., \& Sant'Anna, P. H. (2021).}
  Difference-in-Differences with Multiple Time Periods. \emph{Journal of
  Econometrics}, 225(2), 200-230.
\item
  \textbf{Cameron, A. C., \& Miller, D. L. (2015).} A Practitioner's
  Guide to Cluster-Robust Inference. \emph{Journal of Human Resources},
  50(2), 317-372.
\item
  \textbf{Coase, R. H. (1960).} The Problem of Social Cost.
  \emph{Journal of Law and Economics}, 3, 1-44.
\item
  \textbf{Diamond, P. A. (1965).} National Debt in a Neoclassical Growth
  Model. \emph{The American Economic Review}, 55(5), 1126-1150.
\item
  \textbf{Diamond, D. W., \& Dybvig, P. H. (1983).} Bank Runs, Deposit
  Insurance, and Liquidity. \emph{Journal of Political Economy}, 91(3),
  401-419.
\item
  \textbf{Diamond, P. A., \& Mirrlees, J. A. (1971).} Optimal Taxation
  and Public Production I: Production Efficiency. \emph{The American
  Economic Review}, 61(1), 8-27.
\item
  \textbf{Fama, E. F., \& French, K. R. (1993).} Common Risk Factors in
  the Returns on Stocks and Bonds. \emph{Journal of Financial
  Economics}, 33(1), 3-56.
\item
  \textbf{Fudenberg, D., \& Maskin, E. (1986).} The Folk Theorem in
  Repeated Games with Discounting or with Incomplete Information.
  \emph{Econometrica}, 54(3), 533-554.
\item
  \textbf{Green, E. J., \& Porter, R. H. (1984).} Noncooperative
  Collusion under Imperfect Price Information. \emph{Econometrica},
  52(1), 87-100.
\item
  \textbf{Hausman, J. A. (1978).} Specification Tests in Econometrics.
  \emph{Econometrica}, 46(6), 1251-1271.
\item
  \textbf{Hotelling, H. (1929).} Stability in Competition. \emph{The
  Economic Journal}, 39(153), 41-57.
\item
  \textbf{Kahneman, D., \& Tversky, A. (1979).} Prospect Theory: An
  Analysis of Decision under Risk. \emph{Econometrica}, 47(2), 263-291.
\item
  \textbf{Kaldor, N. (1939).} Welfare Propositions of Economics and
  Inter-Personal Comparisons of Utility. \emph{The Economic Journal},
  49(195), 549-552.
\item
  \textbf{Melitz, M. J. (2003).} The Impact of Trade on Intra-Industry
  Reallocations and Aggregate Industry Productivity.
  \emph{Econometrica}, 71(6), 1695-1725.
\item
  \textbf{Mirrlees, J. A. (1971).} An Exploration in the Theory of
  Optimum Income Taxation. \emph{The Review of Economic Studies}, 38(2),
  175-208.
\item
  \textbf{Modigliani, F., \& Miller, M. H. (1958).} The Cost of Capital,
  Corporation Finance and the Theory of Investment. \emph{The American
  Economic Review}, 48(3), 261-297.
\item
  \textbf{Mortensen, D. T., \& Pissarides, C. A. (1994).} Job Creation
  and Job Destruction in the Theory of Unemployment. \emph{The Review of
  Economic Studies}, 61(3), 397-415.
\item
  \textbf{Oates, W. E. (1972).} \emph{Fiscal Federalism.} Harcourt Brace
  Jovanovich.
\item
  \textbf{Pigou, A. C. (1920).} \emph{The Economics of Welfare.}
  Macmillan and Co.
\item
  \textbf{Ramsey, F. P. (1927).} A Contribution to the Theory of
  Taxation. \emph{The Economic Journal}, 37(145), 47-61.
\item
  \textbf{Rochet, J.-C., \& Tirole, J. (2003).} Platform Competition in
  Two-Sided Markets. \emph{Journal of the European Economic
  Association}, 1(4), 990-1029.
\item
  \textbf{Romer, P. M. (1990).} Endogenous Technological Change.
  \emph{Journal of Political Economy}, 98(5), S71-S102.
\item
  \textbf{Rotemberg, J. J., \& Saloner, G. (1986).} A
  Supergame-Theoretic Model of Price Wars during Booms. \emph{The
  American Economic Review}, 76(3), 390-407.
\item
  \textbf{Saez, E. (2001).} Using Elasticities to Derive Optimal Income
  Tax Rates. \emph{The Review of Economic Studies}, 68(1), 205-229.
\item
  \textbf{Samuelson, P. A. (1954).} The Pure Theory of Public
  Expenditure. \emph{The Review of Economics and Statistics}, 36(4),
  387-389.
\item
  \textbf{Solow, R. M. (1956).} A Contribution to the Theory of Economic
  Growth. \emph{The Quarterly Journal of Economics}, 70(1), 65-94.
\item
  \textbf{Spence, M. (1973).} Job Market Signaling. \emph{The Quarterly
  Journal of Economics}, 87(3), 355-374.
\item
  \textbf{Sun, L., \& Abraham, S. (2021).} Estimating Dynamic Treatment
  Effects in Event Studies with Heterogeneous Treatment Effects.
  \emph{Journal of Econometrics}, 225(2), 175-199.
\item
  \textbf{Tiebout, C. M. (1956).} A Pure Theory of Local Expenditures.
  \emph{Journal of Political Economy}, 64(5), 416-424.
\end{enumerate}

\begin{itemize}
\tightlist
\item
  \textbf{Aghion, P., \& Howitt, P. (1992).} A Model of Growth Through
  Creative Destruction. \emph{Econometrica}, 60(2), 323-351.
\item
  \textbf{Black, F., \& Scholes, M. (1973).} The Pricing of Options and
  Corporate Liabilities. \emph{Journal of Political Economy}, 81(3),
  637-654.
\item
  \textbf{Cournot, A. A. (1838).} \emph{Recherches sur les principes
  mathématiques de la théorie des richesses}.
\item
  \textbf{Gale, D., \& Shapley, L. S. (1962).} College Admissions and
  the Stability of Marriage. \emph{The American Mathematical Monthly},
  69(1), 9-15.
\item
  \textbf{Kreps, D. M., \& Wilson, R. (1982).} Reputation and Imperfect
  Information. \emph{Journal of Economic Theory}, 27(2), 253-279.
\item
  \textbf{Krugman, P. R. (1980).} Scale Economies, Product
  Differentiation, and the Pattern of Trade. \emph{The American Economic
  Review}, 70(5), 950-959.
\item
  \textbf{Maskin, E., \& Tirole, J. (2001).} Markov Perfect Equilibrium:
  I. Observable Actions. \emph{Journal of Economic Theory}, 100(2),
  191-219.
\item
  \textbf{Mokyr, J. (1990).} \emph{The Lever of Riches: Technological
  Creativity and Economic Progress.} Oxford University Press.
\item
  \textbf{Nash, J. F. (1950).} Equilibrium points in n-person games.
  \emph{Proceedings of the National Academy of Sciences}, 36(1), 48-49.
\item
  \textbf{Roth, A. E. (1984).} The Evolution of the Labor Market for
  Medical Interns and Residents: A Case Study in Game Theory.
  \emph{Journal of Political Economy}, 92(6), 991-1016.
\item
  \textbf{Selten, R. (1965).} Spieltheoretische Behandlung eines
  Oligopolmodells mit Nachfrageträgheit. \emph{Zeitschrift für die
  gesamte Staatswissenschaft}, 121(2), 301-324.
\item
  \textbf{Sonnenschein, H. (1972).} Market Excess Demand Functions.
  \emph{Econometrica}, 40(3), 549-563.
\item
  \textbf{Taylor, J. B. (1993).} Discretion versus policy rules in
  practice. \emph{Carnegie-Rochester Conference Series on Public
  Policy}, 39, 195-214.
\item
  \textbf{Vickrey, W. (1961).} Counterspeculation, Auctions, and
  Competitive Sealed Tenders. \emph{The Journal of Finance}, 16(1),
  8-37.
\item
  \textbf{Von Neumann, J., \& Morgenstern, O. (1944).} \emph{Theory of
  Games and Economic Behavior.} Princeton University Press.
\end{itemize}

\end{document}
